% Complete standalone editable source corresponding to the EGA III-2 zh-Hans-CN reader.
% Generated deterministically from the sealed producer sources; no authority scan is included.
\documentclass[UTF8,fontset=none,11pt,oneside,openany]{ctexbook}

\usepackage{iftex}
\usepackage{fontspec}
\ifXeTeX
  \usepackage{xeCJK}
\fi
\defaultfontfeatures{Ligatures=TeX}
\setmainfont{texgyretermes-regular.otf}[
  BoldFont=texgyretermes-bold.otf,
  ItalicFont=texgyretermes-italic.otf,
  BoldItalicFont=texgyretermes-bolditalic.otf
]
\setsansfont{texgyreheros-regular.otf}
\setmonofont{lmmono10-regular.otf}[ItalicFont=lmmono10-italic.otf]
\setCJKmainfont{FandolSong-Regular.otf}[
  BoldFont=FandolSong-Bold.otf,
  ItalicFont=FandolKai-Regular.otf,
  BoldItalicFont=FandolSong-Bold.otf
]
\setCJKsansfont{FandolHei-Regular.otf}
\setCJKmonofont{FandolFang-Regular.otf}
\ifXeTeX
  \xeCJKsetup{PunctStyle=kaiming,CheckSingle=true}
  \XeTeXlinebreaklocale "zh"
  \XeTeXlinebreakskip=0pt plus 1pt
\fi

\usepackage{amsmath}
\ifLuaTeX
  \usepackage{unicode-math}
  \setmathfont{texgyretermes-math.otf}
  \providecommand{\square}{\mdlgwhtsquare}
\else
  \usepackage{amssymb,mathrsfs}
\fi
\ifLuaTeX
  \renewcommand{\leadsto}{\mathrel{\rightsquigarrow}}
\fi
\usepackage{graphicx}
\usepackage{tikz-cd}
\ifLuaTeX
  \usepackage{luatex85}
\fi
\usepackage[all]{xy}
\ifLuaTeX
  % Preserve the established EGA Xy-pic stroke geometry under OpenType math.
  \makeatletter
  \AtBeginDocument{%
    \edef\xP@preclw{0.43799pt}%
    \edef\xP@lw{\xP@dim\xP@preclw}%
  }
  \makeatother
\fi
\usepackage[shortlabels]{enumitem}
\usepackage{marginnote}
\newcommand{\SourceCorrectionNote}[1]{\footnote{#1}}
\ifLuaTeX\else\usepackage{microtype}\fi
\usepackage{fancyhdr}
\usepackage[
  a4paper,
  inner=26mm,
  outer=24mm,
  top=24mm,
  bottom=24mm,
  headheight=14pt,
  headsep=8mm,
  footskip=12mm
]{geometry}
\usepackage[
  unicode=true,
  colorlinks=true,
  linkcolor=black,
  citecolor=black,
  urlcolor=blue,
  hypertexnames=false
]{hyperref}

\hypersetup{
  pdftitle={代数几何原理 III：上同调研究（大陆简体中文工作版，非典范检查点）},
  pdfauthor={Alexander Grothendieck；Jean Dieudonné（协作编写）},
  pdfsubject={依据 NUMDAM 法文印本与法文外交式 TeX 制作的大陆简体中文版本},
  pdfkeywords={代数几何；上同调；EGA；简体中文}
}

\AtBeginDocument{\fontsize{10.95pt}{16.5pt}\selectfont}
\setlength{\parindent}{2em}
\setlength{\parskip}{0pt}
\setlength{\emergencystretch}{2em}
\flushbottom

\pagestyle{fancy}
\fancyhf{}
\fancyhead[L]{\small 代数几何原理 III}
\fancyhead[R]{\small 上同调研究}
\fancyfoot[C]{\thepage}
\renewcommand{\headrulewidth}{0.4pt}
\renewcommand{\footrulewidth}{0pt}
\fancypagestyle{plain}{%
  \fancyhf{}%
  \fancyhead[L]{\small 代数几何原理 III}%
  \fancyhead[R]{\small 上同调研究}%
  \fancyfoot[C]{\thepage}%
  \renewcommand{\headrulewidth}{0.4pt}%
  \renewcommand{\footrulewidth}{0pt}%
}

\renewcommand{\contentsname}{目录}
\renewcommand{\bibname}{参考文献}
\renewcommand{\figurename}{图}
\renewcommand{\tablename}{表}
\renewcommand{\thesection}{\arabic{section}}
\renewcommand{\thesubsection}{\thesection.\arabic{subsection}}

\newenvironment{env}[1][]{%
  \par\medskip\noindent\textbf{(#1)}\ %
}{\par}
\newenvironment{definition}[1][]{%
  \par\medskip\noindent\textbf{定义 (#1).}\ %
}{\par}
\newenvironment{lemma}[1][]{%
  \par\medskip\noindent\textbf{引理 (#1).}\ %
}{\par}
\newenvironment{theorem}[1][]{%
  \par\medskip\noindent\textbf{定理 (#1).}\ %
}{\par}
\newenvironment{proposition}[1][]{%
  \par\medskip\noindent\textbf{命题 (#1).}\ %
}{\par}
\newenvironment{corollary}[1][]{%
  \par\medskip\noindent\textbf{推论 (#1).}\ %
}{\par}
\newenvironment{remark}[1][]{%
  \par\medskip\noindent\textbf{注 (#1).}\ %
}{\par}
\newenvironment{remarks}[1][]{%
  \par\medskip\noindent\textbf{注 (#1).}\ %
}{\par}
\newenvironment{notation}[1][]{%
  \par\medskip\noindent\textbf{记号 (#1).}\ %
}{\par}
\newenvironment{proof}{%
  \par\smallskip\noindent\textit{证明.}\ %
}{\hfill$\square$\par}

\newcommand{\oldpage}[2][]{%
  \hypertarget{ega-source-page-#1-#2}{}%
  \marginnote{\scriptsize\textbf{#1}~|~#2}\ignorespaces}

\let\EGAOriginalLabel\label
\renewcommand{\label}[1]{\phantomsection\EGAOriginalLabel{#1}}

\newenvironment{namedtheorem}[2][]{%
  \par\medskip\noindent\textbf{#2 (#1).}\ %
}{\par}
\newenvironment{scholium}[1][]{%
  \par\medskip\noindent\textbf{附论 (#1).}\ %
}{\par}

\begin{document}
\setcounter{page}{137}
\setcounter{section}{5}
% EGA III-2 mainland Simplified-Chinese producer source.
% Basis: Canon keeper R66 corrected-current/diplomatic pair; source-literal unless an admitted reversible disposition applies.
% This file is a producer draft and is uncertified pending independent replay.
\chapter*{第三章（续）\\[1.2ex]
\Large 相干层的上同调研究}

\section{局部与整体 Tor 函子；Künneth 公式}
\label{section:III.6.fr}

\subsection{引言}
\label{subsection:III.6.1.fr}

\oldpage[III]{137}

\begin{env}[6.1.1]
\label{III.6.1.1-fr}
设 $f:X\to Y$ 是预概形的一个态射，$\mathscr F$ 是拟相干
$\mathscr O_X$-模。在研究“高阶直接像”
$\mathrm R^n f_*(\mathscr F)$ 时，会遇到如下的一般问题：给定一个“基变换”态射
$g:Y'\to Y$，置
\[
  X'=X_{(Y')}=X\times_Y Y',\qquad
  \mathscr F'=\mathscr F\otimes_Y\mathscr O_{Y'},\qquad
  f'=f_{(Y')}:X'\to Y',
\]
并希望得到关于高阶直接像
$\mathrm R^n f'_*(\mathscr F')$ 的信息（假定已知
$\mathrm R^n f_*(\mathscr F)$）。容易看出（例如参见（7.7.2）），归根结底要研究
$\mathrm R^n f_*(\mathscr F\otimes_{\mathscr O_Y}\mathscr G)$ 的变化，其中
$\mathscr O_Y$-模 $\mathscr G$ 是任意拟相干模；换言之，要研究函子
\[
  \mathscr G\longmapsto
  \mathrm R^n f_*(\mathscr F\otimes_{\mathscr O_Y}\mathscr G)
\]
。
若 $\mathscr F$ 在 $Y$ 上平坦，则函子
$\mathscr G\mapsto\mathscr F\otimes_{\mathscr O_Y}\mathscr G$ 是正合的
（（0\textsubscript{I}, 6.7.4）），因而复合函子
$\mathscr G\mapsto\mathrm R^n f_*(\mathscr F\otimes_{\mathscr O_Y}\mathscr G)$
仍是一个上同调函子。但一般情形并非如此；为了能够应用上同调方法，必须用其他函子替代
$\mathscr G\mapsto\mathrm R^n f_*(\mathscr F\otimes_{\mathscr O_Y}\mathscr G)$，
而这些函子始终是上同调函子。这些函子推广了模理论中的“Tor”函子，在第 6.3 至 6.7 节定义；
此外有两种这样的推广，即“局部的”和“整体的”，它们由谱序列联系起来；这些谱序列将在第 6.7 节讨论。
作为谱序列的应用，在一定条件下得到一个“Künneth 公式”，它表示
\[
  \mathrm R^n(f_1\times f_2)_*(\mathscr F_1\otimes_Y\mathscr F_2)
\]
，其中用到高阶直接像 $\mathrm R^p f_{1*}(\mathscr F_1)$ 与
$\mathrm R^q f_{2*}(\mathscr F_2)$。
其他谱序列（6.8）推广了模的“Tor”函子的结合律谱序列；最后，基变换问题本身也导出谱序列（6.9）。
\end{env}

\begin{env}[6.1.2]
\label{III.6.1.2-fr}
此外，由（6.10）可见，如此定义的上同调函子
$\mathscr G\mapsto\mathscr T_\bullet(\mathscr G)$ 在 $Y$ 上局部具有如下形式：
\[
  \mathscr G\longmapsto
  \mathscr H_\bullet(\mathscr L_\bullet\otimes_Y\mathscr G)
\]
，
其中 $\mathscr L_\bullet$ 是局部自由的 $\mathscr O_Y$-模复形（在同伦意义下定义），
而 $\mathscr H_\bullet$ 是同调。于是，有必要忘却产生 $\mathscr T_\bullet$ 的特殊情形，

\oldpage[III]{138}

一般地研究上述形式的函子 $\mathscr G\mapsto\mathscr H_\bullet(\mathscr L_\bullet\otimes_Y\mathscr G)$
（此外，在适当时还要对 $\mathscr L_i$ 或 $\mathscr H_i(\mathscr L_\bullet)$ 作适当的有限性假设）：
这正是第 7 节的对象；原则上，该节的阅读独立于第 6 节。
这些函子最重要的性质涉及分量 $\mathscr T_i$ 的正合性（它是 $\mathscr T_\bullet$ 的一个分量）；
我们将给出若干判据来建立这类性质。作为应用，我们将得到一些条件，使得能够断言（采用（6.1.1）的记号）函子
$\mathscr G\mapsto\mathrm R^n f_*(\mathscr F\otimes_{\mathscr O_Y}\mathscr G)$
是正合的（这称为 $\mathscr F$ 在 $Y$ 上第 $n$ 维上同调平坦）。
对于分量 $\mathscr T_i$，其中
$\mathscr T_\bullet(\mathscr G)=
\mathscr H_\bullet(\mathscr L_\bullet\otimes_Y\mathscr G)$，
另一个重要性质是函数
$y\mapsto\dim_{k(y)}(T_i(k(y)))$ 的半连续性；当 $\mathscr T_i$ 正合时，
该性质被连续性性质取代；反之，当 $Y$ 约化时，由 Grauert 定理（7.8.4）知连续性也推出正合性。
\end{env}

\begin{env}[6.1.3]
\label{III.6.1.3-fr}
在第 6、7 节中，我们始终使用超上同调：处处取层复形而非层作为自变量，
尽管直到本论后续章节才显出这种观点的必要性。借此发展出的上同调形式主义，
在本论专门讨论相干层上同调函子代数（包括对偶形式主义）的章节中会变得更加清晰。
不过，这将需要超出本章范围的展开。
\end{env}

\begin{env}[6.1.4]
\label{III.6.1.4-fr}
为简便起见，设 $K^\bullet$、$K'^\bullet$ 是阿贝尔范畴 $C$ 中的两个复形。
如果复形态射 $f:K^\bullet\to K'^\bullet$ 存在一个复形态射
$g:K'^\bullet\to K^\bullet$，使复合态射 $f\circ g$ 与 $g\circ f$ 都同伦于恒等态射，
则称该态射为一个同伦态射（滥用术语地说，在存在这样的同伦态射时，
也称 $K^\bullet$ 与 $K'^\bullet$ 同伦）。当可以定义加性协变函子 $T$（从阿贝尔范畴
$C$ 到阿贝尔范畴 $C'$）相对于 $C$ 中一个复形的超上同调时（（0，11.4.3）），
显然，复形同伦态射 $K^\bullet\to K'^\bullet$（其复形在 $C$ 中）对超上同调（同上引文）
自然定义一个同构
\[
  \mathrm R^\bullet T(K^\bullet)\xrightarrow{\sim}
  \mathrm R^\bullet T(K'^\bullet)
\]
。
\end{env}

\subsection{预概形上模复形的超上同调}
\label{subsection:III.6.2.fr}

\begin{env}[6.2.1]
设 $X$ 是一个预概形，
$\mathscr K^\bullet=(\mathscr K^i)_{i\in\mathbf Z}$ 是一个
$\mathscr O_X$-模复形，其微分算子次数为 $+1$。回顾一下：对于任意预概形态射
$f:X\to Y$，我们在（0，12.4.1）中定义了超上同调的
$\mathscr O_Y$-模
$\mathscr H^n(f,\mathscr K^\bullet)$（也记作
$\mathscr H_f^n(\mathscr K^\bullet)$ 或
$\mathrm R^n f_*(\mathscr K^\bullet)$），其中 $n\in\mathbf Z$ 任意；
超上同调 $\mathscr H^\bullet(f,\mathscr K^\bullet)$ 是两个谱函子
$'\mathscr E(f,\mathscr K^\bullet)$ 与 $''\mathscr E(f,\mathscr K^\bullet)$ 的极限，
它们的 $\mathscr E_2$ 项为
\begin{align}
  '\mathscr E_2^{pq}
    &=\mathscr H^p\bigl(\mathscr H^q(f,\mathscr K^\bullet)\bigr),
    \tag{6.2.1.1}\label{III.6.2.1.1-fr}\\
  ''\mathscr E_2^{pq}
    &=\mathscr H^p\bigl(f,\mathscr H^q(\mathscr K^\bullet)\bigr)
      =\mathrm R^p f_*\bigl(\mathscr H^q(\mathscr K^\bullet)\bigr).
    \tag{6.2.1.2}\label{III.6.2.1.2-fr}
\end{align}

\oldpage[III]{139}

其中 $\mathscr H^q(f,\mathscr K^\bullet)$ 是这样一个复形：其次数为
$i$ 的分量是
$\mathscr H^q(f,\mathscr K^i)=\mathrm R^q f_*(\mathscr K^i)$（同上引文）。
还要回顾：当 $Y$ 退化为一点时，相应的超上同调记作
$\mathbf H^\bullet(X,\mathscr K^\bullet)$
（它是定义在 $\Gamma(X,\mathscr O_X)$ 上的模，并且与所取的点概形 $Y$ 无关）；
当 $Y=X$ 且 $f=1_X$ 时，有
$\mathscr H^n(f,\mathscr K^\bullet)=\mathscr H^n(\mathscr K^\bullet)$
（复形 $\mathscr K^\bullet$ 的上同调）；当 $\mathscr K^i=0$，除了 $i=i_0$ 以外，均成立时，有
\[
  \mathscr H^n(f,\mathscr K^\bullet)
    =\mathrm R^{n-i_0}f_*(\mathscr K^{i_0}).
\]
谱序列 $'\mathscr E(f,\mathscr K^\bullet)$ 总是正则的；当
$\mathscr K^\bullet$ 下有界时，两个谱序列都是双正则的（0，12.4.1）。

复形的任意一个同伦态射
$h:\mathscr K^\bullet\to\mathscr K'^\bullet$，其中复形由 $\mathscr O_X$-模组成，
都会给出超上同调的同构（6.1.4）
$\mathscr H^\bullet(f,\mathscr K^\bullet)\xrightarrow{\sim}
\mathscr H^\bullet(f,\mathscr K'^\bullet)$。
即使只假设
$\mathscr H^\bullet(h):\mathscr H^\bullet(\mathscr K^\bullet)\to
\mathscr H^\bullet(\mathscr K'^\bullet)$ 是同构，并且
$\mathscr K^\bullet$ 与 $\mathscr K'^\bullet$ 下有界，结论仍然成立；
这是把（0，11.1.5）应用于谱序列（6.2.1.2）以及对
$\mathscr K'^\bullet$ 的相应谱序列后立即得到的。最后，对于任意开覆盖
$\mathfrak U=(U_\alpha)$（它是 $X$ 的开覆盖），我们也在（0，12.4.5）中定义了超上同调
$\mathbf H^\bullet(\mathfrak U,\mathscr K^\bullet)$，它是双复形
$C^\bullet(\mathfrak U,\mathscr K^\bullet)$ 的上同调（其指标为 $(i,j)$ 的分量按定义为
$C^i(\mathfrak U,\mathscr K^j)$）；
$\mathbf H^n(\mathfrak U,\mathscr K^\bullet)$ 仍然是
$\Gamma(X,\mathscr O_X)$ 上的模。
\end{env}

\begin{proposition}[6.2.2]
设 $X$ 是一个概形，$\mathfrak U=(U_\alpha)$ 是由 $X$ 的仿射开集组成的开覆盖。
对于任意复形 $\mathscr K^\bullet$，若它由拟相干 $\mathscr O_X$-模组成，则超上同调模
$\mathbf H^\bullet(X,\mathscr K^\bullet)$ 与
$\mathbf H^\bullet(\mathfrak U,\mathscr K^\bullet)$ 自然同构。
\end{proposition}

\begin{proof}
事实上，任意有限交集 $V$（由覆盖 $\mathfrak U$ 中的开集构成）都是仿射的
（I，5.5.6），因此 $\mathrm H^q(V,\mathscr K^i)=0$ 对任意 $i$ 及任意
$q>0$ 都成立（1.3.1）；所以该命题是（0，12.4.7）的一个特例。
\end{proof}

\begin{proposition}[6.2.3]
设 $f:X\to Y$ 是预概形的一个拟紧且分离的态射。对于任意复形
$\mathscr K^\bullet$，若它由拟相干 $\mathscr O_X$-模组成，则其
$\mathscr O_Y$-模 $\mathscr H^n(f,\mathscr K^\bullet)$ 都是拟相干的。
\end{proposition}

\begin{proof}
由于 $\mathscr H^q(f,\mathscr K^i)=\mathrm R^q f_*(\mathscr K^i)$
是拟相干的 $\mathscr O_Y$-模（1.4.10），所以
$'\mathscr E_2^{pq}$ 也具有这一性质；根据（6.2.1.1），它是
拟相干模态射的一个核的商模，再除以一个这类态射的像
（I，4.1.1）。同理，第一谱序列中的所有
$\mathscr O_Y$-模 $'\mathscr E_r^{pq}$、
$B_k('\mathscr E_2^{pq})$、$Z_k('\mathscr E_2^{pq})$ 都是拟相干的。
谱序列 $'\mathscr E(f,\mathscr K^\bullet)$ 的正则性说明
$Z_\infty('\mathscr E_2^{pq})$ 等于某个
$Z_k('\mathscr E_2^{pq})$，因而是拟相干的；同样，
$B_\infty('\mathscr E_2^{pq})=\varinjlim_k B_k('\mathscr E_2^{pq})$
也是拟相干的（0，11.2.4 及 I，4.1.1）；因此
$'\mathscr E_\infty^{pq}$ 也都是拟相干的。由于上述谱序列是正则的，
$F^p(\mathscr H^n(f,\mathscr K^\bullet))$ 的滤子是离散且穷尽的；
换句话说，$\mathscr O_Y$-模
$\mathscr H^n(f,\mathscr K^\bullet)$ 是一列递增的
$(\mathscr G_k)_{k\geq0}$ 的并，其中各项都是 $\mathscr O_Y$-模，并且
$\mathscr G_0=0$，且每个 $\mathscr G_k/\mathscr G_{k-1}$ 都等于某个
$\mathscr O_Y$-模 $'\mathscr E_\infty^{pq}$，因而是拟相干的。
对 $k$ 归纳，得到每个 $\mathscr G_k$ 都是拟相干的（1.4.17）；又因为
$\mathscr H^n(f,\mathscr K^\bullet)=\varinjlim_k\mathscr G_k$，命题得证
（I，4.1.1）。
\end{proof}
\oldpage[III]{140}

\begin{corollary}[6.2.4]
在（6.2.3）的假设下，对于任意仿射开集 $V$（它是 $Y$ 的开集），规范同态
\[
  \mathbf H^n(f^{-1}(V),\mathscr K^\bullet)
  \longrightarrow
  \Gamma(V,\mathscr H^n(f,\mathscr K^\bullet))
  \tag{6.2.4.1}\label{III.6.2.4.1-fr}
\]
对任意 $n\in\mathbf Z$ 都是双射。
\end{corollary}

\begin{proof}
证明与（1.4.11）的证明相同，只需使用（6.2.2），将其中的
$\mathscr F$ 换成 $\mathscr K^\bullet$，
$\mathscr K^\bullet$ 换成
$f_*(C^\bullet(\mathfrak U,\mathscr K^\bullet))$，
$\mathscr H^\bullet(\mathscr K^\bullet)$ 换成
$\mathscr H^\bullet(f,\mathscr K^\bullet)$，并注意到后者由（6.2.3）是
拟相干的 $\mathscr O_Y$-模。
\end{proof}

\begin{proposition}[6.2.5]
设 $Y$ 是局部诺特预概形，$f:X\to Y$ 是一个固有态射，且
$\mathscr K^\bullet$ 是一个 $\mathscr O_X$-模复形，使得
$\mathscr O_X$-模 $\mathscr H^q(\mathscr K^\bullet)$ 都是凝聚的。
那么 $\mathscr O_Y$-模 $\mathscr H^n(f,\mathscr K^\bullet)$ 都是凝聚的。
\end{proposition}

\begin{proof}
问题在 $Y$ 上是局部的，因而可限于 $Y$ 诺特且仿射的情形；于是根据（6.2.4），
只需证明 $\mathbf H^n(X,\mathscr K^\bullet)$ 是
$\Gamma(Y,\mathscr O_Y)$-有限生成模。此时由（6.2.2）有
$\mathbf H^\bullet(X,\mathscr K^\bullet)= \mathbf H^\bullet(\mathfrak U,\mathscr K^\bullet)$，
其中可设 $\mathfrak U$ 是有限覆盖，因为 $X$ 是拟紧的。每个复形
$C^\bullet(\mathfrak U,\mathscr K^j)$ 的上链按定义都是交错的，因此存在整数
$r>0$，使得 $C^i(\mathfrak U,\mathscr K^j)=0$ 对 $i<0$ 和 $i>r$ 都成立；
由此（0，11.3.3）可知双复形
$C^\bullet(\mathfrak U,\mathscr K^\bullet)$ 的两个谱序列都是双正则的。
由于 $\mathfrak U$ 中各集合的交集是仿射开集（I，5.5.6），函子
$\mathscr F\mapsto C^i(\mathfrak U,\mathscr F)$ 在拟相干
$\mathscr O_X$-模范畴中是正合的；因此
$\mathrm H_I^q(C^i(\mathfrak U,\mathscr K^\bullet))
=C^i(\mathfrak U,\mathscr H^q(\mathscr K^\bullet))$，而第二谱序列的 $E_2$ 项
（该谱序列属于双复形 $C^\bullet(\mathfrak U,\mathscr K^\bullet)$）由（0，11.3.2）给出为
\[
  ''E_2^{pq}
  =\mathrm H^p(C^\bullet(\mathfrak U,\mathscr H^q(\mathscr K^\bullet)))
  =\mathbf H^p(\mathfrak U,\mathscr H^q(\mathscr K^\bullet))
  =\mathrm H^p(X,\mathscr H^q(\mathscr K^\bullet))
\]
根据（1.4.1）；由于 $f$ 是固有态射，这些都是
$\Gamma(Y,\mathscr O_Y)$-有限生成模（3.2.1）。谱序列
$''E(C^\bullet(\mathfrak U,\mathscr K^\bullet))$ 是双正则的，故可推出
$\mathbf H^n(X,\mathscr K^\bullet)$ 是
$\Gamma(Y,\mathscr O_Y)$-有限生成模（0，11.1.8）。

特别地，若 $\mathscr K^\bullet$ 是 $\mathscr O_X$-模的复形且这些模凝聚，则诸
$\mathscr O_Y$-模 $\mathscr H^n(f,\mathscr K^\bullet)$ 在（6.2.5）关于
$Y$ 与 $f$ 的假设下都是凝聚的
（0\textsubscript{I}，5.3.4）。
\end{proof}

\begin{env}[6.2.6]
超上同调 $\mathscr H^\bullet(f,\mathscr K^\bullet)$ 是下有界
$\mathscr O_X$-模复形范畴中的上同调函子（0，12.4.4）。对于所有
$\mathscr O_X$-模复形，它在该范畴中也是上同调函子，只要态射
$f$ 拟紧且 $X$ 的底层空间局部诺特：事实上，这时由
（G，II，3.10.1）可知 $f_*$ 与归纳极限交换（问题在 $Y$ 上是局部的），
因而可以应用（0，11.5.2）。

最后，若 $f$ 是分离的，则
$\mathscr H^\bullet(f,\mathscr K^\bullet)$ 是拟相干
$\mathscr O_X$-模复形范畴中的上同调函子。当 $Y$ 仿射时这立即可见，
因为此时 $X$ 是一个概形；于是根据规范同构（6.2.2），只需说明
$\mathscr K^\bullet\mapsto\mathbf H^\bullet(\mathfrak U,\mathscr K^\bullet)$
是拟相干 $\mathscr O_X$-模复形范畴中的上同调函子，而这显然成立，
因为函子
$\mathscr K^\bullet\mapsto C^\bullet(\mathfrak U,\mathscr K^\bullet)$
在该范畴中是正合的（I，1.3.7）。在一般情形下，对于任意仿射开集 $V$（它是 $Y$ 的开集），

\oldpage[III]{141}

$f^{-1}(V)$ 是一个概形；为了应用上面的结论，只需验证：对于正合序列
\[
  0\longrightarrow\mathscr K'^\bullet
  \longrightarrow\mathscr K^\bullet
  \longrightarrow\mathscr K''^\bullet
  \longrightarrow0
\]
（其中各复形均由拟相干 $\mathscr O_X$-模组成），
同态
\[
  \partial:
  \mathscr H^n(f,\mathscr K''^\bullet)|V
  \longrightarrow
  \mathscr H^{n+1}(f,\mathscr K'^\bullet)|V
\]
不依赖于用于定义它的仿射开覆盖 $\mathfrak U$（它覆盖 $f^{-1}(V)$）。
但这是因为：若 $\mathfrak U'$ 是比 $\mathfrak U$ 更细的仿射开覆盖，
则由规范同构组成的图表
\[
  \xymatrix@C=34pt@R=11pt{
    \mathbf H^\bullet(\mathfrak U,\mathscr K^\bullet)
      \ar[rrd]^-{\sim}\ar[dd]^{\iota}
    && \\
    &&
    \mathbf H^\bullet(f^{-1}(V),\mathscr K^\bullet)
    \\
    \mathbf H^\bullet(\mathfrak U',\mathscr K^\bullet)
      \ar[rru]_-{\sim}
    &&
  }
\]
是交换的；并且下图也交换：
\[
  \xymatrix@C=46pt@R=25pt{
    \mathbf H^n(\mathfrak U,\mathscr K''^\bullet)
      \ar[r]^-{\partial}\ar[d]^{\iota}
    &
    \mathbf H^{n+1}(\mathfrak U,\mathscr K'^\bullet)
      \ar[d]^{\iota}
    \\
    \mathbf H^n(\mathfrak U',\mathscr K''^\bullet)
      \ar[r]_-{\partial}
    &
    \mathbf H^{n+1}(\mathfrak U',\mathscr K'^\bullet).
  }
\]

当上述条件之一成立，且
$\mathscr K^\bullet\to\mathscr K'^\bullet$ 是一个同伦态射（6.1.4）时，
相应的同构
$\mathscr H^\bullet(f,\mathscr K^\bullet)\xrightarrow{\sim}
\mathscr H^\bullet(f,\mathscr K'^\bullet)$
便是 $\partial$-函子的同构（0，11.4.4）。
\end{env}

\begin{env}[6.2.7]
以上一切对链复形 $\mathscr K_\bullet$ 同样自然成立（只需改变记号），其中各项是拟相干
$\mathscr O_X$-模；其微分算子次数为 $-1$。
只需考虑上链复形 $\mathscr K^\bullet=(\mathscr K^i)$，其中
$\mathscr K^i=\mathscr K_{-i}$ 对每个 $i\in\mathbf Z$ 成立。
\end{env}

\subsection{两个模复形的超 Tor}
\label{subsection:III.6.3.fr}

\begin{env}[6.3.1]
设 $A$ 为交换环，$P_\bullet$、$Q_\bullet$ 为两个 $A$-模链复形，
其微分算子次数为 $-1$；设 $L_{\bullet\bullet}$（分别为
$M_{\bullet\bullet}$）是 $P_\bullet$（分别为 $Q_\bullet$）的
Cartan--Eilenberg 投射消解（0，11.6.1）；于是
$L_{\bullet\bullet}\otimes_A M_{\bullet\bullet}$ 是一个双复形（第一次数与
第二次数分别取总和），其微分算子次数为 $-1$；其同调
$H_\bullet(L_{\bullet\bullet}\otimes_A M_{\bullet\bullet})$ 不依赖于所选的
Cartan--Eilenberg 消解 $L_{\bullet\bullet}$、$M_{\bullet\bullet}$，并按定义
称为双函子 $P_\bullet\otimes_A Q_\bullet$ 关于 $P_\bullet$、$Q_\bullet$ 的
超同调（0，11.6.5）。按定义置
\[
  \mathbf{Tor}_n^A(P_\bullet,Q_\bullet)
  =H_n(L_{\bullet\bullet}\otimes_A M_{\bullet\bullet})
  \tag{6.3.1.1}\label{III.6.3.1.1-fr}
\]

\oldpage[III]{142}

并称这个 $A$-模为指标为 $n$ 的、两个复形 $P_\bullet$、$Q_\bullet$ 的超 Tor。
已知，在下有界 $A$-模复形范畴中，$\mathbf{Tor}_n^A(P_\bullet,Q_\bullet)$
关于 $P_\bullet$、$Q_\bullet$ 构成同调双函子（0，11.6.5）。此外：
\end{env}

\begin{proposition}[6.3.2]
双函子 $\mathbf{Tor}_\bullet^A(P_\bullet,Q_\bullet)$ 是两个谱双函子
$'E(P_\bullet,Q_\bullet)$、$''E(P_\bullet,Q_\bullet)$ 的共同收敛项；它们的
$E_2$ 项为
\begin{align}
  'E^2_{pq}
    &=H_p\bigl(\operatorname{Tor}_q^A(P_\bullet,Q_\bullet)\bigr),
    \tag{6.3.2.1}\label{III.6.3.2.1-fr}\\
  ''E^2_{pq}
    &=\bigoplus_{q'+q''=q}
      \operatorname{Tor}_p^A\bigl(H_{q'}(P_\bullet),H_{q''}(Q_\bullet)\bigr).
    \tag{6.3.2.2}\label{III.6.3.2.2-fr}
\end{align}
其中，在（6.3.2.1）中，$\operatorname{Tor}_q^A(P_\bullet,Q_\bullet)$ 表示由
$A$-模 $\operatorname{Tor}_q^A(P_i,Q_j)$ 构成的双复形。谱序列（6.3.2.2）
总是正则的；若 $P_\bullet$ 与 $Q_\bullet$ 下有界，或者 $A$ 的上同调维数有限，
则谱序列（6.3.2.1）与（6.3.2.2）都是双正则的。
\end{proposition}

\begin{proof}
这是（0，11.6.5）的结果，因为当 $A$ 的上同调维数为有限的 $n$ 时，
每个 $A$-模都有长度为 $n$ 的投射消解（M，VI，2.1）。
\end{proof}

\begin{corollary}[6.3.3]
设 $P'_\bullet$、$Q'_\bullet$ 是两个 $A$-模复形，
$u:P_\bullet\to P'_\bullet$、$v:Q_\bullet\to Q'_\bullet$ 是两个复形同态。
若同态
$H_\bullet(u):H_\bullet(P_\bullet)\to H_\bullet(P'_\bullet)$、
$H_\bullet(v):H_\bullet(Q_\bullet)\to H_\bullet(Q'_\bullet)$
（分别由 $u$、$v$ 诱导）都是双射，则同态
$\mathbf{Tor}_\bullet^A(P_\bullet,Q_\bullet)\to
\mathbf{Tor}_\bullet^A(P'_\bullet,Q'_\bullet)$（由 $u$、$v$ 诱导）也是双射。
\end{corollary}

\begin{proof}
事实上，谱序列态射
$''E(P_\bullet,Q_\bullet)\to''E(P'_\bullet,Q'_\bullet)$（由 $u$、$v$ 诱导）
在 $E_2$ 项上是同构；结论由这些谱序列根据（6.3.2）（0，11.1.5）
具有正则性而得。
\end{proof}

\begin{proposition}[6.3.4]
设 $P_\bullet$、$Q_\bullet$ 是下有界的两个 $A$-模复形。设
$L_{\bullet\bullet}$（分别为 $M_{\bullet\bullet}$）是由 $A$-平坦模组成的双复形，
并且对每个 $i$，$L_{i,\bullet}$（分别为 $M_{i,\bullet}$）是 $P_i$
（分别为 $Q_i$）的消解。于是有规范同构
\[
  \mathbf{Tor}_\bullet^A(P_\bullet,Q_\bullet)
  \xrightarrow{\sim}H_\bullet(L_{\bullet\bullet}\otimes_AQ_\bullet)
  \xrightarrow{\sim}H_\bullet(P_\bullet\otimes_AM_{\bullet\bullet})
  \xrightarrow{\sim}H_\bullet(L_{\bullet\bullet}\otimes_AM_{\bullet\bullet})
  \tag{6.3.4.1}\label{III.6.3.4.1-fr}
\]
\end{proposition}

\begin{proof}
这是（0，11.6.5，（ii）与（iii））以及 $A$-平坦模定义的直接结果。
\end{proof}

\begin{namedtheorem}[6.3.5]{注}
(i) 采用（6.3.1）的记号，双复形
$L_{\bullet\bullet}\otimes_A M_{\bullet\bullet}$ 与
$M_{\bullet\bullet}\otimes_A L_{\bullet\bullet}$ 具有规范同构，因而有规范同构
$\mathbf{Tor}_\bullet^A(P_\bullet,Q_\bullet)\xrightarrow{\sim}
\mathbf{Tor}_\bullet^A(Q_\bullet,P_\bullet)$。

(ii) 若 $F$、$G$ 是两个 $A$-模，且 $P_\bullet$、$Q_\bullet$ 是两个由
$A$-模组成的复形，分别在次数 0 取 $F$、$G$ 而在其他次数为零，则两个投射消解
$L_\bullet$、$M_\bullet$（分别为 $F$、$G$ 的投射消解）补以零项后，可视为
$P_\bullet$、$Q_\bullet$ 的
Cartan--Eilenberg 消解。因此在此情形下
$\mathbf{Tor}_\bullet^A(P_\bullet,Q_\bullet)
=\operatorname{Tor}_\bullet^A(F,G)$。
\end{namedtheorem}

\begin{proposition}[6.3.6]
设 $(P_\bullet^\lambda)$、$(Q_\bullet^\mu)$ 是 $A$-模复形的两个滤过归纳系统；
则有规范同构
\[
  \varinjlim_{\lambda,\mu}
  \mathbf{Tor}_\bullet^A(P_\bullet^\lambda,Q_\bullet^\mu)
  \xrightarrow{\sim}
  \mathbf{Tor}_\bullet^A
  \left(\varinjlim_\lambda P_\bullet^\lambda,
        \varinjlim_\mu Q_\bullet^\mu\right).
  \tag{6.3.6.1}\label{III.6.3.6.1-fr}
\]
\end{proposition}

\begin{proof}
置 $P_\bullet=\varinjlim_\lambda P_\bullet^\lambda$、
$Q_\bullet=\varinjlim_\mu Q_\bullet^\mu$；由函子性显然，
$\mathbf{Tor}_\bullet^A(P_\bullet^\lambda,Q_\bullet^\mu)$ 构成一个归纳系统，
且映射
$\mathbf{Tor}_\bullet^A(P_\bullet^\lambda,Q_\bullet^\mu)\to
\mathbf{Tor}_\bullet^A(P_\bullet,Q_\bullet)$（由规范映射
$P_\bullet^\lambda\to P_\bullet$、$Q_\bullet^\mu\to Q_\bullet$ 诱导）
构成同态的归纳系统，从而得到规范同态（6.3.6.1）；更一般地，还得到规范同态
\[
  \varinjlim_{\lambda,\mu}''E(P_\bullet^\lambda,Q_\bullet^\mu)
  \longrightarrow ''E(P_\bullet,Q_\bullet)
\]

\oldpage[III]{143}

其在收敛项上诱导的同态就是（6.3.6.1）。此外，谱序列
$''E(P_\bullet,Q_\bullet)$ 是正则的（6.3.2）；由定义（0，11.1.7）及
（0，11.3.3）的证明，谱序列
$\varinjlim_{\lambda,\mu}''E(P_\bullet^\lambda,Q_\bullet^\mu)$ 也正则。
因此，根据（0，11.1.5），为证明（6.3.6.1）是双射，只需证明同态
\[
  \varinjlim_{\lambda,\mu}''E(P_\bullet^\lambda,Q_\bullet^\mu)
  \longrightarrow ''E(P_\bullet,Q_\bullet)
  \tag{6.3.6.2}\label{III.6.3.6.2-fr}
\]
在 $E^2$ 项上是双射。由于函子 $H_\bullet$ 与模复形的归纳极限交换，
问题最终归结为证明：对于两个滤过归纳系统 $(F^\lambda)$、$(G^\mu)$（由 $A$-模组成），
规范同态
\[
  \varinjlim_{\lambda,\mu}\operatorname{Tor}_\bullet^A(F^\lambda,G^\mu)
  \longrightarrow
  \operatorname{Tor}_\bullet^A
  \left(\varinjlim_\lambda F^\lambda,
        \varinjlim_\mu G^\mu\right)
\]
是双射。为此，对每个 $F^\lambda$ 考虑规范自由消解
\[
  L_\bullet^\lambda:\quad
  \cdots\longrightarrow L_{i+1}^\lambda
  \longrightarrow L_i^\lambda\longrightarrow\cdots
  \longrightarrow L_1^\lambda\longrightarrow L_0^\lambda\longrightarrow0
\]
其中 $L_0^\lambda$ 是 $A$-模，其元素是 $F^\lambda$ 中元素的形式线性组合；
而 $L_{i+1}^\lambda$ 是 $A$-模，其元素是
$\operatorname{Ker}(L_i^\lambda\to L_{i-1}^\lambda)$ 中元素的形式线性组合；
立即可验证 $L_\bullet^\lambda$ 构成复形的归纳系统。置
$F=\varinjlim_\lambda F^\lambda$、$L_i=\varinjlim_\lambda L_i^\lambda$，则由于
诸 $L_i$ 构成消解 $L_\bullet$，它消解 $F$；这是因为归纳极限函子
$\varinjlim$ 是正合的。此外，诸 $L_i$ 作为自由 $A$-模的归纳极限是平坦的
（0\textsubscript{I}，6.1.2）。
对每个 $\mu$，同样考虑规范自由消解 $M_\bullet^\mu$，它消解 $G^\mu$；并且
$M_\bullet=\varinjlim_\mu M_\bullet^\mu$ 是
$G=\varinjlim_\mu G^\mu$ 的平坦消解。于是，由（6.3.5）和（6.3.4）有
$\operatorname{Tor}_\bullet^A(\varinjlim_\lambda F^\lambda,
\varinjlim_\mu G^\mu)=H_\bullet(L_\bullet\otimes_A M_\bullet)$。
但有
\[
  H_\bullet(L_\bullet\otimes_A M_\bullet)
  =\varinjlim_{\lambda,\mu}
   H_\bullet(L_\bullet^\lambda\otimes_A M_\bullet^\mu)
\]
，因为 $H_\bullet$ 与模复形的归纳极限交换；而
$H_\bullet(L_\bullet^\lambda\otimes_A M_\bullet^\mu)
=\operatorname{Tor}_\bullet^A(F^\lambda,G^\mu)$，故证明完成。

若假设存在 $i_0$，使得 $P_i^\lambda=Q_i^\mu=0$ 对 $i<i_0$ 以及任意
$\lambda$、$\mu$ 都成立，则同样可证明规范同态
\[
  \varinjlim_{\lambda,\mu}'E(P_\bullet^\lambda,Q_\bullet^\mu)
  \longrightarrow'E(P_\bullet,Q_\bullet)
  \tag{6.3.6.3}\label{III.6.3.6.3-fr}
\]
 是双射。
\end{proof}

\begin{proposition}[6.3.7]
设 $P_\bullet$ 与 $Q_\bullet$ 都是下有界的。若复形 $P_\bullet$ 由平坦的
$A$-模组成，则有一个规范 $A$-同构，它是 $\partial$-函子关于 $Q_\bullet$ 的同构：
\[
  \mathbf{Tor}_\bullet^A(P_\bullet,Q_\bullet)
  \xrightarrow{\sim}H_\bullet(P_\bullet\otimes_A Q_\bullet).
  \tag{6.3.7.1}\label{III.6.3.7.1-fr}
\]
\end{proposition}

\begin{proof}
事实上，谱序列（6.3.2.1）是双正则且退化的，而同构（6.3.7.1）的存在性由
（0，11.1.6）推出。此外，由 Cartan--Eilenberg 的 $P_\bullet$ 投射消解
（6.3.4）计算超 Tor，立即可见，如此定义的同构是 $\partial$-函子关于
$Q_\bullet$ 的同构。
\end{proof}

\begin{env}[6.3.8]
\oldpage[III]{144}
设 $\rho:A\to A'$ 是环同态。我们要定义一个 $A$-同态；它是 0 次函子性同态，
并由 $\rho$ 规范确定：
\[
  \rho_{P_\bullet,Q_\bullet}:
  \mathbf{Tor}_\bullet^A(P_\bullet,Q_\bullet)
  \longrightarrow
  \mathbf{Tor}_\bullet^{A'}
  (P_\bullet\otimes_A A',Q_\bullet\otimes_A A').
  \tag{6.3.8.1}\label{III.6.3.8.1-fr}
\]

为此，取 Cartan--Eilenberg 投射消解
$L_{\bullet\bullet}$，它消解 $P_\bullet$；另一方面，取 Cartan--Eilenberg 投射消解
$L'_{\bullet\bullet}$，它消解 $P_\bullet\otimes_A A'$。我们将证明可以定义一个
$A'$-复形同态
$L_{\bullet\bullet}\otimes_A A'\to L'_{\bullet\bullet}$，它在同伦意义下确定。
事实上，$L_{\bullet\bullet}$ 的构造完全由如下任意选择确定：对每个 $i$，
取投射消解 $(X_{ij}^{\mathrm B})_{j\geqslant0}$ 来消解 $\mathrm B_i(P_\bullet)$，
并取投射消解 $(X_{ij}^{\mathrm H})_{j\geqslant0}$ 来消解 $\mathrm H_i(P_\bullet)$；
它们分别等于 $\mathrm B_i^{\mathrm I}(L_{\bullet\bullet})$ 和
$\mathrm H_i^{\mathrm I}(L_{\bullet\bullet})$；于是依次得到
\[
  \mathrm Z_i^{\mathrm I}(L_{\bullet\bullet})
  =\mathrm H_i^{\mathrm I}(L_{\bullet\bullet})
   \oplus\mathrm B_i^{\mathrm I}(L_{\bullet\bullet}),
  \qquad
  L_{i,\bullet}=\mathrm Z_i^{\mathrm I}(L_{\bullet\bullet})
   \oplus\mathrm B_{i-1}^{\mathrm I}(L_{\bullet\bullet}).
\]
于是，$X_{i,\bullet}^{\mathrm B}\otimes_A A'$ 通常不再是
$\mathrm B_i(P_\bullet)\otimes_A A'$ 的消解，但仍是由 $A'$-投射模组成的复形，
因而存在一个 $A'$-同态
\[
  X_{i,\bullet}^{\mathrm B}\otimes_A A'
  \longrightarrow \mathrm B_i^{\mathrm I}(L'_{\bullet\bullet})
\]
它与增广相容，并在同伦意义下确定（M，V，1.1）。同样有一个
$A'$-同态
\[
  X_{i,\bullet}^{\mathrm H}\otimes_A A'
  \longrightarrow \mathrm H_i^{\mathrm I}(L'_{\bullet\bullet})
\]
它在同伦意义下确定；由此，利用上面回顾的构造，得到一个
$A'$-同态 $L_{i,\bullet}\otimes_A A'\to L'_{i,\bullet}$，它对每个 $i$ 都定义；这些同态（对
$i\in\mathbf Z$）与 $L_{i,\bullet}\to L_{i-1,\bullet}$ 的微分算子以及
$L'_{\bullet\bullet}$ 中的相应算子相容，这是同一构造的结果；因此它们构成所求的
$A'$-同态
$L_{\bullet\bullet}\otimes_A A'\to L'_{\bullet\bullet}$。

为了定义（6.3.8.1），同样取 Cartan--Eilenberg 投射消解
$M_{\bullet\bullet}$（分别为 $M'_{\bullet\bullet}$），它消解 $Q_\bullet$（分别消解
$Q_\bullet\otimes_A A'$），以及一个
$A'$-同态 $M_{\bullet\bullet}\otimes_A A'\to M'_{\bullet\bullet}$。由这些同态得到一个
$A'$-复形同态
\[
  (L_{\bullet\bullet}\otimes_A A')\otimes_{A'}
  (M_{\bullet\bullet}\otimes_A A')
  \longrightarrow
  L'_{\bullet\bullet}\otimes_{A'}M'_{\bullet\bullet},
\]
继而通过复合得到如下双复形的 $A$-同态：
\[
  L_{\bullet\bullet}\otimes_A M_{\bullet\bullet}
  \longrightarrow
  L'_{\bullet\bullet}\otimes_{A'}M'_{\bullet\bullet},
\]
再取同调便得到（6.3.8.1）；它确实定义良好，因为它来自一个在同伦意义下确定的
复形态射。

若 $\rho':A'\to A''$ 是第二个环同态，且 $\rho'':A\to A''$ 是复合环同态
$\rho'\circ\rho$，则显然有
\[
  \rho''_{P_\bullet,Q_\bullet}
  =\rho'_{P'_\bullet,Q'_\bullet}\circ\rho_{P_\bullet,Q_\bullet},
  \qquad
  P'_\bullet=P_\bullet\otimes_A A',
  \quad Q'_\bullet=Q_\bullet\otimes_A A'.
\]

还要注意，上面考虑的双复形态射
$L_{\bullet\bullet}\otimes_A M_{\bullet\bullet}\to
L'_{\bullet\bullet}\otimes_{A'}M'_{\bullet\bullet}$ 定义了谱序列的（关于 $P_\bullet$ 和
$Q_\bullet$ 的）函子性态射
\[
  'E_{pq}^r(P_\bullet,Q_\bullet)
  \longrightarrow
  'E_{pq}^r(P_\bullet\otimes_A A',Q_\bullet\otimes_A A')
\]
以及
\[
  ''E_{pq}^r(P_\bullet,Q_\bullet)
  \longrightarrow
  ''E_{pq}^r(P_\bullet\otimes_A A',Q_\bullet\otimes_A A'),
\]
它们独立于所取的 Cartan--Eilenberg 消解，并且也具有前述的传递性。
\end{env}

\begin{proposition}[6.3.9]
设 $\rho:A\to A'$ 是环同态，且 $A'$ 是平坦的 $A$-模。则有规范的函子性同构
\[
  \mathbf{Tor}_\bullet^{A'}
  (P_\bullet\otimes_A A',Q_\bullet\otimes_A A')
  \xrightarrow{\sim}
  \mathbf{Tor}_\bullet^A(P_\bullet,Q_\bullet)\otimes_A A'.
  \tag{6.3.9.1}\label{III.6.3.9.1-fr}
\]
\[
  \left\{
  \begin{aligned}
  'E(P_\bullet\otimes_A A',Q_\bullet\otimes_A A')
    &\xrightarrow{\sim}'E(P_\bullet,Q_\bullet)\otimes_A A',\\
  ''E(P_\bullet\otimes_A A',Q_\bullet\otimes_A A')
    &\xrightarrow{\sim}''E(P_\bullet,Q_\bullet)\otimes_A A'.
  \end{aligned}
  \right.
  \tag{6.3.9.2}\label{III.6.3.9.2-fr}
\]
\end{proposition}

\begin{proof}
\oldpage[III]{145}
事实上，由于函子 $M\otimes_A A'$ 关于 $M$ 正合，
$L_{\bullet\bullet}\otimes_A A'$ 和 $M_{\bullet\bullet}\otimes_A A'$ 分别是
$P_\bullet\otimes_A A'$ 和 $Q_\bullet\otimes_A A'$ 的 Cartan--Eilenberg
投射消解，结论随即得出。
\end{proof}

\begin{env}[6.3.10]
设 $\rho:A\to A'$ 是环同态；对于任意复形
$P'_\bullet$（由 $A'$-模组成），$P'_{\bullet[\rho]}$ 是 $A$-模复形；此外，恒等映射
$P'_{\bullet[\rho]}\to P'_\bullet$ 可视为下列规范映射的复合
\[
  P'_{\bullet[\rho]}\longrightarrow
  P'_{\bullet[\rho]}\otimes_A A'
  \xrightarrow{\mu}P'_\bullet,
\]
其中 $\mu$ 是 $A'$-同态 $\mu(x\otimes a')=a'x$。若 $Q'_\bullet$
是另一个 $A'$-模复形，则有次数 0 的规范函子性同态
\[
  \mathbf{Tor}_\bullet^A(P'_{\bullet[\rho]},Q'_{\bullet[\rho]})
  \longrightarrow
  \mathbf{Tor}_\bullet^{A'}
  (P'_{\bullet[\rho]}\otimes_A A',Q'_{\bullet[\rho]}\otimes_A A')
  \longrightarrow
  \mathbf{Tor}_\bullet^{A'}(P'_\bullet,Q'_\bullet),
  \tag{6.3.10.1}\label{III.6.3.10.1-fr}
\]
其中第一箭头是（6.3.8）中定义的 $A$-同态，第二箭头由
$A'$-同态 $P'_{\bullet[\rho]}\otimes_A A'\to P'_\bullet$ 和
$Q'_{\bullet[\rho]}\otimes_A A'\to Q'_\bullet$ 通过函子性得到。对（6.3.2）
的谱序列也有类似的同态，并且具有显然的传递性；其表述留给读者。
\end{env}

\begin{proposition}[6.3.11]
设 $\rho:A\to A'$ 是环同态，使 $A'$ 成为平坦的 $A$-模。对于任一复形
$P'_\bullet$（由 $A'$-模组成）以及任一下有界复形 $Q_\bullet$（由 $A$-模组成），有规范的
函子性同构
\[
  \mathbf{Tor}_\bullet^A(P'_{\bullet[\rho]},Q_\bullet)
  \xrightarrow{\sim}
  \mathbf{Tor}_\bullet^{A'}(P'_\bullet,Q_\bullet\otimes_A A').
  \tag{6.3.11.1}\label{III.6.3.11.1-fr}
\]
\end{proposition}

\begin{proof}
事实上，若 $M_{\bullet\bullet}$ 是 $Q_\bullet$ 的 Cartan--Eilenberg
投射消解，则 $M_{\bullet\bullet}\otimes_A A'$ 是
$Q_\bullet\otimes_A A'$ 的 Cartan--Eilenberg 投射消解，并且在规范同构意义下有
\[
  P'_{\bullet[\rho]}\otimes_A M_{\bullet\bullet}
  =P'_\bullet\otimes_{A'}(M_{\bullet\bullet}\otimes_A A');
\]
结论由（6.3.4）得出。
\end{proof}

\begin{namedtheorem}[6.3.12]{注}
设 $(A^\lambda)$ 是环的滤过归纳系统，$(P_\bullet^\lambda)$、$(Q_\bullet^\lambda)$
是 $(A^\lambda)$-模复形的两个归纳系统；则有推广（6.3.6.1）的规范同构
\[
  \varinjlim_\lambda
  \mathbf{Tor}_\bullet^{A^\lambda}
  (P_\bullet^\lambda,Q_\bullet^\lambda)
  \xrightarrow{\sim}
  \mathbf{Tor}_\bullet^A(P_\bullet,Q_\bullet),
  \tag{6.3.12.1}\label{III.6.3.12.1-fr}
\]
其中 $A=\varinjlim A^\lambda$、$P_\bullet=\varinjlim P_\bullet^\lambda$、
$Q_\bullet=\varinjlim Q_\bullet^\lambda$。一旦借助（6.3.10）定义了同态
\[
  \mathbf{Tor}_\bullet^{A^\lambda}
  (P_\bullet^\lambda,Q_\bullet^\lambda)
  \longrightarrow
  \mathbf{Tor}_\bullet^{A^\mu}
  (P_\bullet^\mu,Q_\bullet^\mu),
\]
（对 $\lambda\leqslant\mu$），
证明即为（6.3.6）的证明。
\end{namedtheorem}

\begin{proposition}[6.3.13]
设 $S$ 是 $A$ 的一个乘法闭子集，$P_\bullet$ 和 $Q_\bullet$ 是两个
$A$-模复形，其中由 $S$ 中元素定义的倍乘映射均为双射；因此，若
$A'=S^{-1}A$，则 $P_\bullet$ 和 $Q_\bullet$ 由 $A'$-模组成。于是有规范同构
\[
  \mathbf{Tor}_\bullet^A(P_\bullet,Q_\bullet)
  \xrightarrow{\sim}
  \mathbf{Tor}_\bullet^{A'}(P_\bullet,Q_\bullet).
\]
\end{proposition}

\begin{proof}
事实上，该假设推出规范同态
$P_\bullet\to P_\bullet\otimes_A A'$、$Q_\bullet\to Q_\bullet\otimes_A A'$
均为双射。另一方面，超 Tor 的函子性表明每个 $s\in S$ 在
$\mathbf{Tor}_\bullet^A(P_\bullet,Q_\bullet)$ 中定义一个双射倍乘，因而
\[
  \mathbf{Tor}_\bullet^A(P_\bullet,Q_\bullet)
  \longrightarrow
  \mathbf{Tor}_\bullet^A(P_\bullet,Q_\bullet)\otimes_A A'
\]
\oldpage[III]{146}
也是双射。由于 $A'$ 是平坦的 $A$-模，结论由（6.3.9）得出；对于谱序列也同样有
规范同构。
\end{proof}

\subsection{拟相干模复形的局部超 Tor 函子：仿射概形情形}

\begin{env}[6.4.1]
设 $S$ 是环为 $A$ 的仿射概形，$X$、$Y$ 是两个 $S$-仿射概形，其环分别为
$B$、$C$，从而 $B$、$C$ 都是 $A$-代数。任意复形
$\mathscr P_\bullet$（相应地，$\mathscr Q_\bullet$），若由拟相干
$\mathscr O_X$-模（相应地，$\mathscr O_Y$-模）组成，都具有
$\widetilde{P_\bullet}$（相应地 $\widetilde{Q_\bullet}$）的形式，其中
$P_\bullet$（相应地 $Q_\bullet$）是 $B$-模（相应地 $C$-模）复形
（I，1.3.7 与 1.3.8）。显然可以把 $P_\bullet$ 和 $Q_\bullet$ 看作
$A$-模复形，并形成 $\mathbf{Tor}_n^A(P_\bullet,Q_\bullet)$；此外，由于
$\mathbf{Tor}_n^A(P_\bullet,Q_\bullet)$ 的双函子性，$A$-代数 $B$、$C$ 在此
$A$-模上作用，而这些作用使其成为一个 $(B,C)$-双模；换言之，它是
$B\otimes_A C=A(X\times_S Y)$ 上的模。于是由此定义了拟相干的
$\mathscr O_{X\times_S Y}$-模
\[
  \mathscr{T}\!or_n^{\mathscr O_S}(\mathscr P_\bullet,\mathscr Q_\bullet)
  =\bigl(\mathbf{Tor}_n^A(P_\bullet,Q_\bullet)\bigr)^{\sim}
  \tag{6.4.1.1}\label{III.6.4.1.1-fr}
\]
称其为指标为 $n$ 的、复形 $\mathscr P_\bullet$ 与 $\mathscr Q_\bullet$ 的
\emph{局部超 Tor}，也记作
$\mathscr{T}\!or_n^S(\mathscr P_\bullet,\mathscr Q_\bullet)$。
\end{env}

\begin{namedtheorem}[6.4.2]{引理}
沿用（6.4.1）的记号，设环 $A$ 具有 $R^{-1}A'$ 的形式，其中 $A'$ 是一个环，
$R$ 是 $A'$ 的一个乘法闭子集。令 $S'=\operatorname{Spec}(A')$，于是可将
$X$、$Y$ 看作 $S'$-预概形，并且有
$X\times_{S'}Y=X\times_S Y$（I，1.6.2 与 3.2.4）。则
\[
  \mathscr{T}\!or_\bullet^{S'}(\mathscr P_\bullet,\mathscr Q_\bullet)
  =\mathscr{T}\!or_\bullet^S(\mathscr P_\bullet,\mathscr Q_\bullet).
\]
\end{namedtheorem}

\begin{proof}
这由公式（6.4.1.1）和（6.3.13）推出。
\end{proof}

\begin{env}[6.4.3]
沿用（6.4.1）的记号与假设，设
$\mathscr F=\widetilde F$ 是拟相干的 $\mathscr O_X$-模，
$\mathscr G=\widetilde G$ 是拟相干的 $\mathscr O_Y$-模；将
$\mathscr F$、$\mathscr G$ 看作模复形，记它们的超 Tor 为
$\mathscr{T}\!or_n^{\mathscr O_S}(\mathscr F,\mathscr G)$ 或
$\mathscr{T}\!or_n^S(\mathscr F,\mathscr G)$，其指标为 $n$；由（6.3.5（ii））有
\[
  \mathscr{T}\!or_n^{\mathscr O_S}(\mathscr F,\mathscr G)
  =\bigl(\operatorname{Tor}_n^A(F,G)\bigr)^{\sim}.
  \tag{6.4.3.1}\label{III.6.4.3.1-fr}
\]

现在回到两个拟相干模复形 $\mathscr P_\bullet$、$\mathscr Q_\bullet$ 的一般情形。
公式（6.4.1.1）和（6.4.3.1），结合命题（6.3.2），表明
$\mathscr{T}\!or_\bullet^S(\mathscr P_\bullet,\mathscr Q_\bullet)$ 是两个谱序列
$'\mathscr E(\mathscr P_\bullet,\mathscr Q_\bullet)$、
$''\mathscr E(\mathscr P_\bullet,\mathscr Q_\bullet)$ 的共同收敛项，它们的
$E_2$ 项由下式给出：
\[
  '\mathscr E_{pq}^2
  =\mathscr H_p\bigl(\mathscr{T}\!or_q^S
    (\mathscr P_\bullet,\mathscr Q_\bullet)\bigr)
  \tag{6.4.3.2}\label{III.6.4.3.2-fr}
\]
\[
  ''\mathscr E_{pq}^2
  =\bigoplus_{q'+q''=q}
  \mathscr{T}\!or_p^S
  \bigl(\mathscr H_{q'}(\mathscr P_\bullet),
        \mathscr H_{q''}(\mathscr Q_\bullet)\bigr)
  \tag{6.4.3.3}\label{III.6.4.3.3-fr}
\]
其中 $\mathscr{T}\!or_q^S(\mathscr P_\bullet,\mathscr Q_\bullet)$ 是由拟相干
$\mathscr O_{X\times_S Y}$-模
$\mathscr{T}\!or_q^S(\mathscr P_i,\mathscr Q_j)$ 组成的双复形。
\end{env}

\begin{env}[6.4.4]
现在再考虑两个仿射概形
$X^{(1)}=\operatorname{Spec}(B^{(1)})$、
$Y^{(1)}=\operatorname{Spec}(C^{(1)})$，其中 $B^{(1)}$ 和 $C^{(1)}$ 是
$A$-代数，并给定两个

\oldpage[III]{147}
S-态射 $u:X^{(1)}\to X$、$v:Y^{(1)}\to Y$，它们分别对应于
$A$-同态 $\varphi:B\to B^{(1)}$、$\psi:C\to C^{(1)}$。
考虑复形
$u^*(\mathscr P_\bullet)=(P_\bullet\otimes_B B^{(1)})^{\sim}$（它是
$\mathscr O_{X^{(1)}}$-模复形）、
$v^*(\mathscr Q_\bullet)=(Q_\bullet\otimes_C C^{(1)})^{\sim}$（它是
$\mathscr O_{Y^{(1)}}$-模复形）。规范的 $A$-同态
\[
  P_\bullet\longrightarrow P_\bullet\otimes_B B^{(1)},
  \qquad
  Q_\bullet\longrightarrow Q_\bullet\otimes_C C^{(1)}
\]
由函子性给出一个 $A$-同态
\[
  \mathbf{Tor}_\bullet^A(P_\bullet,Q_\bullet)
  \longrightarrow
  \mathbf{Tor}_\bullet^A
  (P_\bullet\otimes_B B^{(1)},Q_\bullet\otimes_C C^{(1)});
\]
此外，仍由函子性，这个同态实际上是一个 $(B\otimes_A C)$-模同态。因此，
这样就定义了一个 $(u\times_S v)$-态射
\[
  \theta:\mathscr{T}\!or_\bullet^S(\mathscr P_\bullet,\mathscr Q_\bullet)
  \longrightarrow
  \mathscr{T}\!or_\bullet^S
  (u^*(\mathscr P_\bullet),v^*(\mathscr Q_\bullet))
  \tag{6.4.4.1}\label{III.6.4.4.1-fr}
\]
继而得到一个 $\mathscr O_{X^{(1)}\times_S Y^{(1)}}$-模同态
\[
  \theta^\sharp:(u\times_S v)^*
  \bigl(\mathscr{T}\!or_\bullet^S(\mathscr P_\bullet,\mathscr Q_\bullet)\bigr)
  \longrightarrow
  \mathscr{T}\!or_\bullet^S
  (u^*(\mathscr P_\bullet),v^*(\mathscr Q_\bullet))
  \tag{6.4.4.2}\label{III.6.4.4.2-fr}
\]
这显然是下有界拟相干模范畴中的一个双-$\partial$-函子态射。

同态（6.4.4.2）不一定是双射；但是：
\end{env}

\begin{namedtheorem}[6.4.5]{引理}
沿用（6.4.4）的记号，设 $u$ 和 $v$ 是开浸入；则同态（6.4.4.2）是双射。
\end{namedtheorem}

\begin{proof}
将 $X^{(1)}$（相应地，将 $Y^{(1)}$）识别为 $X$（相应地，$Y$）的开子集；
于是 $X^{(1)}$（相应地，$Y^{(1)}$）是形如
$\mathrm D(f)$（相应地，$\mathrm D(g)$）的开集之并，其中 $f\in B$
（相应地，$g\in C$），且所诱导的预概形 $\mathrm D(\varphi(f))$ 与
$\mathrm D(f)$（相应地，$\mathrm D(\psi(g))$ 与 $\mathrm D(g)$）同构。只需在
$X^{(1)}$（相应地，$Y^{(1)}$）为 $\mathrm D(f)$（相应地，$\mathrm D(g)$）的情形
证明引理；事实上，若这一点成立，回到一般情形时，只需证明
$\theta^\sharp$ 在每个开集 $\mathrm D(f)\times_S\mathrm D(g)$ 上的限制是同构；
而若 $u_1:\mathrm D(f)\to X^{(1)}$、$v_1:\mathrm D(g)\to Y^{(1)}$ 是规范嵌入，
则上述限制正是 $(u_1\times_S v_1)^*(\theta^\sharp)$；但是，由定义（6.4.4）和
（$0_{\mathrm I}$，4.4.8）立即可知，将其与规范同态
\[
  (u_1\times_S v_1)^*
  \bigl(\mathscr{T}\!or_\bullet^S
    (u^*(\mathscr P_\bullet),v^*(\mathscr Q_\bullet))\bigr)
  \longrightarrow
  \mathscr{T}\!or_\bullet^S
  ((u')^*(\mathscr P_\bullet),(v')^*(\mathscr Q_\bullet))
  \tag{6.4.5.1}\label{III.6.4.5.1-fr}
\]
复合，其中 $u'=u\circ u_1$、$v'=v\circ v_1$，就得到规范同态
\[
  (u'\times_S v')^*
  \bigl(\mathscr{T}\!or_\bullet^S(\mathscr P_\bullet,\mathscr Q_\bullet)\bigr)
  \longrightarrow
  \mathscr{T}\!or_\bullet^S
  ((u')^*(\mathscr P_\bullet),(v')^*(\mathscr Q_\bullet))
  \tag{6.4.5.2}\label{III.6.4.5.2-fr}
\]
；若已知（6.4.5.1）和（6.4.5.2）是同构，则 $(u_1\times_S v_1)^*(\theta^\sharp)$
 也必为同构。

因此，设 $X^{(1)}=\mathrm D(f)$、$Y^{(1)}=\mathrm D(g)$，于是
$B^{(1)}=B_f$、$C^{(1)}=C_g$；$u^*(\mathscr P_\bullet)$（相应地，
$v^*(\mathscr Q_\bullet)$）可分别识别为 $((P_\bullet)_f)^{\sim}$（相应地，
$((Q_\bullet)_g)^{\sim}$）；另一方面，$X^{(1)}\times_S Y^{(1)}$ 可识别为开子概形
$\mathrm D(f\otimes g)$；它是 $X\times_S Y=\operatorname{Spec}(B\otimes_A C)$ 的开子概形
（II，4.3.2.4）。因此要证明同态
\[
  \bigl(\mathbf{Tor}_\bullet^A(P_\bullet,Q_\bullet)\bigr)_{f\otimes g}
  \longrightarrow
  \mathbf{Tor}_\bullet^A((P_\bullet)_f,(Q_\bullet)_g)
  \tag{6.4.5.3}\label{III.6.4.5.3-fr}
\]

\oldpage[III]{148}
由规范同态 $P_\bullet\to(P_\bullet)_f$、$Q_\bullet\to(Q_\bullet)_g$ 通过
函子性诱导的同态是双射。由（$0_{\mathrm I}$，1.6.1），可写成
$(P_\bullet)_f=\varinjlim P_\bullet^{(n)}$，其中每个 $P_\bullet^{(n)}$ 都是
$B$-模复形，并且与 $P_\bullet$ 相同；映射
$P_\bullet^{(m)}\to P_\bullet^{(n)}$（对 $m\leqslant n$）是乘以 $f^{n-m}$；对 $Q_\bullet$ 也有类似结果，
只需将 $f$ 换成 $g$。另一方面，显然，同态
\[
  \mathbf{Tor}_\bullet^A(P_\bullet^{(m)},Q_\bullet^{(m)})
  \longrightarrow
  \mathbf{Tor}_\bullet^A(P_\bullet^{(n)},Q_\bullet^{(n)})
\]
对应于映射 $P_\bullet^{(m)}\to P_\bullet^{(n)}$ 和
$Q_\bullet^{(m)}\to Q_\bullet^{(n)}$，按定义就是乘以 $(f\otimes g)^{n-m}$。于是，将（$0_{\mathrm I}$，1.6.1）应用于
（6.4.5.3）的第一个成员，并应用（6.3.6），即可得到结论。
\end{proof}

\begin{env}[6.4.6]
沿用（6.4.4）的记号，同样定义谱函子的规范同态
\[
  \left\{
  \begin{aligned}
  (u\times_S v)^*\bigl('\mathscr E(\mathscr P_\bullet,\mathscr Q_\bullet)\bigr)
    &\longrightarrow
      '\mathscr E(u^*(\mathscr P_\bullet),v^*(\mathscr Q_\bullet)),\\
  (u\times_S v)^*\bigl(''\mathscr E(\mathscr P_\bullet,\mathscr Q_\bullet)\bigr)
    &\longrightarrow
      ''\mathscr E(u^*(\mathscr P_\bullet),v^*(\mathscr Q_\bullet)).
  \end{aligned}
  \right.
  \tag{6.4.6.1}\label{III.6.4.6.1-fr}
\]
而（6.4.5）的论证表明，当 $u$ 和 $v$ 是开浸入时，同态（6.4.6.1）是双射：事实上，
结合（6.3.6.2）和（6.3.6.3），该论证证明它在 $E^2$ 项上是同构，而（6.4.5）
证明它在收敛项上是同构；因此借助（0，11.1.2）和（0，11.2.4）即可得出结论。
\end{env}

\subsection{拟相干模复形的局部超 Tor 函子：一般情形}

\begin{env}[6.5.1]
现在考虑任意预概形 $S$ 和两个任意的 $S$-预概形 $X$、$Y$；设
$\mathscr P_\bullet$（相应地，$\mathscr Q_\bullet$）是一个
$\mathscr O_X$-模（相应地，$\mathscr O_Y$-模）复形，且为拟相干。令
$Z=X\times_S Y$；我们将定义拟相干的 $\mathscr O_Z$-模
$\mathscr{T}\!or_n^S(\mathscr P_\bullet,\mathscr Q_\bullet)$，称为
$\mathscr P_\bullet$ 与 $\mathscr Q_\bullet$ 的\emph{局部超 Tor}；当 $S$、$X$ 和
$Y$ 为仿射概形时，它们将化为（6.4）中已经定义的那些对象。

当 $\mathscr P_\bullet$ 和 $\mathscr Q_\bullet$ 分别退化为它们的 0 次项
$\mathscr F$ 和 $\mathscr G$（其余项均为零）时，我们用
$\mathscr{T}\!or_n^S(\mathscr F,\mathscr G)$ 代替
$\mathscr{T}\!or_n^S(\mathscr P_\bullet,\mathscr Q_\bullet)$。
\end{env}

\begin{env}[6.5.2]
先设 $S$ 为仿射概形，并设 $(X_\lambda)$、$(Y_\mu)$ 分别是由仿射开集构成的
$X$、$Y$ 的覆盖；于是 $Z_{\lambda\mu}=X_\lambda\times_S Y_\mu$ 构成 $Z$ 的
仿射开覆盖。置
$\mathscr P_{\lambda\bullet}=\mathscr P_\bullet|X_\lambda$、
$\mathscr Q_{\mu\bullet}=\mathscr Q_\bullet|Y_\mu$；因此，对于每一对 $(\lambda,\mu)$，
都有一个拟相干的 $\mathscr O_{Z_{\lambda\mu}}$-模
\[
  \mathscr F_{\lambda\mu}
  =\mathscr{T}\!or_n^S
   (\mathscr P_{\lambda\bullet},\mathscr Q_{\mu\bullet}),
\]
还需证明 $\mathscr F_{\lambda\mu}$ 满足拼接条件（$0_{\mathrm I}$，3.3.1）。为此，
只需验证：对于任意仿射开集 $U\subset X_\lambda\cap X_{\lambda'}$（相应地，
$V\subset Y_\mu\cap Y_{\mu'}$），$\mathscr F_{\lambda\mu}$ 和
$\mathscr F_{\lambda'\mu'}$ 在 $U\times_S V$ 上的限制通过规范同构相同；但这立即
由这些限制在 $\mathscr{T}\!or_n^S(\mathscr P_\bullet|U,\mathscr Q_\bullet|V)$
上存在规范同构（6.4.5）推出。此外，由此定义以及（6.4.5）立即可得，如此定义的
$\mathscr O_Z$-模（在同构意义下）不依赖于所考\oldpage[III]{149}虑的开覆盖；
因此我们将其记为
$\mathscr{T}\!or_n^S(\mathscr P_\bullet,\mathscr Q_\bullet)$；最后，由（6.4.5）可知，
对于任意开子集 $U$（相应地，$V$），若它属于 $X$（相应地，$Y$），则
$\mathscr{T}\!or_n^S(\mathscr P_\bullet,\mathscr Q_\bullet)$ 在
$U\times_S V$ 上的限制与
$\mathscr{T}\!or_n^S(\mathscr P_\bullet|U,\mathscr Q_\bullet|V)$ 规范同构。
\end{env}

\begin{env}[6.5.3]
现在进入一般情形，其中 $S$ 是任意的；设 $(S_\alpha)$ 是由仿射开集构成的
$S$ 的覆盖；以 $X_\alpha$（相应地，$Y_\alpha$）表示 $S_\alpha$ 在 $X$
（相应地，$Y$）中的逆像；还需证明层
\[
  \mathscr{T}\!or_n^{S_\alpha}
  (\mathscr P_\bullet|X_\alpha,\mathscr Q_\bullet|Y_\alpha)
  =\mathscr G_\alpha
\]
满足拼接条件。只需对于每个仿射开集 $T$（它包含在 $S_\alpha\cap S_\beta$ 中），定义
$\mathscr G_\alpha$ 和 $\mathscr G_\beta$ 在 $U\times_S V$ 上的限制之间的规范同构，
其中 $U$ 和 $V$ 分别表示 $T$ 在 $X$ 和 $Y$ 中的逆像，并令它们取值于
$\mathscr{T}\!or_n^T(\mathscr P_\bullet|U,\mathscr Q_\bullet|V)$；此外，可以只考虑
$T$ 同时写成 $\mathrm D(f_\alpha)$ 和 $\mathrm D(f_\beta)$ 的情形，其中
$f_\alpha$（相应地，$f_\beta$）是 $\mathscr O_S$ 在 $S_\alpha$
（相应地，$S_\beta$）上的一个截面；但此时，由（6.4.2），
$\mathscr{T}\!or_n^T(\mathscr P_\bullet|U,\mathscr Q_\bullet|V)$ 一方面与
$\mathscr{T}\!or_n^{S_\alpha}(\mathscr P_\bullet|U,\mathscr Q_\bullet|V)$
规范同构，另一方面与
$\mathscr{T}\!or_n^{S_\beta}(\mathscr P_\bullet|U,\mathscr Q_\bullet|V)$
规范同构；由于我们刚刚根据（6.5.2）定义了 $\mathscr G_\alpha$ 到
$\mathscr{T}\!or_n^{S_\alpha}(\mathscr P_\bullet|U,\mathscr Q_\bullet|V)$ 以及
$\mathscr G_\beta$ 到
$\mathscr{T}\!or_n^{S_\beta}(\mathscr P_\bullet|U,\mathscr Q_\bullet|V)$ 的规范同构，
这就完成了 $\mathscr O_Z$-模
$\mathscr{T}\!or_n^S(\mathscr P_\bullet,\mathscr Q_\bullet)$ 的定义。此外，对于任意开子集
$U$（相应地，$V$），若它属于 $X$（相应地，$Y$），则
$\mathscr{T}\!or_n^S(\mathscr P_\bullet|U,\mathscr Q_\bullet|V)$ 与
$\mathscr{T}\!or_n^S(\mathscr P_\bullet,\mathscr Q_\bullet)$ 在
$U\times_S V$ 上的限制规范同构。

这样定义的确是双-$\partial$-函子
$\mathscr{T}\!or_\bullet^S(\mathscr P_\bullet,\mathscr Q_\bullet)$（在下有界拟相干模复形范畴中）；
它取值于 $\mathscr O_Z$-模范畴，
这一点是立即的，因为问题在 $X$、$Y$ 和 $S$ 上都是局部的；这是根据（6.4.5）
以及（6.4.4.2）是双-$\partial$-函子态射这一点得到的。注意，若
$\mathscr P_\bullet$ 和 $\mathscr Q_\bullet$ 分别退化为它们的 0 次项
$\mathscr F$ 和 $\mathscr G$，则根据（6.4.1.1），
$\mathscr{T}\!or_0^S(\mathscr F,\mathscr G)$ 正是（I，9.1.2）中定义的外张量积
$\mathscr F\otimes_S\mathscr G$；事实上，这由（I，9.1.3）推出。
\end{env}

\begin{env}[6.5.4]
由上述构造以及（6.4.6）中的说明可知，
$\mathscr{T}\!or_\bullet^S(\mathscr P_\bullet,\mathscr Q_\bullet)$ 是两个谱函子的
收敛项：
$'\mathscr E(\mathscr P_\bullet,\mathscr Q_\bullet)$、
$''\mathscr E(\mathscr P_\bullet,\mathscr Q_\bullet)$；它们的 $E_2$ 项分别为
\[
  '\mathscr E_{pq}^2
  =\mathscr H_p\bigl(\mathscr{T}\!or_q^S
    (\mathscr P_\bullet,\mathscr Q_\bullet)\bigr)
  \tag{6.5.4.1}\label{III.6.5.4.1-fr}
\]
\[
  ''\mathscr E_{pq}^2
  =\bigoplus_{q'+q''=q}
   \mathscr{T}\!or_p^S
   \bigl(\mathscr H_{q'}(\mathscr P_\bullet),
         \mathscr H_{q''}(\mathscr Q_\bullet)\bigr).
  \tag{6.5.4.2}\label{III.6.5.4.2-fr}
\]
谱序列（6.5.4.2）总是正则的；若 $\mathscr P_\bullet$ 和 $\mathscr Q_\bullet$ 下有界，
则两个谱序列都是双正则的。前述两个谱序列均为双正则的另一种情形如下：
\end{env}

\begin{env}[6.5.5]
我们称拓扑空间 $T$ 上的一个环层 $\mathscr A$ 具有\emph{上同调维数 $\leqslant n$}，
如果对于每个 $t\in T$，环 $\mathscr A_t$ 的上同调维数 $\leqslant n$；此时也称\emph{带环空间}
$(T,\mathscr A)$ 具有\emph{上同调维数 $\leqslant n$}。如果一个环层（相应地，一个带环空间）
存在整数 $n$ 使其上同调维数 $\leqslant n$，则称其上同调维数\emph{有限}。注意，若
$\mathscr A_t$ 是局部（交换）诺特环，

\oldpage[III]{150}
则称它们的上同调维数 $\leqslant n$，意味着它们是正则的且具有 $\leqslant n$ 的
（Krull）维数（$0_{\mathrm{IV}}$，17.3.1）。按照我们将在第 IV 章引入的维数理论术语，
称局部诺特预概形 $T$ 具有上同调维数 $\leqslant n$，等价于称它是正则的（$0_{\mathrm I}$，4.1.4）
且维数 $\leqslant n$；这意味着，对于每个仿射开集 $U$（它是 $T$ 的开集），环
$\Gamma(U,\mathscr O_T)$ 的上同调维数 $\leqslant n$（$0_{\mathrm{IV}}$，17.2.6）。因此，结合
这一最后的说明与（6.3.2）可知，若 $S$ 局部诺特且具有有限上同调维数，则谱序列
$'\mathscr E(\mathscr P_\bullet,\mathscr Q_\bullet)$ 和
$''\mathscr E(\mathscr P_\bullet,\mathscr Q_\bullet)$ 均为双正则的。

显然，$\mathscr{T}\!or_\bullet^S(\mathscr P_\bullet,\mathscr Q_\bullet)$
变为 $\mathscr{T}\!or_\bullet^S(\mathscr Q_\bullet,\mathscr P_\bullet)$（在同构意义下），
所依据的是 $X\times_S Y$ 到 $Y\times_S X$ 的规范同构。
\end{env}
\begin{proposition}[6.5.6]
设 $(\mathscr P_{\alpha\bullet})$ 是由拟相干 $\mathscr O_X$-模复形组成的滤过归纳系统；
则存在规范同构
\[
  \varinjlim_\alpha
  \mathscr{T}\!or_\bullet^S
  (\mathscr P_{\alpha\bullet},\mathscr Q_\bullet)
  \xrightarrow{\sim}
  \mathscr{T}\!or_\bullet^S
  \left(\varinjlim_\alpha\mathscr P_{\alpha\bullet},\mathscr Q_\bullet\right).
  \tag{6.5.6.1}\label{III.6.5.6.1-fr}
\]
\end{proposition}

\begin{proof}
由于问题局部地取决于 $S$、$X$ 和 $Y$，可设 $S$、$X$、$Y$ 都是仿射的；
于是命题化为（6.3.6）。
\end{proof}
\begin{namedtheorem}[6.5.7]{注}
\textnormal{(i)} 特别考虑 $S=X=Y$ 的情形；于是 $\mathscr P_\bullet$ 和
$\mathscr Q_\bullet$ 是两个拟相干的 $\mathscr O_S$-模复形。此时
$\mathscr{T}\!or_n^S(\mathscr P_\bullet,\mathscr Q_\bullet)$ 是拟相干的
$\mathscr O_S$-模；此外，对每个点 $z\in S$，由（6.5.6）得到规范同构
\[
  \bigl(\mathscr{T}\!or_n^S
    (\mathscr P_\bullet,\mathscr Q_\bullet)\bigr)_z
  \xrightarrow{\sim}
  \mathbf{Tor}_n^{\mathscr O_z}
  \bigl((\mathscr P_\bullet)_z,(\mathscr Q_\bullet)_z\bigr)
  \tag{6.5.7.1}\label{III.6.5.7.1-fr}
\]
因为问题是局部的，并且根据（6.4.1.1）归结为模的情形。

\textnormal{(ii)} 可以把超 Tor 的定义推广到两个 $\mathscr O_X$-模复形
$\mathscr P_\bullet$、$\mathscr Q_\bullet$ 的情形；它们位于同一带环空间
$(X,\mathscr O_X)$ 上。事实上，对每个开集 $U$（它是 $X$ 的开集），置
$A(U)=\Gamma(U,\mathscr O_X)$、$P_\bullet(U)=\Gamma(U,\mathscr P_\bullet)$、
$Q_\bullet(U)=\Gamma(U,\mathscr Q_\bullet)$；于是 $A(U)$-模
\[
  \mathbf{Tor}_n^{A(U)}(P_\bullet(U),Q_\bullet(U))
\]
构成 $X$ 上的一个预层；记
$\mathscr{T}\!or_n^X(\mathscr P_\bullet,\mathscr Q_\bullet)$ 为由此预层关联的
$\mathscr O_X$-模。当 $X$ 是预概形时，
由（6.3.12）可知，这个 $\mathscr O_X$-模与上面定义的超 Tor 规范同构。
我们不再进一步展开这一推广。
\end{namedtheorem}
\begin{proposition}[6.5.8]
设 $X$、$Y$ 是两个 $S$-预概形，$\mathscr F$（相应地，$\mathscr G$）是拟相干的
$\mathscr O_X$-模（相应地，$\mathscr O_Y$-模）。若 $\mathscr F$ 或 $\mathscr G$ 在 $S$ 上平坦，
则有 $\mathscr{T}\!or_n^S(\mathscr F,\mathscr G)=0$，只要 $n\neq0$。
\end{proposition}

\begin{proof}
由于问题在 $X$ 和 $Y$ 上是局部的，可设 $X$、$Y$、$S$ 都是仿射的，
其相应的环为 $B$、$C$、$A$，并且
$\mathscr F=\widetilde M$、$\mathscr G=\widetilde N$，其中 $M$（相应地，$N$）
是 $B$-模（相应地，$C$-模）。不妨设 $\mathscr F$ 在 $S$ 上平坦；这意味着对每个
$s\in S$，$M_s$ 是平坦的 $A_s$-模（$0_{\mathrm I}$，6.7.1）；因此 $M$ 是平坦的
$A$-模（$0_{\mathrm I}$，6.3.3），并且已知
$\operatorname{Tor}_n^A(M,N)=0$ 对 $n>0$ 以及任意 $C$-模 $N$ 都成立
（$0_{\mathrm I}$，6.1.1）。由（6.4.1.1）即得结论。
\end{proof}
\begin{namedtheorem}[6.5.9]{推论}
设 $X$、$Y$ 是两个 $S$-预概形，$\mathscr P_\bullet$（相应地，$\mathscr Q_\bullet$）是
下有界的 $\mathscr O_X$-模（相应地，$\mathscr O_Y$-模）复形，并且是拟相干的。假设所有
$\mathscr P_i$ 都在 $S$ 上平坦。则存在规范的 $\partial$-函子同构；该函子关于
$\mathscr Q_\bullet$ 取值：

\oldpage[III]{151}
\[
  \mathscr{T}\!or_\bullet^S(\mathscr P_\bullet,\mathscr Q_\bullet)
  \xrightarrow{\sim}
  \mathscr H_\bullet(\mathscr P_\bullet\otimes_S\mathscr Q_\bullet).
  \tag{6.5.9.1}\label{III.6.5.9.1-fr}
\]
\end{namedtheorem}

\begin{proof}
当 $S$、$X$、$Y$ 都是仿射的时，这不过是（6.3.7）；一般情形则由（6.5.2）和（6.5.3）
的论证推出。
\end{proof}
\begin{namedtheorem}[6.5.10]{推论}
假设 $X$ 在 $S$ 上平坦（$0_{\mathrm I}$，6.7.1），$\mathscr P_\bullet$ 和
$\mathscr Q_\bullet$ 下有界，并且所有 $\mathscr P_i$ 都是局部自由的
$\mathscr O_X$-模（不必是有限型）。则同态（6.5.9.1）是双射。
\end{namedtheorem}

\begin{proof}
事实上，（6.5.9）的假设得到满足，因为按定义平坦性是 $X$ 上的逐点性质，且平坦模的任意
直和仍是平坦模（$0_{\mathrm I}$，6.1.2）。
\end{proof}
\begin{proposition}[6.5.11]
设 $X'$、$Y'$ 是两个 $S$-预概形，$f:X\to X'$、$g:Y\to Y'$ 是两个仿射
$S$-态射。设 $\mathscr P_\bullet$（相应地，$\mathscr Q_\bullet$）是由拟相干
$\mathscr O_X$-模（相应地，$\mathscr O_Y$-模）组成的复形。则存在函子性的规范同构
\[
  (f\times_S g)_*
  \bigl(\mathscr{T}\!or_\bullet^S
    (\mathscr P_\bullet,\mathscr Q_\bullet)\bigr)
  \xrightarrow{\sim}
  \mathscr{T}\!or_\bullet^S
  (f_*(\mathscr P_\bullet),g_*(\mathscr Q_\bullet)).
  \tag{6.5.11.1}\label{III.6.5.11.1-fr}
\]
\end{proposition}

\begin{proof}
由于 $f$ 和 $g$ 都是仿射的，$f_*(\mathscr P_\bullet)$ 和
$g_*(\mathscr Q_\bullet)$ 是拟相干模的复形（II，1.2.6）；若置
$Z'=X'\times_S Y'$，则（6.5.11.1）的两端都是拟相干的 $\mathscr O_{Z'}$-模
（6.5.1）。可容易地归结到 $S$、$X'$ 和 $Y'$ 都是仿射的情形；但此时由假设
$X$ 和 $Y$ 也都是仿射的，验证立即由（6.4.1.1）和（I，1.6.3）得出。
\end{proof}
\begin{namedtheorem}[6.5.12]{注}
设 $X'$、$Y'$ 是两个 $S$-预概形，并假设 $X'$ 在 $S$ 上平坦；设 $X$ 是
$X'$ 的一个闭子预概形，$i:X\to X'$ 是典范嵌入，$\mathscr P_\bullet$
（相应地，$\mathscr Q_\bullet$）是由拟相干 $\mathscr O_X$-模（相应地，
$\mathscr O_{Y'}$-模）组成的下有界复形。最后，设
$\mathscr L'_{\bullet\bullet}$ 是 $i_*(\mathscr P_\bullet)$ 的一个消解，
它由局部自由的 $\mathscr O_{X'}$-模组成，并满足：$X'$ 的每一点都有一个仿射开邻域
$U$，使得 $\mathscr L'_{j,\bullet}|U$ 是
$i_*(\mathscr P_j)|U$ 的一个自由消解，对每个 $j$ 都如此。于是存在规范同构
\[
  (i\times_S 1)_*
  \bigl(\mathscr{T}\!or_\bullet^S
    (\mathscr P_\bullet,\mathscr Q_\bullet)\bigr)
  \xrightarrow{\sim}
  \mathscr H_\bullet
  (\mathscr L'_{\bullet\bullet}\otimes_S\mathscr Q_\bullet).
  \tag{6.5.12.1}\label{III.6.5.12.1-fr}
\]

若 $S$、$X'$、$Y'$ 都是仿射的，并且
$\mathscr L'_{j,\bullet}$ 是 $i_*(\mathscr P_j)$ 的一个自由消解（对每个 $j$），则根据（6.5.11）
可归结到 $X'=X$ 且其在 $S$ 上平坦的情形，此时只需应用（6.3.4）。
在一般情形下，局部定义同构（6.5.12.1），还需验证这一定义确实给出一个整体同构。
为此，必须回到第一个同构（6.3.4.1）的定义：它来自谱序列的一个同构
（0，11.6.5 和 11.5.3），而后者又由双复形态射
$L_{\bullet\bullet}\to L''_{\bullet\bullet}$ 得到，其中 $L_{\bullet\bullet}$ 是
$P_\bullet$ 的给定消解，$L''_{\bullet\bullet}$ 是 $P_\bullet$ 在由

\oldpage[III]{152}
下有界 $A$-模复形组成的范畴中的一个投射消解（参见（0，11.5.2.2））。
我们的断言由以下事实推出：根据任意两个这样的消解之间存在同伦（M，V，1.2），
同构（6.3.4.1）不依赖于所选择的投射消解 $L''_{\bullet\bullet}$。
\end{namedtheorem}
\begin{proposition}[6.5.13]
设 $X$、$Y$ 是两个 $S$-预概形，并假设下列条件之一成立：
\begin{enumerate}
  \item[(i)] $X$ 和 $Z=X\times_S Y$ 局部诺特，并且 $X$ 在 $S$ 上平坦。
  \item[(ii)] $S$ 和 $X$ 局部诺特，并且 $Y$ 在 $S$ 上是有限型的。
\end{enumerate}
设 $\mathscr P_\bullet$（相应地，$\mathscr Q_\bullet$）是由拟相干
$\mathscr O_X$-模（相应地，$\mathscr O_Y$-模）组成的下有界复形。进一步假设，
对每个 $n$，$\mathscr H_n(\mathscr P_\bullet)$（相应地，
$\mathscr H_n(\mathscr Q_\bullet)$）是有限型的 $\mathscr O_X$-模（相应地，
$\mathscr O_Y$-模）。那么
$\mathscr{T}\!or_n^S(\mathscr P_\bullet,\mathscr Q_\bullet)$ 是凝聚的
$\mathscr O_Z$-模。
\end{proposition}

\begin{proof}
由于 $\mathscr P_\bullet$ 和 $\mathscr Q_\bullet$ 下有界，谱序列
$''\mathscr E(\mathscr P_\bullet,\mathscr Q_\bullet)$ 是双正则的（6.5.4）；
又根据（0，11.1.8），由于在（i）、（ii）两种情形下 $Z$ 都局部诺特，只需证明各项
$''\mathscr E_{pq}^2$ 凝聚。关于
$\mathscr H_n(\mathscr P_\bullet)$ 和 $\mathscr H_n(\mathscr Q_\bullet)$ 的假设，
以及 $''\mathscr E_{pq}^2$ 的表达式（6.5.4.2），由此说明本命题等价于
$\mathscr P_\bullet$ 和 $\mathscr Q_\bullet$ 都仅保留次数 0 项的特殊情形，
换言之，即为下述
\end{proof}

\begin{corollary}[6.5.14]
假设（6.5.13）的条件（i）、（ii）之一成立，并设 $\mathscr F$（相应地，
$\mathscr G$）是有限型拟相干 $\mathscr O_X$-模（相应地，$\mathscr O_Y$-模）；
则 $\mathscr{T}\!or_n^S(\mathscr F,\mathscr G)$ 是凝聚的 $\mathscr O_Z$-模。
\end{corollary}

\begin{proof}
由于问题在 $X$ 和 $Y$ 上是局部的，可设 $S$、$X$、$Y$ 都是仿射的。

\textnormal{(i)} 在（i）的假设下，$S$、$X$ 和 $Z$ 都是诺特的。因此存在
一个局部自由消解 $\mathscr L_\bullet$，它消解 $\mathscr F$，并且由有限型的
$\mathscr O_X$-模组成（I，1.3.7）；由于 $X$ 在 $S$ 上平坦，由（6.3.4）得到
\[
  \mathscr{T}\!or_n^S(\mathscr F,\mathscr G)
  =\mathscr H_n(\mathscr L_\bullet\otimes_S\mathscr G);
\]
而 $\mathscr L_i\otimes_S\mathscr G$ 是有限型拟相干 $\mathscr O_Z$-模
（I，9.1.1），因而是凝聚的。由此推出
$\mathscr H_n(\mathscr L_\bullet\otimes_S\mathscr G)$ 是凝聚的
（$0_{\mathrm I}$，5.3.4）。

\textnormal{(ii)} 现在假设条件（ii）成立。由于环 $A(Y)$ 是有限个不定元的
$A(S)$-多项式代数的商（I，6.3.3），$Y$ 是一个闭 $S$-子预概形，它位于仿射
$S$-预概形 $Y'$ 中；后者在 $S$ 上平坦且为有限型。由于 $Y'$ 是诺特的
（I，6.3.7），存在一个局部自由消解 $\mathscr M_\bullet$，它消解
$j_*(\mathscr G)$，并且由有限型的 $\mathscr O_{Y'}$-模组成
（其中 $j:Y\to Y'$ 是典范嵌入）。根据（6.5.12），
$\mathscr{T}\!or_n^S(\mathscr F,\mathscr G)$ 是在 $1\times j$ 下的逆像；被拉回的是
$\mathscr O_{Z'}$-模 $\mathscr H_n(\mathscr F\otimes_S\mathscr M_\bullet)$
（其中 $Z'=X\times_S Y'$）；如同（i）中一样，可见
$\mathscr F\otimes_S\mathscr M_i$ 是凝聚的 $\mathscr O_{Z'}$-模，
再由（$0_{\mathrm I}$，5.3.4）得到结论。
\end{proof}
\begin{env}[6.5.15]
上面对两个复形 $\mathscr P_\bullet$、$\mathscr Q_\bullet$ 所建立的理论（它们是两个
$S$-预概形 $X$、$Y$ 上的拟相干层复形），可以毫无困难地推广到如下情形：
考虑任意有限多个 $S$-预概形 $X^{(i)}$（$1\leqslant i\leqslant m$），并在每个
$X^{(i)}$ 上取一个复形 $\mathscr P_\bullet^{(i)}$，它由拟相干
$\mathscr O_{X^{(i)}}$-模组成；若
$Z=X^{(1)}\times_S X^{(2)}\times_S\cdots\times_S X^{(m)}$，便由此定义一个拟相干
$\mathscr O_Z$-模
$\mathscr{T}\!or_q^S(\mathscr P_\bullet^{(1)},\mathscr P_\bullet^{(2)},
\ldots,\mathscr P_\bullet^{(m)})$。我们把一般情形下理论的展开留给读者；为供以后引用，
这里只写出第二个谱序列（正则谱序列）的 $E_2$ 项，该谱序列收敛到
$\mathscr{T}\!or_\bullet^S(\mathscr P_\bullet^{(1)},
\mathscr P_\bullet^{(2)},\ldots,\mathscr P_\bullet^{(m)})$：

\oldpage[III]{153}

\[
  ''\mathscr E_{pq}^2
  =\bigoplus_{q_1+q_2+\cdots+q_m=q}
   \mathscr{T}\!or_p^S
   \bigl(\mathscr H_{q_1}(\mathscr P_\bullet^{(1)}),\ldots,
         \mathscr H_{q_m}(\mathscr P_\bullet^{(m)})\bigr).
  \tag{6.5.15.1}\label{III.6.5.15.1-fr}
\]
在（6.8）中，我们将研究由任意多个复形的这些超 Tor 函子所产生的结合性谱序列。
\end{env}
\subsection{拟相干模复形的整体超 Tor 函子与 Künneth 谱序列：仿射基底的情形}
\label{subsection:III.6.6.fr}

\begin{env}[6.6.1]
考虑一个仿射概形 $S=\operatorname{Spec}(A)$ 和两个拟紧 $S$-概形
$X^{(i)}$（$i=1,2$）；设 $\mathscr P_\bullet^{(i)}$ 是由拟相干
$\mathscr O_{X^{(i)}}$-模组成的下有界复形，其微分算子的次数为 $-1$
（$i=1,2$）。另一方面，考虑一个由仿射开集组成的有限覆盖
$\mathfrak U^{(i)}=(U_\alpha^{(i)})$，它覆盖 $X^{(i)}$；置
\[
  X=X^{(1)}\times_S X^{(2)},
\]
它是一个拟紧 $S$-概形（I，5.5.1 和 6.6.4）；并令
$\mathfrak U=\mathfrak U^{(1)}\times_S\mathfrak U^{(2)}$ 为 $X$ 的覆盖，它由诸仿射开集
$U_\alpha^{(1)}\times_S U_\beta^{(2)}$ 组成。对每一对整数
$p\leqslant0$、$q\in\mathbf Z$，$(-p)$-次交错上链群
\[
  C^{-p}(\mathfrak U^{(i)},\mathscr P_q^{(i)})
\]
（采用覆盖 $\mathfrak U^{(i)}$ 和系数层 $\mathscr P_q^{(i)}$；见 G，II，5.1）
是一个 $A$-模；当 $p>0$ 时，置
$C^p(\mathfrak U^{(i)},\mathscr P_q^{(i)})=0$；这样便定义了一个双复形
\[
  C^\bullet(\mathfrak U^{(i)},\mathscr P_\bullet^{(i)})
\]
；它由 $A$-模组成，并且两个微分算子的次数均为 $-1$。由定义（$0_{\mathrm I}$，12.1.2）可知，
这个双复形（按通常方式视为关于总次数的单复形）的同调 $A$-模
\[
  \mathrm H_n\bigl(C^\bullet
  (\mathfrak U^{(i)},\mathscr P_\bullet^{(i)})\bigr)
\]
正是超上同调 $A$-模
\[
  \mathbf H^{-n}
  (\mathfrak U^{(i)},\mathscr P^{(i)\bullet}),
\]
其中 $\mathscr P^{(i)\bullet}$ 是微分算子次数为 $+1$ 的复形；它通过把
$\mathscr P_{-q}^{(i)}$ 取作次数 $q$ 的分量而得到。滥用记号，我们把它写作
$\mathbf H^{-n}(\mathfrak U^{(i)},\mathscr P_\bullet^{(i)})$。于是由（6.2.2）可知，
这个 $A$-模规范同构于超上同调 $A$-模
$\mathbf H^{-n}(X^{(i)},\mathscr P^{(i)\bullet})$，我们也把后者写作
$\mathbf H^{-n}(X^{(i)},\mathscr P_\bullet^{(i)})$；因此它不依赖于所选择的有限覆盖
$\mathfrak U^{(i)}$。
\end{env}
\begin{env}[6.6.2]
我们将把关于函子相对于双复形的超同调的一般理论，应用于下列两个 $A$-模双复形
\[
  L_{\bullet\bullet}^{(i)}
  =C^\bullet(\mathfrak U^{(i)},\mathscr P_\bullet^{(i)})
  \qquad(i=1,2)
\]
以及关于这两个双复形的协变双函子
$L_{\bullet\bullet}^{(1)}\otimes_A L_{\bullet\bullet}^{(2)}$（$0_{\mathrm I}$，11.7.4）。
由于所考虑的上链是交错的，
而覆盖 $\mathfrak U^{(i)}$ 是有限的，注意到模
$C^{-p}(\mathfrak U^{(i)},\mathscr P_q^{(i)})$ 满足 $\ne0$ 的情形，只出现于有限多个
（其数目与 $q$ 无关）的 $p$ 值；特别地，每个
$L_{\bullet\bullet}^{(i)}$ 的两个次数都下有界。我们用
\[
  \mathbf{Tor}_n^A(L_{\bullet\bullet}^{(1)},L_{\bullet\bullet}^{(2)})
\]
或
\[
  \mathbf{Tor}_n^S
  (\mathfrak U^{(1)},\mathfrak U^{(2)};
   \mathscr P_\bullet^{(1)},\mathscr P_\bullet^{(2)})
\]
表示第 $n$ 个超同调模，即
$L_{\bullet\bullet}^{(1)}\otimes_A L_{\bullet\bullet}^{(2)}$ 的超同调模，
并称它为指标为 $n$ 的 $\mathscr P_\bullet^{(1)}$ 与 $\mathscr P_\bullet^{(2)}$ 关于覆盖
$\mathfrak U^{(1)}$ 和 $\mathfrak U^{(2)}$ 的\emph{超 Tor}。
当 $\mathscr P_\bullet^{(1)}$ 和 $\mathscr P_\bullet^{(2)}$ 都仅保留次数 0 项
$\mathscr F^{(1)}$ 和 $\mathscr F^{(2)}$ 时，把它们的超 Tor 写作
\[
  \mathbf{Tor}_n^S
  (\mathfrak U^{(1)},\mathfrak U^{(2)};
   \mathscr F^{(1)},\mathscr F^{(2)})
\]
。按照一般约定，我们也用
\[
  \mathbf{Tor}_n^S
  (\mathfrak U^{(1)},\mathfrak U^{(2)};
   \mathscr P_\bullet^{(1)},\mathscr P_\bullet^{(2)}),
\]
表示这样一个双复形：其指标为 $(j,k)$ 的分量是
\[
  \mathbf{Tor}_n^S
  (\mathfrak U^{(1)},\mathfrak U^{(2)};
   \mathscr P_j^{(1)},\mathscr P_k^{(2)}).
\]

由于 $L_{\bullet\bullet}^{(i)}$ 关于 $\mathscr P_\bullet^{(i)}$ 是正合函子，因为覆盖
\oldpage[III]{154}
$\mathfrak U^{(i)}$ 中各开集的交是仿射的（I，5.5.6 和 1.3.11），所以
$\mathbf{Tor}_\bullet^S
(\mathfrak U^{(1)},\mathfrak U^{(2)};
 \mathscr P_\bullet^{(1)},\mathscr P_\bullet^{(2)})$ 是协变双-$\partial$-函子，
它关于 $\mathscr P_\bullet^{(1)}$、$\mathscr P_\bullet^{(2)}$ 而变，
取值于 $A$-模范畴（$0_{\mathrm I}$，11.7.3）。此外，已知（$0_{\mathrm I}$，11.7.2）
这个双函子是六个双正则谱函子的共同收敛项；我们把这些谱函子记作
\[
  {}^{(t)}\mathscr E^S
  (\mathfrak U^{(1)},\mathfrak U^{(2)};
   \mathscr P_\bullet^{(1)},\mathscr P_\bullet^{(2)})
\]
或
\[
  {}^{(t)}\mathscr E
  (\mathfrak U^{(1)},\mathfrak U^{(2)};
   \mathscr P_\bullet^{(1)},\mathscr P_\bullet^{(2)}),
\]
其中 $t$ 应以字母 $a$、$b$、$a'$、$b'$、$c$、$d$ 之一替代；它们的 $E_2$ 项如下：
\[
\begin{aligned}
  {}^{(a)}\mathscr E_{pq}^2
  &=\bigoplus_{q_1+q_2=q}
    \mathbf{Tor}_p^A
    \bigl(\mathrm H_{q_1}(L_{\bullet\bullet}^{(1)}),
          \mathrm H_{q_2}(L_{\bullet\bullet}^{(2)})\bigr),\\
  {}^{(b)}\mathscr E_{pq}^2
  &=\mathrm H_p\bigl(
    \mathbf{Tor}_q^{A,\mathrm{II}}
    (L_{\bullet\bullet}^{(1)},L_{\bullet\bullet}^{(2)})\bigr),\\
  {}^{(a')}\mathscr E_{pq}^2
  &=\bigoplus_{q_1+q_2=q}
    \mathbf{Tor}_p^A
    \bigl(\mathrm H_{q_1}^{\mathrm I}(L_{\bullet\bullet}^{(1)}),
          \mathrm H_{q_2}^{\mathrm I}(L_{\bullet\bullet}^{(2)})\bigr),\\
  {}^{(b')}\mathscr E_{pq}^2
  &=\mathrm H_p\bigl(
    \mathbf{Tor}_q^A
    (L_{\bullet\bullet}^{(1)},L_{\bullet\bullet}^{(2)})\bigr),\\
  {}^{(c)}\mathscr E_{pq}^2
  &=\bigoplus_{q_1+q_2=q}
    \mathbf{Tor}_p^A
    \bigl(\mathrm H_{q_1}^{\mathrm{II}}(L_{\bullet\bullet}^{(1)}),
          \mathrm H_{q_2}^{\mathrm{II}}(L_{\bullet\bullet}^{(2)})\bigr),\\
  {}^{(d)}\mathscr E_{pq}^2
  &=\mathrm H_p\bigl(
    \mathbf{Tor}_q^{A,\mathrm I}
    (L_{\bullet\bullet}^{(1)},L_{\bullet\bullet}^{(2)})\bigr),
\end{aligned}
\]
其中的记号与超同调一般理论中的记号一致。下面我们将进一步明确这些初始项。
\end{env}

\begin{env}[6.6.3]
\emph{谱序列 $(a)$ 与 $(a')$.} 我们已在 (6.6.1) 看到，同调模
$\mathrm H_n(L_{\bullet\bullet}^{(i)})$（这里所指的双复形为
$L_{\bullet\bullet}^{(i)}$）等于
$\mathbf H^{-n}(X^{(i)},\mathscr P_\bullet^{(i)})$；因此
\[
  {}^{(a)}\mathscr E_{pq}^2
  =\bigoplus_{q_1+q_2=q}
   \mathbf{Tor}_p^A
   \bigl(\mathbf H^{-q_1}(X^{(1)},\mathscr P_\bullet^{(1)}),
         \mathbf H^{-q_2}(X^{(2)},\mathscr P_\bullet^{(2)})\bigr).
\]
依定义，复形
$\mathrm H_n^{\mathrm I}(L_{\bullet\bullet}^{(i)})$ 的次数为 $k$ 的项是同调模
$\mathrm H_n(C^\bullet(\mathfrak U^{(i)},\mathscr P_k^{(i)}))$，亦即依定义为上同调模
$\mathbf H^{-n}(\mathfrak U^{(i)},\mathscr P_k^{(i)})$；由 (1.4.1)，已知该模典范同构于
$\mathbf H^{-n}(X^{(i)},\mathscr P_k^{(i)})$；因此
\[
  {}^{(a')}\mathscr E_{pq}^2
  =\bigoplus_{q_1+q_2=q}
   \mathbf{Tor}_p^A
   \bigl(\mathbf H^{-q_1}(X^{(1)},\mathscr P_\bullet^{(1)}),
         \mathbf H^{-q_2}(X^{(2)},\mathscr P_\bullet^{(2)})\bigr).
\]
\end{env}

\begin{env}[6.6.4]
\emph{谱序列 $(b)$ 与 $(b')$.} 依定义，
$\mathbf{Tor}_q^{A,\mathrm{II}}
(L_{\bullet\bullet}^{(1)},L_{\bullet\bullet}^{(2)})$ 是一个双复形，其次数为 $(h,k)$ 的项是
$A$-模
\[
  \mathbf{Tor}_q^A
  \bigl(C^{-h}(\mathfrak U^{(1)},\mathscr P_\bullet^{(1)}),
        C^{-k}(\mathfrak U^{(2)},\mathscr P_\bullet^{(2)})\bigr).
\]
令 $\Phi^{(i)}$ 为 $\mathfrak U^{(i)}$ 的指标集；依定义，模复形
$C^r(\mathfrak U^{(i)},\mathscr P_\bullet^{(i)})$（$r\geqslant0$）是诸复形
$\Gamma(U_\rho^{(i)},\mathscr P_\bullet^{(i)})$ 的直和，其中 $U_\rho^{(i)}$ 是
诸 $U_\xi^{(i)}$（$\xi\in\rho$）的交，而 $\rho$ 遍历
$\mathfrak P(\Phi^{(i)})$；故 $A$-模
\[
  \mathbf{Tor}_q^A
  \bigl(C^{-h}(\mathfrak U^{(1)},\mathscr P_\bullet^{(1)}),
        C^{-k}(\mathfrak U^{(2)},\mathscr P_\bullet^{(2)})\bigr)
\]
是下列 $A$-模的直和：
\[
  \mathbf{Tor}_q^A
  \bigl(\Gamma(U_\sigma^{(1)},\mathscr P_\bullet^{(1)}),
        \Gamma(U_\tau^{(2)},\mathscr P_\bullet^{(2)})\bigr),
\]
这里 $\sigma$（相应地，$\tau$）遍历
$\mathfrak P(\Phi^{(1)})$（相应地，$\mathfrak P(\Phi^{(2)})$）中满足
$\operatorname{Card}(\sigma)=-(h+1)$（相应地，
$\operatorname{Card}(\tau)=-(k+1)$）的元素。由于 $X^{(1)}$ 和 $X^{(2)}$ 是概形，
$U_\rho^{(i)}$ 是仿射的，故由 (6.4.1.1) 有
\[
  \mathbf{Tor}_q^A
  \bigl(\Gamma(U_\sigma^{(1)},\mathscr P_\bullet^{(1)}),
        \Gamma(U_\tau^{(2)},\mathscr P_\bullet^{(2)})\bigr)
  =\Gamma\bigl(U_\sigma^{(1)}\times_S U_\tau^{(2)},
    \mathscr{T}\!or_q^S
    (\mathscr P_\bullet^{(1)},\mathscr P_\bullet^{(2)})\bigr).
\]
\oldpage[III]{155}
因此可见，${}^{(b)}\mathscr E_{pq}^2$ 是第 $(-p)$ 个上同调模，所取复形为
\[
  L^\bullet(\Phi^{(1)},\Phi^{(2)};\mathscr S)
\]
；该复形由 $\Phi^{(1)}$ 与 $\Phi^{(2)}$ 上、取值于下列系数系的
双交错上链组成：
\[
  \mathscr S:(\sigma,\tau)\leadsto
  \Gamma\bigl(U_\sigma^{(1)}\times_S U_\tau^{(2)},
    \mathscr{T}\!or_q^S
    (\mathscr P_\bullet^{(1)},\mathscr P_\bullet^{(2)})\bigr)
\]
（$0_{\mathrm I}$，11.8.4）。于是已知（$0_{\mathrm I}$，11.8.5 和 11.8.6），
该复形的上同调与复形 $C^\bullet(\Phi^{(1)},\Phi^{(2)};\mathscr S)$ 的上同调相同；
这个复形由 $\Phi^{(1)}$ 和 $\Phi^{(2)}$ 上取值于 $\mathscr S$ 的所有上链组成；其上同调也与复形
$P^\bullet(\Phi^{(1)},\Phi^{(2)};\mathscr S)$ 的上同调相同；后者的元素是下列诸元的线性组合：
\[
  \lambda(\sigma,\tau)\in
  \Gamma\bigl(U_\sigma^{(1)}\times_S U_\tau^{(2)},
    \mathscr{T}\!or_q^S
    (\mathscr P_\bullet^{(1)},\mathscr P_\bullet^{(2)})\bigr),
\]
其中 $\sigma=(\alpha_0,\ldots,\alpha_h)$ 与
$\tau=(\beta_0,\ldots,\beta_h)$ 是具有相同项数的序列。但此时有
\[
  U_\sigma^{(1)}\times_S U_\tau^{(2)}
  =(U_{\alpha_0}^{(1)}\times_S U_{\beta_0}^{(2)})\cap\cdots\cap
    (U_{\alpha_h}^{(1)}\times_S U_{\beta_h}^{(2)})
  \qquad\text{(I, 3.2.7)}.
\]
若以 $\mathfrak U$ 表示 $Z=X^{(1)}\times_S X^{(2)}$ 由仿射开集
$U_\alpha^{(1)}\times_S U_\beta^{(2)}$ 构成的覆盖，则最终考虑到
$X^{(1)}\times_S X^{(2)}$ 是概形，由 (1.3.1) 可见
\[
  {}^{(b)}\mathscr E_{pq}^2
  =\mathbf H^{-p}\bigl(X^{(1)}\times_S X^{(2)},
    \mathscr{T}\!or_q^S
    (\mathscr P_\bullet^{(1)},\mathscr P_\bullet^{(2)})\bigr).
\]

其次，
$\mathbf{Tor}_q^A(L_{\bullet\bullet}^{(1)},L_{\bullet\bullet}^{(2)})$ 是一个双复形，
其次数为 $(h,k)$ 的项是下列 $A$-模的直和：
\[
  \mathbf{Tor}_q^A
  \bigl(C^{-h_1}(\mathfrak U^{(1)},\mathscr P_{k_1}^{(1)}),
        C^{-h_2}(\mathfrak U^{(2)},\mathscr P_{k_2}^{(2)})\bigr)
\]
其中 $h_1+h_2=h$ 且 $k_1+k_2=k$；如上展开模
$C^r(\mathfrak U^{(i)},\mathscr P_s^{(i)})$，又可看出该项是下列 $A$-模的直和：
\[
  \Gamma\bigl(U_\sigma^{(1)}\times_S U_\tau^{(2)},
    \mathscr{T}\!or_q^S
    (\mathscr P_{k_1}^{(1)},\mathscr P_{k_2}^{(2)})\bigr),
\]
其中 $k_1+k_2=k$，而 $\sigma$（相应地，$\tau$）遍历
$\mathfrak P(\Phi^{(1)})$（相应地，$\mathfrak P(\Phi^{(2)})$）中满足
$\operatorname{Card}(\sigma)+\operatorname{Card}(\tau)=-h-2$ 的元素。我们所计算的项
${}^{(b')}\mathscr E_{pq}^2$ 是第 $(-p)$ 个上同调模；这里所取的双复形为
$N^{\bullet\bullet}=(N^{hk})$；
这里单复形 $N^{\bullet k}$ 是 $\Phi^{(1)}$ 与 $\Phi^{(2)}$ 上、取值于下列系数系的
双交错上链复形：
\[
  \mathscr S_k:(\sigma,\tau)\leadsto
  \Gamma\left(U_\sigma^{(1)}\times_S U_\tau^{(2)},
    \bigoplus_{k_1+k_2=k}
    \mathscr{T}\!or_q^S
    (\mathscr P_{-k_1}^{(1)},\mathscr P_{-k_2}^{(2)})\right),
\]
这些系数系构成复形 $\mathscr S^\bullet$，其微分来自与双复形
$\mathscr{T}\!or_q^S
(\mathscr P_\bullet^{(1)},\mathscr P_\bullet^{(2)})$ 相伴的单复形的微分。
已知 $N^{\bullet\bullet}$ 的上同调与双复形
$C^\bullet(\Phi^{(1)},\Phi^{(2)};\mathscr S^\bullet)$ 的上同调相同
（$0_{\mathrm I}$，11.8.9），也与双复形
$P^\bullet(\Phi^{(1)},\Phi^{(2)};\mathscr S^\bullet)$ 的上同调相同；其次数为 $(h,k)$ 的元素是
下列诸元的线性组合：
\[
  \lambda(\sigma,\tau)\in
  \Gamma\left(U_\sigma^{(1)}\times_S U_\tau^{(2)},
    \bigoplus_{k_1+k_2=k}
    \mathscr{T}\!or_q^S
    (\mathscr P_{-k_1}^{(1)},\mathscr P_{-k_2}^{(2)})\right),
\]
这里 $\sigma=(\alpha_0,\ldots,\alpha_h)$、
$\tau=(\beta_0,\ldots,\beta_h)$ 是具有相同项数的序列（$0_{\mathrm I}$，11.8.10）。
如上可见，${}^{(b')}\mathscr E_{pq}^2$ 是第 $(-p)$ 个上同调模；所取的双复形为
$C^\bullet(\mathfrak U,\mathscr Q^\bullet)$，其中
$\mathscr Q^\bullet$ 是与双复形
$\mathscr{T}\!or_q^S
(\mathscr P_\bullet^{(1)},\mathscr P_\bullet^{(2)})$ 相伴的单复形，而该双复形由
$\mathscr O_Z$-模组成。因此，按照 (6.6.1) 中的约定，
\[
  {}^{(b')}\mathscr E_{pq}^2
  =\mathbf H^{-p}\bigl(X^{(1)}\times_S X^{(2)},
    \mathscr{T}\!or_q^S
    (\mathscr P_\bullet^{(1)},\mathscr P_\bullet^{(2)})\bigr).
\]
\end{env}

\oldpage[III]{156}
\begin{env}[6.6.5]
\emph{谱序列 $(c)$ 与 $(d)$.} 依定义，复形
$\mathrm H_n^{\mathrm{II}}(L_{\bullet\bullet}^{(i)})$ 的次数为 $h$ 的项是 $A$-模
$C^{-h}(\mathfrak U^{(i)},\mathscr H_n(\mathscr P_\bullet^{(i)}))$；这是由于函子
$C^{-h}$ 的正合性。因此，由两个模关于两个覆盖的超 Tor 定义（6.6.2）有
\[
  {}^{(c)}\mathscr E_{pq}^2
  =\bigoplus_{q_1+q_2=q}
   \mathbf{Tor}_p^S
   \bigl(\mathfrak U^{(1)},\mathfrak U^{(2)};
         \mathscr H_{q_1}(\mathscr P_\bullet^{(1)}),
         \mathscr H_{q_2}(\mathscr P_\bullet^{(2)})\bigr).
\]
最后，依定义，
\[
\mathbf{Tor}_q^{A,\mathrm I}
(L_{\bullet\bullet}^{(1)},L_{\bullet\bullet}^{(2)})
\]
是一个双复形，其次数为 $(h,k)$ 的项是 $A$-模
\[
\mathbf{Tor}_q^S
(\mathfrak U^{(1)},\mathfrak U^{(2)};
 \mathscr P_h^{(1)},\mathscr P_k^{(2)}).
\]
因此
\[
  {}^{(d)}\mathscr E_{pq}^2
  =\mathrm H_p\bigl(
    \mathbf{Tor}_q^S
    (\mathfrak U^{(1)},\mathfrak U^{(2)};
     \mathscr P_\bullet^{(1)},\mathscr P_\bullet^{(2)})\bigr).
\]
\end{env}

\begin{env}[6.6.6]
双复形函子的超同调理论（$0_{\mathrm I}$，11.7.3）如同在 (6.3.4) 中那样表明：
对于任一 Cartan--Eilenberg \emph{平坦}消解 $M_{\bullet\bullet}^{(i)}$，其所消解的复形为
$L_{\bullet\bullet}^{(i)}$（在下有界模复形范畴中）（$i=1,2$），存在双-$\partial$-函子的典范同构
\[
\begin{aligned}
  \mathbf{Tor}_\bullet^S
  (\mathfrak U^{(1)},\mathfrak U^{(2)};
   \mathscr P_\bullet^{(1)},\mathscr P_\bullet^{(2)})
  &\xrightarrow{\sim}
   \mathrm H_\bullet
   (M_{\bullet\bullet}^{(1)}\otimes_A M_{\bullet\bullet}^{(2)})\\
  &\xrightarrow{\sim}
   \mathrm H_\bullet
   (M_{\bullet\bullet}^{(1)}\otimes_A L_{\bullet\bullet}^{(2)})
   \xrightarrow{\sim}
   \mathrm H_\bullet
   (L_{\bullet\bullet}^{(1)}\otimes_A M_{\bullet\bullet}^{(2)}).
\end{aligned}
\tag{6.6.6.1}\label{III.6.6.6.1.fr}
\]
\end{env}

\begin{env}[6.6.7]
现在证明，(6.6.2) 中定义的整体超 Tor 以及相应的六个谱序列，在典范同构的意义下，
并不依赖于定义它们时所用的有限仿射开覆盖 $\mathfrak U^{(i)}$。为此，只须证明如下断言：
若 $\mathfrak V^{(i)}$ 是另外两个同类覆盖，并且 $\mathfrak V^{(i)}$ 比
$\mathfrak U^{(i)}$ \emph{更细}（$i=1,2$），则存在谱函子的\emph{典范同构}
\[
  {}^{(t)}\mathscr E
  (\mathfrak U^{(1)},\mathfrak U^{(2)};
   \mathscr P_\bullet^{(1)},\mathscr P_\bullet^{(2)})
  \xrightarrow{\sim}
  {}^{(t)}\mathscr E
  (\mathfrak V^{(1)},\mathfrak V^{(2)};
   \mathscr P_\bullet^{(1)},\mathscr P_\bullet^{(2)})
  \tag{6.6.7, $t$}\label{III.6.6.7.t.fr}
\]
其中 $t$ 分别取 $a$、$b$、$a'$、$b'$、$c$ 或 $d$。

另一方面，对 $i=1,2$，有双复形同态
\[
  C^\bullet(\mathfrak U^{(i)},\mathscr P_\bullet^{(i)})
  \longrightarrow
  C^\bullet(\mathfrak V^{(i)},\mathscr P_\bullet^{(i)})
\]
它们在同伦意义下有良定义（G，II，5.7.1）；由此已得到典范确定的同态
(6.6.7, $t$)，并且这些同态与收敛项中的边界算子相容（$0_{\mathrm I}$，11.3.2）。
此外，对这些谱序列的 $E_2$ 项作计算，即对 $(a)$、$(b)$、$(a')$、$(b')$ 的计算表明，
对于这些谱序列，同态 (6.6.7, $t$) 在 $E_2$ 项上是同构；由于这些谱序列都是双正则的，
可见 (6.6.7, $t$) 对这四个谱函子都是同构，从而在它们的共同收敛项上给出
\emph{双-$\partial$-函子的同构}（$0_{\mathrm I}$，11.1.5）。

特别地，对于拟相干 $\mathscr O_{X^{(i)}}$-模 $\mathscr F^{(i)}$（$i=1,2$），典范同态
\[
  \mathbf{Tor}_\bullet^S
  (\mathfrak U^{(1)},\mathfrak U^{(2)};
   \mathscr F^{(1)},\mathscr F^{(2)})
  \longrightarrow
  \mathbf{Tor}_\bullet^S
  (\mathfrak V^{(1)},\mathfrak V^{(2)};
   \mathscr F^{(1)},\mathscr F^{(2)})
\]
是双射；由 (6.6.5) 的计算可见，同态 (6.6.7, $t$) 在 $E_2$ 项上也为同构，其中 $t=c$ 或 $t=d$。
同上可得，(6.6.7, $t$) 也是谱序列的同构，其中 $t=c$ 或 $t=d$。
\oldpage[III]{157}
可以认为，同构 (6.6.7, $t$) 在有限仿射开覆盖对
$(\mathfrak U^{(1)},\mathfrak U^{(2)})$ 所成的滤过集合上定义了谱函子的归纳系统；
这里两个覆盖分别覆盖 $X^{(1)}$ 与 $X^{(2)}$。
我们用
\[
  {}^{(t)}\mathscr E
  (X^{(1)},X^{(2)};
   \mathscr P_\bullet^{(1)},\mathscr P_\bullet^{(2)})
  \quad\text{或}\quad
  {}^{(t)}\mathscr E^S
  (X^{(1)},X^{(2)};
   \mathscr P_\bullet^{(1)},\mathscr P_\bullet^{(2)})
\]
表示这个系统的归纳极限，用
$\mathbf{Tor}_\bullet^S
(X^{(1)},X^{(2)};\mathscr P_\bullet^{(1)},\mathscr P_\bullet^{(2)})$
表示这个谱函子的收敛项，并称之为两个复形 $\mathscr P_\bullet^{(1)}$ 与
$\mathscr P_\bullet^{(2)}$ 的\emph{整体超 Tor}；当 $\mathscr P_\bullet^{(1)}$ 与
$\mathscr P_\bullet^{(2)}$ 都仅保留其 0 次项 $\mathscr F^{(1)}$ 与 $\mathscr F^{(2)}$ 时，记作
\[
  \mathbf{Tor}_\bullet^S
  (X^{(1)},X^{(2)};\mathscr F^{(1)},\mathscr F^{(2)}),
\]
并且依照一般约定，
\[
  \mathbf{Tor}_q^S
  (X^{(1)},X^{(2)};\mathscr P_\bullet^{(1)},\mathscr P_\bullet^{(2)})
\]
因而就是由下列各分量组成的双复形：
\[
  \mathbf{Tor}_q^S
  (X^{(1)},X^{(2)};\mathscr P_h^{(1)},\mathscr P_k^{(2)}).
\]
\end{env}

\begin{env}[6.6.8]
仍在 (6.6.1) 的假设下，现在考虑两个 $S$-态射 $f_i:X^{(i)}\to Y^{(i)}$，其中
$Y^{(i)}=\operatorname{Spec}(B_i)$ 是\emph{仿射} $S$-概形，因而 $B_i$ 是 $A$-代数（$i=1,2$）；
这便定义了 $A$-同态 $B_i\to\Gamma(X^{(i)},\mathscr O_{X^{(i)}})$（I，2.2.4），
从而 (6.6.2) 中定义的每个 $L_{\bullet\bullet}^{(i)}$ 都是 $B_i$-模双复形；于是
$L_{\bullet\bullet}^{(1)}\otimes_A L_{\bullet\bullet}^{(2)}$ 是
$(B_1\otimes_A B_2)$-模四重复形，而且它的六个超同调谱函子可视为取值于
$(B_1\otimes_A B_2)$-模谱序列的范畴。若置
$Y=Y^{(1)}\times_S Y^{(2)}=\operatorname{Spec}(B_1\otimes_A B_2)$，
则可考虑与这些模相伴的拟相干 $\mathscr O_Y$-模（I，1.3.4）；我们用
${}^{(t)}\mathscr E
(f_1,f_2;\mathscr P_\bullet^{(1)},\mathscr P_\bullet^{(2)})$
（其中 $t=a$、$b$、$a'$、$b'$、$c$ 或 $d$）表示下列六个谱序列
\[
  \bigl({}^{(t)}\mathscr E
  (X^{(1)},X^{(2)};
   \mathscr P_\bullet^{(1)},\mathscr P_\bullet^{(2)})\bigr)^{\sim}
\]
它们是 $\mathscr O_Y$-模的谱序列；并用
$\mathfrak{Tor}_\bullet^S
(f_1,f_2;\mathscr P_\bullet^{(1)},\mathscr P_\bullet^{(2)})$ 表示它们的共同收敛项
\[
  \bigl(\mathbf{Tor}_\bullet^S
  (X^{(1)},X^{(2)};
   \mathscr P_\bullet^{(1)},\mathscr P_\bullet^{(2)})\bigr)^{\sim}.
\]
当记作
$\mathfrak{Tor}_\bullet^S(f_1,f_2;\mathscr F^{(1)},\mathscr F^{(2)})$
时，指的是 $\mathscr P_\bullet^{(i)}$ 仅保留其 0 次项 $\mathscr F^{(i)}$（$i=1,2$）的情形。
\end{env}

\subsection{拟相干模复形的整体超 Tor 函子与 Künneth 谱序列：一般情形}

\begin{env}[6.7.1]
现在把 (6.6.8) 的定义推广到如下情形：$S$ 是任意预概形，$X^{(i)}$、$Y^{(i)}$
是 $S$-预概形，而 $f_i:X^{(i)}\to Y^{(i)}$ 是\emph{分离且拟紧}的态射。
此时，对于任意一对下有界复形 $\mathscr P_\bullet^{(i)}$，其各项为拟相干
$\mathscr O_{X^{(i)}}$-模（$i=1,2$），要对每个 $n$ 定义拟相干 $\mathscr O_Y$-模
$\mathfrak{Tor}_n^S
(f_1,f_2;\mathscr P_\bullet^{(1)},\mathscr P_\bullet^{(2)})$ 以及六个谱函子；
当 $S$、$Y^{(1)}$ 和 $Y^{(2)}$ 都是仿射的时，这些定义应退化为 (6.6.8) 的定义
（置 $Y=Y^{(1)}\times_S Y^{(2)}$）。

先假定 $S=\operatorname{Spec}(A)$ 是\emph{仿射}的，而 $Y^{(1)}$ 和 $Y^{(2)}$ 任意；
设 $W^{(i)}$ 是 $Y^{(i)}$ 的一个仿射开集；于是 $f_i^{-1}(W^{(i)})$ 是拟紧
$S$-概形，$W=W^{(1)}\times_S W^{(2)}$ 是 $Y$ 的一个仿射开集；设
$f'_i:f_i^{-1}(W^{(i)})\to W^{(i)}$ 是 $f_i$ 的限制，并设
$\mathscr P'_\bullet{}^{(i)}$ 是限制
$\mathscr P_\bullet^{(i)}|f_i^{-1}(W^{(i)})$（$i=1,2$）。由 (6.6.8)，便得到谱序列
${}^{(t)}\mathscr E
(f'_1,f'_2;\mathscr P'_\bullet{}^{(1)},\mathscr P'_\bullet{}^{(2)})$；
它们由拟相干 $(\mathscr O_Y|W)$-模组成；
所需验证的是它们满足拼接条件（$0_{\mathrm I}$，3.3.1）。
问题立即归结到如下情形：$Y^{(i)}=\operatorname{Spec}(B_i)$ 是仿射的，并且
$W^{(i)}=D(g_i)$，其中 $g_i\in B_i$，从而
$W=D(g_1\otimes g_2)$
\oldpage[III]{158}
是 $Y=\operatorname{Spec}(B_1\otimes_A B_2)$ 中的开集（II，4.3.2.1）；若
$X'^{(i)}=f_i^{-1}(W^{(i)})$，则须建立谱函子的典范同构
\[
  {}^{(t)}\mathscr E
  (X'^{(1)},X'^{(2)};
   \mathscr P'_\bullet{}^{(1)},\mathscr P'_\bullet{}^{(2)})
  \xrightarrow{\sim}
  {}^{(t)}\mathscr E
  (X^{(1)},X^{(2)};
   \mathscr P_\bullet^{(1)},\mathscr P_\bullet^{(2)})
  \otimes_B B_g
  \tag{6.7.1.1}\label{III.6.7.1.1.fr}
\]
这里置 $B=B_1\otimes_A B_2$ 和 $g=g_1\otimes g_2$。为此，取有限仿射开覆盖
$\mathfrak U^{(i)}$，覆盖 $X^{(i)}$（$i=1,2$），并令 $\mathfrak U'^{(i)}$ 为
$\mathfrak U^{(i)}$ 在 $X'^{(i)}$ 上的迹；它仍由仿射开集组成（I，5.5.10）；更确切地，
\[
  C^\bullet(\mathfrak U'^{(i)},\mathscr P'_\bullet{}^{(i)})
  =C^\bullet(\mathfrak U^{(i)},\mathscr P_\bullet^{(i)})
   \otimes_{B_i}(B_i)_{g_i}.
\]
若置
$L'_{\bullet\bullet}{}^{(i)}
=C^\bullet(\mathfrak U'^{(i)},\mathscr P'_\bullet{}^{(i)})$，则有
\[
  L'_{\bullet\bullet}{}^{(1)}\otimes_A L'_{\bullet\bullet}{}^{(2)}
  =\bigl(L_{\bullet\bullet}^{(1)}\otimes_{B_1}(B_1)_{g_1}\bigr)
   \otimes_A
   \bigl(L_{\bullet\bullet}^{(2)}\otimes_{B_2}(B_2)_{g_2}\bigr);
\]
由于在典范同构意义下 $B_g=(B_1)_{g_1}\otimes_A(B_2)_{g_2}$，仍在典范同构意义下有
\[
  L'_{\bullet\bullet}{}^{(1)}\otimes_A L'_{\bullet\bullet}{}^{(2)}
  =(L_{\bullet\bullet}^{(1)}\otimes_A L_{\bullet\bullet}^{(2)})
    \otimes_B B_g.
\]
若 $M_{\bullet\bullet\bullet}^{(i)}$ 是 $L_{\bullet\bullet}^{(i)}$ 的 Cartan--Eilenberg 投射消解，
并可设它由 $B_i$-模组成，则由 $(B_i)_{g_i}$ 在 $B_i$ 上\emph{平坦}可知，
$M'_{\bullet\bullet\bullet}{}^{(i)}
=M_{\bullet\bullet\bullet}^{(i)}\otimes_{B_i}(B_i)_{g_i}$ 是双复形
$L'_{\bullet\bullet}{}^{(i)}$ 的 Cartan--Eilenberg 投射消解；此外有
\[
  M'_{\bullet\bullet\bullet}{}^{(1)}\otimes_A M'_{\bullet\bullet\bullet}{}^{(2)}
  =(M_{\bullet\bullet\bullet}^{(1)}\otimes_A M_{\bullet\bullet\bullet}^{(2)})
    \otimes_B B_g.
\]
于是，由双复形超同调的定义以及函子 $G\otimes_B B_g$ 关于 $B$-模 $G$ 的正合性，
立即得到所求同构 (6.7.1.1)。
\end{env}

\begin{env}[6.7.2]
现在设 $S$ 任意，并设 $u_i:Y^{(i)}\to S$ 是结构态射（$i=1,2$）。设 $(S_\alpha)$ 是
$S$ 的仿射开覆盖，置 $Y_\alpha^{(i)}=u_i^{-1}(S_\alpha)$、
$X_\alpha^{(i)}=f_i^{-1}(Y_\alpha^{(i)})$，并设
$f_{i\alpha}:X_\alpha^{(i)}\to Y_\alpha^{(i)}$ 是 $f_i$ 的限制；这是一个
\emph{分离且拟紧}的态射。诸
$Y_\alpha=Y_\alpha^{(1)}\times_{S_\alpha}Y_\alpha^{(2)}$ 构成 $Y$ 的开覆盖，
而 (6.7.1) 在每个 $Y_\alpha$ 上定义了谱函子
\[
  {}^{(t)}\mathscr E^{S_\alpha}
  \bigl(f_{1\alpha},f_{2\alpha};
    \mathscr P_\bullet^{(1)}|X_\alpha^{(1)},
    \mathscr P_\bullet^{(2)}|X_\alpha^{(2)}\bigr);
\]
仍须证明这些函子满足拼接条件。

问题立即归结到如下情形：$S=\operatorname{Spec}(A)$ 是仿射的，$S'=D(h)$，其中 $h\in A$，
并且 $u_i(Y^{(i)})\subset S'$；还可设 $Y^{(i)}=\operatorname{Spec}(B_i)$ 是仿射的；
所需定义的是典范同构
\[
  {}^{(t)}\mathscr E^S
  (X^{(1)},X^{(2)};
   \mathscr P_\bullet^{(1)},\mathscr P_\bullet^{(2)})
  \xrightarrow{\sim}
  {}^{(t)}\mathscr E^{S'}
  (X^{(1)},X^{(2)};
   \mathscr P_\bullet^{(1)},\mathscr P_\bullet^{(2)}).
  \tag{6.7.2.1}\label{III.6.7.2.1.fr}
\]
而在 (6.6.2) 的记号下，$L_{\bullet\bullet}^{(i)}$ 由 $A_h$-模组成，故在典范同构意义下有
$L_{\bullet\bullet}^{(1)}\otimes_{A_h}L_{\bullet\bullet}^{(2)}
=L_{\bullet\bullet}^{(1)}\otimes_A L_{\bullet\bullet}^{(2)}$；由于可以取 Cartan--Eilenberg 投射消解
$M_{\bullet\bullet\bullet}^{(i)}$，它消解 $L_{\bullet\bullet}^{(i)}$ 且由 $A_h$-模组成，
这便立即给出所求的典范同构。
\end{env}

综上，我们已经证明如下定理。

\begin{theorem}[6.7.3]
设 $S$ 是预概形，$f_i:X^{(i)}\to Y^{(i)}$ 是 $S$-预概形之间的分离且拟紧的
$S$-态射，$\mathscr P_\bullet^{(i)}$ 是由拟相干 $\mathscr O_{X^{(i)}}$-模组成的下有界复形
（$i=1,2$）；以 $u_i:Y^{(i)}\to S$ 表示结构态射，并置
$Y=Y^{(1)}\times_S Y^{(2)}$。存在取值于拟相干模范畴的双-$\partial$-函子
  $\mathfrak{Tor}_\bullet^S
(f_1,f_2;\mathscr P_\bullet^{(1)},\mathscr P_\bullet^{(2)})$，该范畴中的模是
$\mathscr O_Y$-模，而且它具有如下性质：若 $S_\alpha$ 是 $S$ 的仿射开集，
$V^{(i)}$ 是 $Y_\alpha^{(i)}=u_i^{-1}(S_\alpha)$ 的仿射开集（$i=1,2$），并且
$V=V^{(1)}\times_{S_\alpha}V^{(2)}$，则
\[
  \mathfrak{Tor}_\bullet^S
  (f_1,f_2;\mathscr P_\bullet^{(1)},\mathscr P_\bullet^{(2)})|V
  =\Bigl(\mathbf{Tor}_\bullet^{S_\alpha}
    \bigl(f_1^{-1}(V^{(1)}),f_2^{-1}(V^{(2)});
      \mathscr P_\bullet^{(1)}|f_1^{-1}(V^{(1)}),
      \mathscr P_\bullet^{(2)}|f_2^{-1}(V^{(2)})\bigr)\Bigr)^{\sim}.
\]
\oldpage[III]{159}
这个双函子是六个双正则谱函子的共同收敛项
\[
  {}^{(t)}\mathscr E
  (f_1,f_2;\mathscr P_\bullet^{(1)},\mathscr P_\bullet^{(2)})
  \qquad(t=a,b,a',b',c,d)
\]
其 $E_2$ 项由下式给出：
\[
\begin{aligned}
  {}^{(a)}\mathscr E_{pq}^2
  &=\bigoplus_{q_1+q_2=q}
    \mathscr{T}\!or_p^S
    \bigl(\mathscr H^{-q_1}(f_1,\mathscr P_\bullet^{(1)}),
          \mathscr H^{-q_2}(f_2,\mathscr P_\bullet^{(2)})\bigr),\\
  {}^{(b)}\mathscr E_{pq}^2
  &=\mathscr H^{-p}\bigl(f_1\times_S f_2,
    \mathfrak{Tor}_q^S
    (\mathscr P_\bullet^{(1)},\mathscr P_\bullet^{(2)})\bigr),\\
  {}^{(a')}\mathscr E_{pq}^2
  &=\bigoplus_{q_1+q_2=q}
    \mathfrak{Tor}_p^S
    \bigl(\mathscr H^{-q_1}(f_1,\mathscr P_\bullet^{(1)}),
          \mathscr H^{-q_2}(f_2,\mathscr P_\bullet^{(2)})\bigr),\\
  {}^{(b')}\mathscr E_{pq}^2
  &=\mathscr H^{-p}\bigl(f_1\times_S f_2,
    \mathscr{T}\!or_q^S
    (\mathscr P_\bullet^{(1)},\mathscr P_\bullet^{(2)})\bigr),\\
  {}^{(c)}\mathscr E_{pq}^2
  &=\bigoplus_{q_1+q_2=q}
    \mathfrak{Tor}_p^S
    \bigl(f_1,f_2;
      \mathscr H_{q_1}(\mathscr P_\bullet^{(1)}),
      \mathscr H_{q_2}(\mathscr P_\bullet^{(2)})\bigr),\\
  {}^{(d)}\mathscr E_{pq}^2
  &=\mathscr H_p\bigl(\mathfrak{Tor}_q^S
    (f_1,f_2;\mathscr P_\bullet^{(1)},\mathscr P_\bullet^{(2)})\bigr).
\end{aligned}
\]
谱序列 $(a)$ 和 $(b)$ 称为\emph{Künneth 谱序列}。

注意，谱序列 $(a)$ 与 $(a')$（相应地，$(b)$ 与 $(b')$）在
$\mathscr P_\bullet^{(1)}$ 和 $\mathscr P_\bullet^{(2)}$ 都仅保留其 0 次项时相同；在这种情形下，
谱序列 $(c)$ 与 $(d)$ 退化，因而没有意义。
\end{theorem}

\begin{remark}[6.7.4]
上面定义的整体超 Tor 同时以 (6.2.1) 中定义的\emph{超上同调模}和
(6.5.3) 中定义的\emph{局部超 Tor} 为特殊情形。我们来证明：对于任意\emph{拟紧且分离}的
态射 $f:X\to Y$ 和任意复形 $\mathscr P_\bullet$（由拟相干 $\mathscr O_X$-模组成且下有界），
存在典范的 $\partial$-函子同构（以 $\mathscr P_\bullet$ 为变元）
\[
  \mathfrak{Tor}_n^Y(f,1_Y;\mathscr P_\bullet,\mathscr O_Y)
  \xrightarrow{\sim}
  \mathscr H^{-n}(f,\mathscr P_\bullet)
  \qquad\text{（任意 $n\in\mathbf Z$）}.
  \tag{6.7.4.1}\label{III.6.7.4.1.fr}
\]
事实上，(6.7.2) 的拼接方法立即把问题化归到 $Y$ 为仿射的情形；于是根据 (6.2.2)，
可以借助同一个有限覆盖 $\mathfrak U$ 来计算 (6.7.4.1) 的两端；该覆盖由 $X$ 的仿射开集组成，
而计算左端时还取仅由 $Y$ 本身构成的 $Y$ 的覆盖。在 (6.6.2) 的记号下，双复形
$L_{\bullet\bullet}^{(2)}$ 此时仅保留双次数为 $(0,0)$ 的项，该项等于 $A$，
结论遂由 ($0_{\mathrm I}$, 11.7.5) 得出。关于这个结果的推广，见 (6.7.7)；但应注意，
若在 (6.7.4.1) 左端把 $\mathscr O_Y$ 换成\emph{任意}拟相干 $\mathscr O_Y$-模 $\mathscr F$，
则一般不再有与
$\mathscr H^{-n}(f,\mathscr P_\bullet\otimes_{\mathscr O_Y}\mathscr F)$ 的同构，尽管在上述计算中，
双复形 $L_{\bullet\bullet}^{(1)}\otimes_A L_{\bullet\bullet}^{(2)}$ 仍可认同于双复形
$C^\bullet(\mathfrak U,\mathscr P_\bullet\otimes_Y\mathscr F)$。

另一方面，存在典范的双-$\partial$-函子同构
\[
  \mathfrak{Tor}_\bullet^S
  (1_{X^{(1)}},1_{X^{(2)}};
   \mathscr P_\bullet^{(1)},\mathscr P_\bullet^{(2)})
  \xrightarrow{\sim}
  \mathscr{T}\!or_\bullet^S
  (\mathscr P_\bullet^{(1)},\mathscr P_\bullet^{(2)}).
  \tag{6.7.4.2}\label{III.6.7.4.2.fr}
\]
事实上，利用 (6.7.1) 和 (6.7.2)，问题仍可化归到 $S$ 和诸 $X^{(i)}$ 都是仿射的情形；
在 (6.7.4.2) 左端的计算中，可以把覆盖 $\mathfrak U^{(i)}$ 取为仅含 $X^{(i)}$ 的族，
从而在 (6.6.2) 的记号下，
\oldpage[III]{160}
$L_{\bullet\bullet}^{(i)}$ 化为 $\Gamma(X^{(i)},\mathscr P_\bullet^{(i)})$
（视为第一次数 $\ne0$ 的项均为零的双复形），而 (6.7.4.2) 两端的相等性由 (6.4.1.1) 和 (6.3.1) 得出。
\end{remark}

\begin{proposition}[6.7.5]
设 $u:\mathscr P_\bullet^{(1)}\to\mathscr Q_\bullet^{(1)}$ 是由拟相干
$\mathscr O_{X^{(1)}}$-模组成的下有界复形之间的同态，使得同态
\[
  \mathscr H_\bullet(u):
  \mathscr H_\bullet(\mathscr P_\bullet^{(1)})
  \longrightarrow
  \mathscr H_\bullet(\mathscr Q_\bullet^{(1)})
\]
（由 $u$ 导出）为同构。于是，同态
\[
  {}^{(t)}\mathscr E
  (f_1,f_2;\mathscr P_\bullet^{(1)},\mathscr P_\bullet^{(2)})
  \longrightarrow
  {}^{(t)}\mathscr E
  (f_1,f_2;\mathscr Q_\bullet^{(1)},\mathscr P_\bullet^{(2)})
\]
（由 $u$ 导出）在 $t=a$、$t=b$ 和 $t=c$ 时都是同构。
\end{proposition}

\begin{proof}
关于谱序列 $(c)$ 的断言来自如下事实：这个谱序列是双正则的，并且按假设，所考虑的同态在
$E_2$ 项上是同构（$0_{\mathrm I}$, 11.1.5）。这已经表明
\[
  \mathfrak{Tor}_\bullet^S
  (f_1,f_2;\mathscr P_\bullet^{(1)},\mathscr P_\bullet^{(2)})
  \longrightarrow
  \mathfrak{Tor}_\bullet^S
  (f_1,f_2;\mathscr Q_\bullet^{(1)},\mathscr P_\bullet^{(2)})
\]
是同构。应用关系式 (6.7.4.1) 和 (6.7.4.2)，首先可见同态
\[
  \mathscr H^{-n}(f_1,\mathscr P_\bullet^{(1)})
  \longrightarrow
  \mathscr H^{-n}(f_1,\mathscr Q_\bullet^{(1)})
  \quad\text{和}\quad
  \mathscr{T}\!or_n^S
  (\mathscr P_\bullet^{(1)},\mathscr P_\bullet^{(2)})
  \longrightarrow
  \mathscr{T}\!or_n^S
  (\mathscr Q_\bullet^{(1)},\mathscr P_\bullet^{(2)})
\]
都是由 $u$ 导出的同构。于是，关于谱序列 $(a)$ 和 $(b)$ 的断言来自如下事实：
这些谱序列是双正则的（6.7.3），而且所考虑的同态在 $E_2$ 项上是双射
（$0_{\mathrm I}$, 11.1.5）。

另一方面，注意到若 $u:\mathscr P_\bullet^{(1)}\to\mathscr Q_\bullet^{(1)}$ 是一个
\emph{同伦态射}，则对于六个谱序列都由此得到典范同构
\[
  {}^{(t)}\mathscr E
  (f_1,f_2;\mathscr P_\bullet^{(1)},\mathscr P_\bullet^{(2)})
  \longrightarrow
  {}^{(t)}\mathscr E
  (f_1,f_2;\mathscr Q_\bullet^{(1)},\mathscr P_\bullet^{(2)})
\]
。事实上，若 $S$ 和诸 $Y^{(i)}$ 都是仿射的，则由 $u$ 得到双复形之间的同伦态射
\[
  C^\bullet(\mathfrak U^{(1)},\mathscr P_\bullet^{(1)})
  \longrightarrow
  C^\bullet(\mathfrak U^{(1)},\mathscr Q_\bullet^{(1)}),
\]
而命题随即由超同调的一般理论推出（$0_{\mathrm I}$, 11.3.2）；一般情形通过拼接得到，
这里使用了如下事实：复形之间的同伦态射会导出这些复形的 Cartan--Eilenberg 投射消解之间的
同伦态射（M, XVII, 1.2）。
\end{proof}

\begin{proposition}[6.7.6]
假设复形 $\mathscr P_\bullet^{(1)}$ 或复形 $\mathscr P_\bullet^{(2)}$
由 $S$-平坦模组成（两个复形均下有界）。则有双-$\partial$-函子的典范同构
\[
  \mathfrak{Tor}_n^S
  (f_1,f_2;\mathscr P_\bullet^{(1)},\mathscr P_\bullet^{(2)})
  \xrightarrow{\sim}
  \mathscr H^{-n}\bigl(f_1\times_S f_2,
    \mathscr P_\bullet^{(1)}\otimes_S\mathscr P_\bullet^{(2)}\bigr).
  \tag{6.7.6.1}\label{III.6.7.6.1.fr}
\]
\end{proposition}

\begin{proof}
先假设 $S$、$Y^{(1)}$ 和 $Y^{(2)}$ 都是仿射的，从而处于（6.6.2）的情形，
并沿用其记号。例如，假设 $\mathscr P_\bullet^{(1)}$ 由 $S$-平坦模组成，
利用评注（6.6.6）计算超 Tor；它就是
$L_{\bullet\bullet}^{(1)}\otimes_A M_{\bullet\bullet\bullet}^{(2)}$ 的同调，其中
$M_{\bullet\bullet\bullet}^{(2)}$ 是 $L_{\bullet\bullet}^{(2)}$ 的
Cartan--Eilenberg 投射消解，其意义如（$0_{\mathrm I}$，11.7.1）。另一方面，
由假设（I，4.15.1），诸模 $L_{ij}^{(1)}$ 在 $A$ 上是\emph{平坦的}；
于是由（$0_{\mathrm I}$，11.7.5）得到典范同构
\[
  \mathbf{Tor}_\bullet^S
  (\mathfrak U^{(1)},\mathfrak U^{(2)};
   \mathscr P_\bullet^{(1)},\mathscr P_\bullet^{(2)})
  \xrightarrow{\sim}
  \mathscr H_\bullet
  (L_{\bullet\bullet}^{(1)}\otimes_A L_{\bullet\bullet}^{(2)}).
  \tag{6.7.6.2}\label{III.6.7.6.2.fr}
\]
另一方面，有一个从
$L_{\bullet\bullet}^{(1)}\otimes_A L_{\bullet\bullet}^{(2)}$ 到
$C^\bullet(\mathfrak U,\mathscr Q_\bullet)$ 的双复形自然同态；这里
$\mathfrak U$ 是 $Z=X^{(1)}\times_S X^{(2)}$ 由诸仿射开集
\oldpage[III]{161}
$U_\alpha^{(1)}\times_S U_\beta^{(2)}$ 组成的覆盖，而
$\mathscr Q_\bullet
=\mathscr P_\bullet^{(1)}\otimes_S\mathscr P_\bullet^{(2)}$
被视为关于总次数的单复形。事实上，这一同态的定义实质上已在（6.6.4）中
谱序列 $(b')$ 的计算过程中对 $q=0$ 给出；只须注意（仍采用（6.6.4）的记号），
一方面有从复形 $N^{\bullet k}$ 到复形
$C^\bullet(\Phi^{(1)},\Phi^{(2)};\mathscr S_k)$ 的自然同态
（$0_{\mathrm I}$，11.8.5），另一方面有从后一个复形到复形
$P^\bullet(\Phi^{(1)},\Phi^{(2)};\mathscr S_k)$ 的自然同态
（$0_{\mathrm I}$，11.8.6），最后还有从后一个上链复形到交错上链子复形的自然同态
（$0_{\mathrm I}$，11.8.7）。此外，如（6.6.4）所见，如此定义的双复形同态诱导
同调上的同构；故与（6.7.6.2）复合，得到同构
\[
  \mathbf{Tor}_n^S
  (\mathfrak U^{(1)},\mathfrak U^{(2)};
   \mathscr P_\bullet^{(1)},\mathscr P_\bullet^{(2)})
  \xrightarrow{\sim}
  \mathbf H^{-n}
  (\mathfrak U,\mathscr P_\bullet^{(1)}\otimes_S\mathscr P_\bullet^{(2)}).
  \tag{6.7.6.3}\label{III.6.7.6.3.fr}
\]
其次须证明，如此定义的同构与所选开覆盖无关（由（6.2.2），（6.7.6.3）的右端
典范同构于
$\mathbf H^{-n}
(X^{(1)}\times_S X^{(2)},
 \mathscr P_\bullet^{(1)}\otimes_S\mathscr P_\bullet^{(2)})$）；
利用（6.6.7）即可做到这一点：采用（6.6.7）的记号，注意到有一个在同伦意义下交换的图
\[
\xymatrix@C=18pt@R=28pt{
  C^\bullet(\mathfrak U^{(1)},\mathscr P_\bullet^{(1)})
  \otimes_A
  C^\bullet(\mathfrak U^{(2)},\mathscr P_\bullet^{(2)})
  \ar[r]\ar[d]
  & C^\bullet(\mathfrak U,\mathscr Q_\bullet)\ar[d]
  \\
  C^\bullet(\mathfrak V^{(1)},\mathscr P_\bullet^{(1)})
  \otimes_A
  C^\bullet(\mathfrak V^{(2)},\mathscr P_\bullet^{(2)})
  \ar[r]
  & C^\bullet(\mathfrak V,\mathscr Q_\bullet)
}
\]
其中水平箭头就是上面定义的同态。最后，须通过拼接过渡到一般情形；这如（6.7.1）和
（6.7.2）一样毫无困难，细节留给读者。
\end{proof}

\begin{proposition}[6.7.7]
假设 $\mathscr P_\bullet^{(1)}$ 和 $\mathscr P_\bullet^{(2)}$ 均下有界，
并且诸模 $\mathscr H^{-n}(f_1,\mathscr P_\bullet^{(1)})$ 或诸模
$\mathscr H^{-n}(f_2,\mathscr P_\bullet^{(2)})$ 中的一族全部为 $S$-平坦模。
则有双-$\partial$-函子的典范同构（$n$ 遍历 $\mathbf Z$）
\[
  \mathfrak{Tor}_n^S
  (f_1,f_2;\mathscr P_\bullet^{(1)},\mathscr P_\bullet^{(2)})
  \xrightarrow{\sim}
  \bigoplus_{q_1+q_2=n}
  \mathscr H^{-q_1}(f_1,\mathscr P_\bullet^{(1)})
  \otimes_S
  \mathscr H^{-q_2}(f_2,\mathscr P_\bullet^{(2)}).
  \tag{6.7.7.1}\label{III.6.7.7.1.fr}
\]
\end{proposition}

\begin{proof}
事实上，由（6.5.8），（6.7.3）的谱序列 $(a)$ 退化；又因该谱序列是双正则的
（6.7.3），命题立即由（$0_{\mathrm I}$，11.1.6）得出。
\end{proof}

\begin{theorem}[6.7.8]
假设：$1^\circ$ 复形 $\mathscr P_\bullet^{(1)}$ 和
$\mathscr P_\bullet^{(2)}$ 均下有界；$2^\circ$ 复形
$\mathscr P_\bullet^{(1)}$ 或复形 $\mathscr P_\bullet^{(2)}$ 由 $S$-平坦模组成；
$3^\circ$ 下列两族
\oldpage[III]{162}
模 $\mathscr H^{-n}(f_1,\mathscr P_\bullet^{(1)})$ 或
$\mathscr H^{-n}(f_2,\mathscr P_\bullet^{(2)})$ 中的一族全部为 $S$-平坦模。
则有双-$\partial$-函子的典范同构（$n$ 遍历 $\mathbf Z$）
\[
  \mathscr H^n\bigl(f_1\times_S f_2,
    \mathscr P_\bullet^{(1)}\otimes_S\mathscr P_\bullet^{(2)}\bigr)
  \xrightarrow{\sim}
  \bigoplus_{n_1+n_2=n}
  \mathscr H^{n_1}(f_1,\mathscr P_\bullet^{(1)})
  \otimes_S
  \mathscr H^{n_2}(f_2,\mathscr P_\bullet^{(2)}).
  \tag{6.7.8.1}\label{III.6.7.8.1.fr}
\]
（“Künneth 公式”）。
\end{theorem}

\begin{proof}
这由（6.7.6）和（6.7.7）得出。

当 $S$、$Y^{(1)}$ 和 $Y^{(2)}$ 都是仿射的，（6.7.8.1）的逆同构可由
双复形同态
\[
  C^\bullet(\mathfrak U^{(1)},\mathscr P_\bullet^{(1)})
  \otimes_A
  C^\bullet(\mathfrak U^{(2)},\mathscr P_\bullet^{(2)})
  \longrightarrow
  C^\bullet(\mathfrak U,\mathscr Q_\bullet)
\]
按照（G，I，2.7）中定义的方法导出（沿用（6.7.6）的记号），这由（G，I，5.5）可知。
\end{proof}

\begin{proposition}[6.7.9]
假设下列三个条件成立：$1^\circ$ $S$、$Y^{(1)}$ 和 $Y^{(2)}$ 均为局部
Noether 概形，$f_1$ 和 $f_2$ 均为固有态射，且 $Y^{(1)}$ 或 $Y^{(2)}$
在 $S$ 上为有限型；$2^\circ$ $\mathscr P_\bullet^{(1)}$ 和
$\mathscr P_\bullet^{(2)}$ 均下有界；$3^\circ$ 对每个 $n\in\mathbf Z$，
$\mathscr H_n(\mathscr P_\bullet^{(i)})$ 都是凝聚模（$i=1,2$）。在这些条件下，
$\mathfrak{Tor}_n^S
(f_1,f_2;\mathscr P_\bullet^{(1)},\mathscr P_\bullet^{(2)})$ 是凝聚
$\mathscr O_Y$-模（其中 $Y=Y^{(1)}\times_S Y^{(2)}$）。
\end{proposition}

\begin{proof}
由（6.5.13），诸局部超 Tor
$\mathscr{T}\!or_n^S(\mathscr P_\bullet^{(1)},\mathscr P_\bullet^{(2)})$
是凝聚 $\mathscr O_X$-模（$X=X^{(1)}\times_S X^{(2)}$ 是局部 Noether 概形，
因为按假设某个 $X^{(i)}$ 在 $S$ 上为有限型（I，6.3.4 和 6.3.8））。
由于 $Y$ 是局部 Noether 概形，且 $f_1\times_S f_2$ 是固有态射（II，5.4.2），
由（6.2.5）可知（6.7.3）的诸项 ${}^{(b)}\mathscr E_{pq}^2$ 是凝聚
$\mathscr O_Y$-模。由假设 $2^\circ$，（6.7.3）的所有谱序列都是双正则的，
故由（$0_{\mathrm I}$，11.1.8）得出结论。
\end{proof}

\begin{env}[6.7.10]
现设 $Y'^{(i)}$ 是两个 $S$-预概形，
$v_i:Y'^{(i)}\to Y^{(i)}$ 是两个 $S$-态射（$i=1,2$），
$v=v_1\times_S v_2$ 是它们的积；这是一个 $S$-态射 $Y'\to Y$，其中置
$Y'=Y'^{(1)}\times_S Y'^{(2)}$。另一方面，对 $i=1,2$，考虑一个
$S$-预概形 $X'^{(i)}$ 以及两个 $S$-态射
$u_i:X'^{(i)}\to X^{(i)}$、$f'_i:X'^{(i)}\to Y'^{(i)}$，使诸图
\[
\xymatrix@C=36pt@R=28pt{
  X'^{(i)}\ar[r]^{u_i}\ar[d]_{f'_i}
  & X^{(i)}\ar[d]^{f_i}
  \\
  Y'^{(i)}\ar[r]_{v_i}
  & Y^{(i)}
}
\tag{6.7.10.1}\label{III.6.7.10.1.fr}
\]
交换，并设诸态射 $f'_i$ \emph{分离且拟紧}。于是有谱函子的典范
$\mathscr O_{Y'}$-同态
\[
  v^*\bigl({}^{(t)}\mathscr E
  (f_1,f_2;\mathscr P_\bullet^{(1)},\mathscr P_\bullet^{(2)})\bigr)
  \longrightarrow
  {}^{(t)}\mathscr E
  (f'_1,f'_2;u_1^*(\mathscr P_\bullet^{(1)}),
                 u_2^*(\mathscr P_\bullet^{(2)}))
  \tag{6.7.10.2}\label{III.6.7.10.2.fr}
\]
其中 $t=a$、$a'$、$b$、$b'$、$c$、$d$。为定义这些同态，先假设
$S=\operatorname{Spec}(A)$、
$Y^{(i)}=\operatorname{Spec}(B_i)$、
$Y'^{(i)}=\operatorname{Spec}(B'_i)$ 均为仿射的；此时 $X^{(i)}$ 和 $X'^{(i)}$
是拟紧概形。为了计算谱序列
${}^{(t)}\mathscr E
(f_1,f_2;\mathscr P_\bullet^{(1)},\mathscr P_\bullet^{(2)})$，如（6.6.1）那样，
考虑有限覆盖 $\mathfrak U^{(i)}$，它由 $X^{(i)}$ 的仿射开集组成（$i=1,2$）；
为了计算
${}^{(t)}\mathscr E
(f'_1,f'_2;u_1^*(\mathscr P_\bullet^{(1)}),
u_2^*(\mathscr P_\bullet^{(2)}))$，考虑有限覆盖 $\mathfrak U'^{(i)}$

\oldpage[III]{163}

由 $X'^{(i)}$ 的仿射开集组成，并分别\emph{细于}覆盖
$u_i^{-1}(\mathfrak U^{(i)})$（$i=1,2$）。存在从双复形
$C^\bullet(\mathfrak U^{(i)},\mathscr P_\bullet^{(i)})
=L_{\bullet\bullet}^{(i)}$ 到
$C^\bullet(u_i^{-1}(\mathfrak U^{(i)}),u_i^*(\mathscr P_\bullet^{(i)}))$
的典范同态（$0_{\mathrm I}$，4.4.3.2）；此外，选取一个从
$\mathfrak U'^{(i)}$ 到 $u_i^{-1}(\mathfrak U^{(i)})$ 的单纯映射（G，II，5.7），
便定义双复形同态
\[
  C^\bullet(u_i^{-1}(\mathfrak U^{(i)}),u_i^*(\mathscr P_\bullet^{(i)}))
  \longrightarrow
  C^\bullet(\mathfrak U'^{(i)},u_i^*(\mathscr P_\bullet^{(i)}))
\]
从而通过复合得到双复形同态
\[
  L_{\bullet\bullet}^{(i)}
  \longrightarrow
  L_{\bullet\bullet}^{\prime(i)}
  =C^\bullet(\mathfrak U'^{(i)},u_i^*(\mathscr P_\bullet^{(i)})).
\]
而且，改变单纯映射时，这一同态被一个与之同伦的同态替代（G，II，5.7.1）；
于是得到良定义的谱函子同态：
\[
  {}^{(t)}\mathscr E
  (\mathfrak U^{(1)},\mathfrak U^{(2)};
   \mathscr P_\bullet^{(1)},\mathscr P_\bullet^{(2)})
  \longrightarrow
  {}^{(t)}\mathscr E
  (\mathfrak U'^{(1)},\mathfrak U'^{(2)};
   u_1^*(\mathscr P_\bullet^{(1)}),u_2^*(\mathscr P_\bullet^{(2)})).
  \tag{6.7.10.3}\label{III.6.7.10.3.fr}
\]
立即可验证：若 $\mathfrak V^{(i)}$ 是 $X^{(i)}$ 的一个细于
$\mathfrak U^{(i)}$ 的有限仿射覆盖，而 $\mathfrak V'^{(i)}$ 是
$X'^{(i)}$ 的一个同时细于 $u_i^{-1}(\mathfrak V^{(i)})$ 和
$\mathfrak U'^{(i)}$ 的有限仿射覆盖，则图
\[
\xymatrix@C=34pt@R=30pt{
  C^\bullet(\mathfrak U^{(i)},\mathscr P_\bullet^{(i)})
  \ar[r]\ar[d]
  & C^\bullet(\mathfrak V^{(i)},\mathscr P_\bullet^{(i)})\ar[d]
  \\
  C^\bullet(\mathfrak U'^{(i)},u_i^*(\mathscr P_\bullet^{(i)}))
  \ar[r]
  & C^\bullet(\mathfrak V'^{(i)},u_i^*(\mathscr P_\bullet^{(i)}))
}
\]
交换。这意味着同态（6.7.10.3）在本质上与所考虑的覆盖
$\mathfrak U^{(i)}$ 和 $\mathfrak U'^{(i)}$ 无关。因此，我们实际定义了
$A$-模的同态
\[
  {}^{(t)}E
  (X^{(1)},X^{(2)};\mathscr P_\bullet^{(1)},\mathscr P_\bullet^{(2)})
  \longrightarrow
  {}^{(t)}E
  (X'^{(1)},X'^{(2)};
   u_1^*(\mathscr P_\bullet^{(1)}),u_2^*(\mathscr P_\bullet^{(2)})).
  \tag{6.7.10.4}\label{III.6.7.10.4.fr}
\]
然而，由 $u_i^*(\mathscr P_\bullet^{(i)})$ 的定义及（6.7.10.1）的交换性可见，
这一同态也是 $(B_1\otimes_A B_2)$-模的同态；由于（6.7.10.4）的右端由
$(B'_1\otimes_A B'_2)$-模组成，便从（6.7.10.4）典范得到一个
$(B'_1\otimes_A B'_2)$-模同态
\[
  {}^{(t)}E
  (X^{(1)},X^{(2)};\mathscr P_\bullet^{(1)},\mathscr P_\bullet^{(2)})
  \otimes_{B_1\otimes_A B_2}(B'_1\otimes_A B'_2)
  \longrightarrow
  {}^{(t)}E
  (X'^{(1)},X'^{(2)};
   u_1^*(\mathscr P_\bullet^{(1)}),u_2^*(\mathscr P_\bullet^{(2)}))
  \tag{6.7.10.5}\label{III.6.7.10.5.fr}
\]
计及（I，1.6.5），这在所考虑的特殊情形中恰是所求同态（6.7.10.2）。

还须依照（6.7.1）和（6.7.2）的拼接过渡到一般情形；第二个过渡是立即的。
对于第一个过渡，如（6.7.1）那样考虑元素 $g_i\in B_i$ 及其像
$g'_i\in B'_i$，考虑张量积 $g=g_1\otimes g_2$（它位于
$B=B_1\otimes_A B_2$ 中）及其像 $g'=g'_1\otimes g'_2$（它位于
$B'=B'_1\otimes_A B'_2$ 中）；一切归结为使用

\oldpage[III]{164}

典范同构
\[
  (M\otimes_B B')_{g'}
  \xrightarrow{\sim}
  M_g\otimes_{B_g}B'_{g'}
\]
（$0_{\mathrm I}$，1.5.4）；细节留给读者。
\end{env}

\begin{env}[6.7.11]
以上为两个 $S$-态射 $X^{(i)}\to Y^{(i)}$ 和两个下有界拟相干模复形
$\mathscr P_\bullet^{(i)}$ 所建立的整体超 Tor 理论，可立即推广到如下的一般情形：
给定一个预概形 $S$、一个 $S$-预概形的有限族 $Y^{(i)}$（$i\in I$）、
一个分离且拟紧的 $S$-态射的有限族
$f_i:X^{(i)}\to Y^{(i)}$，并且对每个 $i$，给定一个下有界拟相干
$\mathscr O_{X^{(i)}}$-模复形 $\mathscr P_\bullet^{(i)}$。若 $Y$ 是诸
$S$-预概形 $Y^{(i)}$ 的积，则对每个整数 $n\in\mathbf Z$，定义一个拟相干
$\mathscr O_Y$-模
$\mathfrak{Tor}_n^S((f_i)_{i\in I};
(\mathscr P_\bullet^{(i)})_{i\in I})$；这些模构成协变的 $\partial$-函子，其每个复形变元为
$\mathscr P_\bullet^{(i)}$；此外，这个函子是六个谱函子
${}^{(t)}\mathscr E((f_i)_{i\in I};(\mathscr P_\bullet^{(i)})_{i\in I})$
的共同收敛项。我们把针对 $I=\{1,2\}$ 在上面作出的定义和论证推广到这一
一般情形的工作留给读者。仅须指出，当 $I$ 只有一个元素时，便回到（6.2.7）中定义的
超上同调 $\mathscr H^\bullet(f,\mathscr P_\bullet)$（这一点在（6.7.4）中已经指出）。
当 $I$ 是区间 $1\leq i\leq m$（作为 $\mathbf N$ 的子集）时，我们写
\[
  \mathfrak{Tor}_n^S
  (f_1,\ldots,f_m;
   \mathscr P_\bullet^{(1)},\ldots,\mathscr P_\bullet^{(m)})
  \quad\text{pour}\quad
  \mathfrak{Tor}_n^S
  ((f_i)_{i\in I};(\mathscr P_\bullet^{(i)})_{i\in I}).
\]
\end{env}

\begin{proposition}[6.7.12]
沿用（6.7.11）的记号，设 $J$ 是 $I$ 的一个子集，使得对
$i\in I-J$，有 $X^{(i)}=Y^{(i)}=S$，$f_i$ 为恒等态射，并且
$\mathscr P_\bullet^{(i)}$ 等于只在次数 $0$ 有项且该项为
$\mathscr O_S$ 的复形。于是存在 $\partial$-函子的典范同构
\[
  \mathfrak{Tor}_\bullet^S
  ((f_i)_{i\in I};(\mathscr P_\bullet^{(i)})_{i\in I})
  \xrightarrow{\sim}
  \mathfrak{Tor}_\bullet^S
  ((f_i)_{i\in J};(\mathscr P_\bullet^{(i)})_{i\in J}).
  \tag{6.7.12.1}\label{III.6.7.12.1.fr}
\]
\end{proposition}

\begin{proof}
只须在 $S$ 和诸 $Y^{(i)}$ 均为仿射时定义这个同构，拼接照常进行。
对于 $i\in I-J$，可以取覆盖 $\mathfrak U^{(i)}$，它只由集合
$S$ 组成；此时
$L_{\bullet\bullet}^{(i)}=C^\bullet(\mathfrak U^{(i)},
\mathscr P_\bullet^{(i)})$ 只剩次数为 $(0,0)$ 的一项，且该项等于
$\Gamma(S,\mathscr O_S)=A(S)$；同构（6.7.12.1）遂为显然。
\end{proof}

\begin{remark}[6.7.13]
沿用（6.7.3）的记号，考虑典范的 $S$-同构
$Y^{(1)}\times_S Y^{(2)}\to Y^{(2)}\times_S Y^{(1)}$（I，3.3.5）；
则
$\mathfrak{Tor}_\bullet^S(f_1,f_2;
\mathscr P_\bullet^{(1)},\mathscr P_\bullet^{(2)})$ 在这个同构下的像是
$\mathfrak{Tor}_\bullet^S(f_2,f_1;
\mathscr P_\bullet^{(2)},\mathscr P_\bullet^{(1)})$。由于问题是局部的，
可归结为（6.6.2）中考虑的情形；若以
$M_{\bullet\bullet\bullet}^{(i)}$ 表示 $L_{\bullet\bullet}^{(i)}$ 的一个
Cartan--Eilenberg 投射消解（$i=1,2$），则所考虑的同构把
$M_{\bullet\bullet\bullet}^{(1)}\otimes_A
M_{\bullet\bullet\bullet}^{(2)}$ 变成
$M_{\bullet\bullet\bullet}^{(2)}\otimes_A
M_{\bullet\bullet\bullet}^{(1)}$；考虑这些三重复形所关联单复形的同调，便得到所述断言。
\end{remark}

\subsection{整体超 Tor 的结合性谱序列}

\begin{env}[6.8.1]
沿用（6.7.11）的假设和记号（特别地，假定诸
$\mathscr P_\bullet^{(i)}$ 下有界），设给定一个分割
$(I_j)_{j\in J}$，它把指标集 $I$ 分开；我们要在超 Tor
$\mathfrak{Tor}_n^S((f_i)_{i\in I};(\mathscr P_\bullet^{(i)})_{i\in I})$
与每一个“部分”超 Tor
\[
  \mathcal T_{\bullet j}
  =\mathfrak{Tor}_\bullet^S
  ((f_i)_{i\in I_j};(\mathscr P_\bullet^{(i)})_{i\in I_j}).
\]

\oldpage[III]{165}
之间建立一种“结合性”关系。为简化书写，我们只考虑 $I$ 是区间
$1\leq i\leq m$，且分割 $(I_j)$ 由两个区间
$\{1,2,\ldots,r\}$ 和 $\{r+1,\ldots,m\}$ 组成的情形。
\end{env}

\begin{proposition}[6.8.2]
存在一个典范的双正则谱函子（称为“结合性谱函子”），记作
\[
  {}^{(e)}\mathscr E^S
  (f_1,\ldots,f_m;
   \mathscr P_\bullet^{(1)},\ldots,\mathscr P_\bullet^{(m)})
\]
（或简记作 ${}^{(e)}\mathscr E(f_1,\ldots,f_m;
\mathscr P_\bullet^{(1)},\ldots,\mathscr P_\bullet^{(m)})$），其
收敛项为
$\mathfrak{Tor}_\bullet^S(f_1,\ldots,f_m;
\mathscr P_\bullet^{(1)},\ldots,\mathscr P_\bullet^{(m)})$，而其
$E_2$ 项由下式给出：
\[
\begin{split}
  {}^{(e)}\mathscr E_{pq}^2
  =\bigoplus_{q_1+q_2=q}
  \mathscr{T}\!or_p^S\Bigl(
  &\mathfrak{Tor}_{q_1}^S
  (f_1,\ldots,f_r;
   \mathscr P_\bullet^{(1)},\ldots,\mathscr P_\bullet^{(r)}),\\
  &\mathfrak{Tor}_{q_2}^S
  (f_{r+1},\ldots,f_m;
   \mathscr P_\bullet^{(r+1)},\ldots,\mathscr P_\bullet^{(m)})
  \Bigr).
\end{split}
\]
\end{proposition}

\begin{proof}
在这一陈述中，我们把 $Y$ 典范地等同于积
$Z^{(1)}\times_S Z^{(2)}$，其中
$Z^{(1)}=Y^{(1)}\times_S Y^{(2)}\times_S\cdots\times_S Y^{(r)}$，而
$Z^{(2)}=Y^{(r+1)}\times_S\cdots\times_S Y^{(m)}$。我们只考虑
$S$ 与诸 $Y^{(i)}$ 均为仿射概形的情形；借助（6.7.1）和（6.7.2）中
建立的方法，可以由这一特殊情形过渡到一般情形，推理的细节（并无困难）留给读者。
因此，我们将证明下面的推论。
\end{proof}

\begin{corollary}[6.8.3]
设 $A$ 是一个环，$S=\operatorname{Spec}(A)$，$X^{(i)}$
（$1\leq i\leq m$）是拟紧 $S$-概形，并且对每个 $i$，设
$\mathscr P_\bullet^{(i)}$ 是下有界的拟相干
$\mathscr O_{X^{(i)}}$-模复形。存在一个典范的双正则谱函子，其收敛项为
\[
  \mathbf{Tor}_\bullet^S
  (X^{(1)},\ldots,X^{(m)};
   \mathscr P_\bullet^{(1)},\ldots,\mathscr P_\bullet^{(m)})
\]
而其 $E_2$ 项由下式给出：
\[
\begin{split}
  {}^{(e)}E_{pq}^2
  =\bigoplus_{q_1+q_2=q}
  \operatorname{Tor}_p^A\Bigl(
  &\mathbf{Tor}_{q_1}^S
  (X^{(1)},\ldots,X^{(r)};
   \mathscr P_\bullet^{(1)},\ldots,\mathscr P_\bullet^{(r)}),\\
  &\mathbf{Tor}_{q_2}^S
  (X^{(r+1)},\ldots,X^{(m)};
   \mathscr P_\bullet^{(r+1)},\ldots,\mathscr P_\bullet^{(m)})
  \Bigr).
\end{split}
\]
\end{corollary}

\begin{proof}
按照（6.6.2）给出的定义，所考虑的超 Tor 的计算如下进行：对每个 $i$，取
一个有限仿射开覆盖 $\mathfrak U^{(i)}$ 来覆盖 $X^{(i)}$，考虑双复形
$L_{\bullet\bullet}^{(i)}=C^\bullet(\mathfrak U^{(i)},
\mathscr P_\bullet^{(i)})$，取每一个这类双复形的一个 Cartan--Eilenberg
投射消解 $M_{\bullet\bullet\bullet}^{(i)}$（其意义如（0，11.7.1）），形成这些
三重复形的张量积
\[
  M_{\bullet\bullet\bullet}
  =\bigotimes_{i=1}^m M_{\bullet\bullet\bullet}^{(i)}
\]
并取 $M_{\bullet\bullet\bullet}$ 的同调。把 $M_{\bullet\bullet\bullet}$ 看作
单复形 $N_\bullet$，它是以下两个单复形的张量积：
\[
  N'_\bullet=\bigotimes_{i=1}^r M_{\bullet\bullet\bullet}^{(i)},
  \qquad
  N''_\bullet=\bigotimes_{i=r+1}^m M_{\bullet\bullet\bullet}^{(i)},
\]
其中 $N'_\bullet$ 和 $N''_\bullet$ 的分次由诸
$M_{\bullet\bullet\bullet}^{(i)}$ 的总次数之和给出。此外，诸 $A$-模
（即复形 $N'_\bullet$ 和 $N''_\bullet$ 的各项）都是投射模，故由（6.5.9）可得
\[
  H_\bullet(M_{\bullet\bullet\bullet})
  =\mathbf{Tor}_\bullet^A(N'_\bullet,N''_\bullet);
\]
于是，所求谱序列正是把（6.3.2.2）的谱序列应用于复形 $N'_\bullet$ 和
$N''_\bullet$ 所得的谱序列；这里还要利用前文对这些复形的同调模所作的解释
（即把开头的说明分别应用于两个部分积
$X^{(1)}\times_S\cdots\times_S X^{(r)}$ 和
$X^{(r+1)}\times_S\cdots\times_S X^{(m)}$）。最后，正则性来自（6.3.2）以及
如下事实：诸 $\mathscr P_\bullet^{(i)}$ 下有界，因而诸
$M_{\bullet\bullet\bullet}^{(i)}$ 也下有界；所以 $N'_\bullet$ 和
$N''_\bullet$ 都下有界。
\end{proof}

\oldpage[III]{166}
\subsection{整体超 Tor 中的基变换谱序列}

\begin{env}[6.9.1]
仍沿用（6.7.11）的假设和记号（特别地，仍假定诸
$\mathscr P_\bullet^{(i)}$ 下有界），考虑预概形的一个态射 $g:S'\to S$，并令
$Y'^{(i)}=Y_{(S')}^{(i)}$，$X'^{(i)}=X_{(S')}^{(i)}$ 以及
$\mathscr P_\bullet'^{(i)}=
\mathscr P_\bullet^{(i)}\otimes_{\mathscr O_S}\mathscr O_{S'}$；因而
$\mathscr P_\bullet'^{(i)}$ 是拟相干 $\mathscr O_{X'^{(i)}}$-模复形。令
$f'_i=(f_i)_{(S')}:X'^{(i)}\to Y'^{(i)}$；这是一个分离且拟紧的态射
（I，5.5.1 和 6.6.4）。我们要研究拟相干 $\mathscr O_{Y'}$-模
$\mathfrak{Tor}_n^{S'}((f'_i)_{i\in I};
(\mathscr P_\bullet'^{(i)})_{i\in I})$ 与
$\mathfrak{Tor}_n^S((f_i)_{i\in I};
(\mathscr P_\bullet^{(i)})_{i\in I})
\otimes_{\mathscr O_S}\mathscr O_{S'}$ 之间的关系，其中
$Y'=Y\times_S S'=Y_{(S')}$。以下是一个特别简单的情形；当 $I$ 只有一个元素且
$\mathscr P_\bullet$ 只含一个模时，它归结为（1.4.15）：
\end{env}

\begin{proposition}[6.9.2]
若态射 $g:S'\to S$ 平坦，则有 $\partial$-函子的典范同构
（以诸 $\mathscr P_\bullet^{(i)}$ 为变元）：
\[
  \mathfrak{Tor}_\bullet^S
  ((f_i)_{i\in I};(\mathscr P_\bullet^{(i)})_{i\in I})
  \otimes_{\mathscr O_S}\mathscr O_{S'}
  \xrightarrow{\sim}
  \mathfrak{Tor}_\bullet^{S'}
  ((f'_i)_{i\in I};(\mathscr P_\bullet'^{(i)})_{i\in I}).
  \tag{6.9.2.1}\label{III.6.9.2.1.fr}
\]
\end{proposition}

\begin{proof}
仍可只考虑 $S$、$S'$ 和诸 $Y^{(i)}$ 均为仿射的情形，拼接按（6.7.1）和
（6.7.2）的方法进行。设 $S=\operatorname{Spec}(A)$，
$S'=\operatorname{Spec}(A')$，并对每个 $i$ 取一个仿射开覆盖
$\mathfrak U^{(i)}$ 来覆盖 $X^{(i)}$；若 $u_i:X'^{(i)}\to X^{(i)}$ 是典范投影，则
$u_i^{-1}(\mathfrak U^{(i)})$ 是 $X'^{(i)}$ 的一个仿射开覆盖（II，1.5.5），
记作 $\mathfrak U'^{(i)}$；于是显然
\[
  C^\bullet(\mathfrak U'^{(i)},\mathscr P_\bullet'^{(i)})
  =C^\bullet(\mathfrak U^{(i)},\mathscr P_\bullet^{(i)})\otimes_A A',
\]
从而同构（6.9.2.1）的存在性立得。事实上，若
$M_{\bullet\bullet\bullet}^{(i)}$ 是
$L_{\bullet\bullet}^{(i)}=C^\bullet(\mathfrak U^{(i)},
\mathscr P_\bullet^{(i)})$ 的 Cartan--Eilenberg 投射消解（其意义如
（0，11.7.1）），且由 $A$-模组成，则
$M_{\bullet\bullet\bullet}^{(i)}\otimes_A A'$ 是
$L_{\bullet\bullet}^{(i)}\otimes_A A'$ 的、由 $A'$-模组成的
Cartan--Eilenberg 投射消解；这是因为 $A'$ 是平坦 $A$-模。同一假设还表明
\[
  H_\bullet\left(
    \bigotimes_{i=1}^m
    (M_{\bullet\bullet\bullet}^{(i)}\otimes_A A')
  \right)
  =H_\bullet\left(
    \bigotimes_{i=1}^m M_{\bullet\bullet\bullet}^{(i)}
  \right)\otimes_A A'.
\]
\end{proof}

\begin{env}[6.9.2.2]
应注意，当 $I$ 只含元素 $1$ 时，公式（6.9.2.1）可由（6.7.7）直接推出：在其中取
$Y_2=X_2=S'$、$f_2=1_{S'}$，并令复形 $\mathscr P_\bullet^{(2)}$
只在次数 $0$ 有一项，且该项等于 $\mathscr O_{S'}$；此时我们知道，超上同调
$\mathscr H^n(1_{S'},\mathscr O_{S'})$ 对每个 $n\ne0$ 都为零，并化为
$\mathscr O_{S'}$（当 $n=0$ 时）（6.2.1）。
\end{env}

在一般情形下，我们不再使用 $\mathscr O_{S'}$，而引入一个复形
$\mathscr Q'_\bullet$，它是下有界的拟相干 $\mathscr O_{S'}$-模复形；于是，为简化起见，
若取 $I=\{1,2,\ldots,m\}$，便可以考虑 $\partial$-函子
\[
  \mathfrak{Tor}_\bullet^S
  (f_1,\ldots,f_m,1_{S'};
   \mathscr P_\bullet^{(1)},\ldots,\mathscr P_\bullet^{(m)},
   \mathscr Q'_\bullet).
\]

\begin{proposition}[6.9.3]
存在三个典范的双正则谱函子，记为
${}^{(t)}\mathscr E(f_1,\ldots,f_m;
\mathscr P_\bullet^{(1)},\ldots,\mathscr P_\bullet^{(m)},
\mathscr Q'_\bullet)$（其中 $t=e$、$f$ 或 $f'$）；它们具有共同的收敛项
$\mathfrak{Tor}_\bullet^S(f_1,\ldots,f_m,1_{S'};
\mathscr P_\bullet^{(1)},\ldots,\mathscr P_\bullet^{(m)},
\mathscr Q'_\bullet)$，并且它们的 $E_2$ 项分别为

\oldpage[III]{167}
\[
\begin{split}
  {}^{(e)}\mathscr E_{pq}^2
  =\bigoplus_{q'+q''=q}
  \mathscr{T}\!or_p^S\Bigl(
  &\mathfrak{Tor}_{q'}^S
  (f_1,\ldots,f_m;
   \mathscr P_\bullet^{(1)},\ldots,\mathscr P_\bullet^{(m)}),\\
  &\mathscr H_{q''}(\mathscr Q'_\bullet)
  \Bigr),
\end{split}
\]
\[
\begin{split}
  {}^{(f)}\mathscr E_{pq}^2
  =\bigoplus_{q_1+q_2+\cdots+q_{m+1}=q}
  \mathscr{T}\!or_p^{S'}\Bigl(
  &\mathfrak{Tor}_{q_1}^S
  (f_1,1_{S'};\mathscr P_\bullet^{(1)},\mathscr O_{S'}),\ldots,\\
  &\mathfrak{Tor}_{q_m}^S
  (f_m,1_{S'};\mathscr P_\bullet^{(m)},\mathscr O_{S'}),
  \mathscr H_{q_{m+1}}(\mathscr Q'_\bullet)
  \Bigr),
\end{split}
\]
\[
\begin{split}
  {}^{(f')}\mathscr E_{pq}^2
  =\bigoplus_{q_1+\cdots+q_m=q}
  \mathfrak{Tor}_p^{S'}\Bigl(
  &f'_1,\ldots,f'_m,1_{S'};\\
  &\mathscr{T}\!or_{q_1}^S
  (\mathscr P_\bullet^{(1)},\mathscr O_{S'}),\ldots,
  \mathscr{T}\!or_{q_m}^S
  (\mathscr P_\bullet^{(m)},\mathscr O_{S'}),
  \mathscr Q'_\bullet
  \Bigr).
\end{split}
\]
\end{proposition}

\begin{proof}
谱序列 $(e)$ 正是（6.8.2）中取 $r=m$ 时的结合性谱序列。为定义另外两个谱序列，
我们仍只考虑 $S$、$S'$ 和诸 $Y^{(i)}$ 均为仿射的情形；从这一情形过渡到一般情形，
可用（6.7.1）和（6.7.2）的方法完成，其细节留给读者。因此，我们将证明以下推论。
\end{proof}

\begin{corollary}[6.9.4]
设 $A$ 是一个环，$A'$ 是一个 $A$-代数，
$S=\operatorname{Spec}(A)$，$S'=\operatorname{Spec}(A')$；设
$X^{(i)}$（$1\leq i\leq m$）是拟紧 $S$-概形，并且对每个
$i$，令 $X'^{(i)}=X_{(S')}^{(i)}$，这是一个拟紧 $S'$-概形。对每个
$i$，设 $\mathscr P_\bullet^{(i)}$ 是拟相干
$\mathscr O_{X^{(i)}}$-模复形；最后，设 $Q'_\bullet$ 是一个
$A'$-模复形，并假定这些复形都下有界。存在三个以诸
$\mathscr P_\bullet^{(i)}$ 和 $Q'_\bullet$ 为变元的双正则谱函子；它们具有共同的收敛项
\[
  \mathbf{Tor}_\bullet^S
  (X^{(1)},\ldots,X^{(m)},S';
   \mathscr P_\bullet^{(1)},\ldots,\mathscr P_\bullet^{(m)},
   \widetilde{Q'_\bullet})
\]
并且它们的 $E_2$ 项分别为
\[
\begin{split}
  {}^{(e)}E_{pq}^2
  =\bigoplus_{q'+q''=q}
  \operatorname{Tor}_p^A\Bigl(
  &\mathbf{Tor}_{q'}^S
  (X^{(1)},\ldots,X^{(m)};
   \mathscr P_\bullet^{(1)},\ldots,\mathscr P_\bullet^{(m)}),\\
  &H_{q''}(Q'_\bullet)
  \Bigr),
\end{split}
\]
\[
\begin{split}
  {}^{(f)}E_{pq}^2
  =\bigoplus_{q_1+q_2+\cdots+q_{m+1}=q}
  \operatorname{Tor}_p^{A'}\Bigl(
  &\mathbf{Tor}_{q_1}^S
  (X^{(1)},S';\mathscr P_\bullet^{(1)},\mathscr O_{S'}),\ldots,\\
  &\mathbf{Tor}_{q_m}^S
  (X^{(m)},S';\mathscr P_\bullet^{(m)},\mathscr O_{S'}),
  H_{q_{m+1}}(Q'_\bullet)
  \Bigr),
\end{split}
\]
\[
\begin{split}
  {}^{(f')}E_{pq}^2
  =\bigoplus_{q_1+\cdots+q_m=q}
  \mathbf{Tor}_p^{S'}\Bigl(
  &X'^{(1)},\ldots,X'^{(m)},S';\\
  &\mathscr{T}\!or_{q_1}^S
  (\mathscr P_\bullet^{(1)},\mathscr O_{S'}),\ldots,
  \mathscr{T}\!or_{q_m}^S
  (\mathscr P_\bullet^{(m)},\mathscr O_{S'}),
  \widetilde{Q'_\bullet}
  \Bigr).
\end{split}
\]
\end{corollary}

\begin{proof}
对于这些谱函子中的第一个，我们不再赘述，因为它已经在（6.8.3）中处理过，这里列出它
只是为了备忘。为定义另外两个谱函子，对每个 $i$，取有限仿射开覆盖
$\mathfrak U^{(i)}$ 来覆盖 $X^{(i)}$；若 $u_i:X'^{(i)}\to X^{(i)}$ 是典范投影，则考虑相应的有限仿射开覆盖
$\mathfrak U'^{(i)}=u_i^{-1}(\mathfrak U^{(i)})$。

由（6.6.6），为了得到
$\mathbf{Tor}_\bullet^S(X^{(1)},\ldots,X^{(m)},S';
\mathscr P_\bullet^{(1)},\ldots,\mathscr P_\bullet^{(m)},
\widetilde{Q'_\bullet})$，可对 $1\leq i\leq m$ 取一个 Cartan--Eilenberg 投射消解
$M_{\bullet\bullet\bullet}^{(i)}$；它所消解的双复形为
$L_{\bullet\bullet}^{(i)}=C^\bullet(\mathfrak U^{(i)},
\mathscr P_\bullet^{(i)})$（其意义如（0，11.7.1））。再考虑三重复形
\[
  M_{\bullet\bullet\bullet}
  =M_{\bullet\bullet\bullet}^{(1)}\otimes_A
   M_{\bullet\bullet\bullet}^{(2)}\otimes_A\cdots\otimes_A
   M_{\bullet\bullet\bullet}^{(m)}\otimes_A Q'_\bullet
\]
（这里把 $Q'_\bullet$ 看作一个三重复形，其后两个次数均集中于次数 $0$），并取其同调。
若令 $M_{\bullet\bullet\bullet}'^{(i)}=
M_{\bullet\bullet\bullet}^{(i)}\otimes_A A'$，则由于 $Q'_\bullet$ 是一个
$A'$-模复形，有
\[
  M_{\bullet\bullet\bullet}
  =M_{\bullet\bullet\bullet}'^{(1)}\otimes_{A'}
   M_{\bullet\bullet\bullet}'^{(2)}\otimes_{A'}\cdots\otimes_{A'}
   M_{\bullet\bullet\bullet}'^{(m)}\otimes_{A'} Q'_\bullet.
\]

现在把每个复形 $M_{\bullet\bullet\bullet}'^{(i)}$ 按其总次数看成单复形；并注意这个复形由
$A'$-投射模组成。由（6.3.7）（推广到任意多个复形）可知，
$H_\bullet(M_{\bullet\bullet\bullet})$ 也等于
\[
  \operatorname{Tor}_\bullet^{A'}
  (M_{\bullet\bullet\bullet}'^{(1)},\ldots,
   M_{\bullet\bullet\bullet}'^{(m)},Q'_\bullet);
\]
所以由（6.5.15），它是一个具有所需正则性的谱序列的收敛项（
$M_{\bullet\bullet\bullet}^{(i)}$ 的三个次数，在
$\mathscr P_\bullet^{(i)}$ 下有界时，都下有界），
而该谱序列的 $E_2$ 项为
\[
  E_{pq}^2
  =\bigoplus_{q_1+\cdots+q_{m+1}=q}
  \operatorname{Tor}_p^{A'}
  (H_{q_1}(M_{\bullet\bullet\bullet}'^{(1)}),\ldots,
   H_{q_m}(M_{\bullet\bullet\bullet}'^{(m)}),H_{q_{m+1}}(Q'_\bullet)).
\]

\oldpage[III]{168}
此外，由整体超 Tor 的定义，有
\[
  H_{q_i}(M_{\bullet\bullet\bullet}'^{(i)})
  =H_{q_i}(M_{\bullet\bullet\bullet}^{(i)}\otimes_A A')
  =\mathbf{Tor}_{q_i}^S
  (X^{(i)},S';\mathscr P_\bullet^{(i)},\mathscr O_{S'})
\]
，这便给出所求的谱序列 $(f)$。另一方面，可以把
$M_{\bullet\bullet\bullet}'^{(i)}$ 看成一个双复形，其中第一个次数是三重复形
$M_{\bullet\bullet\bullet}^{(i)}$ 的第一、第二个次数之和，而第二个次数是该三重复形的第三个次数；
由于构成诸 $M_{\bullet\bullet\bullet}'^{(i)}$ 的模都是 $A'$-投射模，超同调的一般理论表明，
双复形
\[
  M_{\bullet\bullet\bullet}'^{(1)}\otimes_{A'}
  M_{\bullet\bullet\bullet}'^{(2)}\otimes_{A'}\cdots\otimes_{A'}
  M_{\bullet\bullet\bullet}'^{(m)}\otimes_{A'} Q'_\bullet
\]
的同调典范同构于它的超同调（0，11.6.5）；所以它是一个谱序列的收敛项，而该谱序列的
$E_2$ 项为
\[
  E_{pq}^2
  =\bigoplus_{q_1+\cdots+q_{m+1}=q}
  \operatorname{Tor}_p^{A'}
  (H_{q_1}^{\mathrm{II}}(M_{\bullet\bullet\bullet}'^{(1)}),\ldots,
   H_{q_m}^{\mathrm{II}}(M_{\bullet\bullet\bullet}'^{(m)}),
   H_{q_{m+1}}^{\mathrm{II}}(Q'_\bullet)).
\]
由于 $Q'_\bullet$ 的第二个次数只集中在 $0$，故有
$H_n^{\mathrm{II}}(Q'_\bullet)=0$（$n\ne0$），而
$H_0^{\mathrm{II}}(Q'_\bullet)=Q'_\bullet$；上式也可写成
\[
  E_{pq}^2
  =\bigoplus_{q_1+\cdots+q_m=q}
  \operatorname{Tor}_p^{A'}
  (H_{q_1}^{\mathrm{II}}(M_{\bullet\bullet\bullet}'^{(1)}),\ldots,
   H_{q_m}^{\mathrm{II}}(M_{\bullet\bullet\bullet}'^{(m)}),Q'_\bullet).
\]
此外，由（6.3.4），有
\[
  H_{q_i}^{\mathrm{II}}(M_{\bullet\bullet\bullet}'^{(i)})
  =H_{q_i}^{\mathrm{II}}(M_{\bullet\bullet\bullet}^{(i)}\otimes_A A')
  =\operatorname{Tor}_{q_i}^A(L_{\bullet\bullet}^{(i)},A')
\]
但是，$L_{-j,k}^{(i)}=C^j(\mathfrak U^{(i)},\mathscr P_k^{(i)})$ 是诸
$\Gamma(V,\mathscr P_k^{(i)})$ 的直和，其中 $V$ 遍历由
$j+1$ 个覆盖集组成的（仿射）交，这些覆盖集取自 $\mathfrak U^{(i)}$；若
$V'=u_i^{-1}(V)$，则 $V'$ 是 $X'^{(i)}$ 中的仿射开集，而且由（6.4.1.1）有
\[
  \Gamma\bigl(V',\mathscr{T}\!or_{q_i}^S
  (\mathscr P_k^{(i)},\mathscr O_{S'})\bigr)
  =\operatorname{Tor}_{q_i}^A(\Gamma(V,\mathscr P_k^{(i)}),A')
\]
因而，双复形 $H_{q_i}^{\mathrm{II}}(M_{\bullet\bullet\bullet}'^{(i)})$ 可表示为
\[
  C^\bullet\bigl(\mathfrak U'^{(i)},
  \mathscr{T}\!or_{q_i}^S(\mathscr P_\bullet^{(i)},\mathscr O_{S'})\bigr),
\]
这最终给出了 $E_2$ 项的所求表达式；所对应的是谱序列 $(f')$。在所述条件下，该谱序列的双正则性
也照常验证；须注意，当 $\mathscr P_\bullet^{(i)}$ 下有界时，
$M_{\bullet\bullet\bullet}'^{(i)}$ 的所有次数都下有界。
\end{proof}

\begin{remark}[6.9.5]
与（6.7.6）相同，可以看出，用分别与
$\mathscr P_\bullet^{(i)}$ 和 $\mathscr Q'_\bullet$ 同伦的复形替换它们，
并不改变谱序列 $(e)$、$(f)$ 与 $(f')$，亦即它们在典范同构意义下不变。此外，
对于谱序列 $(f)$，若复形同态
$\mathscr P_\bullet^{(i)}\to\mathscr R_\bullet^{(i)}$、
$\mathscr Q'_\bullet\to\mathscr T'_\bullet$ 在同调上给出同构
$\mathscr H_\bullet(\mathscr P_\bullet^{(i)})\xrightarrow{\sim}
\mathscr H_\bullet(\mathscr R_\bullet^{(i)})$、
$\mathscr H_\bullet(\mathscr Q'_\bullet)\xrightarrow{\sim}
\mathscr H_\bullet(\mathscr T'_\bullet)$，则它们给出谱序列的同构
\[
\begin{split}
  {}^{(f)}\mathscr E
  (f_1,\ldots,f_m;
   \mathscr P_\bullet^{(1)},\ldots,\mathscr P_\bullet^{(m)},
   \mathscr Q'_\bullet)
  \xrightarrow{\sim}\;&
  {}^{(f)}\mathscr E
  (f_1,\ldots,f_m;\\
  &\mathscr R_\bullet^{(1)},\ldots,\mathscr R_\bullet^{(m)},
   \mathscr T'_\bullet);
\end{split}
\]
证明与（6.7.6）的相同，只需考虑（6.7.6）的结果以及谱序列 $(f)$ 的正则性。
\end{remark}

\begin{corollary}[6.9.6]
在（6.9.1）的条件下，假设：
\begin{enumerate}
  \item[$1^\circ$] 诸复形 $\mathscr P_\bullet^{(i)}$ 均由
  $S$-平坦模组成，而且下列 $\mathscr O_Y$-模
  \[
    \mathfrak{Tor}_n^S
    (f_1,\ldots,f_m;
     \mathscr P_\bullet^{(1)},\ldots,\mathscr P_\bullet^{(m)})
  \]
  也都是 $S$-平坦的。
  \item[$2^\circ$] 诸 $\mathscr P_\bullet^{(i)}$ 以及
  $\mathscr Q'_\bullet$ 都下有界。
\end{enumerate}

\oldpage[III]{169}
令
$\mathscr P_\bullet'^{(i)}=
\mathscr P_\bullet^{(i)}\otimes_{\mathscr O_S}\mathscr O_{S'}$，则有函子性的典范同构
\[
\begin{split}
  &\mathfrak{Tor}_n^{S'}
  (f'_1,\ldots,f'_m,1_{S'};
   \mathscr P_\bullet'^{(1)},\ldots,
   \mathscr P_\bullet'^{(m)},\mathscr Q'_\bullet)
  \xrightarrow{\sim}\\
  &\qquad
  \bigoplus_{n'+n''=n}
  \mathfrak{Tor}_{n'}^S
  (f_1,\ldots,f_m;
   \mathscr P_\bullet^{(1)},\ldots,
   \mathscr P_\bullet^{(m)})
  \otimes_{\mathscr O_S}\mathscr H_{n''}(\mathscr Q'_\bullet).
\end{split}
\tag{6.9.6.1}\label{III.6.9.6.1.fr}
\]
特别地，当 $\mathscr Q'_\bullet$ 只含单项 $\mathscr F'$，且该项位于次数
$0$ 时，有函子性的典范同构
\[
\begin{split}
  &\mathfrak{Tor}_n^{S'}
  (f'_1,\ldots,f'_m,1_{S'};
   \mathscr P_\bullet'^{(1)},\ldots,
   \mathscr P_\bullet'^{(m)},\mathscr F')
  \xrightarrow{\sim}\\
  &\qquad
  \mathfrak{Tor}_n^S
  (f_1,\ldots,f_m;
   \mathscr P_\bullet^{(1)},\ldots,
   \mathscr P_\bullet^{(m)})
  \otimes_{\mathscr O_S}\mathscr F'
\end{split}
\tag{6.9.6.2}\label{III.6.9.6.2.fr}
\]
更特别地，当 $\mathscr F'=\mathscr O_{S'}$ 时，有
\[
\begin{split}
  &\mathfrak{Tor}_n^{S'}
  (f'_1,\ldots,f'_m;
   \mathscr P_\bullet'^{(1)},\ldots,
   \mathscr P_\bullet'^{(m)})
  \xrightarrow{\sim}\\
  &\qquad
  \mathfrak{Tor}_n^S
  (f_1,\ldots,f_m;
   \mathscr P_\bullet^{(1)},\ldots,
   \mathscr P_\bullet^{(m)})
  \otimes_{\mathscr O_S}\mathscr O_{S'}.
\end{split}
\tag{6.9.6.3}\label{III.6.9.6.3.fr}
\]
\end{corollary}

\begin{proof}
关于组成诸 $\mathscr P_\bullet^{(i)}$ 的模的平坦性假设说明，复形
$\mathscr{T}\!or_{q_i}^S(\mathscr P_\bullet^{(i)},\mathscr O_{S'})$ 在
$q_i\ne0$ 时为零（6.5.8）。因而谱序列 $(f')$ 退化；此外，假设
$2^\circ$ 还说明它是双正则的（6.9.3），所以边缘同态
\[
\begin{split}
  &\mathfrak{Tor}_n^S
  (f_1,\ldots,f_m,1_{S'};
   \mathscr P_\bullet^{(1)},\ldots,
   \mathscr P_\bullet^{(m)},\mathscr Q'_\bullet)
  \longrightarrow {}^{(f')}\mathscr E_{n0}^2\\
  &\qquad=
  \mathfrak{Tor}_n^{S'}
  (f'_1,\ldots,f'_m,1_{S'};
   \mathscr P_\bullet'^{(1)},\ldots,
   \mathscr P_\bullet'^{(m)},\mathscr Q'_\bullet)
\end{split}
\tag{6.9.6.4}\label{III.6.9.6.4.fr}
\]
是双射（0，11.1.6）。关于下列诸模的平坦性假设
\[
  \mathfrak{Tor}_n^S
  (f_1,\ldots,f_m;
   \mathscr P_\bullet^{(1)},\ldots,\mathscr P_\bullet^{(m)})
\]
说明 ${}^{(e)}\mathscr E_{pq}^2=0$（$p\ne0$）（6.5.8）。因而谱序列
$(e)$ 也退化；又因它是双正则的，边缘同态
\[
\begin{split}
  &\mathfrak{Tor}_n^S
  (f_1,\ldots,f_m,1_{S'};
   \mathscr P_\bullet^{(1)},\ldots,
   \mathscr P_\bullet^{(m)},\mathscr Q'_\bullet)
  \longrightarrow {}^{(e)}\mathscr E_{0n}^2\\
  &\qquad=
  \bigoplus_{n'+n''=n}
  \mathfrak{Tor}_{n'}^S
  (f_1,\ldots,f_m;
   \mathscr P_\bullet^{(1)},\ldots,
   \mathscr P_\bullet^{(m)})
  \otimes_{\mathscr O_S}\mathscr H_{n''}(\mathscr Q'_\bullet)
\end{split}
\tag{6.9.6.5}\label{III.6.9.6.5.fr}
\]
也是双射（0，11.1.6）；因此，结合前述两个同构便得到同构（6.9.6.1）。
同构（6.9.6.2）由此立即推出，因为这时
$\mathscr H_q(\mathscr Q'_\bullet)=0$（$q\ne0$），而且
$\mathscr H_0(\mathscr Q'_\bullet)=\mathscr F'$。最后，在（6.9.6.2）中取
$\mathscr F'=\mathscr O_{S'}$，再考虑（6.7.12），便得到同构（6.9.6.3）。
\end{proof}

\begin{corollary}[6.9.7]
在（6.9.1）的条件下，假设 $S$ 和 $S'$ 都是仿射的，并且对每个
$i$ 给定一个整数 $d_i$（$1\leq i\leq m$）。那么存在一个整数 $N$，
它只依赖于 $S$、诸 $X^{(i)}$ 和诸 $d_i$，并具有如下性质：对每个整数
$n_0$，当 $n\leq n_0$ 时，对于满足下列条件的每一族复形
$\mathscr P_\bullet^{(i)}$，都有典范同构（6.9.6.3）：
\begin{enumerate}
  \item[$1^\circ$] $\mathscr P_k^{(i)}=0$（$k<d_i$）；
  \item[$2^\circ$] $\mathscr P_k^{(i)}$ 在 $S$ 上平坦（$k<n_0+N$）；
  \item[$3^\circ$]
  $\mathfrak{Tor}_q^S(f_1,\ldots,f_m;
  \mathscr P_\bullet^{(1)},\ldots,\mathscr P_\bullet^{(m)})$ 在
  $S$ 上平坦（$q<n_0+N$）。
\end{enumerate}
\end{corollary}

\begin{proof}
假设 $\mathscr P_k^{(i)}$ 在 $S$ 上平坦（$k<r$）；那么
$\mathscr{T}\!or_{q_i}^S(\mathscr P_k^{(i)},\mathscr O_{S'})=0$
（$k<r$ 且 $q_i\ne0$）。我们用
\[
\begin{split}
  \mathfrak{Tor}_p^{S'}\Bigl(
  &f'_1,\ldots,f'_m;
  \mathscr{T}\!or_{q_1}^S(\mathscr P_\bullet^{(1)},\mathscr O_{S'}),
  \ldots,\\
  &\mathscr{T}\!or_{q_m}^S(\mathscr P_\bullet^{(m)},\mathscr O_{S'})
  \Bigr).
\end{split}
\tag{6.9.7.1}\label{III.6.9.7.1.fr}
\]

\oldpage[III]{170}
按照（6.6.2）的方法计算它，其中使用一个固定仿射覆盖
$\mathfrak U^{(i)}$（它覆盖 $X^{(i)}$，$1\leq i\leq m$）的逆像；这个覆盖与 $S'$ 及诸
$\mathscr P_\bullet^{(i)}$ 无关。诸项 $L_{jk}^{(i)}$ 在 $j<-N_i$ 时为零，
其中这个界只依赖于 $\mathfrak U^{(i)}$。如果某个 $q_i$ 不为零，则以
（6.9.7.1）为其 $p$ 次同调的单复形，在所有满足下式的次数上，其项均为零：
\[
  <r+\sum_{j\ne i}d_j-\sum_{i=1}^m N_i,
\]
所以，当 $p<r-N$ 时（6.9.7.1）为零，其中 $N$ 表示下列诸数中的最大者：
\[
  \sum_{i=1}^m N_i-\sum_{j\ne i}d_j.
\]
由此可知，${}^{(f')}\mathscr E_{pq}^2=0$（$q\ne0$ 且
$p<r-N$）；另一方面，${}^{(f')}\mathscr E_{pq}^2=0$（$q<0$），
故边缘同态（6.9.6.4）在 $n<r-N$
时是双射（M，XV，5.6）（这里 $\mathscr Q'_\bullet=\mathscr O_{S'}$）。其次，若当
下列模
\[
  \mathfrak{Tor}_q^S
  (f_1,\ldots,f_m;
   \mathscr P_\bullet^{(1)},\ldots,\mathscr P_\bullet^{(m)})
\]
在 $S$ 上平坦（$q<r$），则 ${}^{(e)}\mathscr E_{pq}^2=0$
（$p\ne0$ 且 $q<r$）；此外，${}^{(e)}\mathscr E_{pq}^2=0$
（$p<0$），所以边缘同态（6.9.6.5）在 $n<r$ 时是双射，
这就完成了证明。
\end{proof}

在应用中，（6.9.3）最重要的情形是 $m=1$，此时
$\mathscr Q'_\bullet$ 只含单项 $\mathscr F'$，且该项位于次数 $0$；为了便于以后引用，
我们在这一情形下重述该结果\footnote{（6.9.8）所处理的情形，尤其是谱序列
（6.9.8.3）和（6.9.8.4），是 J.-P. Serre 于 1957 年向我们指出的。}：

\begin{proposition}[6.9.8]
设 $S$ 是一个预概形，$g:S'\to S$ 是一个态射，
$f:X\to Y$ 是 $S$-预概形之间的一个分离拟紧 $S$-态射，
$\mathscr P_\bullet$ 是由拟相干 $\mathscr O_X$-模组成的下有界复形，
$\mathscr F'$ 是拟相干 $\mathscr O_{S'}$-模。存在两个以
$\mathscr P_\bullet$ 和 $\mathscr F'$ 为变元、取值于拟相干
$\mathscr O_{Y_{(S')}}$-模范畴的双正则谱函子；它们有共同收敛项
$\mathfrak{Tor}_\bullet^S(f,1_{S'};\mathscr P_\bullet,\mathscr F')$，而其
$E_2$ 项为
\[
  {}'\mathscr E_{pq}^2
  =\mathscr{T}\!or_p^S
  (\mathscr H^{-q}(f,\mathscr P_\bullet),\mathscr F')
  \tag{6.9.8.1}\label{III.6.9.8.1.fr}
\]
和
\[
  {}''\mathscr E_{pq}^2
  =\mathscr H^{-p}
  (f',\mathscr{T}\!or_q^S(\mathscr P_\bullet,\mathscr F')),
  \tag{6.9.8.2}\label{III.6.9.8.2.fr}
\]
其中 $f'=f_{(S')}:X_{(S')}\to Y_{(S')}$。
\end{proposition}

\begin{proof}
这些谱序列也可以不从（6.9.3）出发，而从（6.7.3）的谱序列 $(a)$ 和
$(b')$ 得到，只须取 $X^{(1)}=X$、$Y^{(1)}=Y$、
$X^{(2)}=Y^{(2)}=S'$、$f_1=f$、$f_2=1_{S'}$。当
$S=S'=Y$ 且 $Y$ 为仿射时，得到两个谱序列，其 $E_2$ 项为
\[
  {}'\mathscr E_{pq}^2
  =\mathscr{T}\!or_p^Y
  (\mathscr H^{-q}(f,\mathscr P_\bullet),\mathscr F)
  \tag{6.9.8.3}\label{III.6.9.8.3.fr}
\]
\[
  {}''\mathscr E_{pq}^2
  =\mathscr H^{-p}
  (f,\mathscr{T}\!or_q^Y(\mathscr P_\bullet,\mathscr F))
  \tag{6.9.8.4}\label{III.6.9.8.4.fr}
\]
它们（由（6.7.6））以超上同调
$\mathscr H^\bullet(f,\mathscr P_\bullet\otimes_Y\mathscr F)$ 为收敛项；
这是函子 $f_*$ 关于 $\mathscr P_\bullet\otimes_Y\mathscr F$ 这一
$\mathscr O_X$-模复形的超上同调。这里，对每个拟相干 $\mathscr O_Y$-模，
若它在 $Y$ 上平坦并记作 $\mathscr F$，上述结论成立（或者，对每个拟相干
$\mathscr O_Y$-模 $\mathscr F$，当 $\mathscr P_\bullet$ 由
$\mathscr O_X$-模组成且这些模在 $Y$ 上平坦时，上述结论成立）。这两个谱序列
不同于（6.2.1）中的那些谱序列。
\end{proof}

\begin{corollary}[6.9.9]
在（6.9.8）的条件下，假设复形
$\mathscr P_\bullet$ 下有界，由在
$S$ 上平坦的模组成，并且诸 $\mathscr O_Y$-模
$\mathscr H^n(f,\mathscr P_\bullet)$ 在 $S$ 上平坦。

\oldpage[III]{171}
那么有函子性的典范同构
\[
  \mathfrak{Tor}_n^{S'}
  (f',1_{S'};\mathscr P'_\bullet,\mathscr F')
  \xrightarrow{\sim}
  \mathscr H^{-n}(f,\mathscr P_\bullet)
  \otimes_{\mathscr O_S}\mathscr F'
  \tag{6.9.9.1}\label{III.6.9.9.1.fr}
\]
其中
$\mathscr P'_\bullet=\mathscr P_\bullet\otimes_{\mathscr O_S}
\mathscr O_{S'}$；特别地，当 $\mathscr F'=\mathscr O_{S'}$ 时，有
函子性的典范同构
\[
  \mathscr H^n(f',\mathscr P'_\bullet)
  \xrightarrow{\sim}
  \mathscr H^n(f,\mathscr P_\bullet)
  \otimes_{\mathscr O_S}\mathscr O_{S'}.
  \tag{6.9.9.2}\label{III.6.9.9.2.fr}
\]
这是（6.9.6）中 $m=1$ 的特殊情形。更特别地：
\end{corollary}

\begin{corollary}[6.9.10]
设 $S$ 为预概形，$f:X\to Y$ 为 $S$-预概形之间的一个分离拟紧
$S$-态射，$\mathscr P_\bullet$ 为下有界复形，由拟相干
$\mathscr O_X$-模组成，并且这些模在 $S$ 上平坦。另假设诸
$\mathscr O_Y$-模 $\mathscr H^n(f,\mathscr P_\bullet)$ 在 $S$ 上平坦。
对每个 $s\in S$，以 $X_s$ 和 $Y_s$ 表示纤维
$X\otimes_S k(s)$、$Y\otimes_S k(s)$，以
$f_s:X_s\to Y_s$ 表示态射 $f\times_S1$，并以
$\mathscr P_\bullet^s$ 表示复形
$\mathscr P_\bullet\otimes_{\mathscr O_S}k(s)$，它由
$\mathscr O_{X_s}$-模组成。那么有函子性的典范同构
\[
  \mathscr H^n(f_s,\mathscr P_\bullet^s)
  \xrightarrow{\sim}
  \mathscr H^n(f,\mathscr P_\bullet)\otimes_{\mathscr O_S}k(s).
  \tag{6.9.10.1}\label{III.6.9.10.1.fr}
\]
\end{corollary}

所以，在适当的平坦性假设下，导出函子 $R^nf_*(\mathscr F)$ 的形成
“与取纤维交换”；我们将在第~7~节用另一种方法再次得到这一情形。

\subsection{某些上同调函子的局部结构。}

\begin{proposition}[6.10.1]
设 $S=\operatorname{Spec}(A)$ 为仿射概形，
$Y^{(i)}$（$1\leq i\leq n$）为有限族在 $S$ 上平坦的仿射
$S$-概形；对每个 $i$，设 $f_i:X^{(i)}\to Y^{(i)}$ 为分离拟紧的
$S$-态射，并设 $\mathscr P_\bullet^{(i)}$ 为由拟相干
$\mathscr O_{X^{(i)}}$-模组成的下有界复形。设 $Y$ 为诸 $S$-概形
$Y^{(i)}$ 的积。存在一个复形 $\mathscr R_\bullet$，它由拟相干
$\mathscr O_Y$-模组成且在 $S$ 上平坦，并具有如下性质：对每个仿射
$S$-概形 $S'$ 以及每个由拟相干 $\mathscr O_{S'}$-模组成的下有界复形
$\mathscr Q'_\bullet$，有同构
\[
  \mathbf{Tor}_\bullet^S
  (f_1,\ldots,f_n,1_{S'};
   \mathscr P_\bullet^{(1)},\ldots,\mathscr P_\bullet^{(n)},
   \mathscr Q'_\bullet)
  \xrightarrow{\sim}
  \mathscr H_\bullet
  (\mathscr R_\bullet\otimes_{\mathscr O_S}\mathscr Q'_\bullet)
  \tag{6.10.1.1}\label{III.6.10.1.1.fr}
\]
这是一个 $\partial$-函子同构，其变元为 $\mathscr Q'_\bullet$。此外，
对每个 $S$-态射 $u:S''\to S'$，其中两端是仿射 $S$-概形，图
\[
\xymatrix@C=12pt@R=30pt{
  {\begin{gathered}
    \mathbf{Tor}_\bullet^S
    (f_1,\ldots,f_n,1_{S'};\\
    \mathscr P_\bullet^{(1)},\ldots,\mathscr P_\bullet^{(n)},
    \mathscr Q'_\bullet)
  \end{gathered}}
  \ar[r]^-{\sim}\ar[d]
  &
  {\mathscr H_\bullet
    (\mathscr R_\bullet\otimes_{\mathscr O_S}\mathscr Q'_\bullet)}
  \ar[d]
  \\
  {\begin{gathered}
    \mathbf{Tor}_\bullet^S
    (f_1,\ldots,f_n,1_{S''};\\
    \mathscr P_\bullet^{(1)},\ldots,\mathscr P_\bullet^{(n)},
    u^*(\mathscr Q'_\bullet))
  \end{gathered}}
  \ar[r]^-{\sim}
  &
  {\mathscr H_\bullet
    (\mathscr R_\bullet\otimes_{\mathscr O_S}u^*(\mathscr Q'_\bullet))}
}
\tag{6.10.1.2}\label{III.6.10.1.2.fr}
\]
是交换的（其中竖直箭头是（6.7.10）中定义的典范
$(1_Y\times_Su)$-态射）。
\end{proposition}

\begin{proof}
\oldpage[III]{172}
考虑评注（6.6.6），用（6.6.2）的方法计算超 Tor；沿用（6.6.2）的记号，
每个 $L_{\bullet\bullet}^{(i)}$ 都是 $A_i$-模双复形，这里 $A_i$ 表示
$Y^{(i)}$ 的环；所以它有一个 Cartan--Eilenberg 投射消解
$M_{\bullet\bullet\bullet}^{(i)}$（其意义如（0，11.7.1）），该消解由
$A_i$-模组成，
而由（6.6.6），（6.10.1.1）的左端典范同构于
\[
  \mathscr H_\bullet
  (M_{\bullet\bullet\bullet}\otimes_A Q'_\bullet),
\]
其中
\[
  M_{\bullet\bullet\bullet}
  =
  M_{\bullet\bullet\bullet}^{(1)}
  \otimes_A M_{\bullet\bullet\bullet}^{(2)}
  \otimes_A\cdots\otimes_A M_{\bullet\bullet\bullet}^{(n)}
\]
且 $\mathscr Q'_\bullet=\widetilde{Q'_\bullet}$。根据假设，诸环
$A_i$ 是平坦 $A$-模，所以 $M_{\bullet\bullet\bullet}^{(i)}$ 是由平坦
$A$-模组成的三重复形（$0_{\mathrm I}$，6.2.1），
$M_{\bullet\bullet\bullet}$ 也是如此；此外，若 $B$ 是 $Y$ 的环，即诸
$A_i$ 的张量积，则 $M_{\bullet\bullet\bullet}$ 是 $B$-模三重复形。因此，复形
\[
  \mathscr R_\bullet=(M_{\bullet\bullet\bullet})^\sim
\]
是一个 $\mathscr O_Y$-模复形（这里把 $M_{\bullet\bullet\bullet}$ 看作
单复形），并满足要求；这由（6.7.10）立即得出。
\end{proof}

\begin{corollary}[6.10.2]
在（6.10.1）的陈述中，可以假设 $\mathscr R_\bullet$ 下有界。当诸
$\mathscr P_\bullet^{(i)}$ 上有界且诸 $Y^{(i)}$ 的上同调维数有限时，
可以假设 $\mathscr R_\bullet$ 上有界。
\end{corollary}

\begin{proof}
第一项断言来自如下事实：每个 $M_{\bullet\bullet\bullet}^{(i)}$ 的三个次数都
下有界；另一方面，若诸环 $A_i$ 的上同调维数有限，则每个
$M_{\bullet\bullet\bullet}^{(i)}$ 的第三次数只能取有限多个值，而由其构造，
它的第一次数也是如此（6.6.2）；如果 $\mathscr P_\bullet^{(i)}$ 的次数上有界，
则它的第二次数也上有界（6.6.2），第二项断言立即由此得出。
\end{proof}

\begin{remarks}[6.10.3]
\textnormal{(i)} 沿用（6.10.1）的记号，
$\mathscr H_\bullet(\mathscr R_\bullet\otimes_S\mathscr Q'_\bullet)$ 同构于
$\mathbf{Tor}_\bullet(\mathscr R_\bullet,\mathscr Q'_\bullet)$，因为
$\mathscr R_\bullet$ 由 $\mathscr O_Y$-模组成且这些模在 $S$ 上平坦
（6.5.9）；所以它是（6.5.4）一个正则谱序列的收敛项，该谱序列的
$E_2$ 项为
\[
  \mathscr E_{pq}^2
  =
  \bigoplus_{q'+q''=q}
  \mathbf{Tor}_p^S
  \bigl(
    \mathscr H_{q'}(\mathscr R_\bullet),
    \mathscr H_{q''}(\mathscr Q'_\bullet)
  \bigr)
  \tag{6.10.3.1}\label{III.6.10.3.1.fr}
\]
这正是（6.9.3）中基变换的谱序列 $(e)$。

\medskip
\textnormal{(ii)} 设 $\mathscr R'_\bullet$ 是另一个由拟相干
$\mathscr O_Y$-模组成且在 $S$ 上平坦的复形，并设
$g:\mathscr R'_\bullet\to\mathscr R_\bullet$ 是复形同态，使得
\[
  \mathscr H_\bullet(g):
  \mathscr H_\bullet(\mathscr R'_\bullet)
  \longrightarrow
  \mathscr H_\bullet(\mathscr R_\bullet)
\]
为同构。于是由（6.3.3）和（6.5.9），从 $g$ 得到一个
$\partial$-函子同构，其变元为 $\mathscr Q'_\bullet$：
\[
  \mathscr H_\bullet
  (\mathscr R'_\bullet\otimes_S\mathscr Q'_\bullet)
  \xrightarrow{\sim}
  \mathscr H_\bullet
  (\mathscr R_\bullet\otimes_S\mathscr Q'_\bullet)
\]
并且图
\[
\xymatrix@C=28pt@R=30pt{
  \mathscr H_\bullet
  (\mathscr R'_\bullet\otimes_S\mathscr Q'_\bullet)
  \ar[r]^-{\sim}\ar[d]
  &
  \mathscr H_\bullet
  (\mathscr R_\bullet\otimes_S\mathscr Q'_\bullet)
  \ar[d]
  \\
  \mathscr H_\bullet
  (\mathscr R'_\bullet\otimes_Su^*(\mathscr Q'_\bullet))
  \ar[r]^-{\sim}
  &
  \mathscr H_\bullet
  (\mathscr R_\bullet\otimes_Su^*(\mathscr Q'_\bullet))
}
\]
\par\medskip\noindent
是交换的。这就证明复形 $\mathscr R_\bullet$ 并不完全由（6.10.1）的性质所确定。

\medskip
\oldpage[III]{173}
\textnormal{(iii)} 在（6.10.1）的证明中，可以假设
$M_{\bullet\bullet\bullet}^{(i)}$ 由自由 $A_i$-模组成（这由
（0，11.5.2.1）的“对偶化”证明容易得出）；于是
\[
  M_{\bullet\bullet\bullet}^{(i)}\otimes_{A_i}B
\]
由自由 $B$-模组成；又因为 $M_{\bullet\bullet\bullet}$ 等于它们在
$B$ 上的张量积，故在（6.10.1）中还可以假设 $\mathscr R_\bullet$
与一个自由 $B$-模复形相伴。此外，由（M，XVII，1.2），三重复形
$M_{\bullet\bullet\bullet}$ 函子性地依赖于每个双复形
$L_{\bullet\bullet}^{(i)}$（因而也依赖于每个
$\mathscr P_\bullet^{(i)}$，这里固定每个 $X^{(i)}$ 的一个有限覆盖）；
此处三重复形的“态射”应理解为按同伦关系所得的同态类。另一方面，用
$X^{(i)}$ 的一个更细覆盖替换原覆盖时，会给出
$L_{\bullet\bullet}^{(i)}$ 之间恰好只在同伦意义下定义的同态（6.6.8）；
所以最终可见，在上述态射约定下，三重复形
$M_{\bullet\bullet\bullet}$ 关于每个 $\mathscr P_\bullet^{(i)}$ 都是函子。
我们将在论述上同调函子一般代数的那一章中（该章已在（6.1.3）中提到），
详细说明这种函子性依赖，尤其说明 $\mathscr R_\bullet$ 对复形正合列
\[
  0\longrightarrow\mathscr P_\bullet^{\prime(i)}
  \longrightarrow\mathscr P_\bullet^{(i)}
  \longrightarrow\mathscr P_\bullet^{\prime\prime(i)}
  \longrightarrow0
\]
的行为。
\end{remarks}

\begin{scholium}[6.10.4]
$\mathscr R_\bullet$ 由 $\mathscr O_Y$-模组成，且这些模在 $S$ 上平坦；由此
容易推出
\[
  \mathscr H_\bullet
  (\mathscr R_\bullet\otimes_S\mathscr Q'_\bullet)
\]
是关于 $\mathscr Q'_\bullet$ 的同调函子（见（7.7.1）的论证）。正如
（6.1.1）中提到的，正是这一性质促使我们引入超 Tor。事实上，令
\[
  X=X^{(1)}\times_SX^{(2)}\times_S\cdots\times_SX^{(n)},
  \qquad
  f=f_1\times_Sf_2\times_S\cdots\times_Sf_n,
\]
\[
  \mathscr P_\bullet
  =
  \mathscr P_\bullet^{(1)}
  \otimes_S\mathscr P_\bullet^{(2)}
  \otimes_S\cdots\otimes_S\mathscr P_\bullet^{(n)},
\]
\[
  X'=X\times_SS',
  \qquad
  Y'=Y\times_SS',
  \qquad
  f'=f\times_S1_{S'}.
\]
基变换问题使我们研究超上同调
\[
  \mathscr H^\bullet_{f'}
  (\mathscr P_\bullet\otimes_S\mathscr N')
\]
将其视为关于拟相干 $\mathscr O_{S'}$-模 $\mathscr N'$ 的函子，或者研究
超上同调
\[
  \mathscr H^\bullet_f
  (\mathscr P_\bullet\otimes_S\mathscr N)
\]
将其视为关于拟相干 $\mathscr O_S$-模 $\mathscr N$ 的函子。当诸
$\mathscr P_\bullet^{(i)}$（因而 $\mathscr P_\bullet$ 也）在 $S$ 上平坦时，
由上述结果及（6.7.6），这个函子确实是关于 $\mathscr N$ 的上同调函子；
但是，当不再对诸 $\mathscr P_\bullet^{(i)}$ 作平坦性假设时，结论便不再成立，
而且这时不能再用同调代数的通常方法来研究
\[
  \mathscr H^\bullet_f
  (\mathscr P_\bullet\otimes_S\mathscr N)
\]
。

不过，我们主要要用的是如下情形：$n=1$、$Y=S$，并且
$\mathscr P_\bullet$ 由 $\mathscr O_X$-模组成，而这些模在 $Y$ 上平坦。
在这种情形下，有下列定理。
\end{scholium}

\begin{theorem}[6.10.5]
设 $Y=\operatorname{Spec}(A)$ 是诺特仿射概形，$f:X\to Y$ 是固有态射，
而 $\mathscr P_\bullet$ 是由凝聚 $\mathscr O_X$-模组成、在 $Y$ 上平坦的
下有界复形。那么存在一个下有界复形 $\mathscr L_\bullet$，它由
$\mathscr O_Y$-模组成；其各项 $\mathscr L_i$ 都是
$\mathscr O_Y$-模，且形如 $\mathscr O_Y^{m_i}$；此外有一个同构
\[
  \mathscr H^\bullet
  (f,\mathscr P_\bullet\otimes_Y\mathscr Q_\bullet)
  \xrightarrow{\sim}
  \mathscr H_\bullet
  (\mathscr L_\bullet\otimes_Y\mathscr Q_\bullet)
  \tag{6.10.5.1}\label{III.6.10.5.1.fr}
\]
；这是 $\partial$-函子关于下有界复形 $\mathscr Q_\bullet$ 的同构，其中
该复形由拟相干 $\mathscr O_Y$-模组成。此外，对于任意态射
$u:Y'\to Y$，令
\[
  X'=X_{(Y')},
  \qquad
  f'=f_{(Y')},
  \qquad
  \mathscr P'_\bullet
  =
  \mathscr P_\bullet\otimes_Y\mathscr O_{Y'},
  \qquad
  \mathscr L'_\bullet=u^*(\mathscr L_\bullet),
\]
\oldpage[III]{174}
（后者是一个由有限型局部自由 $\mathscr O_{Y'}$-模组成的复形），则有同构
\[
  \mathscr H^\bullet
  (f',\mathscr P'_\bullet\otimes_{Y'}\mathscr Q'_\bullet)
  \xrightarrow{\sim}
  \mathscr H_\bullet
  (\mathscr L'_\bullet\otimes_{Y'}\mathscr Q'_\bullet)
  \tag{6.10.5.2}\label{III.6.10.5.2.fr}
\]
；这是 $\partial$-函子关于下有界复形 $\mathscr Q'_\bullet$ 的同构，其中
该复形由拟相干 $\mathscr O_{Y'}$-模组成，并且可以使图
\[
\xymatrix@C=25pt@R=30pt{
  \mathscr H^\bullet
  (f,\mathscr P_\bullet\otimes_Y\mathscr Q_\bullet)
  \ar[r]^-{\sim}\ar[d]
  &
  \mathscr H_\bullet
  (\mathscr L_\bullet\otimes_Y\mathscr Q_\bullet)
  \ar[d]
  \\
  \mathscr H^\bullet
  (f',\mathscr P'_\bullet\otimes_{Y'}u^*(\mathscr Q_\bullet))
  \ar[r]^-{\sim}
  &
  \mathscr H_\bullet
  (\mathscr L'_\bullet\otimes_{Y'}u^*(\mathscr Q_\bullet))
}
\tag{6.10.5.3}\label{III.6.10.5.3.fr}
\]
交换。
\end{theorem}

\begin{proof}
应用（6.10.1），首先得到一个下有界复形（6.10.2）
$\mathscr R_\bullet$；它由拟相干 $\mathscr O_Y$-模组成，并且这些模在
$Y$ 上自由（6.10.3，(iii)），同时得到一个同构
\[
  \mathscr H^\bullet
  (f,\mathscr P_\bullet\otimes_Y\mathscr Q_\bullet)
  \xrightarrow{\sim}
  \mathscr H_\bullet
  (\mathscr R_\bullet\otimes_Y\mathscr Q_\bullet)
  \tag{6.10.5.4}\label{III.6.10.5.4.fr}
\]
；这是 $\partial$-函子关于 $\mathscr Q_\bullet$ 的同构。不过，
\emph{先验地}，$\mathscr R_\bullet$ 的各项不一定是有限型
$\mathscr O_Y$-模。现在把（6.10.5.4）应用于如下情形：
$\mathscr Q_\bullet$ 是仅有一项 $\mathscr O_Y$ 的复形；于是可见
$\mathscr H_\bullet(\mathscr R_\bullet)$ 与
$\mathscr H^\bullet(f,\mathscr P_\bullet)$ 同构，因而由凝聚
$\mathscr O_Y$-模组成（6.2.5）。这时由（0，11.9.2）可知，存在一个
下有界复形 $\mathscr L_\bullet$，它由 $\mathscr O_Y$-模组成，而这些模
与有限型自由 $A$-模相伴；并有一个同态
\[
  \mathscr L_\bullet\longrightarrow\mathscr R_\bullet,
\]
使相应的同调同态
\[
  \mathscr H_\bullet(\mathscr L_\bullet)
  \longrightarrow
  \mathscr H_\bullet(\mathscr R_\bullet)
\]
为双射；由（6.10.3，(ii)）即得同构（6.10.5.1）。

当 $Y'$ 仿射时，（6.10.5）的其他断言由（6.10.1）和（6.10.3，(ii)）
得到；在一般情形，只需验证如下事实：用仿射开集 $(V_\alpha)$ 覆盖
$Y'$，并在每个 $V_\alpha$ 上取相应的同构（6.10.5.2）时，把这些同构
限制到仿射开集 $W\subset V_\alpha\cap V_\beta$ 后，关于
$V_\alpha$ 和 $V_\beta$ 的两个限制都与关于 $W$ 的同构一致。
这是把交换图（6.10.1.2）应用于典范嵌入
$W\to V_\alpha$ 和 $W\to V_\beta$ 的结果。
\end{proof}

\begin{remark}[6.10.6]
在后续各章中，我们主要把（6.10.5）应用于如下情形：
$\mathscr P_\bullet$ 约化为单个凝聚 $\mathscr O_X$-模
$\mathscr F$，并且它在 $Y$ 上平坦。由于这时
\[
  \mathscr H_n(\mathscr L_\bullet)
  =
  \mathscr H^{-n}(f,\mathscr F)
  =
  R^{-n}f_*(\mathscr F)
  \qquad(6.2.1),
\]
可见诸 $\mathscr H_n(\mathscr L_\bullet)$ 对 $n>0$ 为零；以后我们将证明
（7.7.12，(i)），此时可以假设 $\mathscr L_\bullet$ 只有次数
$\leqslant0$ 的项（因而只有有限项），但须把“诸 $\mathscr L_i$ 都与
有限型自由 $A$-模相伴”这一假设换成“诸 $\mathscr L_i$ 都是有限型
局部自由模”。

复形 $\mathscr L_\bullet$ 对应于这样的 $\mathscr O_X$-模 $\mathscr F$；
除上述次数限制外，它看来并没有什么特殊性质。
\oldpage[III]{175}
因此可以反过来问：给定一个下有界复形 $\mathscr L_\bullet$，它由
$\mathscr O_Y$-模组成，这些模与有限型投射 $A$-模相伴，并且其次数
$>0$ 的项均为零，是否存在一个 $Y$-概形 $X$，它在 $Y$ 上投射且平坦，
以及一个局部自由 $\mathscr O_X$-模 $\mathscr F$，使得存在同构
\[
  \mathscr H^\bullet
  (f,\mathscr F\otimes_Y\mathscr Q_\bullet)
  \xrightarrow{\sim}
  \mathscr H_\bullet
  (\mathscr L_\bullet\otimes_Y\mathscr Q_\bullet)
\]
并且关于 $\mathscr Q_\bullet$ 具有函子性。这样的结果将把在 $Y$ 上平坦的
凝聚模在固有 $Y$-概形上的上同调理论，完全归结为有限型投射 $A$-模复形
在诺特环 $A$ 上的“在同伦意义下”理论。
\end{remark}

% EGA III-2 mainland Simplified-Chinese producer source.
% Basis: Canon keeper R71 corrected-current/diplomatic pair; source-literal unless an admitted reversible disposition applies.
% This file is a producer draft and is uncertified pending independent replay.

\section{模的协变同调函子中的基变换研究}
\label{section:III.7.fr}

\subsection[A-模函子]{$A$-模函子}
\label{subsection:III.7.1.fr}

\begin{env}[7.1.1]
给定一个环 $A$（不必交换），记 $\mathbf{Ab}_A$ 为左 $A$-模范畴，
并简记 $\mathbf{Ab}$ 为 $\mathbf Z$-模范畴，也就是交换群范畴。设
\[
  T:\mathbf{Ab}_A\longrightarrow\mathbf{Ab}
\]
是协变加性函子，并设 $M$ 是一个 $(A,A)$-双模；于是 $T(M)$ 自然地
带有右 $A$-模结构。事实上，对每个 $a\in A$，记 $h_{a,M}$（或简记
$h_a$）为自同态 $x\mapsto xa$；这里它作用在左 $A$-模 $M$ 上。依假设，$T(h_a)$
是 $\mathbf Z$-模 $T(M)$ 的一个自同态；又由于 $T$ 是协变加性函子，
对 $a\in A$、$b\in A$ 有
\[
  T(h_{ab})
  =
  T(h_b\circ h_a)
  =
  T(h_b)\circ T(h_a),
  \qquad\text{且}\qquad
  T(h_{a+b})
  =
  T(h_a+h_b)
  =
  T(h_a)+T(h_b);
\]
这证明映射
\[
  (y,a)\longmapsto T(h_a)(y)
\]
定义了一个右 $A$-模结构，其底层群为 $T(M)$。\SourceCorrectionNote{法文印本及外交式转录在此把自变量次序写成
“标量、元素”，但右模标量乘法的定义域应先取元素、再取标量。本译依 Canon
校订现行法文读法“元素、标量”，并在本注保留原印次序。}特别地，
$T(A_s)$ 是一个右 $A$-模。
\end{env}

\begin{env}[7.1.2]
当 $A$ 是交换环时，由（7.1.1）可知，对每个 $A$-模 $M$，$T(M)$
自然地带有 $A$-模结构；此外，若 $u:M\to N$ 是 $A$-模同态，则对每个
$a\in A$ 有
\[
  u\circ h_{a,M}=h_{a,N}\circ u,
\]
从而
\[
  T(u)\circ T(h_{a,M})
  =
  T(h_{a,N})\circ T(u),
\]
这证明 $T(u):T(M)\to T(N)$ 是 $A$-模同态；所以可以把 $T$ 看作从
$\mathbf{Ab}_A$ 到其自身的协变加性函子。更确切地说，我们由此定义了
协变加性函子 $\mathbf{Ab}_A\to\mathbf{Ab}$ 所成范畴与协变 $A$-线性函子
\[
  T:\mathbf{Ab}_A\longrightarrow\mathbf{Ab}_A,
\]
所成范畴之间的一个典范等价；这里的线性条件是
\[
  T(h_{a,M})=h_{a,T(M)}
  \qquad\text{对每个 }a\in A.
\]
由于包含函子
\[
  I:\mathbf{Ab}_A\longrightarrow\mathbf{Ab},
\]
把每个 $A$-模对应到其底层 $\mathbf Z$-模，并且是正合且忠实的，
所以上述等价对应的两个函子具有相同的正合性。
\end{env}

\begin{env}[7.1.3]
仍设环 $A$ 交换，设 $B$ 是一个 $A$-代数（不必交换），并设
$\rho:A\to B$ 是与这一代数结构相应的环同态；该同态定义了一个协变
\oldpage[III]{176}
加性函子
\[
  \rho_*:M\longmapsto M_{[\rho]}
\]
，它从范畴 $\mathbf{Ab}_B$（左 $B$-模范畴）到范畴
$\mathbf{Ab}_A$（左 $A$-模范畴）。复合以后得到函子
\[
  T_{(B)}:
  \mathbf{Ab}_B
  \xrightarrow{\rho_*}
  \mathbf{Ab}_A
  \xrightarrow{T}
  \mathbf{Ab},
\]
它显然是协变且加性的；由于排印上的原因，我们也把它记作
$T^{(B)}$ 或 $T\otimes_A B$，并称它是由 $T$ 经从 $A$ 到 $B$ 的
标量扩张所得的函子。当然，如果 $B$ 交换，就可以把 $T_{(B)}$ 看作从
$\mathbf{Ab}_B$ 到其自身的函子（7.1.2）。当 $B$ 交换且 $C$ 是一个
$B$-代数时，立即可见
\[
  T_{(C)}=(T_{(B)})_{(C)};
\]
标量扩张关于 $T$ 显然具有函子性和加性；此外，当 $T$ 与归纳极限或
直和交换（相应地，当它左正合、右正合或正合）时，$T_{(B)}$ 也具有
相同性质：事实上，$\rho_*$ 是正合的，并且与归纳极限及直和交换。
\end{env}

\begin{env}[7.1.4]
仍设 $A$ 交换，并设
$T:\mathbf{Ab}_A\to\mathbf{Ab}_A$ 是一个与归纳极限交换的协变加性
$A$-线性函子。那么，对任意乘法子集 $S$（取自 $A$）和任意 $A$-模 $M$，
都有一个具有函子性的典范 $A$-模同构
\[
  T(S^{-1}M)\xrightarrow{\sim}S^{-1}T(M).
  \tag{7.1.4.1}\label{III.7.1.4.1.fr}
\]
先设 $S$ 是各次幂 $f^n$ 所成的集合，其中 $(n\geqslant0)$，且 $f\in A$。
我们知道
\[
  M_f=\varinjlim M_n,
\]
其中 $(M_n,\varphi_{nm})$ 是 $A$-模的归纳系，各 $M_n=M$，且
\[
  \varphi_{nm}:z\longmapsto f^{\,n-m}z
  \qquad(0_{\mathrm I}, 1.6.1);
\]
于是由关于 $T$ 的假设，在这种情形得到同构（7.1.4.1）。若 $S$ 任意，
则 $S^{-1}M$ 是诸 $M_f$（$f\in S$）的归纳极限（$0_{\mathrm I}$，
1.4.5），同样即可得出结论。此外，同构（7.1.4.1）的函子性表明它是
$S^{-1}A$-模同构，因而在典范同构意义下可以写成
\[
  T_{(S^{-1}A)}(S^{-1}M)
  =
  S^{-1}T(M)
  =
  T(S^{-1}M).
  \tag{7.1.4.2}\label{III.7.1.4.2.fr}
\]
当 $S=A-\mathfrak p$ 是素理想 $\mathfrak p$（属于环 $A$）的补集时，以
$T_{\mathfrak p}$ 代替 $T_{(A_{\mathfrak p})}$。
\end{env}

\begin{proposition}[7.1.5]
在（7.1.4）的假设下，若 $T_{\mathfrak m}$ 左正合（相应地，右正合或正合），
这里 $\mathfrak m$ 遍历 $A$ 的全部极大理想，则 $T$ 左正合
（相应地，右正合或正合）。
\end{proposition}

\begin{proof}
事实上，已知两个子模 $N$、$P$ 若同为一个 $A$-模 $M$ 的子模，并且
对每个极大理想都满足 $N_{\mathfrak m}=P_{\mathfrak m}$，其中
$\mathfrak m$ 遍历 $A$ 的全部极大理想，则
$N=P$（Bourbaki，\emph{Alg. comm.}，第 II 章，\S~3，第 3 号，
定理 1）。
\end{proof}

\subsection{张量积函子的刻画}
\label{subsection:III.7.2.fr}

\begin{env}[7.2.1]
设 $A$ 是一个环（不必交换），$M$（相应地，$N$）是一个左（相应地，
右）$A$-模，$P$ 是一个 $\mathbf Z$-模。回顾一下，给定一个
$\mathbf Z$-同态
\[
  v:N\otimes_A M\longrightarrow P
\]
等价于给定一个 $\mathbf Z$-双线性映射
\[
  u:N\times M\longrightarrow P
\]
使得
\[
  u(ta,x)=u(t,ax)
  \qquad\text{对 }a\in A,\ t\in N,\ x\in M,
\]
而这两个映射由 $v(t\otimes x)=u(t,x)$ 联系起来。另一方面，给定 $u$
等价于给定一个
\oldpage[III]{177}
$\mathbf Z$-同态 $x\mapsto f_x$，它从 $M$ 映到
$\operatorname{Hom}_{\mathbf Z}(N,P)$，且满足
\[
  f_{ax}(t)=f_x(ta)
  \qquad\text{对 }a\in A,\ t\in N,\ x\in M,
\]
而这两个映射由 $u(t,x)=f_x(t)$ 联系起来。
\end{env}

\begin{env}[7.2.2]
设 $T:\mathbf{Ab}_A\to\mathbf{Ab}$ 是协变加性函子。我们将对每个左
$A$-模 $M$ 定义一个关于 $M$ 具有函子性的典范 $\mathbf Z$-模同态
\[
  t_M:T(A_s)\otimes_A M\longrightarrow T(M).
  \tag{7.2.2.1}\label{III.7.2.2.1.fr}
\]
依（7.2.1），为此只须定义一个 $\mathbf Z$-同态 $x\mapsto t'_M(x)$，
它从 $M$ 映到 $\operatorname{Hom}_{\mathbf Z}(T(A_s),T(M))$，使得
\[
  t'_M(ax)(y)=t'_M(x)(ya)
  \qquad\text{对 }a\in A,\ x\in M,\ y\in T(A_s).
\]
为此注意到，$\operatorname{Hom}_{\mathbf Z}(T(A_s),T(M))$ 典范地带有
一个左 $A$-模结构；这个结构来自 $A$-模 $T(A_s)$ 的右模结构，其外乘法则满足：
若 $a\in A$，$v\in\operatorname{Hom}_{\mathbf Z}(T(A_s),T(M))$，则
\[
  (a\cdot v)(y)=v(ya)
  \qquad\text{对 }y\in T(A_s).
\]
现在，把 $t'_M$ 定义为以下两个典范同态的复合：
\[
  M\xrightarrow{\sim}\operatorname{Hom}_A(A_s,M)
  \xrightarrow{T}
  \operatorname{Hom}_{\mathbf Z}(T(A_s),T(M)),
\]
第二个箭头是映射 $u\mapsto T(u)$；第一个箭头是 $A$-模的典范同构
$x\mapsto\theta_x$，它满足 $\theta_x(\xi)=\xi x$，其中
$\xi\in A$、$x\in M$。又有 $\theta_{ax}=\theta_x\circ h_a$，所以
\[
  T(\theta_{ax})
  =
  T(\theta_x\circ h_a)
  =
  T(\theta_x)\circ T(h_a)
\]
因而对 $y\in T(A_s)$，
\[
  T(\theta_{ax})(y)
  =
  T(\theta_x)(T(h_a)(y))
  =
  T(\theta_x)(ya)
\]
根据 $T(A_s)$ 上外乘法则的定义，这证明了 $t_M$ 的存在性。容易验证该
同态关于 $M$ 具有函子性；也就是说，对任意同态 $w:M\to M'$（它是左
$A$-模同态），图表
\[
\xymatrix@C=30pt@R=30pt{
  T(A_s)\otimes_A M
  \ar[r]^-{t_M}\ar[d]_{1\otimes w}
  &
  T(M)\ar[d]^{T(w)}
  \\
  T(A_s)\otimes_A M'
  \ar[r]_-{t_{M'}}
  &
  T(M')
}
\tag{7.2.2.2}\label{III.7.2.2.2.fr}
\]
交换。

同态（7.2.2.1）的函子性表明，当 $A$ 交换时，它是一个 $A$-模同态
（参见（7.1.2））。
\end{env}

\begin{env}[7.2.3]
当 $A$ 交换时，可以更一般地定义典范 $A$-模同态
\[
  T(N)\otimes_A M\longrightarrow T(N\otimes_A M)
  \tag{7.2.3.1}\label{III.7.2.3.1.fr}
\]
（这里取任意 $A$-模 $N$）。只须在（7.2.2）的构造中，把同态
$\theta_x$ 换成如下 $A$-模同态：
\[
  N\longrightarrow N\otimes_A M
\]
它把每个 $y\in N$ 映到 $y\otimes x$。显然，该同态关于 $M$ 和 $N$
都具有函子性。

\oldpage[III]{178}
特别地，若 $B$ 是一个 $A$-代数（不必交换），则有关于 $M$ 具有函子性的
同态
\[
  (T(M))_{(B)}
  =
  T(M)\otimes_A B
  \longrightarrow
  T(M\otimes_A B)
  =
  T_{(B)}(M_{(B)})
  \tag{7.2.3.2}\label{III.7.2.3.2.fr}
\]
由（7.2.3.1）关于 $M$ 的函子性可知，这是一个 $B$-模同态。

此外，还有交换图表
\[
\xymatrix@C=30pt@R=30pt{
  T(A)\otimes_A M
  \ar[r]^-{t_M}\ar[d]
  &
  T(M)\ar[d]
  \\
  T_{(B)}(B_s)\otimes_B M_{(B)}
  \ar[r]_-{t_{M_{(B)}}}
  &
  T_{(B)}(M_{(B)})
}
\tag{7.2.3.3}\label{III.7.2.3.3.fr}
\]
其中右边的竖直箭头是（7.2.3.2）与典范同态的复合同态
\[
  T(M)
  \longrightarrow
  T(M)\otimes_A B
  \longrightarrow
  T(M\otimes_A B)
  =
  T_{(B)}(M_{(B)})
\]
至于（7.2.3.3）左边的竖直箭头，则是同态
\[
  T(A)\otimes_A M
  \longrightarrow
  T_{(B)}(B_s)\otimes_B(B\otimes_A M)
  =
  T_{(B)}(B_s)\otimes_A M
\]
这里
\[
  T(A)\longrightarrow T_{(B)}(B_s)=T(B)
\]
是 $T(\rho)$，而 $\rho$ 被看作 $A$-模同态 $A\to B_{[\rho]}$。
\end{env}

\begin{lemma}[7.2.4]
若 $T$ 是从 $\mathbf{Ab}_A$ 到 $\mathbf{Ab}$、与直和交换的协变加性函子，
则典范同态 $t_L$（7.2.2.1）对每个自由 $A$-模 $L$ 都是同构。
\end{lemma}

\begin{proof}
事实上，
\[
  L=\bigoplus_{\alpha\in I}L_\alpha
\]
其中 $L_\alpha$ 同构于 $A_s$，这里 $\alpha\in I$；（7.2.2）中
$t_M$ 的定义表明
\[
  t_L=\bigoplus_{\alpha\in I}t_{L_\alpha},
\]
这是因为
\[
  T:
  \operatorname{Hom}_A(A_s,L)
  \longrightarrow
  \operatorname{Hom}_{\mathbf Z}(T(A_s),T(L))
\]
是诸 $\mathbf Z$-线性映射
\[
  T_\alpha:
  \operatorname{Hom}_A(A_s,L_\alpha)
  \longrightarrow
  \operatorname{Hom}_{\mathbf Z}(T(A_s),T(L_\alpha))
\]
的直和（依据关于 $T$ 的假设）。因此只须对 $L=A_s$ 证明引理；但这时
$t_L$ 正是如下典范同构：
\[
  T(A_s)\otimes_A A_s\xrightarrow{\sim}T(A_s)
\]
该同构对每个右 $A$-模都成立。
\end{proof}

\begin{proposition}[7.2.5]
设 $T$ 是从 $\mathbf{Ab}_A$ 到 $\mathbf{Ab}$、与直和交换的协变加性函子。
以下条件等价：
\begin{enumerate}
  \item[a)] $T$ 右正合。
  \item[b)] 典范同态 $t_M$（7.2.2.1）对每个左 $A$-模 $M$ 都是同构。
  \item[$b'$)] $T$ 半正合，并且同态 $t_M$ 对每个左 $A$-模 $M$ 都满射。
  \item[c)] $T$ 同构于一个关于 $M$ 的函子，这个函子形如
    $N\otimes_A M$，其中 $N$ 是一个右 $A$-模。
\end{enumerate}
\end{proposition}

\begin{proof}
显然，$b)$ 蕴涵 $c)$，而 $c)$ 蕴涵 $a)$；现在证明 $a)$ 蕴涵 $b)$。
\oldpage[III]{179}
置
\[
  T'(M)=T(A_s)\otimes_A M
\]
（其中取任意左 $A$-模 $M$）。
存在正合列
\[
  L'\longrightarrow L\longrightarrow M\longrightarrow0,
\]
其中 $L$ 和 $L'$ 是两个自由左 $A$-模；由于 $T$ 和 $T'$ 都右正合，
遂有交换图表
\[
\xymatrix@C=22pt@R=25pt{
  T'(L')\ar[r]\ar[d]^{t_{L'}}
  &
  T'(L)\ar[r]\ar[d]^{t_L}
  &
  T'(M)\ar[r]\ar[d]^{t_M}
  &
  0
  \\
  T(L')\ar[r]
  &
  T(L)\ar[r]
  &
  T(M)\ar[r]
  &
  0
}
\]
其中两行都正合；由（7.2.4），$t_L$ 和 $t_{L'}$ 都是同构，故由五引理，
$t_M$ 也是同构。最后，$b)$ 显然蕴涵 $b')$。为证明 $b')$ 蕴涵 $a)$，
只须证明以下引理。
\end{proof}

\begin{lemma}[7.2.5.1]
设 $\mathcal K$、$\mathcal K'$ 是两个阿贝尔范畴，$F$、$G$ 是从
$\mathcal K$ 到 $\mathcal K'$ 的两个协变加性函子，$f:F\to G$ 是一个
函子态射（T, I, 1.2），并且对每个对象 $E$（属于范畴 $\mathcal K$），
\[
  f_E:F(E)\longrightarrow G(E)
\]
都是满态射。若 $F$ 右正合且 $G$ 半正合，则 $G$ 右正合。
\end{lemma}

\begin{proof}
事实上，只须证明：对每个满态射 $v:E'\to E$（在 $\mathcal K$ 中），
\[
  G(v):G(E')\longrightarrow G(E)
\]
都是满态射。现有交换图表
\[
\xymatrix@C=35pt@R=30pt{
  F(E')\ar[r]^{F(v)}\ar[d]_{f_{E'}}
  &
  F(E)\ar[d]^{f_E}
  \\
  G(E')\ar[r]_{G(v)}
  &
  G(E)
}
\]
其中 $F(v)$、$f_{E'}$ 和 $f_E$ 都是满态射；因此 $G(v)$ 也是满态射。
\end{proof}

\begin{remark}[7.2.6]
对每个右 $A$-模 $N$，置
\[
  T_N(M)=N\otimes_A M
\]
（其中取任意左 $A$-模 $M$），于是 $T_N$ 是从 $\mathbf{Ab}_A$ 到
$\mathbf{Ab}$、右正合且与直和交换的协变加性函子。把 $T_N(A_s)$
典范地等同于 $N$，立即可以验证，相应的同态（7.2.2.1）变成恒等映射。
由此可知，其中的右 $A$-模 $N$（7.2.5，$c)$）在唯一同构意义下确定，
并且典范同构于 $T(A_s)$。还可以说，函子态射
\[
  T\longmapsto T(A_s),
  \qquad
  N\longmapsto T_N
\]
构成右 $A$-模范畴与如下函子范畴之间的一个等价（T, I, 1.2）：后者的对象
是从 $\mathbf{Ab}_A\to\mathbf{Ab}$、右正合且与直和交换的协变加性函子。
\end{remark}

\begin{proposition}[7.2.7]
设 $A$ 是一个左 Artin 环，其关于根基 $\mathfrak m$ 的商环是一个域
$k$。设 $T$ 是从 $\mathbf{Ab}_A$ 到 $\mathbf{Ab}$、与直和交换的协变
加性函子。此时，（7.2.5）的条件还等价于
\begin{enumerate}
  \item[d)] $T$ 半正合，并且同态
    \[
      T(\varepsilon):T(A_s)\longrightarrow T(k)
    \]
    满射；该同态由典范同态 $\varepsilon:A_s\to k$ 导出。
\end{enumerate}
\end{proposition}

\begin{proof}
\oldpage[III]{180}
显然，（7.2.5）的条件 $b')$ 蕴涵 $d)$；现在证明 $d)$ 蕴涵 $b')$。

存在一个整数 $n$ 使得 $\mathfrak m^n=0$；对每个 $A$-模 $M$，置
$M_h=\mathfrak m^hM$。我们将对 $h$ 作降序归纳，证明 $t_{M_h}$ 满射。
当 $h=n$ 时命题显然成立；当 $h<n$ 时，有正合列
\[
  0\longrightarrow M_{h+1}
  \longrightarrow M_h
  \longrightarrow M_h/M_{h+1}
  \longrightarrow 0
\]
而归纳假设表明 $t_{M_{h+1}}$ 满射。另一方面，$M_h/M_{h+1}$ 被
$\mathfrak m$ 零化，因而是一个 $(A/\mathfrak m)$-模；换言之，它是诸
$A$-模（每一个都同构于 $k$）的直和。为证明 $t_{M_h/M_{h+1}}$ 满射，
只须证明 $t_k$ 满射，因为 $T$ 与直和交换。由交换图表
\[
\xymatrix@C=35pt@R=30pt{
  T(A_s)\otimes_A A_s
  \ar[r]^-{t_{A_s}}\ar[d]_{1\otimes\varepsilon}
  &
  T(A_s)\ar[d]^{T(\varepsilon)}
  \\
  T(A_s)\otimes_A k
  \ar[r]_-{t_k}
  &
  T(k)
}
\]
及（7.2.4），假设 $d)$ 确实蕴涵 $t_k$ 满射。为完成证明，只须证明：
若有正合列
\[
  0\longrightarrow M'
  \longrightarrow M
  \xrightarrow{v}M''
  \longrightarrow 0
\]
（它由 $A$-模组成），并且 $t_{M'}$ 与 $t_{M''}$ 都满射，则 $t_M$
满射。现有交换图表
\[
\xymatrix@C=27pt@R=30pt{
  T'(M')\ar[r]\ar[d]_{t_{M'}}
  &
  T'(M)\ar[r]\ar[d]_{t_M}
  &
  T'(M'')\ar[r]\ar[d]_{t_{M''}}
  &
  0\ar[d]
  \\
  T(M')\ar[r]
  &
  T(M)\ar[r]_-{T(v)}
  &
  T(M'')\ar[r]
  &
  \operatorname{Coker}(T(v))
}
\]
其中两行因 $T$ 半正合而正合。由归纳假设，$t_{M'}$ 与 $t_{M''}$ 都是
满态射，而最后一条竖直箭头是单态射，所以五引理（M, I, 1.1）表明
$t_M$ 是满态射。
\end{proof}

\subsection{模的同调函子的正合性判据}
\label{subsection:III.7.3.fr}

\begin{proposition}[7.3.1]
设 $A$ 是一个环（不必交换），$T_\bullet$ 是从范畴 $\mathbf{Ab}_A$
到范畴 $\mathbf{Ab}$、与直和交换的协变同调函子（T，II，2.1）。设
$p$ 是一个整数，使 $T_p$ 与 $T_{p-1}$ 均有定义。下列条件等价：
\begin{enumerate}
  \item[a)] $T_p$ 右正合。
  \item[b)] $T_{p-1}$ 左正合。
  \oldpage[III]{181}
  \item[c)] 对每个左 $A$-模 $M$，典范的函子性同态（7.2.2.1）
    \[
      T_p(A_s)\otimes_A M\longrightarrow T_p(M)
      \tag{7.3.1.1}\label{III.7.3.1.1.fr}
    \]
    是同构。
  \item[d)] 对每个左 $A$-模 $M$，同态（7.3.1.1）是满态射。
  \item[e)] $T_p$ 同构于形如
    $M\mapsto N\otimes_A M$ 的函子，其中 $N$ 是一个右 $A$-模。
\end{enumerate}
此外，若 $A$ 与 $\mathfrak m$ 满足（7.2.7）的条件，则上述条件还等价于
\begin{enumerate}
  \item[f)] 典范同态
    \[
      T_p(\varepsilon):T_p(A_s)\longrightarrow T_p(k)
    \]
    是满态射。
\end{enumerate}
\end{proposition}

\begin{proof}
按同调函子的定义，对每个使 $T_i$ 有定义的 $i$，函子 $T_i$ 都半正合；
并且对任意正合列
\[
  0\longrightarrow M'\xrightarrow{u}M\xrightarrow{v}M''\longrightarrow0,
\]
都有
\[
  \operatorname{Ker}(T_{i-1}(u))
  =
  \operatorname{Coker}(T_i(v)),
\]
因此 $a)$ 与 $b)$ 的等价性是显然的，其余断言由（7.2.5）和（7.2.7）
立即得到。
\end{proof}

\begin{corollary}[7.3.2]
设 $A$ 是交换环。沿用（7.3.1）的记号，假设 $T_p$ 右正合。若
$f\in A$ 不零化任何 $\ne0$ 的元素（这里这些元素取自某个 $A$-模
$M$），则 $f$ 也不零化任何 $\ne0$ 的元素（这里这些元素取自
$T_{p-1}(M)$）。特别地，若 $A$ 是整环，则 $A$-模
$T_{p-1}(A)$ 无挠。
\end{corollary}

\begin{proof}
事实上，若以 $h_f$ 表示位似 $x\mapsto fx$（定义在 $M$ 上），则假设
意味着 $h_f$ 是单射；$T_{p-1}(h_f)$ 也为单射，这是（7.3.1）的条件
$b)$ 所给出的。
\end{proof}

\begin{proposition}[7.3.3]
设 $A$ 是一个环，$T_\bullet$ 是从 $\mathbf{Ab}_A$ 到 $\mathbf{Ab}$、
与直和交换的协变同调函子。设 $p$ 是一个整数，使 $T_{p-1}$、$T_p$
与 $T_{p+1}$ 均有定义。下列条件等价：
\begin{enumerate}
  \item[a)] $T_p$ 正合。
  \item[b)] $T_{p+1}$ 与 $T_p$ 右正合。
  \item[c)] $T_p$ 与 $T_{p-1}$ 左正合。
  \item[d)] $T_{p+1}$ 右正合且 $T_{p-1}$ 左正合。
  \item[e)] 对每个 $A$-模 $M$，当 $i=p$ 与 $i=p+1$ 时，典范同态
    \[
      T_i(A_s)\otimes_A M\longrightarrow T_i(M)
      \tag{7.3.3.1}\label{III.7.3.3.1.fr}
    \]
    都是同构。
  \item[$e'$)] 对每个 $A$-模 $M$，当 $i=p$ 与 $i=p+1$ 时，典范同态
    （7.3.3.1）都是满态射。
  \item[f)] 对每个 $A$-模 $M$，当 $i=p$ 时同态（7.3.3.1）是同构，
    并且 $T_p(A_s)$ 是平坦右 $A$-模。
  \item[$f'$)] 对每个 $A$-模 $M$，当 $i=p$ 时同态（7.3.3.1）是满态射，
    并且 $T_p(A_s)$ 是平坦右 $A$-模。
\end{enumerate}
\end{proposition}

\begin{proof}
条件 $a)$、$b)$、$c)$、$d)$ 的等价性来自（7.3.1）中条件 $a)$ 与
$b)$ 的等价性。$b)$、$e)$、$e')$ 的等价性来自（7.3.1）中 $a)$、
$c)$、$d)$ 的等价性。最后，说 $T_p(A_s)$ 平坦，就是说函子
\[
  M\mapsto T_p(A_s)\otimes_A M
\]
左正合；$a)$、$f)$、$f')$ 的等价性仍由（7.3.1）中 $a)$、$c)$、
$d)$ 的等价性得到。
\end{proof}

\oldpage[III]{182}

\begin{corollary}[7.3.4]
假设 $A$ 交换、$T_p$ 正合，并且 $T_p(A)$ 是有限表示 $A$-模。则函数
\[
  x\mapsto\operatorname{rang}_{\kappa(x)}(T_p(\kappa(x)))
\]
在 $X=\operatorname{Spec}(A)$ 上局部常值；因而当
$\operatorname{Spec}(A)$ 连通时，它为常值。
\end{corollary}

\begin{proof}
事实上，因为 $T_p(A)$ 是平坦 $A$-模（由（7.3.3，$f)$），它是有限生成投射模，
故 $\bigl(T_p(A)\bigr)^{\sim}$ 是局部自由 $\mathscr O_X$-模
（Bourbaki，\emph{Alg. comm.}，第 II 章，\S 5，第 2 号，定理 1）。此外，
\[
  T_p(\kappa(x))=T_p(A)\otimes_A\kappa(x)
\]
（7.3.3，$e)$），而且已知在点 $x$ 处，$A$-模 $T_p(A)$ 的秩局部常值
（同上），于是得到推论。
\end{proof}

\begin{proposition}[7.3.5]
假设 $A$ 是一个左 Artin 环，其关于根基 $\mathfrak m$ 的商环是一个域
$k$。则（7.3.3）的条件还分别等价于下列每一个条件：
\begin{enumerate}
  \item[g)] 当 $i=p$ 与 $i=p+1$ 时，典范同态
    \[
      T_i(\varepsilon):T_i(A_s)\longrightarrow T_i(k)
    \]
    是满态射。
  \item[h)] $T_p(\varepsilon)$ 是满态射，并且 $T_p(A_s)$ 是平坦右
    $A$-模（或者，等价地（Bourbaki，\emph{Alg. comm.}，第 II 章，
    \S 3，第 2 号，命题 5 的推论 2），是自由 $A$-模）。
\end{enumerate}
再假设 $A$ 交换，且 $A$-模 $T_p(k)$ 的长度有限并等于 $d$。则上述条件
还分别等价于下列每一个条件：
\begin{enumerate}
  \item[i)] 对每个有限长度 $A$-模 $M$，有
    \[
      \operatorname{long}(T_p(M))=d\mathbin{.}\operatorname{long}(M).
      \tag{7.3.5.1}\label{III.7.3.5.1.fr}
    \]
  \item[j)] 有
    \[
      \operatorname{long}(T_p(A))=d\mathbin{.}\operatorname{long}(A).
      \tag{7.3.5.2}\label{III.7.3.5.2.fr}
    \]
\end{enumerate}
\end{proposition}

$g)$、$h)$ 与（7.3.3）中诸条件的等价性立即由（7.2.7）得到。为证明
其余断言，我们将使用下列引理：

\begin{lemma}[7.3.5.3]
设 $\mathcal K$、$\mathcal K'$ 是两个阿贝尔范畴，
$F:\mathcal K\to\mathcal K'$ 是协变加性函子；假设 $F$ 半正合，并且对
每个单对象 $S$（属于 $\mathcal K$），$F(S)$ 在 $\mathcal K'$ 中都是
有限长度对象。则对每个有限长度对象 $E$（属于 $\mathcal K$），$F(E)$ 在
$\mathcal K'$ 中也具有有限长度。对 $\mathcal K$ 中有限长度对象的任意
正合列
\[
  0\longrightarrow E'\xrightarrow{u}E\xrightarrow{v}E''\longrightarrow0
\]
都有
\[
  \operatorname{long}F(E)
  \leq
  \operatorname{long}F(E')+\operatorname{long}F(E'')
  \tag{7.3.5.4}\label{III.7.3.5.4.fr}
\]
并且（7.3.5.4）两边相等的充要条件是序列
\[
  0\longrightarrow F(E')\longrightarrow F(E)\longrightarrow F(E'')
  \longrightarrow0
\]
正合。
\end{lemma}

\begin{proof}
事实上，序列
\[
  F(E')\xrightarrow{F(u)}F(E)\xrightarrow{F(v)}F(E'')
\]
依假设正合；若假设 $F(E')$ 与 $F(E'')$ 具有有限长度，则
$\operatorname{Im}(F(u))$ 与 $\operatorname{Im}(F(v))$ 也具有有限长度。
又因为 $\operatorname{Ker}(F(v))=\operatorname{Im}(F(u))$，所以 $F(E)$
具有有限长度，并且
\[
  \operatorname{long}F(E)
  =
  \operatorname{long}\operatorname{Im}(F(u))
  +
  \operatorname{long}\operatorname{Im}(F(v))
  \leq
  \operatorname{long}F(E')+\operatorname{long}F(E'').
  \tag{7.3.5.5}\label{III.7.3.5.5.fr}
\]
\oldpage[III]{183}
对 $E$ 的长度作归纳，这已经证明了第一个断言。此外，（7.3.5.5）
两边相等，只能发生在
\[
  \operatorname{long}\operatorname{Im}(F(u))=\operatorname{long}F(E')
\]
（这等价于
$\operatorname{long}\operatorname{Ker}(F(u))=0$，亦即
$\operatorname{Ker}(F(u))=0$）以及
\[
  \operatorname{long}\operatorname{Im}(F(v))=\operatorname{long}F(E'')
\]
（这等价于
$\operatorname{long}\operatorname{Coker}(F(v))=0$，亦即
$\operatorname{Coker}(F(v))=0$）同时成立之时。
\end{proof}

现在注意，若 $M$ 是有限长度 $A$-模（其中 $A$ 交换），则 $M$ 的一个
Jordan--Hölder 列的各个商都必同构于 $A$-模 $k$。因此，对 $M$ 的长度
作归纳，由（7.3.5.4）得到
\[
  \operatorname{long}T_p(M)\leq d\mathbin{.}\operatorname{long}(M).
  \tag{7.3.5.6}\label{III.7.3.5.6.fr}
\]
此外，由（7.3.5.3），若 $T_p$ 正合，则（7.3.5.1）取等号；所以
（7.3.3）的条件 $a)$ 蕴涵 $i)$；显然 $i)$ 蕴涵 $j)$，余下只须证明
下列引理。

\begin{lemma}[7.3.5.7]
关系
\[
  \operatorname{long}T_p(A)=d\mathbin{.}\operatorname{long}A
\]
蕴涵 $T_p(\varepsilon)$ 是满态射，并且 $T_p(A)$ 是平坦 $A$-模。
\end{lemma}

\begin{proof}
事实上，从正合列
\[
  0\longrightarrow\mathfrak m\longrightarrow A\longrightarrow k
  \longrightarrow0,
\]
出发，由（7.3.5.4）和（7.3.5.6）有
\[
  \operatorname{long}T_p(A)
  \leq
  \operatorname{long}T_p(\mathfrak m)+\operatorname{long}T_p(k)
  \leq
  d\bigl(\operatorname{long}\mathfrak m+\operatorname{long}k\bigr)
  =d\mathbin{.}\operatorname{long}A
\]
而且等号成立时，由（7.3.5.3），序列
\[
  0\longrightarrow T_p(\mathfrak m)\longrightarrow T_p(A)
  \longrightarrow T_p(k)\longrightarrow0
  \tag{7.3.5.8}\label{III.7.3.5.8.fr}
\]
必正合。由（7.2.7）和（7.2.5），$T_p$ 同构于一个函子
$M\mapsto N\otimes_A M$；序列（7.3.5.8）的正合性以及
$\operatorname{Tor}$ 正合列表明
$\operatorname{Tor}_1^A(N,k)=0$。由此得出 $N=T_p(A)$ 是平坦 $A$-模
（0，10.1.3）。
\end{proof}

\begin{lemma}[7.3.6]
设 $A$ 是一个环，$T_\bullet$ 是从 $\mathbf{Ab}_A$ 到 $\mathbf{Ab}$、
与直和交换的协变同调函子。假设 $T_p$ 与 $T_{p+1}$ 均有定义，且
$T_p$ 左正合。则 $T_{p+1}$ 正合的充要条件是 $T_{p+1}(A_s)$ 为平坦右
$A$-模。
\end{lemma}

\begin{proof}
事实上，由（7.3.1）已知，典范同态
\[
  T_{p+1}(A_s)\otimes_A M\longrightarrow T_{p+1}(M)
\]
是函子的同构；应用平坦 $A$-模的定义即可。
\end{proof}

\begin{proposition}[7.3.7]
设 $A$ 是一个环，$T_\bullet$ 是从 $\mathbf{Ab}_A$ 到 $\mathbf{Ab}$、
与直和交换的协变同调函子。假设存在 $i_0$，使 $T_i$ 对 $i\leq i_0$
正合。则对每个整数 $p>i_0$，下列条件等价：
\begin{enumerate}
  \item[a)] $T_q$ 对 $q\leq p$ 正合。
  \item[b)] $T_q(A_s)$ 是平坦右 $A$-模，其中 $q\leq p$。
  \item[c)] 对每个 $A$-模 $M$，当 $q\leq p+1$ 时，典范同态
    \[
      T_q(A_s)\otimes_A M\longrightarrow T_q(M)
    \]
    是满射。
\end{enumerate}
\end{proposition}

\begin{proof}
$a)$ 与 $b)$ 的等价性由（7.3.6）对 $q$ 作归纳得到，因为按假设
$T_{i_0}$ 正合；$a)$ 与 $c)$ 的等价性由（7.3.3）中条件 $a)$ 与
$e')$ 的等价性得到。
\end{proof}

\oldpage[III]{184}

\begin{env}[7.3.8]
若 $A$ 是交换环，$B$ 是一个 $A$-代数（不必交换），$T_\bullet$ 是从
$\mathbf{Ab}_A$ 到 $\mathbf{Ab}$ 的协变同调函子，则由定义（7.1.3）可知，
从 $\mathbf{Ab}_B$ 到 $\mathbf{Ab}$ 的函子（它由把标量从 $A$ 扩张到
$B$ 得到）
（记作
\[
  T_\bullet^{(B)}=(T_i^{(B)}),
\]
）仍是一个同调函子。
\end{env}

\begin{corollary}[7.3.9]
假设 $T_\bullet$ 满足（7.3.7）的一般条件并与归纳极限交换；再假设
$A$ 是整环，且所有 $T_n(A)$ 都是有限表示 $A$-模。则对每个整数 $N$，
存在 $f\in A-\{0\}$，使函子
\[
  T_p^{(A_f)}:\mathbf{Ab}_{A_f}\longrightarrow\mathbf{Ab}
\]
对 $p\leq N$ 正合。
\end{corollary}

\begin{proof}
按假设，$T_i$ 对 $i\leq i_0$ 正合，所以 $T_i(A)$ 对这些 $i$ 平坦。
由（7.3.7，$b)$），只须选取 $f$，使得
\[
  T_p^{(A_f)}(A_f)=T_p(A_f)
\]
是自由 $A_f$-模，其中 $i_0<p\leq N$。又有
$T_p(A_f)=(T_p(A))_f$，这是因为 $T_p$ 与归纳极限交换（7.1.4）。若
$x$ 是 $\operatorname{Spec}(A)$ 的泛点，
则 $(T_p(A))_x$ 是 $A$ 的分式域上的有限维向量空间。由于每个
$T_p(A)$ 都是有限表示模，确实存在具有所需性质的 $f$
（Bourbaki，\emph{Alg. comm.}，第 II 章，\S 5，第 1 号，命题 2 的推论）。

应当注意，若只有有限多个指标 $i$ 满足 $T_i\ne0$，则存在
$f\in A-\{0\}$，使所有 $T_p^{(A_f)}$ 都正合。
\end{proof}

\begin{corollary}[7.3.10]
假设 $T_\bullet$ 满足（7.3.7）的一般条件并与归纳极限交换，又假设
$A$ 是交换 Noether 环，且 $T_n(A)$ 都是有限生成 $A$-模。则对每个整数
$N$，存在一个稠密开集 $U$（属于 $\operatorname{Spec}(A)$），使对每个
$p\leq N$，函数
\[
  x\mapsto\operatorname{rang}_{\kappa(x)}(T_p(\kappa(x)))
\]
在 $U$ 上为常值。
\end{corollary}

\begin{proof}
设 $\mathfrak p$ 是 $A$ 的一个极小素理想；按假设，环
$B=A/\mathfrak p$ 是整环，并且 $\operatorname{Spec}(B)$ 可与拓扑空间
$\operatorname{Spec}(A)$ 的一个不可约分支等同。我们对 $p\leq N$ 作
归纳，证明存在 $f_p\in B-\{0\}$，使得若置 $B'=B_{f_p}$，则
$T_i^{(B')}$ 正合，$T_i(B')$ 是有限生成 $B'$-模，其中 $i\leq p$。
当 $p\leq i_0$ 时，该命题由假设成立，取 $f_p=1$（于是
$B'=B=A/\mathfrak p$）即可；因为此时 $T_p$ 正合，$T_p(B)$ 同构于
$T_p(A)/T_p(\mathfrak p)$，故是有限生成 $A$-模，从而更是有限生成
$B$-模。现在对 $p$ 作归纳推理；$f_p$ 是在 $B$ 中的典范像，它来自
某个元素 $g_p\in A$。若置 $A'=A_{g_p}$，则 $B'=A'/\mathfrak p'$，其中
$\mathfrak p'$ 是 $A'$ 的一个极小素理想，等于 $\mathfrak p_{g_p}$。由于
\[
  T_i(A_{g_p})=(T_i(A))_{g_p},
\]
诸 $T_i(A')$ 都是有限生成 $A'$-模，因此函子 $T_\bullet^{(A')}$ 满足
与 $T_\bullet$ 相同的假设，只须把 $i_0$ 替换为 $p$。所以可归约到
$A'=A$ 且 $T_p$ 正合的情形；于是正合列
\[
  0\longrightarrow\mathfrak p\longrightarrow A\longrightarrow A/\mathfrak p
  \longrightarrow0
\]
给出正合列
\[
  T_{p+1}(A)\longrightarrow T_{p+1}(A/\mathfrak p)
  \xrightarrow{\partial}T_p(\mathfrak p)\longrightarrow T_p(A),
\]
又因 $T_p$ 正合，最右边的箭头为单射，所以
$T_{p+1}(A/\mathfrak p)$ 是 $T_{p+1}(A)$ 的商，从而有限生成。现在注意，
（7.3.9）的论证只使用了 $T_p(A)$ 在 $p\leq N$ 时有限生成这一事实；
因此可将该论证应用于整环 $B$、函子 $T_\bullet^{(B)}$ 以及 $N=p+1$，
从而完成归纳。这样，存在 $f_N\in B-\{0\}$，使
$\operatorname{Spec}(B_{f_N})$ 成为一个开集 $V$；它在
$\operatorname{Spec}(B)$ 中稠密，并且不与 $\operatorname{Spec}(A)$ 的
任何其他不可约分支
相交。
\oldpage[III]{185}
若命题已对 $B_{f_N}$ 证明，则存在开集 $W$；它在 $V$ 中稠密，且题述
函数在其中为常值，因为有 $A_x=(B_{f_N})_x$，这里 $x\in W$。对
$\operatorname{Spec}(A)$ 的每个不可约分支作同样
的论证，即得推论。因此可归约到 $A$ 为整环的情形；这时（7.3.9）的
论证证明存在 $f\in A-\{0\}$，使 $T_p(A_f)$ 都是有限生成自由
$A_f$-模，其中 $p\leq N$；由（7.3.4）即得（7.3.10）的结论。
\end{proof}

\begin{proposition}[7.3.11]
设 $A$ 是交换局部环，$k$ 是其剩余域，$T_\bullet$ 是从
$\mathbf{Ab}_A$ 到 $\mathbf{Ab}$、与直和交换的协变同调函子。假设存在
$i_0$，使 $T_i$ 对 $i\leq i_0$ 正合，并且所有 $T_n(A)$ 都是有限表示
$A$-模。则（7.3.7）的等价条件 $a)$、$b)$、$c)$ 蕴涵下列两个条件；
当该环还约化时，它们也与下列条件等价：
\begin{enumerate}
  \item[d)] 对每个 $x\in\operatorname{Spec}(A)$，有
    \[
      \operatorname{rang}_{\kappa(x)}T_q(\kappa(x))
      =
      \operatorname{rang}_kT_q(k)
      \quad\text{对 }q\leq p.
    \]
  \item[$d'$)] 对每个泛点 $x_j$（它属于 $\operatorname{Spec}(A)$ 的
    某个不可约分支），有
    \[
      \operatorname{rang}_{\kappa(x_j)}T_q(\kappa(x_j))
      =
      \operatorname{rang}_kT_q(k)
      \quad\text{对 }q\leq p.
    \]
\end{enumerate}
\end{proposition}

\begin{proof}
因为 $T_q(A)$ 是有限表示 $A$-模，（7.3.7）的条件 $b)$ 等价于说
$T_q(A)$ 是自由 $A$-模，其中 $q\leq p$（Bourbaki，\emph{Alg. comm.}，
第 II 章，\S 3，第 2 号，命题 5 的推论 2）；条件 $c)$ 蕴涵
\[
  T_q(\kappa(x))=T_q(A)\otimes_A\kappa(x)
  \quad\text{对 }q\leq p,
\]
所以（7.3.7）的等价条件蕴涵 $d)$，而 $d)$ 蕴涵 $d')$ 是显然的。
余下须证明 $d')$ 蕴涵 $a)$（在 $A$ 约化时）。我们对 $q\leq p$ 作
归纳，因为 $T_q$ 对 $q\leq i_0$ 正合。因而假设 $T_k$ 对 $k\leq q<p$
正合，并证明 $T_{q+1}(A)$ 是自由 $A$-模。由归纳假设，
$T_{q+1}(A)\otimes_A M$ 同构于 $T_{q+1}(M)$；对每个 $A$-模 $M$
都如此，这是由（7.3.7）的条件 $c)$ 和（7.3.3）得到的。把这一性质
应用于 $M=\kappa(x_j)$ 和 $M=k$，
由假设 $d')$ 得到
\[
  \operatorname{rang}_{\kappa(x_j)}
  \bigl(T_{q+1}(A)\otimes_A\kappa(x_j)\bigr)
  =
  \operatorname{rang}_kT_{q+1}(k)
\]
对每个 $j$ 都成立\SourceCorrectionNote{法文印本及外交式来源在此写作“对每个
$i$”；本工作版依 Canon corrected-current 改为与前文泛点 $x_j$ 一致的
指标 $j$，并保留本说明以便精确逆转。}；但这蕴涵 $T_{q+1}(A)$ 自由
（Bourbaki，\emph{Alg. comm.}，第 II 章，\S 3，第 2 号，命题 7），
从而完成证明。
\end{proof}

对我们将在（7.4）中研究的特殊类型的同调函子，上述结果将得到显著
加强；事实上，我们会得到只涉及单个 $T_p$ 的正合性判据。

\subsection[函子正合性判据]{函子 $H_\bullet(P_\bullet\otimes_A M)$ 的正合性判据}
\label{subsection:III.7.4.fr}

\begin{env}[7.4.1]
设 $A$ 是一个环（不必交换），$P_\bullet$ 是由平坦右 $A$-模组成的复形。
由于函子 $M\mapsto P_k\otimes_A M$ 在 $\mathbf{Ab}_A$ 中对每个 $k$
都正合，故 $\partial$-函子
\[
  T_\bullet(M)=H_\bullet(P_\bullet\otimes_A M)
  \tag{7.4.1.1}\label{III.7.4.1.1.fr}
\]
\oldpage[III]{186}
是从 $\mathbf{Ab}_A$ 到 $\mathbf{Ab}$ 的同调函子；当 $A$ 交换时，
它显然是 $A$-线性的（7.1.2），并且与归纳极限交换。

若 $A$ 交换，则对每个 $A$-代数 $B$，同调函子 $T_\bullet^{(B)}$
（7.3.8）依定义由下式给出：
\[
  T_\bullet^{(B)}(N)
  =
  H_\bullet(P_\bullet\otimes_A N_{[\rho]})
  \tag{7.4.1.2}\label{III.7.4.1.2.fr}
\]
其中 $\rho:A\to B$ 是定义 $B$ 的代数结构的同态；又因为也可写成
\[
  P_\bullet\otimes_A N_{[\rho]}
  =P_\bullet\otimes_A(B\otimes_B N)_{[\rho]}
  =(P_\bullet\otimes_A B)\otimes_B N,
\]
所以对每个 $B$-模 $N$ 都有
\[
  T_\bullet^{(B)}(N)=H_\bullet(P'_\bullet\otimes_B N)
  \tag{7.4.1.3}\label{III.7.4.1.3.fr}
\]
这里 $P'_\bullet$ 是复形 $P_\bullet\otimes_A B$，它由平坦 $B$-模
组成（0\textsubscript{I}，6.2.1）。
\end{env}

\begin{proposition}[7.4.2]
在（7.4.1）的一般条件下，对给定整数 $p\in\mathbf Z$，下列性质等价：
\begin{enumerate}
  \item[a)] $T_p$ 左正合（或者等价地，$T_{p+1}$ 右正合）。
  \item[b)]
    \[
      Z'_p(P_\bullet)
      =
      \operatorname{Coker}(P_{p+1}\longrightarrow P_p)
    \]
    是平坦右 $A$-模。
  \item[c)] 存在一个复形 $P'_\bullet$，它由平坦右 $A$-模组成，使微分
    \[
      d'_{p+1}:P'_{p+1}\longrightarrow P'_p
    \]
    为零，并且存在从 $H_\bullet(P_\bullet\otimes_A M)$ 到
    $H_\bullet(P'_\bullet\otimes_A M)$ 的同调函子同构。
\end{enumerate}
\end{proposition}

\begin{proof}
依定义，有一个关于 $M$ 函子性的正合列
\[
  0\longrightarrow T_p(M)
  \longrightarrow Z'_p(P_\bullet\otimes M)
  \longrightarrow P_{p-1}\otimes M
\]
其中
\[
  Z'_p(P_\bullet\otimes M)
  =
  \operatorname{Coker}(P_{p+1}\otimes M\longrightarrow P_p\otimes M)
  =
  Z'_p(P_\bullet)\otimes M
\]
这是张量积的右正合性所给出的。因而对每个同态 $f:M\to N$，有交换图
\[
\xymatrix@C=30pt@R=38pt{
  0\ar[r]
  &T_p(M)\ar[r]\ar[d]_u
  &Z'_p(P_\bullet)\otimes M\ar[r]\ar[d]_v
  &P_{p-1}\otimes M\ar[d]_w
  \\
  0\ar[r]
  &T_p(N)\ar[r]
  &Z'_p(P_\bullet)\otimes N\ar[r]
  &P_{p-1}\otimes N
}
\tag{7.4.2.1}\label{III.7.4.2.1.fr}
\]
其两行正合。若 $f$ 是单态射，则 $w$ 也是单态射，因为 $P_{p-1}$ 平坦；
若 $T_p$ 左正合，则 $u$ 也是单态射，从而推出 $v$ 是单态射，故
$Z'_p(P_\bullet)$ 平坦。反之，若它平坦，则 $v$ 对每个单态射
$f:M\to N$ 都是单态射，于是图（7.4.2.1）表明 $u$ 是单态射，
从而 $T_p$（它已经半正合）左正合。因此 $a)$ 与 $b)$ 等价。显然
$c)$ 蕴涵 $a)$，因为若 $d'_{p+1}:P'_{p+1}\to P'_p$ 为零，而
$0\to M'\to M\to M''\to0$ 是 $A$-模正合列，则正合列
\[
  H_{p+1}(P'_\bullet\otimes M'')
  \xrightarrow{\partial}H_p(P'_\bullet\otimes M')
  \longrightarrow H_p(P'_\bullet\otimes M)
\]
\oldpage[III]{187}
中的连接同态依定义为零（M，IV，1），故 $T_p$ 左正合。

反过来证明 $b)$ 蕴涵 $c)$。置
\SourceCorrectionNote{法文印本及外交式来源在此证明中使用了一个非正合的短列，
并把循环、边界与余循环在张量后对所有次数作了无依据的整体识别。
Canon disposition EGA3-2-CANON-DISP-0007 保持命题及其全部原假设不变，
以以下两个在任意左 $A$-模张量后仍正合的截断列及连接同态为零的论证
恢复原强度证明；本工作版采用该 corrected-current 证明并保留本说明以便逆转。}
\[
 C=Z'_p(P_\bullet),\qquad
 B=B_p(P_\bullet)=\operatorname{Im}(d_{p+1}),\qquad
 Z=Z_{p+1}(P_\bullet)=\operatorname{Ker}(d_{p+1}).
\]
有两个正合列
\[
 0\longrightarrow B\longrightarrow P_p\longrightarrow C\longrightarrow0,
 \qquad
 0\longrightarrow Z\longrightarrow P_{p+1}\longrightarrow B\longrightarrow0.
\]
由于 $C$ 与 $P_p$ 平坦，$B$ 平坦；由于 $B$ 与 $P_{p+1}$ 平坦，
$Z$ 平坦。这两个列在与任意左 $A$-模 $M$ 张量后仍正合
（$0_{\mathrm I}$，6.1.2）。特别地，典范同态把 $B\otimes_A M$
与 $d_{p+1}\otimes1_M$ 的像等同，并把 $Z\otimes_A M$ 与其核等同。

定义复形 $U_\bullet$ 如下：置 $U_i=P_i$（$i\geq p+2$），
$U_{p+1}=Z$，并置 $U_i=0$（$i\leq p$）；次数 $p+2$ 处的微分是
$d_{p+2}$，其值落在 $Z$ 中，其余非零微分继承自 $P_\bullet$。同样定义
$L_\bullet$：置 $L_i=P_i$（$i\leq p-1$），$L_p=C$，并置
$L_i=0$（$i\geq p+1$）；它在次数 $p$ 处的微分由 $d_p$ 在商模
$C$ 上诱导。于是 $P'_\bullet=U_\bullet\oplus L_\bullet$ 是由平坦模
组成的复形，并且 $P'_i=P_i$（$i\ne p,p+1$），$P'_p=C$，
$P'_{p+1}=Z$ 且 $d'_{p+1}=0$。

次数 $p+1$ 处的包含定义复形态射 $a:U_\bullet\to P_\bullet$，
次数 $p$ 处取商定义复形态射 $b:P_\bullet\to L_\bullet$。与 $M$
张量后，前者在所有次数 $i\geq p+1$ 上诱导同调同构：次数 $p+1$
的情形来自上述核的等同；次数 $p+2$ 的情形来自
$Z\otimes_A M\to P_{p+1}\otimes_A M$ 的单射性；更高次数不变。
后者在所有次数 $i\leq p$ 上诱导同构：在次数 $p$ 处取商掉像
$B\otimes_A M$；在次数 $p-1$ 处使用
$P_p\otimes_A M\to C\otimes_A M$ 的满射性；更低次数不变。
将前一组同构取逆并保留后一组同构，便对每个 $i$ 得到函子性同构
\[
 H_i(P_\bullet\otimes_A M)\xrightarrow{\sim}
 H_i(P'_\bullet\otimes_A M).
\]
它们与任意正合列 $0\to M'\to M\to M''\to0$ 所联系的连接同态
交换：在次数 $i\geq p+2$ 和 $i\leq p$，这来自复形态射 $a$ 与 $b$
的函子性。在次数 $p+1$，$P_\bullet\otimes_A M''$ 的每个循环都来自
$Z\otimes_A M''$，并通过满射 $Z\otimes_A M\to Z\otimes_A M''$
提升为 $P_\bullet\otimes_A M$ 的循环；故连接同态为零。$P'_\bullet$
的连接同态同样为零，因为 $d'_{p+1}=0$。于是得到了所需的
$\partial$-函子同构。
\end{proof}

应当注意，（7.4.2）的条件还蕴涵 $B_p(P_\bullet)$ 平坦，因为有正合列
\[
  0\longrightarrow B_p(P_\bullet)\longrightarrow P_p
  \longrightarrow Z'_p(P_\bullet)\longrightarrow0,
\]
其中 $P_p$ 与 $Z'_p(P_\bullet)$ 都平坦（$0_{\mathrm I}$，6.1.2）。

\begin{corollary}[7.4.3]
假设 $A$ 是维数 $1$ 的正则 Noether 环（换言之，是 Dedekind 环的乘积
（$0_{\mathrm{IV}}$，17.1.3 与 17.3.7）；例如主理想环）。则 $T_p$
左正合的充要条件是 $T_p(A)$ 为平坦 $A$-模；$T_p$ 正合的充要条件是
$T_p(A)$ 与 $T_{p-1}(A)$ 都为平坦 $A$-模。
\end{corollary}

\begin{proof}
回忆，对 Dedekind 环上的模 $M$，说 $M$ 平坦等价于说它无挠（$0_{\mathrm I}$，
6.3.3 与 6.3.4）；所以在（7.4.3）的假设下，平坦 $A$-模的每个子模
都平坦。

（7.4.3）的第二个断言由第一个断言得到，因为说 $T_p$ 正合就等于说
$T_p$ 与 $T_{p-1}$ 都左正合。为证明第一个断言，注意有正合列
\[
  0\longrightarrow H_p(P_\bullet)\longrightarrow Z'_p(P_\bullet)
  \longrightarrow B_{p-1}(P_\bullet)\longrightarrow0
\]
其中 $B_{p-1}(P_\bullet)$ 是平坦 $A$-模，因为它是平坦 $A$-模
$P_{p-1}$ 的子模。
所以 $H_p(P_\bullet)$ 平坦等价于 $Z'_p(P_\bullet)$ 平坦
（$0_{\mathrm I}$，6.1.2）。
\end{proof}

（7.4.2）最重要的应用如下。

\begin{proposition}[7.4.4]
设 $A$ 是 Noether 环，$P_\bullet$ 是由平坦 $A$-模组成的复形；假设或者
诸 $P_i$ 都为有限型，或者诸 $H_i(P_\bullet)$ 都为有限型 $A$-模且存在
$i_0$，使 $H_i(P_\bullet)=0$（$i<i_0$）。设 $T$ 是由（7.4.1.1）定义的
同调函子。则集合 $U$——其元素 $y\in\operatorname{Spec}(A)$ 满足
$(T_p)_y$（7.1.4）右正合（相应地，左正合、正合）——在
$\operatorname{Spec}(A)$ 中是开集。
\end{proposition}

\begin{proof}
在关于 $P_\bullet$ 的第二种假设下，可以把 $P_\bullet$ 替换为复形
$P'_\bullet$，
\oldpage[III]{188}
它由有限型自由 $A$-模组成，并使函子
$H_\bullet(P'_\bullet\otimes_A M)$ 作为 $\partial$-函子同构于
$T_\bullet(M)$（0，11.9.3）。因而总可归约到第一种假设；这时诸
$Z'_i(P_\bullet)$ 都为有限型。此外，只须证明关于左正合性的断言
（参见（7.4.2，$a)$）。现在设 $x\in U$；因为函子 $M\mapsto M_x$
正合，所以
\[
  (Z'_p(P_\bullet))_x=Z'_p((P_\bullet)_x),
\]
而结合（7.4.1.3），由假设及（7.4.2，$b)$）可知
$(Z'_p(P_\bullet))_x$ 是平坦 $A_x$-模；它又为有限型，并且 $A_x$
是 Noether 局部环，故它自由（Bourbaki，\emph{Alg. comm.}，第 II 章，
\S 3，第 2 号，命题 5 的推论 2）。因此存在 $f\in A$，使
$(Z'_p(P_\bullet))_f$ 是自由 $A_f$-模（Bourbaki，\emph{Alg. comm.}，
第 II 章，\S 5，第 1 号，命题 2 的推论），从而更有
$(Z'_p(P_\bullet))_y$ 都是自由 $A_y$-模（对每个 $y\in D(f)$）；由
（7.4.2，$b)$）即完成证明。
\end{proof}

\begin{corollary}[7.4.5]
在（7.4.4）的假设下，再假设 $A$ 是整环。则集合 $U$——其元素
$x\in\operatorname{Spec}(A)$ 满足 $(T_p)_x$ 正合——是非空开集。
\end{corollary}

\begin{proof}
事实上，只须证明 $(T_p)_x$ 对泛点 $x$（属于 $\operatorname{Spec}(A)$）
正合；这是显然的，因为 $A_x$ 是域，所以 $\mathbf{Ab}_{A_x}$ 上每个
加性函子都正合。
\end{proof}

\begin{proposition}[7.4.6]
在（7.4.4）的一般假设下，（7.4.2）的条件 $a)$、$b)$、$c)$ 还等价于：
\begin{enumerate}
  \item[d)] 存在 $A$-模 $Q$ 以及函子性同构
    \[
      T_p(M)\xrightarrow{\sim}\operatorname{Hom}_A(Q,M).
      \tag{7.4.6.1}\label{III.7.4.6.1.fr}
    \]
    此外，该性质在相差唯一同构的意义下决定 $A$-模 $Q$，并且它为有限型。
\end{enumerate}
\end{proposition}

\begin{proof}
$Q$ 的唯一性是可表函子的表示对象的唯一性（0，8.1.5）的特例。显然，
（7.4.6.1）的右端左正合。反过来，为证明 $Q$ 的存在，当 $T_p$ 左正合时，
先可像（7.4.4）中那样，归约到诸 $P_i$ 平坦且有限型的情形；因此，由于
$A$ Noether，它们是有限型投射模（Bourbaki，\emph{Alg. comm.}，
第 II 章，\S 5，第 2 号，定理 1 的推论）。此时对偶 $\check P_i$
（即 $P_i$ 的对偶）也是有限型投射 $A$-模，$P_i$ 典范同构于 $\check P_i$
的对偶，并且典范同态
\[
  P_i\otimes_A M\longrightarrow\operatorname{Hom}_A(\check P_i,M)
\]
为双射（Bourbaki，\emph{Alg.}，第 II 章，第 3 版，\S 4，第 2 号，
命题 2）。另一方面，由（7.4.2，$c)$）可假设
$d_{p+1}:P_{p+1}\to P_p$ 为零，故有正合列
\[
  0\longrightarrow T_p(M)\xrightarrow{u}P_p\otimes M
  \xrightarrow{v}P_{p-1}\otimes M
\]
其中 $v=d_p\otimes1$。置 $Q'=\operatorname{Ker}(d_p)$，于是有正合列
\[
  0\longrightarrow Q'\xrightarrow{w}P_p\xrightarrow{d_p}P_{p-1},
\]
转置后得到正合列
\[
  \check P_{p-1}\xrightarrow{{}^t d_p}\check P_p
  \xrightarrow{{}^t w}\check Q'\longrightarrow0.
\]
我们来说明
\[
  Q=\check Q'=\operatorname{Coker}({}^t d_p)
\]
满足要求。事实上，有正合列
\[
  0\longrightarrow\operatorname{Hom}_A(Q,M)
  \longrightarrow\operatorname{Hom}_A(\check P_p,M)
  \xrightarrow{v'}\operatorname{Hom}_A(\check P_{p-1},M)
\]
\oldpage[III]{189}
其中 $v'=\operatorname{Hom}({}^t d_p,1)$；在典范地把
$P_i\otimes M$ 与 $\operatorname{Hom}(\check P_i,M)$ 等同后，$v'$
也就与 $v=d_p\otimes1$ 等同，因而得到所求函子性同构
\[
  T_p(M)\xrightarrow{\sim}\operatorname{Hom}_A(Q,M)
\]
此外，$Q$ 是 $\check P_p$ 的商模，故为有限型。
\end{proof}

\begin{proposition}[7.4.7]
假设（7.4.4）的一般条件成立。则对每个有限型 $A$-模 $M$：
\begin{enumerate}
  \item[(i)] 诸 $T_i(M)$ 都是有限型 $A$-模。
  \item[(ii)] 对每个理想 $\mathfrak m$（属于 $A$），典范同态
    \[
      \widehat{T_i(M)}
      \longrightarrow
      \varprojlim_n T_i\bigl(M\otimes_A(A/\mathfrak m^{n+1})\bigr)
      \tag{7.4.7.1}\label{III.7.4.7.1.fr}
    \]
    是双射，其中左端是 $T_i(M)$ 对 $\mathfrak m$-预进制拓扑的分离完备化。
\end{enumerate}
\end{proposition}

\begin{proof}
像（7.4.4）中一样，先归约到诸 $P_i$ 都为有限型的情形；由于 $A$
Noether，$P_i\otimes_A M$ 的子模都为有限型，故断言（i）显然成立。
至于断言（ii），它由下面这个更一般的引理得到。
\end{proof}

\begin{lemma}[7.4.7.2]
设 $A$ 是 Noether 环，$u:E\to F$ 是有限型 $A$-模的同态。对每个有限型
$A$-模 $M$，置
\[
  K(M)=\operatorname{Ker}(u\otimes1_M),
  \qquad
  C(M)=\operatorname{Coker}(u\otimes1_M);
\]
则典范同态
\[
  \widehat{K(M)}\longrightarrow\varprojlim_n K(M_n),
  \qquad
  \widehat{C(M)}\longrightarrow\varprojlim_n C(M_n)
  \tag{7.4.7.3}\label{III.7.4.7.3.fr}
\]
都是双射，其中置
$M_n=M\otimes_A(A/\mathfrak m^{n+1})=M/\mathfrak m^{n+1}M$；这对每个
理想 $\mathfrak m$（属于 $A$）都成立。
\end{lemma}

\begin{proof}
由于 $E\otimes M$ 与 $F\otimes M$ 都为有限型，并且函子
$M\mapsto\widehat M$ 在有限型 $A$-模范畴中正合（$0_{\mathrm I}$，
7.3.3），故 $\widehat{K(M)}$ 与 $\widehat{C(M)}$ 分别是
\[
  \widehat{(u\otimes1)}:
  \widehat{(E\otimes M)}\longrightarrow\widehat{(F\otimes M)}.
\]
的核与余核。函子 $\varprojlim$ 的左正合性因而表明
$\widehat{K(M)}=\varprojlim_n K(M_n)$；另一方面，张量积的右正合性证明
$C(M_n)=C(M)\otimes_A(A/\mathfrak m^{n+1})$，故依定义
$\widehat{C(M)}=\varprojlim_n C(M_n)$。
\end{proof}

\begin{remark}[7.4.8]
考虑到（6.10.5）与（6.10.6），可以看出，在加上一个平坦性假设后，
（7.4.7）重新给出（4.1.7.1）为同构这一事实，即固有态射的“第一比较
定理”的实质；此外，该命题不再只适用于一个凝聚 $\mathscr O_X$-模，
而适用于由这样的模组成的复形。得到一个同时包含（7.4.7）和
（4.1.7.1）作为特例的命题将很有意义。还应注意，当诸 $P_i$ 为有限型时，
（7.4.7）的证明并未使用它们是平坦模这一事实；值得研究的是，如果既不
假设诸 $P_i$ 平坦也不假设其为有限型，而只假设诸 $H_i(P_\bullet)$ 对
每个 $i$ 都为有限型并且在 $i<i_0$ 时为零，（7.4.7）的结论是否仍成立。
这时能否把 $P_\bullet$ 替换为复形 $P'_\bullet$，后者由有限型 $A$-模
组成，并使函子 $H_\bullet(P_\bullet\otimes M)$ 与
$H_\bullet(P'_\bullet\otimes M)$（它们不再是同调函子）仍然同构？
\end{remark}

\oldpage[III]{190}

\subsection{Noether 局部环的情形}
\label{subsection:III.7.5.fr}

\begin{env}[7.5.1]
设 $A$ 是 Noether 局部环，$\mathfrak m$ 是它的极大理想；对每个 $A$-模
$M$，以 $\widehat M$ 表示该模对 $\mathfrak m$-预进制拓扑的分离完备化，
它同构于
\[
  \varprojlim_n\bigl(M\otimes_A(A/\mathfrak m^{n+1})\bigr)
  =
  \varprojlim_n\bigl(M/\mathfrak m^{n+1}M\bigr).
\]
设 $T$ 是从 $\mathbf{Ab}_A$ 到 $\mathbf{Ab}$ 的协变加性函子；典范同态
（7.2.3.1）
\[
  T(M)\otimes_A(A/\mathfrak m^{n+1})
  \longrightarrow
  T\bigl(M\otimes_A(A/\mathfrak m^{n+1})\bigr)
\]
显然组成一个由 $A$-同态构成的逆系，因而在取极限后给出一个
$\widehat A$-同态，它关于 $M$ 函子性：
\[
  \widehat{T(M)}
  \longrightarrow
  \varprojlim_n T(M_n)
  \tag{7.5.1.1}\label{III.7.5.1.1.fr}
\]
其中置
\[
  M_n=M\otimes_A(A/\mathfrak m^{n+1}),
  \qquad
  A_n=A/\mathfrak m^{n+1}.
\]
\end{env}

\begin{proposition}[7.5.2]
设 $A$ 是极大理想为 $\mathfrak m$ 的 Noether 局部环，$k=A/\mathfrak m$
是它的剩余域；设 $T$ 是从 $\mathbf{Ab}_A$ 到 $\mathbf{Ab}$ 的协变加性
函子，它半正合并且与归纳极限交换。再假设，对每个有限型 $A$-模 $M$，
$T(M)$ 都是有限型 $A$-模，并且典范同态（7.5.1.1）是同构。在这些条件下，
下列性质等价：
\begin{enumerate}
  \item[a)] $T$ 右正合。
  \item[b)] 对每个 $n$，函子 $N\mapsto T(N)$ 在有限型 $A_n$-模范畴中
    右正合（这等价于说，$T$ 在有限长度 $A$-模范畴中右正合）。
  \item[c)] 典范同态
    \[
      T(\varepsilon):T(A)\longrightarrow T(k)
    \]
    是满射。
  \item[d)] 对一切充分大的 $n$，典范同态
    \[
      T(A_n)\longrightarrow T(k)
    \]
    是满射。
\end{enumerate}
\end{proposition}

\begin{proof}
显然，a) 蕴涵 b)。来证明 b) 蕴涵 a)；也就是说，若
$u:M\to N$ 是 $A$-模的满同态，则 $T(u)$ 是满同态。由于 $T$ 与归纳
极限交换，并且在模范畴中函子 $\varinjlim$ 对滤过指标集是正合的，故可
归约到 $M$ 与 $N$ 都为有限型的情形。这时 $T(M)$ 与 $T(N)$ 都为有限型；
又因 $A$ 是 Noether 局部环，只须证明
\[
  \widehat{T(u)}:
  \widehat{T(M)}
  \longrightarrow
  \widehat{T(N)}
\]
为满射（$0_{\mathrm I}$，7.3.5 及 $0_{\mathrm I}$，6.4.1）。由假设，
$\widehat{T(M)}$ 与 $\widehat{T(N)}$ 分别为 $\varprojlim_n T(M_n)$ 与
$\varprojlim_n T(N_n)$，所以 $\widehat{T(u)}$ 是同态逆系
\[
  T(u\otimes1_{A_n}):T(M_n)\longrightarrow T(N_n).
\]
的极限。而 b) 正是说这些同态都是满射；此外，$T(M_n)$ 是有限型
$A_n$-模，且按假设 $A_n$ 是 Artin 环；由此可知 $\widehat{T(u)}$ 是
满射（$0_{\mathrm I}$，13.1.2 及 13.2.2）。显然，a) 蕴涵 c)；又因
$T(\varepsilon)$ 分解为
\[
  T(A)\longrightarrow T(A_n)\longrightarrow T(k),
\]
故 c) 蕴涵 d)；最后，由（7.2.7）可知 b) 与 d) 等价，因为 $T$
在 $\mathbf{Ab}_{A_n}$ 中半正合。这就完成了证明。
\end{proof}

\begin{corollary}[7.5.3]
在（7.5.2）的一般假设下，若 $T(k)=0$，则 $T(M)=0$ 对每个 $A$-模
$M$ 都成立。
\end{corollary}

\begin{proof}
\oldpage[III]{191}
由于 $k$ 是唯一的单 $A$-模，由（7.3.5.4）可知，$T(E)=0$ 对每个有限
长度 $A$-模 $E$ 都成立。现在若 $M$ 为有限型，则
$\widehat{T(M)}$ 同构于 $\varprojlim_n T(M_n)$；而诸 $M_n$ 都具有有限
长度，故 $\widehat{T(M)}=0$。按假设 $T(M)$ 为有限型，并且它同构于
$\widehat{T(M)}$ 的一个子模（$0_{\mathrm I}$，7.3.5），所以
$T(M)=0$。最后，对任意 $A$-模 $M$，$T(M)$ 是诸 $T(N_\alpha)$ 的
归纳极限，其中 $N_\alpha$ 遍历 $M$ 的有限型子模；证明至此完成。
\end{proof}

\begin{proposition}[7.5.4]
设 $A$ 是极大理想为 $\mathfrak m$ 的 Noether 局部环，$k=A/\mathfrak m$
是它的剩余域；设 $T_\bullet$ 是从 $\mathbf{Ab}_A$ 到 $\mathbf{Ab}$ 的
同调函子，并且与归纳极限交换。再假设，对每个 $i$ 及每个有限型
$A$-模 $M$，$T_i(M)$ 都是有限型的，并且典范同态
\[
  \widehat{T_i(M)}
  \longrightarrow
  \varprojlim_n T_i(M_n)
\]
为双射。对给定的整数 $p$，下列条件等价：
\begin{enumerate}
  \item[a)] $T_p$ 正合。
  \item[b)] $T_p$ 右正合，并且 $T_p(A)$ 是自由 $A$-模。
  \item[c)] 典范同态
    \[
      T_{p+1}(A)\longrightarrow T_{p+1}(k)
      \quad\text{及}\quad
      T_p(A)\longrightarrow T_p(k)
    \]
    都是满射。
  \item[d)] 对每个 $n$，典范同态
    \[
      T_{p+1}(A_n)\longrightarrow T_{p+1}(k)
      \quad\text{及}\quad
      T_p(A_n)\longrightarrow T_p(k)
    \]
    都是满射。
  \item[e)] 对每个 $n$，函子 $N\mapsto T_p(N)$ 在有限型 $A_n$-模
    范畴中正合。
\end{enumerate}
\end{proposition}

\begin{proof}
由（7.3.3）知道，a) 等价于 $T_{p+1}$ 与 $T_p$ 都右正合；又因为
$T_\bullet$ 是 $\mathbf{Ab}_{A_n}$ 范畴中的同调函子，与（7.3.1）中
相同的论证表明，e) 等价于 $T_p$ 与 $T_{p+1}$ 在有限型 $A_n$-模范畴中
都右正合。因此由（7.5.2）可知 a) 与 e) 等价；a)、c)、d) 的
等价性同样由（7.5.2）得到。最后，已知局部环 $A$ 上每个有限型平坦模
都是自由模（Bourbaki，\emph{Alg. comm.}，第 II 章，\S 3，第 2 号，
命题 5 的推论 2）；于是 a) 与 b) 的等价性由（7.3.1）及（7.3.3）
得到。
\end{proof}

\begin{corollary}[7.5.5]
假设（7.5.4）的一般条件成立。
\begin{enumerate}
  \item[(i)] 若 $T_p(k)=0$，则 $T_p=0$，$T_{p+1}$ 右正合，而
    $T_{p-1}$ 左正合。
  \item[(ii)] 若 $T_{p-1}(k)=T_{p+1}(k)=0$，则 $T_p$ 正合，典范同态
    \[
      T_p(A)\otimes_A M\longrightarrow T_p(M)
    \]
    为双射，并且 $T_p(A)$ 是自由 $A$-模。
\end{enumerate}
\end{corollary}

\begin{proof}
由于 $T_p$ 半正合，断言（i）立即由（7.5.3）得到；其最后一个断言则
来自同调函子的定义。考虑到（7.3.3），由（i）立即得到（ii）的前两个
断言；$T_p(A)$ 为自由模则由（7.5.4）得到。
\end{proof}

\begin{corollary}[7.5.6]
假设（7.5.4）的一般假设成立，并且还假设 $A$ 是离散赋值环。
\begin{enumerate}
  \item[(i)] $T_p$ 右正合，当且仅当 $T_{p-1}(A)$ 是自由 $A$-模。
  \item[(ii)] $T_p$ 正合，当且仅当 $T_p(A)$ 与 $T_{p-1}(A)$ 都是
    自由 $A$-模。
\end{enumerate}
\end{corollary}

\begin{proof}
\oldpage[III]{192}
显然，（ii）蕴涵（i）（参见（7.3.3））。\SourceCorrectionNote{法文印本及
外交式来源在此写作“（i）蕴涵（ii）”；当前校勘法文改作“（ii）蕴涵
（i）”，此处采用该更正读法。} 为证明（i），注意，若 $f$ 是 $A$ 的极大
理想的一个生成元（$A$ 的“一致化元”），则有限型 $A$-模 $M$ 是自由模
（或等价地，是平坦模），当且仅当位似 $h_f:x\mapsto fx$ 在 $M$ 上是
单射，因为在这里这等价于说 $M$ 无挠（$0_{\mathrm I}$，6.3.4）。现考虑
正合列
\[
  0\longrightarrow A\xrightarrow{h_f}A\longrightarrow k\longrightarrow0,
\]
它给出同调正合列
\[
  T_p(A)
  \longrightarrow T_p(k)
  \longrightarrow T_{p-1}(A)
  \xrightarrow{h_f}T_{p-1}(A).
\]
由此可见，$T_{p-1}(A)$ 为自由模，当且仅当
$T_p(A)\to T_p(k)$ 为满射；结论遂由（7.5.2）得到。
\end{proof}

\begin{remark}[7.5.7]
应注意，由（7.4.7），（7.4.4）的一般假设蕴涵：由（7.4.1.1）定义的
同调函子 $T_\bullet$ 满足（7.5.4）的一般假设。在这种情形下，
（7.5.6）因而已包含在（7.4.3）中。
\end{remark}

\subsection{正合性质的下降、半连续性定理与 Grauert 正合性判据}
\label{subsection:III.7.6.fr}

\begin{proposition}[7.6.1]
在（7.4.1）的条件下，设 $B$ 是交换 $A$-代数。若 $T_p$ 右正合（相应地，
左正合、正合），则 $T_p^{(B)}$ 也具有同一性质；当 $B$ 是忠实平坦
$A$-模时，逆命题成立。
\end{proposition}

\begin{proof}
第一个断言是（7.1.3）中一个平凡断言的特例。反过来，先假设 $B$ 是平坦
$A$-模。于是对每个 $A$-模 $M$ 都有
\[
  H_\bullet\bigl(P_\bullet\otimes_A(M\otimes_A B)\bigr)
  =
  \bigl(H_\bullet(P_\bullet\otimes_A M)\bigr)\otimes_A B,
\]
这也可以对每个 $p$ 写成
\[
  T_p(M)\otimes_A B
  =
  T_p^{(B)}(M_{(B)})
  \tag{7.6.1.1}\label{III.7.6.1.1.fr}
\]
其中等号理解为相差一个典范同构。假设 $T_p^{(B)}$ 右正合（相应地，
左正合、正合）；由于 $M\mapsto M_{(B)}$ 是正合函子，（7.6.1.1）的
左端作为 $M$ 的函子右正合（相应地，左正合、正合）；现在若 $B$ 在
$A$ 上忠实平坦，便可推出 $T_p$ 具有同样的正合性质（$0_{\mathrm I}$，
6.4.1）。
\end{proof}

\begin{proposition}[7.6.2]
在（7.4.1）的条件下，再假设 $A$ 是约化 Noether 环，诸 $P_i$ 都是
有限型 $A$-模。则 $T_p$ 右正合（相应地，左正合、正合），当且仅当对
每个本身为离散赋值环的 $A$-代数 $B$，$T_p^{(B)}$ 具有相应性质。
\end{proposition}

\begin{proof}
由（7.3.1）与（7.3.3），只须考虑右正合性；自然也只须证明条件
（7.6.1）的充分性。由（7.4.2），只须证明 $Z'_{p-1}(P_\bullet)$ 是
平坦 $A$-模。因为 $P_{p-1}$ 为有限型，$Z'_{p-1}(P_\bullet)$ 也为
有限型；判据（$0$，10.2.8）表明，只须
$Z'_{p-1}(P_\bullet)\otimes_A B$ 是平坦 $B$-模，这对每个本身为离散
赋值环的 $A$-代数 $B$ 都成立。而因为
$P_\bullet$ 是平坦 $A$-模的复形，有
\[
  Z'_{p-1}(P_\bullet)\otimes_A B
  =
  Z'_{p-1}(P_\bullet\otimes_A B);
\]
\oldpage[III]{193}
$P_\bullet\otimes_A B$ 是平坦 $B$-模的复形（$0_{\mathrm I}$，6.2.1），
并且对每个 $B$-模 $N$ 有
\[
  H_\bullet(P_\bullet\otimes_A N)
  =
  H_\bullet\bigl((P_\bullet\otimes_A B)\otimes_B N\bigr),
\]
所以
\[
  T_p^{(B)}(N)
  =
  H_p\bigl((P_\bullet\otimes_A B)\otimes_B N\bigr);
\]
把（7.4.2）应用于 $T_p^{(B)}$，可见 $T_p^{(B)}$ 右正合这一假设等价于
$Z'_{p-1}(P_\bullet\otimes_A B)$ 是平坦 $B$-模。

前述判据促使我们更细致地研究离散赋值环的情形：
\end{proof}

\begin{proposition}[7.6.3]
在（7.4.1）的条件下，假设 $A$ 是维数为 1 的正则 Noether 环（换言之，
$A$ 为 Noether 环，并且对每个 $x\in\operatorname{Spec}(A)$，$A_x$ 是域
或离散赋值环）。则对每个整数 $p$ 及每个 $A$-模 $M$，有一个关于 $M$
函子性的典范正合列
\[
  0
  \longrightarrow T_p(A)\otimes_A M
  \xrightarrow{t_M}T_p(M)
  \longrightarrow \operatorname{Tor}_1^A\bigl(T_{p-1}(A),M\bigr)
  \longrightarrow0.
  \tag{7.6.3.1}\label{III.7.6.3.1.fr}
\]
\end{proposition}

\begin{proof}
在下文中，为简化记号，我们将省去复形 $P_\bullet$ 在通常同调记号
$H_p(P_\bullet)$、$B_p(P_\bullet)$、$Z_p(P_\bullet)$ 与
$Z'_p(P_\bullet)$ 中的出现。有三个正合列
\begin{gather*}
  0\longrightarrow H_p\longrightarrow Z'_p\longrightarrow B_{p-1}
  \longrightarrow0,\\
  0\longrightarrow B_{p-1}\longrightarrow Z_{p-1}\longrightarrow H_{p-1}
  \longrightarrow0,\\
  0\longrightarrow Z_{p-1}\longrightarrow P_{p-1}\longrightarrow B_{p-2}
  \longrightarrow0.
\end{gather*}
由于 $P_{p-1}$ 与 $P_{p-2}$ 平坦，它们的相应子模 $B_{p-1}$、
$Z_{p-1}$ 与 $B_{p-2}$ 也平坦，因为平坦 $A_x$-模与无挠 $A_x$-模等同
（对每个 $x\in\operatorname{Spec}(A)$）；
因此与 $M$ 作张量积后，得到正合列
\[
  0=\operatorname{Tor}_1^A(B_{p-1},M)
  \longrightarrow H_p\otimes M
  \longrightarrow Z'_p\otimes M
  \xrightarrow{u}B_{p-1}\otimes M
  \longrightarrow0
  \tag{7.6.3.2}\label{III.7.6.3.2.fr}
\]
\[
  0=\operatorname{Tor}_1^A(Z_{p-1},M)
  \longrightarrow\operatorname{Tor}_1^A(H_{p-1},M)
  \longrightarrow B_{p-1}\otimes M
  \xrightarrow{v}Z_{p-1}\otimes M
  \tag{7.6.3.3}\label{III.7.6.3.3.fr}
\]
\[
  0=\operatorname{Tor}_1^A(B_{p-2},M)
  \longrightarrow Z_{p-1}\otimes M
  \xrightarrow{w}P_{p-1}\otimes M.
  \tag{7.6.3.4}\label{III.7.6.3.4.fr}
\]
依定义，
$T_p(M)=\operatorname{Ker}(d_p\otimes1)/\operatorname{Im}(d_{p+1}\otimes1)$；
所以它是同态
\[
  (P_p\otimes M)/\operatorname{Im}(d_{p+1}\otimes1)
  \longrightarrow P_{p-1}\otimes M
\]
的核；该同态由 $d_p\otimes1$ 取商得到，也可写成
$Z'_p\otimes M\to P_{p-1}\otimes M$，因为依定义 $Z'_p=P_p/B_p$；而这个
同态又可看成复合
\[
  Z'_p\otimes M
  \xrightarrow{u}B_{p-1}\otimes M
  \xrightarrow{v}Z_{p-1}\otimes M
  \xrightarrow{w}P_{p-1}\otimes M.
\]
由（7.6.3.4），$w$ 是单射，故有正合列
\[
  0\longrightarrow\operatorname{Ker}u
  \longrightarrow T_p(M)
  \longrightarrow\operatorname{Ker}v
  \longrightarrow0,
\]
考虑到（7.6.3.2）与（7.6.3.3），并注意依定义 $H_p=T_p(A)$，这正是
（7.6.3.1）。
\end{proof}

\begin{remarks}[7.6.4]
\begin{enumerate}
  \item[(i)] $H_\bullet(P_\bullet\otimes_A M)$ 是双复形
    $P_\bullet\otimes_A M$ 的同调，其中把 $M$ 看成只在次数 0 有一项的
    复形；因而由（6.3.6）与（6.3.2），它是一个正则谱序列的收敛项，
    该谱序列的 $E_2$ 项为
    \[
      E^2_{pq}
      =
      \operatorname{Tor}_p^A\bigl(H_q(P_\bullet),M\bigr)
      =
      \operatorname{Tor}_p^A\bigl(T_q(A),M\bigr).
    \]
    \oldpage[III]{194}
    而关于环 $A$ 的假设蕴涵 $\operatorname{Tor}_p^A(E,F)=0$，其中
    $p\geq2$，两个模都是任意 $A$-模（$0_{\mathrm{IV}}$，17.2.2）；已知
    （M，XV）这蕴涵正合列
    \[
      0\longrightarrow E^2_{0,q}
      \longrightarrow H_q(P_\bullet\otimes_A M)
      \longrightarrow E^2_{1,q-1}
      \longrightarrow0
    \]
    而这正是（7.6.3.1）。
  \item[(ii)] 考虑到（7.3.1），正合列（7.6.3.1）把（7.4.3）的结果
    作为特例重新给出。
\end{enumerate}
\end{remarks}

\begin{corollary}[7.6.5]
在（7.4.1）的条件下，假设 $A$ 是离散赋值环，其分式域为 $K$、剩余域
为 $k$，并且诸 $T_i(A)$ 都是有限型 $A$-模。则
\[
  \operatorname{rang}_k T_p(k)
  \geq
  \operatorname{rang}_k\bigl(T_p(A)\otimes_A k\bigr)
  \geq
  \operatorname{rang}_A T_p(A)
  =
  \operatorname{rang}_K T_p(K).
  \tag{7.6.5.1}\label{III.7.6.5.1.fr}
\]
此外，上式两端的项相等，当且仅当 $T_p$ 正合，也当且仅当 $T_p(A)$ 与
$T_{p-1}(A)$ 都是自由 $A$-模。
\end{corollary}

\begin{proof}
事实上，在正合列（7.6.3.1）中取 $M=k$；因为这里涉及的是 $k$-向量
空间，得到
\[
  \operatorname{rang}_k T_p(k)
  =
  \operatorname{rang}_k\bigl(T_p(A)\otimes_A k\bigr)
  +
  \operatorname{rang}_k\bigl(\operatorname{Tor}_1^A(T_{p-1}(A),k)\bigr).
\]
另一方面，因为 $T_p(A)$ 是局部整环 $A$ 上的有限型模，有
（Bourbaki，\emph{Alg. comm.}，第 II 章，\S 3，第 2 号，命题 4 的
推论 1）
\[
  \operatorname{rang}_k\bigl(T_p(A)\otimes_A k\bigr)
  \geq
  \operatorname{rang}_A T_p(A)
  =
  \operatorname{rang}_K\bigl(T_p(A)\otimes_A K\bigr)
  \tag{7.6.5.2}\label{III.7.6.5.2.fr}
\]
而且（7.6.5.2）的两端相等，当且仅当 $T_p(A)$ 是自由 $A$-模
（\emph{同上}，命题 7）。还应注意，因为 $K$ 是平坦 $A$-模，依定义有
\[
  T_p(A)\otimes_A K
  =
  H_p(P_\bullet)\otimes_A K
  =
  H_p(P_\bullet\otimes_A K)
  =
  T_p(K).
\]
因此确有不等式（7.6.5.1），并且还可看出等号成立必须满足：
1\textsuperscript{o} $T_p(A)$ 自由；
2\textsuperscript{o} $\operatorname{Tor}_1^A(T_{p-1}(A),k)=0$，而已知
（$0$，10.1.3）这一条件等价于 $T_{p-1}(A)$ 是自由 $A$-模。最后，因为诸
$T_i(A)$ 都是有限型 $A$-模，说它们平坦或自由是同一回事
（Bourbaki，\emph{Alg. comm.}，第 II 章，\S 3，第 2 号，命题 5 的
推论 2）；结论遂由（7.4.3）得到。
\end{proof}

\begin{env}[7.6.6]
仍在（7.4.1）的假设下，对每个 $x\in\operatorname{Spec}(A)$，置
\[
  d_p(x)=d_p^T(x)=\operatorname{rang}_{\kappa(x)}T_p(\kappa(x)).
  \tag{7.6.6.1}\label{III.7.6.6.1.fr}
\]
\end{env}

\begin{lemma}[7.6.7]
设 $\varphi:A\to A'$ 是环同态，并设
\[
  f={}^a\!\varphi:
  \operatorname{Spec}(A')\longrightarrow\operatorname{Spec}(A)
\]
为相应映射（I，1.2.1）。若置 $T'_\bullet=T_\bullet^{(A')}$（7.1.3），则
\[
  d_p^{T'}=d_p^T\circ f.
  \tag{7.6.7.1}\label{III.7.6.7.1.fr}
\]
\end{lemma}

\begin{proof}
\oldpage[III]{195}
事实上，对每个 $x'\in\operatorname{Spec}(A')$，置 $x=f(x')$，则
\[
  H_\bullet(P_\bullet\otimes_A\kappa(x'))
  =
  H_\bullet\bigl((P_\bullet\otimes_A\kappa(x))
    \otimes_{\kappa(x)}\kappa(x')\bigr)
  =
  H_\bullet(P_\bullet\otimes_A\kappa(x))
    \otimes_{\kappa(x)}\kappa(x'),
\]
因为 $\kappa(x')$ 在 $\kappa(x)$ 上平坦；由此得到关系（7.6.7.1）。
\end{proof}

\begin{lemma}[7.6.8]
若环 $A$ 为 Noether 环，而复形 $P_\bullet$ 由有限型 $A$-模组成，则函数
$x\mapsto d_p^T(x)$ 在 $\operatorname{Spec}(A)$ 上是可构造的。
\end{lemma}

\begin{proof}
须证明，对每个不可约闭子集 $Y$（属于 $X=\operatorname{Spec}(A)$），
存在一个非空开集 $U$（属于 $Y$），使 $d_p$ 在其中为常值（$0$，9.2.2）。由于
$Y=\operatorname{Spec}(A/\mathfrak a)$，其中 $\mathfrak a$ 是 $A$ 的一个
理想，并且 $A/\mathfrak a$ 约化，故由（7.6.7）可归约到 $Y=X$ 且 $A$
为 Noether 整环的情形；这时断言由（7.4.5）得到。
\end{proof}

\begin{theorem}[7.6.9]
设 $A$ 是 Noether 环，$P_\bullet$ 是由有限型平坦 $A$-模组成的复形，
$T_\bullet(M)=H_\bullet(P_\bullet\otimes_A M)$ 是由 $P_\bullet$ 定义的
同调函子；对每个 $x\in\operatorname{Spec}(A)$，置
\[
  d_p(x)=\operatorname{rang}_{\kappa(x)}T_p(\kappa(x)).
\]
则：
\begin{enumerate}
  \item[(i)] 函数 $d_p$ 在 $\operatorname{Spec}(A)$ 中可构造且上半连续。
  \item[(ii)] 若 $T_p$ 正合，则 $d_p$ 在 $\operatorname{Spec}(A)$ 中连续
    （因而局部常值）；当环 $A$ 约化时，逆命题成立。
\end{enumerate}
\end{theorem}

\begin{proof}
\begin{enumerate}
  \item[(i)] 第一个断言已在（7.6.8）中证明。为证明第二个断言，只须
    （$0$，9.2.4）证明：若 $x'\ne x$ 是 $x$ 在
    $\operatorname{Spec}(A)$ 中的一个一般化，则 $d_p(x')\leq d_p(x)$。
    这时存在一个离散赋值环
    $B$ 及态射
    \[
      f:\operatorname{Spec}(B)\longrightarrow\operatorname{Spec}(A)
    \]
    满足：若 $a$ 表示 $\operatorname{Spec}(B)$ 的闭点、$b$ 表示其泛点，
    则 $f(a)=x$ 且 $f(b)=x'$（II，7.1.9）。\SourceCorrectionNote{法文印本
    及外交式来源在此把第二个等式写成 $f(b)=y$；当前校勘法文改为与本段
    既定一般化点一致的 $f(b)=x'$，此处采用该更正读法。} 由公式
    （7.6.7.1），可见问题归约为证明不等式 $d_p(a)\geq d_p(b)$ 在
    $\operatorname{Spec}(B)$ 中成立；而这正是不等式（7.6.5.1）\footnote{把（i）的
    证明归约到离散赋值环情形这一思路，是 Hironaka 口头告知我们的。}。
  \item[(ii)] 第一个断言已经证明（7.3.4）。为证明逆命题，使用赋值判据
    （7.6.2）；考虑到公式（7.6.7.1），问题便归约到 $A$ 是离散赋值环的
    情形；但此时 $\operatorname{Spec}(A)$ 只有两个点，所以 $d_p$ 为常值
    这一假设由（7.6.5）确实蕴涵 $T_p$ 正合。
\end{enumerate}
\end{proof}

\begin{corollary}[7.6.10]
设 $A$ 是 Noether 环，$\mathfrak p_i$ $(1\leq i\leq r)$ 是它的极小素
理想，$k_i$ 是 $A_{\mathfrak p_i}$ 的剩余域 $(1\leq i\leq r)$。
\begin{enumerate}
  \item[(i)] 对每个 $x\in\operatorname{Spec}(A)$，存在一个指标 $i$，使
    \[
      d_p(x)\geq\operatorname{rang}_{k_i}T_p(k_i).
      \tag{7.6.10.1}\label{III.7.6.10.1.fr}
    \]
\end{enumerate}
特别地，若 $A$ 是整环，$K$ 是它的分式域，则对每个
$x\in\operatorname{Spec}(A)$ 都有
\[
  d_p(x)\geq\operatorname{rang}_K T_p(K)
  \tag{7.6.10.2}\label{III.7.6.10.2.fr}
\]
\oldpage[III]{196}
\begin{enumerate}
  \item[(ii)] 再假设 $A$ 局部且约化，并设 $k$ 是它的剩余域。则 $T_p$
    正合，当且仅当
    \[
      \operatorname{rang}_k T_p(k)
      =
      \operatorname{rang}_{k_i}T_p(k_i)
      \qquad\text{对 }1\leq i\leq r.
      \tag{7.6.10.3}\label{III.7.6.10.3.fr}
    \]
\end{enumerate}
\end{corollary}

\begin{proof}
（i）是立即的，因为 $x$ 的每个邻域都含有某个 $\mathfrak p_i$，只须应用
半连续性的定义。另一方面，若 $A$ 局部，则在 $\operatorname{Spec}(A)$ 中
极大理想 $\mathfrak m$ 的唯一邻域是整个 $\operatorname{Spec}(A)$，
所以对每个 $x\in\operatorname{Spec}(A)$ 都有
\[
  d_p(x)\leq\operatorname{rang}_k T_p(k)
\]
此关系结合（i）表明，条件（7.6.10.3）蕴涵 $d_p(x)$ 在
$\operatorname{Spec}(A)$ 中为常值，从而由（7.6.9，（ii））可知 $T_p$
正合；逆命题由（7.6.9，（ii））显然成立。
\end{proof}

\begin{remark}[7.6.11]
可以问，（7.6.9，（i））的断言能否加强为不等式
\[
  \operatorname{rang}_{\kappa(x)}T_p(\kappa(x))
  \geq
  \operatorname{rang}_{\kappa(x)}
    \bigl(T_p(A)\otimes_A\kappa(x)\bigr)
  \tag{7.6.11.1}\label{III.7.6.11.1.fr}
\]
对每个 $x\in\operatorname{Spec}(A)$ 都成立；当 $A$ 是离散赋值环且 $x$
是其极大理想时，它确实成立（7.6.5）。我们只考察 $A$ 为 Noether 局部
环的情形；设其极大理想为 $\mathfrak m$、剩余域为 $k$。这时，下列条件
等价：
\begin{enumerate}
  \item[a)] 对每个复形 $P_\bullet$，若它由有限型平坦 $A$-模组成，则都有
    \[
      \operatorname{rang}_k T_p(k)
      \geq
      \operatorname{rang}_k\bigl(T_p(A)\otimes_A k\bigr)
      \quad\text{对每个整数 }p.
      \tag{7.6.11.2}\label{III.7.6.11.2.fr}
    \]
  \item[b)] 对每个有限型 $A$-模 $M$，都有
    \[
      \operatorname{rang}_k(M\otimes_A k)
      \geq
      \operatorname{rang}_k(\check M\otimes_A k).
      \tag{7.6.11.3}\label{III.7.6.11.3.fr}
    \]
  \item[c)] 对每个有限型 $A$-模 $N$，都有
    \[
      \operatorname{rang}_k\bigl(\operatorname{Tor}_1^A(N,k)\bigr)
      \geq
      \operatorname{rang}_k\bigl(\operatorname{Tor}_2^A(N,k)\bigr).
      \tag{7.6.11.4}\label{III.7.6.11.4.fr}
    \]
\end{enumerate}
应注意，由移位（M，V，7.2），这等价于说，对每个 $i\geq1$ 都有
\[
  \operatorname{rang}_k\bigl(\operatorname{Tor}_i^A(N,k)\bigr)
  \geq
  \operatorname{rang}_k\bigl(\operatorname{Tor}_{i+1}^A(N,k)\bigr).
  \tag{7.6.11.5}\label{III.7.6.11.5.fr}
\]
下面迅速给出证明的一些说明。为看出 a) 蕴涵 b)，考虑正合列
\[
  L_1\xrightarrow{d}L_0\longrightarrow M\longrightarrow0
\]
其中 $L_0$ 与 $L_1$ 是有限型自由模，并把 a) 应用于复形
\[
  P_1\xrightarrow{{}^t d}P_0,
  \qquad
  P_0=\check L_1,
  \quad
  P_1=\check L_0,
\]
而其他各项均为零；此时 $T_1(A)=\check M$，并且
\[
  T_1(k)
  =
  \operatorname{Hom}_A(M,k)
  =
  \operatorname{Hom}_k(M/\mathfrak mM,k),
\]
换言之，$T_1(k)$ 是向量空间 $M\otimes_A k$ 的对偶，因而与后者同秩。
为证明 b) 蕴涵 c)，先证明下述引理：

\begin{lemma}[7.6.11.6]
给定一个由平坦 $A$-模组成的复形
\[
  \cdots\longrightarrow0\longrightarrow
  P_1\xrightarrow{d}P_0
  \longrightarrow0\longrightarrow\cdots
\]
则有正合列
\[
  0\longrightarrow\operatorname{Tor}_2^A(Z'_0,k)
  \longrightarrow T_1(A)\otimes_A k
  \longrightarrow T_1(k)
  \longrightarrow\operatorname{Tor}_1^A(Z'_0,k)
  \longrightarrow0.
  \tag{7.6.11.7}\label{III.7.6.11.7.fr}
\]
\end{lemma}

\begin{proof}
\oldpage[III]{197}
事实上，从正合列 $0\to Z_1\to P_1\to B_0\to0$ 出发，得到正合列
\[
  0\longrightarrow\operatorname{Tor}_1^A(B_0,k)
  \longrightarrow Z_1\otimes k
  \longrightarrow P_1\otimes k
  \xrightarrow{u}B_0\otimes k
  \longrightarrow0.
\]
由正合列 $0\to B_0\to Z_0\to Z'_0\to0$，又因 $Z_0=P_0$ 平坦，得到
$\operatorname{Tor}_1^A(B_0,k)=\operatorname{Tor}_2^A(Z'_0,k)$；依定义有
$Z_1=T_1(A)$；最后，$T_1(k)=\operatorname{Ker}(d\otimes1)$，而
$d\otimes1$ 分解为
\[
  P_1\otimes k
  \xrightarrow{u}B_0\otimes k
  \xrightarrow{v}Z_0\otimes k
  =P_0\otimes k;
\]
有 $T_1(k)=u^{-1}(R)$，其中 $R=\operatorname{Ker}v$；又因为 $u$ 满射，
$R=u(T_1(k))$；最后，$R=\operatorname{Tor}_1^A(Z'_0,k)$，这是由
$v$ 的定义得到的。这就证明了正合列（7.6.11.7）。
\end{proof}

为由 b) 推出 c)，考虑正合列
\[
  L_1\xrightarrow{d}L_0\longrightarrow N\longrightarrow0,
\]
其中 $L_0$ 与 $L_1$ 是有限型自由模；考虑函子 $T$，它与由 $L_1$ 和
$L_0$ 组成的复形相伴。因为 $L_0$ 与 $L_1$ 都自由，它们等同于自己的双对偶；
所以若
\[
  M=\operatorname{Coker}({}^t d),
  \qquad
  T_1(A)=\operatorname{Ker}(d)=\check M;
\]
另一方面，$M\otimes_A k=\operatorname{Coker}({}^t d\otimes1_k)$ 在 $k$
上的秩与 $\operatorname{Ker}(d\otimes1_k)$ 相同。因此由假设 b) 可知
\[
  \operatorname{rang}_k\bigl(T_1(A)\otimes_A k\bigr)
  \leq
  \operatorname{rang}_k\bigl(T_1(k)\bigr);
\]
又因 $Z'_0=N$，故不等式（7.6.11.4）由正合列（7.6.11.7）得到。

最后，为证明 c) 蕴涵 a)，把（7.6.11.6）应用于二项复形
$P_p\longrightarrow P_{p-1}$；把假设 c) 应用于模 $Z'_{p-1}$，得到
$\operatorname{rang}_kR\geq\operatorname{rang}_kS$，其中
\SourceCorrectionNote{法文印本及外交式来源在此说把 $P_0,P_1$ 替换为
$P_p,P_{p-1}$，并把假设 c) 应用于 $Z'_p$；当前校勘法文改为直接考察
二项复形 $P_p\to P_{p-1}$ 并应用于 $Z'_{p-1}$，此处采用该更正读法。}
\[
  R=\operatorname{Ker}(P_p\otimes k\longrightarrow P_{p-1}\otimes k)
  \quad\text{且}\quad
  S=Z_p\otimes k.
\]
若把 $d_{p+1}:P_{p+1}\to P_p$ 分解为
\[
  P_{p+1}\xrightarrow{v}Z_p\xrightarrow{j}P_p,
\]
则
\[
  \operatorname{Im}(d_{p+1}\otimes1)
  =
  (j\otimes1)\bigl(\operatorname{Im}(v\otimes1)\bigr).
\]
由于
\[
  T_p(A)\otimes_A k
  =(Z_p/B_p)\otimes_A k
  =(Z_p\otimes k)/\operatorname{Im}(v\otimes1),
\]
并且
\[
  T_p(k)=R/\operatorname{Im}(d_{p+1}\otimes1),
\]
便确实得到不等式（7.6.11.2）。

现在，假设局部环 $A$ 正则且维数为 $n$；已知 [17]，$A$-模
$\operatorname{Tor}_i^A(k,k)$ 同构于外幂
\[
  \bigwedge\nolimits^i(\mathfrak m/\mathfrak m^2);
\]
因而可见，条件（7.6.11.4）对 $N=k$ 不成立，只要 $n\geq4$。

反过来，若局部整环 $A$ 的每个有限型自反 $A$-模都自由（$A$ 为维数 2
的正则环时就是如此），则条件（7.6.11.3）成立：事实上，已知对偶
$\check M$（这里 $A$-模 $M$ 为有限型）是自反的，因而自由；所以
\[
  \operatorname{rang}_k(\check M\otimes_A k)
  =
  \operatorname{rang}_K(\check M)
  =
  \operatorname{rang}_K(M),
\]
其中 $K$ 是 $A$ 的分式域；另一方面，已知在 $k$ 上，$M\otimes_A k$ 的
每组基都由 $M$ 的一个生成系的像组成（Bourbaki，\emph{Alg. comm.}，
第 II 章，\S 3，第 2 号，命题 4 的推论 2），所以
$\operatorname{rang}_K(M)\leq\operatorname{rang}_k(M\otimes_A k)$，这就
证明了断言。
\end{remark}

\subsection{对固有态射的应用：I. 交换性质}
\label{subsection:III.7.7.fr}

以下三个小节实质上是把前几小节的结果改写为预概形态射的语言。

\begin{env}[7.7.1]
设 $f:X\to Y$ 是预概形的一个拟紧且分离的态射，且 $\mathscr P_\bullet$
是由拟凝聚 $\mathscr O_X$-模组成的复形，其微分算子次数为 $-1$；再假设
$\mathscr O_X$-模 $\mathscr P_i$ 都在 $Y$ 上平坦（$0_{\mathrm I}$，6.7.1）。
\oldpage[III]{198}
我们考察 $\partial$-函子
$\mathscr M\mapsto\mathscr T_\bullet(\mathscr M)$（也记作
$\mathscr T_\bullet(\mathscr P_\bullet,\mathscr M)$）；它定义在拟凝聚
$\mathscr O_Y$-模范畴上，其值仍在拟凝聚 $\mathscr O_Y$-模范畴中
（由（6.2.3）），并由下式给出：
\[
  \mathscr T_n(\mathscr P_\bullet,\mathscr M)
  =\mathscr T_n(\mathscr M)
  =\mathscr H^{-n}
    \bigl(f,\mathscr P^\bullet\otimes_{\mathscr O_Y}\mathscr M\bigr)
  \qquad\text{对 }n\in\mathbf Z,
  \tag{7.7.1.1}\label{III.7.7.1.1.fr}
\]
其中 $\mathscr P^\bullet$ 的次数 $j$ 项为 $\mathscr P_{-j}$，故其微分算子
次数为 $+1$。如此定义的 $\mathscr T_\bullet$ 关于 $\mathscr M$ 是一个
\emph{同调函子}（6.2.6）。
\end{env}

\begin{env}[7.7.2]
设 $g:Y'\to Y$ 是一个态射，并置
\[
  X'=X_{(Y')}=X\times_Y Y',
  \qquad
  f'=f_{(Y')}:X'\longrightarrow Y',
\]
后者仍是拟紧且分离的态射；另一方面，令
\[
  \mathscr P'_\bullet
  =\mathscr P_\bullet\otimes_{\mathscr O_Y}\mathscr O_{Y'};
\]
这是由拟凝聚 $\mathscr O_{X'}$-模组成的复形；由（I，9.1.12），这些模
在 $Y'$ 上平坦，再由（$0_{\mathrm I}$，6.2.1）可作如下定义。我们置
（次数约定同上）
\[
  \mathscr T^{Y'}_\bullet(\mathscr M')
  =\mathscr H^\bullet
    \bigl(f',\mathscr P'^\bullet\otimes_{\mathscr O_{Y'}}\mathscr M'\bigr)
  =\mathscr H^\bullet
    \bigl(f',\mathscr P^\bullet\otimes_{\mathscr O_Y}\mathscr M'\bigr)
  \tag{7.7.2.1}\label{III.7.7.2.1.fr}
\]
它关于拟凝聚 $\mathscr O_{Y'}$-模 $\mathscr M'$ 是同调函子。当 $Y'$ 是
环 $A'$ 的仿射概形时，用 $\mathscr T^{A'}_\bullet$ 代替
$\mathscr T^{Y'}_\bullet$；于是对每个 $A'$-模 $M'$，有
\[
  \mathscr T^{A'}_\bullet(\widetilde{M'})
  =\bigl(\Gamma(Y',\mathscr T^{Y'}_\bullet(\widetilde{M'}))\bigr)^\sim;
\]
我们置
\[
  T^{A'}_\bullet(M')
  =\Gamma(Y',\mathscr T^{Y'}_\bullet(\widetilde{M'})),
\]
这是一个从 $A'$-模到 $A'$-模范畴的同调函子。

注意，若 $Y=\operatorname{Spec}(A)$ 也为仿射，则 $A'$-模函子
$T^{A'}_\bullet$ 与经由 $A$ 到 $A'$ 的标量扩张所得的函子相同；后者
从 $A$-模同调函子 $T^A_\bullet$（7.1.3）出发。事实上，设 $g:Y'\to Y$ 是与
环同态 $A\to A'$ 对应的态射，$g':X'\to X$ 是相应的仿射态射
（II，1.6.2）；若 $\mathfrak U$ 是 $X$ 的一个仿射开覆盖，则
$\mathfrak U'=g'^{-1}(\mathfrak U)$ 是 $X'$ 的一个仿射开覆盖。由（6.2.2），
一切归结为验证
\[
  C^\bullet
    \bigl(\mathfrak U,
      \mathscr P_\bullet\otimes_{\mathscr O_Y}g_*(\mathscr M')\bigr)
  =C^\bullet
    \bigl(\mathfrak U',
      \mathscr P_\bullet\otimes_{\mathscr O_Y}\mathscr M'\bigr),
\]
并最终验证：对每个仿射开集 $U$（它是 $X$ 的开集），若置
$U'=g'^{-1}(U)$，则
\[
  \Gamma
    \bigl(U,\mathscr P_\bullet\otimes_{\mathscr O_Y}g_*(\mathscr M')\bigr)
  =\Gamma
    \bigl(U',\mathscr P_\bullet\otimes_{\mathscr O_Y}\mathscr M'\bigr),
\]
这由（I，1.3 和 3.2）显然成立。

特别地，若 $U$ 是 $Y$ 的开集，则对每个拟凝聚 $\mathscr O_Y$-模
$\mathscr M$，有
\[
  \mathscr T^U_\bullet(\mathscr M|U)
  =\bigl(\mathscr T_\bullet(\mathscr M)\bigr)|U.
  \tag{7.7.2.2}\label{III.7.7.2.2.fr}
\]
\end{env}

\begin{env}[7.7.3]
对每个拟凝聚 $\mathscr O_Y$-模 $\mathscr M$，都有一个关于 $\mathscr M$
函子性的典范同态：
\[
  \mathscr T_p(\mathscr O_Y)\otimes_{\mathscr O_Y}\mathscr M
  \longrightarrow
  \mathscr T_p(\mathscr M).
  \tag{7.7.3.1}\label{III.7.7.3.1.fr}
\]
事实上，若 $Y$ 为仿射，此同态已在（7.2.2）中定义；该定义可毫无困难地
推广到一般情形。只需注意，若 $U$、$V$ 是 $Y$ 的两个仿射开集且
$V\subset U$，则图表
\oldpage[III]{199}
\[
\xymatrix@C=18pt@R=32pt{
  {\begin{gathered}
    \bigl(\mathscr T_p(\mathscr O_Y)
      \otimes_{\mathscr O_Y}\mathscr M\bigr)|U\\
    =\mathscr T^U_p(\mathscr O_Y|U)
      \otimes_{\mathscr O_Y|U}(\mathscr M|U)
  \end{gathered}}
  \ar[r]\ar[d]
  &
  {\begin{gathered}
    \mathscr T^U_p(\mathscr M|U)\\
    =\bigl(\mathscr T_p(\mathscr M)\bigr)|U
  \end{gathered}}
  \ar[d]
  \\
  {\begin{gathered}
    \bigl(\mathscr T_p(\mathscr O_Y)
      \otimes_{\mathscr O_Y}\mathscr M\bigr)|V\\
    =\mathscr T^V_p(\mathscr O_Y|V)
      \otimes_{\mathscr O_Y|V}(\mathscr M|V)
  \end{gathered}}
  \ar[r]
  &
  {\begin{gathered}
    \mathscr T^V_p(\mathscr M|V)\\
    =\bigl(\mathscr T_p(\mathscr M)\bigr)|V
  \end{gathered}}
}
\]
由（7.2.3.3）可交换。

对每个态射 $g:Y'\to Y$，都有一个典范同态
\[
  \mathscr T_p(\mathscr O_Y)\otimes_{\mathscr O_Y}\mathscr O_{Y'}
  \longrightarrow
  \mathscr T^{Y'}_p(\mathscr O_{Y'}),
  \tag{7.7.3.2}\label{III.7.7.3.2.fr}
\]
它正是（6.7.11.2）在 $S=Y$、$v_1=f$、$v_2=1_Y$ 且
$\mathscr P^{(2)}_\bullet$ 仅由次数 0 的项 $\mathscr M$ 组成时的特例
（就收敛目标而言）。

当 $Y=\operatorname{Spec}(A)$、$Y'=\operatorname{Spec}(A')$ 都是仿射时，
（7.7.3.2）就是与（7.2.2）中定义的典范 $A'$-模同态
\[
  T^A_p(A)\otimes_A A'
  \longrightarrow
  T^{A'}_p(A')=T^A_p(A')
\]
对应的层同态；这很容易由（6.7.11）得到（因为在这里可于（6.7.11）中
取 $\mathfrak U'^{(i)}=u_i^{-1}(\mathfrak U^{(i)})$）。
\end{env}

\begin{env}[7.7.4]
当 $f$ 是\emph{真}态射、$Y=\operatorname{Spec}(A)$ 是仿射 Noether
概形，而 $\mathscr P_\bullet$ 是由凝聚 $\mathscr O_X$-模组成、在 $Y$
上平坦的\emph{下有界}复形时，我们已在（6.10.5）中看到，
在相差一个同构的意义下可写成
\[
  \mathscr T_p(\mathscr M)
  =\mathscr H_p
    \bigl(\mathscr L_\bullet\otimes_{\mathscr O_Y}\mathscr M\bigr),
  \qquad
  \mathscr L_\bullet=\widetilde{L_\bullet},
\]
其中 $L_\bullet$ 是一个由\emph{有限型自由 $A$-模组成的下有界复形}；
因此，函子 $\mathscr T_\bullet$ 正是（7.4）和（7.6）中已详细研究的
类型。下面把该研究的结果改写出来：
\end{env}

\begin{theorem}[7.7.5]
设 $Y$ 是局部 Noether 预概形，$(U_\alpha)$ 是由仿射开集组成的 $Y$ 的
一个覆盖，$f:X\to Y$ 是固有态射，而 $\mathscr P_\bullet$ 是由凝聚
$\mathscr O_X$-模组成、在 $Y$ 上平坦的下有界复形。由（7.7.1.1）定义的同调
函子 $\mathscr T_\bullet(\mathscr M)$ 具有下列性质：

\medskip
\noindent
\textnormal{I)（半连续性质）。}\footnote{该定理的一个特例已经见于
Chow--Igusa 的注记 [3]。半连续性质由 Kodaira--Spencer 在解析空间的框架中
（并在相当特殊的假设下）发现，见 “On the variations of almost-complex
structures”，\emph{Algebraic Geometry and Topology, A Symposium in honor of
S. Lefschetz}，Princeton Series 第 12 号，第 139--150 页，Princeton，1957；
一般版本由 Grauert [5] 证明。}
\emph{函数}
\[
  y\longmapsto
  d_p(y)
  =
  \operatorname{rang}_{\kappa(y)}
    \mathscr T_p^{\kappa(y)}(\kappa(y))
  \tag{7.7.5.1}\label{III.7.7.5.1.fr}
\]
\emph{是上半连续的。}

\oldpage[III]{200}
\medskip
\noindent
\textnormal{II)（交换性质）。}
\emph{对给定的整数 $p$，下列条件等价：}
\begin{enumerate}
  \item[a)] $\mathscr T_p$ 右正合。

  \item[a$'$)] $\mathscr T_p(\mathscr M)$ 同构于形如
    $\mathscr N\otimes_{\mathscr O_Y}\mathscr M$ 的函子，其中 $\mathscr N$
    必然同构于
    \[
      \mathscr T_p(\mathscr O_Y)
      =
      \mathscr H^{-p}(f,\mathscr P^\bullet).
    \]

  \item[a$''$)] 函子性的典范同态（7.7.3.1）是同构。

  \item[b)] $\mathscr T_{p-1}$ 左正合。

  \item[b$'$)] 存在一个 $\mathscr O_Y$-模 $\mathscr Q$（它必然凝聚，且
    在相差唯一同构的意义下确定）以及函子同构
    \[
      \mathscr T_{p-1}(\mathscr M)
      \xrightarrow{\sim}
      \mathscr{H}\!om_{\mathscr O_Y}(\mathscr Q,\mathscr M).
      \tag{7.7.5.2}\label{III.7.7.5.2.fr}
    \]

  \item[c)] 以 $A_\alpha$ 表示仿射开集 $U_\alpha$ 的环；对每个指标
    $\alpha$，$A_\alpha$-模函子 $T_p^{A_\alpha}$ 都右正合。

  \item[d)] 对每个态射 $g:Y'\to Y$，典范同态
    \[
      \mathscr T_p(\mathscr O_Y)
      \otimes_{\mathscr O_Y}\mathscr O_{Y'}
      \longrightarrow
      \mathscr T_p^{Y'}(\mathscr O_{Y'})
      \tag{7.7.5.3}\label{III.7.7.5.3.fr}
    \]
    都是同构。
\end{enumerate}
\end{theorem}

\begin{proof}
半连续性质在 $Y$ 上是局部的，因而由注（7.7.4）和（7.6.9）得到。显然，
a$''$) 蕴涵 a$'$)，而 a$'$) 蕴涵 a)。a)、a$''$)、b) 与 b$'$) 的
等价性，当 $Y$ 为仿射时已在（7.3.1）和（7.4.6）中证明，其中要计及注
（7.7.4）。为过渡到一般情形，先证明 a) 等价于 c)，这也证明性质 a)
在 $Y$ 上具有局部性；同一论证也适用于证明 a$''$) 和 b) 的局部性。
因为 c) 蕴涵 a) 是显然的，只须证明其逆命题。显然，只要说明：对每个
仿射开集 $U$（它是 $Y$ 的开集）以及由拟凝聚 $(\mathscr O_Y|U)$-模组成的
每个正合列
\[
  0\longrightarrow\mathscr F'
  \longrightarrow\mathscr F
  \longrightarrow\mathscr F''
  \longrightarrow0
\]
都存在一个由拟凝聚 $\mathscr O_Y$-模组成的正合列
\[
  0\longrightarrow\mathscr G'
  \longrightarrow\mathscr G
  \longrightarrow\mathscr G''
  \longrightarrow0
\]
使得
\[
  \mathscr F'=\mathscr G'|U,
  \qquad
  \mathscr F=\mathscr G|U,
  \qquad
  \mathscr F''=\mathscr G''|U;
\]
但这立即由 $Y$ 局部 Noether 的假设和（I，9.4.2）得到：事实上，把
$\mathscr F$ 延拓为一个拟凝聚 $\mathscr O_Y$-模 $\mathscr G$，再把
$\mathscr F'$ 延拓为一个 $\mathscr O_Y$-子模 $\mathscr G'$，它是
$\mathscr G$ 的子模，并取 $\mathscr G''=\mathscr G/\mathscr G'$ 即可。

为证明一般情形下 b) 与 b$'$) 的等价性，注意当 $Y$ 为仿射时，已知
$\mathscr Q$ 在相差唯一同构的意义下确定；若 $U$ 是仿射概形 $Y$ 的一个
仿射开集，便得到函子同构
\[
  \mathscr T_{p-1}^{U}(\mathscr M|U)
  \xrightarrow{\sim}
  \mathscr{H}\!om_{\mathscr O_Y|U}
    (\mathscr Q|U,\mathscr M|U).
\]
在一般情形下，对每个仿射开集 $U$（它是 $Y$ 的开集），存在一个凝聚
$(\mathscr O_Y|U)$-模 $\mathscr Q_U$ 以及函子同构
\[
  \mathscr T_{p-1}^{U}(\mathscr M|U)
  \xrightarrow{\sim}
  \mathscr{H}\!om_{\mathscr O_Y|U}
    (\mathscr Q_U,\mathscr M|U);
\]
上述说明表明，若 $V$ 是包含于 $U$ 的仿射开集，则
$\mathscr Q_U|V=\mathscr Q_V$；于是得到满足（7.7.5.2）的
$\mathscr O_Y$-模 $\mathscr Q$ 的存在性与唯一性。

最后还须证明 a) 与 d) 等价。显然，d) 在 $Y$ 上具有局部性，而上面已
看到 a) 也如此；此外，d) 还在 $Y'$ 上具有局部性。当
$Y=\operatorname{Spec}(A)$、$Y'=\operatorname{Spec}(A')$ 时，我们已经
看到 $T_\bullet^{A'}$ 是由 $T_\bullet^A$ 经到 $A'$ 的标量扩张所得的函子；

\oldpage[III]{201}
因此 a$'$) 显然蕴涵（7.7.5.3）是同构。反过来，仍设
$Y=\operatorname{Spec}(A)$ 为仿射，并令 $A'$ 为 $A$-代数 $A\oplus M$，
其中 $M$ 是任意 $A$-模，且 $A'$ 中的乘法定义为
\[
  (a_1,m_1)(a_2,m_2)
  =
  (a_1a_2,a_1m_2+a_2m_1);
\]
于是
\[
  T_p^{A'}(A')
  =
  T_p(A\oplus M)
  =
  T_p(A)\oplus T_p(M),
\]
而（7.7.5.3）为双射的假设蕴涵典范映射
\[
  T_p(A)\otimes_A M\longrightarrow T_p(M),
\]
也为双射；换言之，d) 蕴涵 a$''$)，证明完毕。
\end{proof}

\begin{theorem}[7.7.6]
设 $Y$ 是局部 Noether 预概形，$f:X\to Y$ 是固有态射，$\mathscr F$ 是
凝聚 $\mathscr O_X$-模，并在 $Y$ 上平坦。则存在一个凝聚 $\mathscr O_Y$-模
$\mathscr Q$（在相差唯一同构的意义下确定），并且关于拟凝聚
$\mathscr O_Y$-模 $\mathscr M$ 有函子同构
\[
  f_*(\mathscr F\otimes_{\mathscr O_Y}\mathscr M)
  \xrightarrow{\sim}
  \mathscr{H}\!om_{\mathscr O_Y}(\mathscr Q,\mathscr M)
  \tag{7.7.6.1}\label{III.7.7.6.1.fr}
\]
（从而有函子同构
\[
  \Gamma(X,\mathscr F\otimes_{\mathscr O_Y}\mathscr M)
  \xrightarrow{\sim}
  \operatorname{Hom}_{\mathscr O_Y}(\mathscr Q,\mathscr M).
  \tag{7.7.6.2}\label{III.7.7.6.2.fr}
\]
\SourceCorrectionNote{法文印本及外交式来源在公式（7.7.6.2）末尾多出一个
孤立的右括号；当前校勘法文删除该括号，此处采用现行读法，并在本注保留
外交式异文。}
\end{theorem}

\begin{proof}
事实上，$\mathscr M\mapsto\mathscr F\otimes_{\mathscr O_Y}\mathscr M$
正合（0\textsubscript{I}，6.7.4），而 $f_*$ 左正合，所以函子
$\mathscr M\mapsto f_*(\mathscr F\otimes_{\mathscr O_Y}\mathscr M)$ 左正合。
只须应用（7.7.5，b)）与（7.7.5，b$'$)）的等价性，其中 $p=1$。
\end{proof}

\begin{corollary}[7.7.7]
设 $Y$ 是局部 Noether 预概形，$f:X\to Y$ 是固有态射，$\mathscr F$、
$\mathscr F'$ 是两个凝聚 $\mathscr O_X$-模，并在 $Y$ 上平坦；再设
$u:\mathscr F\to\mathscr F'$ 是同态。对拟凝聚 $\mathscr O_Y$-模
$\mathscr M$，考察两个函子
\begin{align*}
  \mathscr T(\mathscr M)
  &=
  \operatorname{Ker}
  \bigl(
    f_*(\mathscr F\otimes_{\mathscr O_Y}\mathscr M)
    \longrightarrow
    f_*(\mathscr F'\otimes_{\mathscr O_Y}\mathscr M)
  \bigr),\\
  T(\mathscr M)
  &=
  \Gamma(Y,\mathscr T(\mathscr M))\\
  &=
  \operatorname{Ker}
  \bigl(
    \Gamma(X,\mathscr F\otimes_{\mathscr O_Y}\mathscr M)
    \longrightarrow
    \Gamma(X,\mathscr F'\otimes_{\mathscr O_Y}\mathscr M)
  \bigr).
\end{align*}
则存在一个凝聚 $\mathscr O_Y$-模 $\mathscr R$（在相差唯一同构的意义下
确定）以及函子同构
\[
  \mathscr T(\mathscr M)
  \xrightarrow{\sim}
  \mathscr{H}\!om_{\mathscr O_Y}(\mathscr R,\mathscr M)
  \tag{7.7.7.1}\label{III.7.7.7.1.fr}
\]
\[
  T(\mathscr M)
  \xrightarrow{\sim}
  \operatorname{Hom}_{\mathscr O_Y}(\mathscr R,\mathscr M).
  \tag{7.7.7.2}\label{III.7.7.7.2.fr}
\]
\end{corollary}

\begin{proof}
只须证明（7.7.7.2）；事实上，这将在 $Y$ 为仿射时证明（7.7.7.1），
再利用可表函子的表示对象在相差唯一同构的意义下的唯一性（0，8.1.8），
像证明（7.7.5，b) 与 b$'$)）等价时那样过渡到一般情形。由（7.7.6），
存在两个凝聚 $\mathscr O_Y$-模 $\mathscr Q$、$\mathscr Q'$，给出函子同构
\[
  \Gamma(X,\mathscr F\otimes_{\mathscr O_Y}\mathscr M)
  \xrightarrow{\sim}
  \operatorname{Hom}_{\mathscr O_Y}(\mathscr Q,\mathscr M),
  \qquad
  \Gamma(X,\mathscr F'\otimes_{\mathscr O_Y}\mathscr M)
  \xrightarrow{\sim}
  \operatorname{Hom}_{\mathscr O_Y}(\mathscr Q',\mathscr M).
\]
而 $u:\mathscr F\to\mathscr F'$ 典范地定义了一个函子态射
\[
  \Gamma(X,\mathscr F\otimes_{\mathscr O_Y}\mathscr M)
  \longrightarrow
  \Gamma(X,\mathscr F'\otimes_{\mathscr O_Y}\mathscr M);
\]

\oldpage[III]{202}
与后者对应，存在唯一的同态 $v:\mathscr Q'\to\mathscr Q$，它是
$\mathscr O_Y$-模的同态，并使图表
\[
\xymatrix@C=42pt@R=28pt{
  \Gamma(X,\mathscr F\otimes_{\mathscr O_Y}\mathscr M)
    \ar[r]\ar[d]_-{\sim}
  &
  \Gamma(X,\mathscr F'\otimes_{\mathscr O_Y}\mathscr M)
    \ar[d]^-{\sim}
  \\
  \operatorname{Hom}_{\mathscr O_Y}(\mathscr Q,\mathscr M)
    \ar[r]
  &
  \operatorname{Hom}_{\mathscr O_Y}(\mathscr Q',\mathscr M)
}
\]
可交换（0，8.1.4）。因为反变函子
$\mathscr N\mapsto\operatorname{Hom}_{\mathscr O_Y}(\mathscr N,\mathscr M)$
在 $\mathscr O_Y$-模范畴中左正合，只须取
$\mathscr R=\operatorname{Coker}(v)$，便得到所求的同构（7.7.7.2）。
\end{proof}

\begin{corollary}[7.7.8]
在（7.7.6）关于 $X$、$Y$ 和 $f$ 的假设下，设 $\mathscr F$、
$\mathscr G$ 是两个满足下列条件的凝聚 $\mathscr O_X$-模：
（i）$\mathscr F$ 在 $Y$ 上平坦；（ii）$\mathscr G$ 同构于有限型局部自由
$\mathscr O_X$-模同态 $\mathscr E_1\to\mathscr E_0$ 的余核。对拟凝聚
$\mathscr O_Y$-模 $\mathscr M$，考察两个函子
\begin{align*}
  \mathscr T(\mathscr M)
  &=
  f_*
  \bigl(
    \mathscr{H}\!om_{\mathscr O_X}
      (\mathscr G,\mathscr F\otimes_{\mathscr O_Y}\mathscr M)
  \bigr),\\
  T(\mathscr M)
  &=
  \Gamma(Y,\mathscr T(\mathscr M))
  =
  \operatorname{Hom}_{\mathscr O_X}
    (\mathscr G,\mathscr F\otimes_{\mathscr O_Y}\mathscr M).
\end{align*}
则存在一个凝聚 $\mathscr O_Y$-模 $\mathscr N$（在相差唯一同构的意义下
确定）以及函子同构
\[
  \mathscr T(\mathscr M)
  \xrightarrow{\sim}
  \mathscr{H}\!om_{\mathscr O_Y}(\mathscr N,\mathscr M)
  \tag{7.7.8.1}\label{III.7.7.8.1.fr}
\]
\[
  T(\mathscr M)
  \xrightarrow{\sim}
  \operatorname{Hom}_{\mathscr O_Y}(\mathscr N,\mathscr M).
  \tag{7.7.8.2}\label{III.7.7.8.2.fr}
\]
\end{corollary}

\begin{proof}
由函子同构（$0_{\mathrm I}$，5.4.2.1），对 $\mathscr M$ 有函子同构
\begin{align*}
  \mathscr{H}\!om_{\mathscr O_X}
    (\mathscr E_i,\mathscr F\otimes_{\mathscr O_Y}\mathscr M)
  &\xrightarrow{\sim}
  \check{\mathscr E}_i
    \otimes_{\mathscr O_X}
    (\mathscr F\otimes_{\mathscr O_Y}\mathscr M)\\
  &\xrightarrow{\sim}
  (\check{\mathscr E}_i\otimes_{\mathscr O_X}\mathscr F)
    \otimes_{\mathscr O_Y}\mathscr M\\
  &\xrightarrow{\sim}
  \mathscr{H}\!om_{\mathscr O_X}(\mathscr E_i,\mathscr F)
    \otimes_{\mathscr O_Y}\mathscr M
\end{align*}
其中 $i=0,1$。置
$\mathscr F_i=\mathscr{H}\!om_{\mathscr O_X}(\mathscr E_i,\mathscr F)$，其中
$i=0,1$；
这些是凝聚 $\mathscr O_X$-模（$0_{\mathrm I}$，5.3.5），并在 $Y$ 上平坦
（$0_{\mathrm I}$，5.4.2）。令
$u=\mathscr{H}\!om(v,1_{\mathscr F}):\mathscr F_0\to\mathscr F_1$。由函子
$\mathscr H\mapsto\mathscr{H}\!om_{\mathscr O_X}
(\mathscr H,\mathscr F\otimes_{\mathscr O_Y}\mathscr M)$ 的左正合性，对
$\mathscr M$ 有函子同构
\begin{align*}
  \mathscr{H}\!om_{\mathscr O_X}
    (\mathscr G,\mathscr F\otimes_{\mathscr O_Y}\mathscr M)
  &\xrightarrow{\sim}
  \operatorname{Ker}
  \bigl(
    \mathscr{H}\!om_{\mathscr O_X}
      (\mathscr E_0,\mathscr F\otimes_{\mathscr O_Y}\mathscr M)
    \longrightarrow
    \mathscr{H}\!om_{\mathscr O_X}
      (\mathscr E_1,\mathscr F\otimes_{\mathscr O_Y}\mathscr M)
  \bigr)\\
  &\xrightarrow{\sim}
  \operatorname{Ker}
  \bigl(
    \mathscr F_0\otimes_{\mathscr O_Y}\mathscr M
    \longrightarrow
    \mathscr F_1\otimes_{\mathscr O_Y}\mathscr M
  \bigr).
\end{align*}
因为 $f_*$ 左正合，由此得到函子同构
\[
  f_*
  \bigl(
    \mathscr{H}\!om_{\mathscr O_X}
      (\mathscr G,\mathscr F\otimes_{\mathscr O_Y}\mathscr M)
  \bigr)
  \xrightarrow{\sim}
  \operatorname{Ker}
  \bigl(
    f_*(\mathscr F_0\otimes_{\mathscr O_Y}\mathscr M)
    \longrightarrow
    f_*(\mathscr F_1\otimes_{\mathscr O_Y}\mathscr M)
  \bigr)
\]
于是只须应用（7.7.7）。
\end{proof}

\oldpage[III]{203}
\begin{remarks}[7.7.9]
\textnormal{(i)} 在（7.7.6）、（7.7.7）、（7.7.8）中，$\mathscr O_Y$-模
$\mathscr Q$、$\mathscr R$、$\mathscr N$ 的构造都\emph{与基变换相容}。
例如，在（7.7.6）的情形中（沿用（7.7.2）的记号），对每个拟凝聚
$\mathscr O_{Y'}$-模 $\mathscr M'$，有同构
\[
  f'_*(\mathscr F\otimes_{\mathscr O_Y}\mathscr M')
  \xrightarrow{\sim}
  \mathscr{H}\!om_{\mathscr O_{Y'}}(g^*(\mathscr Q),\mathscr M')
\]
因为由（7.7.2）中的说明，只须验证
\[
  \operatorname{Hom}_{\mathscr O_Y}(\mathscr Q,g_*(\mathscr M'))
  =
  \operatorname{Hom}_{\mathscr O_{Y'}}(g^*(\mathscr Q),\mathscr M')
\]
而这正是（$0_{\mathrm I}$，4.4.3.1）。同样，在（7.7.7）中把 $Y$、$f$、
$\mathscr M$、$\mathscr F$、$\mathscr F'$ 分别替换为 $Y'$、$f'$、
$\mathscr M'$、$\mathscr F\otimes_{\mathscr O_X}\mathscr O_{X'}$、
$\mathscr F'\otimes_{\mathscr O_X}\mathscr O_{X'}$ 时，应把 $\mathscr R$
替换为 $g^*(\mathscr R)$。最后，在（7.7.8）中把 $X$、$Y$、$f$、
$\mathscr M$、$\mathscr F$ 分别替换为 $X'$、$Y'$、$f'$、$\mathscr M'$、
$\mathscr F\otimes_{\mathscr O_X}\mathscr O_{X'}$，并把 $\mathscr G$ 替换为
$\mathscr G'=\mathscr G\otimes_{\mathscr O_X}\mathscr O_{X'}$
\SourceCorrectionNote{法文印本及外交式来源在此把基环写成
$\mathscr O_Y$；当前校勘法文依 $\mathscr G$ 为 $\mathscr O_X$-模且
基变换落在 $X'$ 上，改为 $\mathscr O_X$，此处采用该更正读法。}时，应把
$\mathscr N$ 替换为 $g^*(\mathscr N)$。这是因为仍有正合列
\[
  \mathscr E'_1\longrightarrow
  \mathscr E'_0\longrightarrow
  \mathscr G'\longrightarrow0
  \qquad\text{且}\qquad
  \mathscr E'_i=\mathscr E_i\otimes_{\mathscr O_X}\mathscr O_{X'}
\]
\SourceCorrectionNote{同一处法文印本及外交式来源把上式张量积的基环写成
$\mathscr O_Y$；当前校勘法文改为 $\mathscr O_X$，此处采用该更正读法。}
以及等式
\[
  \check{\mathscr E}'_i
  \otimes_{\mathscr O_{X'}}
  (\mathscr F\otimes_{\mathscr O_Y}\mathscr M')
  =
  \check{\mathscr E}_i
  \otimes_{\mathscr O_X}
  (\mathscr F\otimes_{\mathscr O_Y}\mathscr M')
  \qquad(i=0,1).
\]

\medskip
\textnormal{(ii)} （7.7.8）的陈述中关于 $\mathscr G$ 的条件（ii），对任意凝聚
$\mathscr O_X$-模 $\mathscr G$ 总是成立，只要存在一个可逆 $\mathscr O_X$-模，
它在 $Y$ 上丰沛；例如，当 $Y$ 为仿射且
$f:X\to Y$ 是射影态射时就是如此。只须注意（并计及（II，5.5.1）），
存在有限型局部自由 $\mathscr O_X$-模 $\mathscr E_0$，使 $\mathscr G$
同构于 $\mathscr E_0$ 的一个商（II，2.7.10）。因为 $\mathscr E_0$ 与
$\mathscr G$ 都凝聚，其核 $\mathscr G_1$（这里的映射是
$\mathscr E_0\to\mathscr G$）
也凝聚；再对 $\mathscr G_1$ 应用同一结果，便得到正合列
$\mathscr E_1\to\mathscr E_0\to\mathscr G\to0$，其中 $\mathscr E_0$ 与
$\mathscr E_1$ 都是有限型局部自由模。

\medskip
\textnormal{(iii)} 我们将在第 V 章证明，在（7.7.8）中，限制性假设
（ii）是\emph{多余的}。
\end{remarks}

\begin{proposition}[7.7.10]
\textnormal{（交换性质的局部判据）。}在（7.7.5）的一般条件下，设 $y$
是 $Y$ 的一点，$p$ 是整数。下列性质等价：
\begin{enumerate}
  \item[a)] 函子 $T_p^{\mathscr O_y}$ 右正合。

  \item[b)] 典范同态
    \[
      T_p^{\mathscr O_y}(\mathscr O_y)
      \longrightarrow
      T_p^{\mathscr O_y}(\kappa(y))
    \]
    是满射。

  \item[c)] 对每个整数 $n$，典范同态
    \[
      T_p^{\mathscr O_y}
      (\mathscr O_y/\mathfrak m_y^{n+1})
      \longrightarrow
      T_p^{\mathscr O_y}(\kappa(y))
    \]
    是满射。
\end{enumerate}
此外，满足这些条件的 $y\in Y$ 所成的集合，是一个开集 $U$；它是 $Y$
中使 $\mathscr T_p^U$ 右正合的最大开集。
\end{proposition}

\begin{proof}
计及（7.7.4），a)、b)、c) 的等价性由（7.4.7）和（7.5.2）得到。
集合 $U$（即 $T_p^{\mathscr O_y}$ 右正合的点集）为开集，也由（7.4.4）
得到；反过来，若 $\mathscr T_p^V$ 右正合，则由（7.7.5，c)）和（7.6.1），
$T_p^{\mathscr O_y}$ 对每个 $y\in V$ 也右正合。
\end{proof}

\begin{corollary}[7.7.11]
若 $\mathscr T_p$ 右正合（相应地，左正合），则对每个态射 $g:Y'\to Y$，
$\mathscr T_p^{Y'}$ 也右正合（相应地，左正合）。当态射 $g$ 忠实平坦时，
逆命题成立。
\end{corollary}

\oldpage[III]{204}
\begin{proof}
第一项断言立即由（7.6.1）以及问题在 $Y$ 和 $Y'$ 上具有局部性得到；
后一局部性依（7.7.5，c) 与 b)）。为证明第二项断言，由（7.7.10），只须
验证对每个 $y\in Y$，$T_p^{\mathscr O_y}$ 右正合（相应地，左正合）。
但由假设，存在 $y'\in Y'$ 使 $g(y')=y$，并且 $\mathscr O_{y'}$ 是忠实
平坦的 $\mathscr O_y$-模；结论遂由假设和（7.6.1）得到。
\end{proof}

\begin{remarks}[7.7.12]
\textnormal{(i)} 在（7.7.4）的假设下，再假设 $\mathscr P_\bullet$ 是一个
\emph{有限}复形；则由（7.7.1）（因为可取一个\emph{有限}仿射开覆盖
$\mathfrak U$，它覆盖 $X$），双复形
$C^\bullet(\mathfrak U,\mathscr P^\bullet\otimes_{\mathscr O_Y}\mathscr M)$
也是有限的。更确切地说，存在有限二元组集合 $E$，它\emph{与 $\mathscr M$
无关}，使得
\[
  C^h(\mathfrak U,\mathscr P^k\otimes_{\mathscr O_Y}\mathscr M)=0
  \qquad\text{对所有 }(h,k)\notin E.
\]
由此可知，存在 $i_1$，使得对 $i\geq i_1$，都有
$\mathscr T_i(\mathscr M)=0$；这里任取拟凝聚 $\mathscr O_Y$-模
$\mathscr M$。特别地，$\mathscr T_i$ 关于 $\mathscr M$ 显然是正合函子，
对上述 $i$ 都如此；因而由（7.4.1），$Z'_i(L_\bullet)$ 是有限型
\emph{平坦} $A$-模（所以是有限型\emph{投射}模，因为 $A$ Noether），
对这些 $i$ 都如此。现考察复形 $(L'_\bullet)$，它由 $A$-模组成：令
$L'_i=L_i$（当 $i<i_1$），令 $L'_{i_1}=Z'_{i_1}(L_\bullet)$，令
$L'_i=0$（当 $i>i_1$）；再令
$\mathscr L'_\bullet=\widetilde{L'_\bullet}$。显然，
\[
  \mathscr H_i(\mathscr L'_\bullet\otimes_{\mathscr O_Y}\mathscr M)
  =
  \mathscr H_i(\mathscr L_\bullet\otimes_{\mathscr O_Y}\mathscr M)
\]
当 $i<i_1-1$ 时成立，而当 $i\geq i_1$ 时也成立（此时两边都为零）。
最后，因为依定义有
$\operatorname{Im}(Z'_{i_1}\otimes_A M)=\operatorname{Im}(L_{i_1}\otimes_A M)$，
还得到
$\mathscr H_i(\mathscr L'_\bullet\otimes_{\mathscr O_Y}\mathscr M)
=\mathscr H_i(\mathscr L_\bullet\otimes_{\mathscr O_Y}\mathscr M)$，此处
$i=i_1-1$。
因此可见，在（7.7.4）中可再假设 $\mathscr L_\bullet$ 是一个\emph{有限}
复形，只须把对 $\mathscr L_i$ 的要求改为：它们是\emph{局部自由}
$\mathscr O_Y$-模（与有限型投射 $A$-模相伴）。

上述论证特别适用于 $\mathscr P_\bullet$ 仅由次数 0 的\emph{一个}非零项
$\mathscr F\ne0$ 组成的情形（这时
$\mathscr T_n(\mathscr M)=\mathrm R^{-n}f_*
(\mathscr F\otimes_{\mathscr O_Y}\mathscr M)$）；此时可假设 $\mathscr L_i$
均为零，只要 $i>0$。此时最好采用上同调记号，因而写
$\mathscr T^{-p}$ 而不写 $\mathscr T_p$。

\medskip
\textnormal{(ii)} 若在（7.7.5）的陈述中不再假设 $\mathscr P_i$ 在 $Y$
上平坦，则各结论仍成立，只须改为令
\[
  \mathscr T_p(\mathscr M)
  =
  \operatorname{Tor}^Y_p
    (f,1_Y;\mathscr P_\bullet,\mathscr M).
  \tag{7.7.12.1}\label{III.7.7.12.1.fr}
\]
事实上，由（6.7.9），
$\operatorname{Tor}^Y_n(f,1_Y;\mathscr P_\bullet,\mathscr O_Y)$ 此时是一个
\emph{凝聚} $\mathscr O_Y$-模。计及（6.10.1），（6.10.5）的证明无需改变
仍可应用，并且仍表明：当 $Y=\operatorname{Spec}(A)$ 为仿射时，有
\[
  \mathscr T_p(\mathscr M)
  =
  \mathscr H_p(\mathscr L_\bullet\otimes_{\mathscr O_Y}\mathscr M),
  \qquad
  \mathscr L_\bullet=\widetilde{L_\bullet},
\]
其中 $L_\bullet$ 是由有限型自由 $A$-模组成的复形；这就证明了断言。
\end{remarks}

\subsection{对固有态射的应用：II. 上同调平坦性判据}
\label{subsection:III.7.8.fr}

\begin{definition}[7.8.1]
设 $X$、$Y$ 是两个预概形，$f:X\to Y$ 是拟紧分离态射，
$\mathscr P_\bullet$ 是由拟凝聚 $\mathscr O_X$-模组成且在 $Y$ 上平坦的复形，
$\mathscr T_\bullet$ 是由（7.7.1.1）定义在拟凝聚 $\mathscr O_Y$-模上的
同调函子，并设 $y$ 是 $Y$ 的一点。称 $\mathscr P_\bullet$

\oldpage[III]{205}
\emph{相对于 $Y$ 在点 $y$ 处、维数 $p$ 上同调平坦}（或
\emph{相对于 $Y$ 在点 $y$ 处、维数 $-p$ 上上同调平坦}），如果存在
开邻域 $U$，它是 $y$ 在 $Y$ 中的邻域，并使
$\mathscr T_p^U=\mathscr T_U^{-p}$ 正合。称 $\mathscr P_\bullet$
\emph{在维数 $p$ 上相对于 $Y$ 同调平坦}（或
\emph{在维数 $-p$ 上相对于 $Y$ 上同调平坦}），如果它相对于 $Y$
在每一点 $y\in Y$ 处、维数 $p$ 上同调平坦。

当 $\mathscr P_\bullet$ 相对于 $Y$（相应地，相对于 $Y$ 在点 $y$ 处）
对\emph{每个维数 $p$} 都同调平坦时，就简称 $\mathscr P_\bullet$
\emph{相对于 $Y$ 同调平坦}（相应地，相对于 $Y$ 在点 $y$ 处同调平坦），
或\emph{相对于 $Y$ 上同调平坦}（相应地，相对于 $Y$ 在点 $y$ 处
上同调平坦）。
\end{definition}

\begin{env}[7.8.2]
按定义，相对于 $Y$ 的同调平坦性概念在 $Y$ 上是\emph{局部的}。若 $Y$
是局部 Noether 预概形，或是一个\emph{概形}，则 $\mathscr P_\bullet$
相对于 $Y$ 在维数 $p$ 上同调平坦，当且仅当函子 $\mathscr T_p$
\emph{正合}。在 $Y$ 局部 Noether 的情形中，这一点已在（7.7.5）的证明
过程中得到；当 $Y$ 是概形时，论证相同，它以把（I，9.4.2）应用于一个
拟紧概形中的仿射开集为基础。
\end{env}

\begin{proposition}[7.8.3]
沿用（7.8.1）的记号与假设，下列条件等价：
\begin{enumerate}
  \item[a)] $\mathscr P_\bullet$ 相对于 $Y$ 在点 $y$ 处、维数 $p$ 上
    同调平坦。
  \item[b)] 存在一个开邻域 $U$，它是 $y$ 在 $Y$ 中的邻域，并使
    $\mathscr T_p^U$ 与 $\mathscr T_{p+1}^U$ 右正合。
  \item[c)] 存在一个开邻域 $U$，它是 $y$ 在 $Y$ 中的邻域，并使
    $\mathscr T_p^U$ 与 $\mathscr T_{p-1}^U$ 左正合。
  \item[d)] 存在一个开邻域 $U$，它是 $y$ 在 $Y$ 中的邻域，并使
    $\mathscr T_{p+1}^U$ 右正合而 $\mathscr T_{p-1}^U$ 左正合。
\end{enumerate}
\end{proposition}

\begin{proof}
计及对 $\mathscr T_p$ 在 $Y$ 为仿射时的解释，这不过是（7.3.3）中一部分
内容的改述。
\end{proof}

\begin{proposition}[7.8.4]
设 $Y$ 是局部 Noether 预概形，$f:X\to Y$ 是固有态射，
$\mathscr P_\bullet$ 是由凝聚 $\mathscr O_X$-模组成且在 $Y$ 上平坦的
下有界复形，$\mathscr T_\bullet$ 是由（7.7.1.1）定义的函子。对每个
$y\in Y$，下列条件等价：
\begin{enumerate}
  \item[a)] $\mathscr P_\bullet$ 相对于 $Y$ 在点 $y$ 处、维数 $p$ 上
    同调平坦。
  \item[b)] 函子 $T_p^{\mathscr O_y}$ 正合。
  \item[c)] 存在整数 $n_0$，使得当 $n\geq n_0$ 时，
    \[
      \operatorname{long}T_p^{\mathscr O_y}
        (\mathscr O_y/\mathfrak m_y^{n+1})
      =
      \operatorname{long}T_p^{\mathscr O_y}(\kappa(y))
      \cdot
      \operatorname{long}(\mathscr O_y/\mathfrak m_y^{n+1})
      \tag{7.8.4.1}\label{III.7.8.4.1.fr}
    \]
    （这里的长度都是 $\mathscr O_y$-模的长度）。
  \item[d)] 存在一个开邻域 $U$（属于 $y$），使
    $(\mathscr H^{-p}(f,\mathscr P^\bullet))|U$ 是一个
    $(\mathscr O_Y|U)$-模且同构于 $(\mathscr O_Y|U)^m$，并且对每个拟凝聚
    $(\mathscr O_Y|U)$-模 $\mathscr M$，典范同态
    \[
      ((\mathscr H^{-p}(f,\mathscr P^\bullet))|U)
      \otimes_{\mathscr O_Y|U}\mathscr M
      \longrightarrow
      \mathscr H^{-p}
      \bigl(
        f,(\mathscr P^\bullet|U)
          \otimes_{\mathscr O_Y|U}\mathscr M
      \bigr)
      \tag{7.8.4.2}\label{III.7.8.4.2.fr}
    \]
    是双射。
\end{enumerate}
当这些条件成立时，还具有下列性质：
\begin{enumerate}
  \item[e)] 存在 $y$ 的一个邻域，使函数 $z\mapsto d_p(z)$
    （定义见（7.7.5.1））在该邻域中为常数。
\end{enumerate}
此外，若 $Y$ 在点 $y$ 处是约化的（$0_{\mathrm I}$，4.1.4），则 e)
与其余条件等价。
\end{proposition}

\oldpage[III]{206}
\begin{proof}
事实上，条件 b) 等价于 $T_p^{\mathscr O_y}$ 与
$T_{p+1}^{\mathscr O_y}$ 都右正合（7.3.3）；于是 a) 与 b) 的等价性由
（7.7.10）和（7.8.3）得到。因为 $\mathscr O_y/\mathfrak m_y^{n+1}$
是 Artin 环，而
$T_p^{\mathscr O_y}(\mathscr O_y/\mathfrak m_y^{n+1})$ 与
$T_p^{\mathscr O_y}(\kappa(y))$ 都是有限型
$(\mathscr O_y/\mathfrak m_y^{n+1})$-模（7.7.4），因而长度有限，所以
b) 与 c) 的等价性再次由（7.7.10）和（7.3.5.7）得到。a) 蕴含 e)，
并且当 $\mathscr O_y$ 约化时二者等价，这是（7.6.9）的推论。最后，
由定义（7.8.1）和（7.7.5），a) 蕴含（7.8.4.2）是双射；另一方面，
由（7.3.3，f)），a) 蕴含
$(\mathscr H^{-p}(f,\mathscr P^\bullet))_y$ 是平坦 $\mathscr O_y$-模，
因而为自由模（$0_{\mathrm I}$，10.1.3），因为它由（7.7.4）为有限型
$\mathscr O_y$-模。
再因 $\mathscr H^{-p}(f,\mathscr P^\bullet)$ 是凝聚 $\mathscr O_Y$-模
（7.7.4），它在 $y$ 的某个邻域中局部自由（$0_{\mathrm I}$，5.2.7）。
反过来，由函子 $\mathscr T_p^U$ 的定义（7.7.2.2），d) 显然蕴含 a)。
\end{proof}

\begin{proposition}[7.8.5]
在（7.8.4）的假设下，下列条件等价：
\begin{enumerate}
  \item[a)] $\mathscr P_\bullet$ 相对于 $Y$ 在每个维数 $i\leq p$ 上
    同调平坦。
  \item[b)] 对 $i\leq p+1$，函子 $\mathscr T_i$ 右正合。
  \item[c)] 对 $i\geq-p$，$\mathscr O_Y$-模
    $\mathscr H^i(f,\mathscr P^\bullet)$ 局部自由。
\end{enumerate}
\end{proposition}

\begin{proof}
a) 与 b) 的等价性由（7.8.3）立即得到，而 a) 由（7.8.4）蕴含 c)。
反过来，设 c) 成立；还应注意 $\mathscr L_i=0$ 对 $i<i_0$ 成立
（7.7.4），因而 $\mathscr T_i=0$ 对 $i\leq i_0$ 成立。因此，对每一点
$y\in Y$，都有一个仿射邻域 $U=\operatorname{Spec}(A)$，使
$T_i^A(A)$ 对 $A$ 为自由模，只要 $i\leq p$。由（7.3.7）遂知，
$\mathscr T_i^U=T_i^A$ 正合，只要 $i\leq p$。
\end{proof}

我们主要把上同调平坦性判据应用于下列情形：复形
$\mathscr P_\bullet$ 仅由\emph{一个}非零项组成，这一项是凝聚
$\mathscr O_X$-模 $\mathscr F$，且在 $Y$ 上\emph{平坦}，并取作
$\mathscr P_0$。回忆此时
$\mathscr T_p(\mathscr M)=\mathrm R^{-p}f_*
(\mathscr F\otimes_{\mathscr O_Y}\mathscr M)$。

\begin{proposition}[7.8.6]
设 $Y$ 是局部 Noether 预概形，$f:X\to Y$ 是固有平坦态射，$y$ 是
$Y$ 的一点；用 $X_y$ 表示纤维
$f^{-1}(y)=X\otimes_Y\kappa(y)$。假设
$\Gamma(X_y,\mathscr O_{X_y})=R$ 是一个可分 $\kappa(y)$-代数
（Bourbaki，\emph{Alg.}，第 VIII 章，\S~7，第 5 号），亦即有限多个
$\kappa(y)$ 的有限次可分扩域的直和。则 $\mathscr O_X$ 相对于 $Y$
在点 $y$ 处、维数 0 上上同调平坦。
\end{proposition}

\begin{proof}
由（7.8.4），可归约到 $Y$ 是局部环 $A=\mathscr O_y$ 的谱这一情形。
$f$ 平坦这一假设蕴含 $\mathscr T_{-1}=0$，所以 $T_0^A$ 已经左正合，
只须证明它右正合；由（7.7.10，c)），甚至可进一步归约到
$A=\mathscr O_y$ 为 Artin 环的情形。设 $k'$ 是 $\kappa(y)$ 的一个
使 $R$ 分裂的有限扩域，从而
$R\otimes_{\kappa(y)}k'$ 是有限多个与 $k'$ 同构的域的直和。已知存在
从 $A$ 到某个局部环 $A'$ 的局部同态，使 $A'$ 是 $A$-代数且作为
$A$-模有限自由，
且 $A'$ 的剩余域同构于 $k'$（$0_{\mathrm I}$，10.3.2）。由（7.6.1），
只须证明 $T_0^{A'}$ 右正合；换言之，可以假设 $R$ 是 $m$ 个与
$\kappa(y)$ 同构的域的直和。现记下如下初等引理：

\begin{lemma}[7.8.6.1]
设 $Z$ 是一个局部赋环空间；$Z$ 连通，

\oldpage[III]{207}
当且仅当环 $\Gamma(Z,\mathscr O_Z)$ 不是两个非零环的积。
\end{lemma}

\begin{proof}
事实上，若 $Z$ 是两个互不相交的非空开集之并，显然
$\Gamma(Z,\mathscr O_Z)$ 同构于两个非零环
$\Gamma(Z_1,\mathscr O_Z)$ 与 $\Gamma(Z_2,\mathscr O_Z)$ 的积。
反过来，$\Gamma(Z,\mathscr O_Z)$ 是这样的积，等价于在
$\Gamma(Z,\mathscr O_Z)$ 中存在一个
既不等于 0 也不等于 1 的幂等元 $s$。对每个 $z\in Z$，芽 $s_z$
于是是局部环 $\mathscr O_z$ 中的幂等元，因而等于 0 或 1。所有 $z$
（满足 $s_z=0$）所成的集合显然是开集；另一方面，若 $s_z=1$，则按定义
$s(z)\ne0$，所以所有 $z$（满足 $s_z=1$）所成的集合也是开集
（$0_{\mathrm I}$，5.5.2）；结论随即得到。
\end{proof}

由该引理，$X_y$ 恰有 $m$ 个连通分支 $X'_i$，并且
$\Gamma(X'_i,\mathscr O_{X'_i})=\kappa(y)$ 对每个 $i$ 都成立。因为
已经假设 $A$ 是 Artin
局部环，它的谱仅由一点组成，所以 $X$ 与 $X_y$ 具有相同的底空间；
因而 $X$ 有 $m$ 个连通分支 $X_i$，满足
$X'_i=X_i\otimes_Y\kappa(y)$。于是最终归约到 $R=\kappa(y)$ 的情形。
由（7.7.10，b)），只须证明典范同态
\[
  \Gamma(X,\mathscr O_X)
  \longrightarrow
  \Gamma(X_y,\mathscr O_{X_y})
\]
为满射；但这是显然的，因为复合同态
\[
  \Gamma(Y,\mathscr O_Y)=A
  \longrightarrow
  \Gamma(X,\mathscr O_X)
  \longrightarrow
  \Gamma(X_y,\mathscr O_{X_y})=\kappa(y)
\]
已经是满射。
\end{proof}

\begin{corollary}[7.8.7]
在（7.8.6）的假设下，存在一个开邻域 $U$（属于 $y$），使得：
\begin{enumerate}
  \item[(i)] $f_*(\mathscr O_X)|U$ 是一个 $(\mathscr O_Y|U)$-模，
    并同构于 $(\mathscr O_Y|U)^m$；
  \item[(ii)] 对每个 $z\in U$，典范同态
    \[
      (f_*(\mathscr O_X))_z
      \otimes_{\mathscr O_z}\kappa(z)
      \longrightarrow
      \Gamma(X_z,\mathscr O_{X_z})
    \]
    是双射。
\end{enumerate}
\end{corollary}

\begin{proof}
(i) 由（7.8.6）与（7.8.4）得到。

(ii) 由 $\mathscr T_0^U$ 的正合性（适当选取 $U$）以及（7.7.5.3）得到。
\end{proof}

\begin{corollary}[7.8.8]
设（7.8.6）的条件成立，并且还有
$\Gamma(X_y,\mathscr O_{X_y})=\kappa(y)$。则存在一个开邻域
$U$（属于 $y$），使典范同态
$\mathscr O_Y|U\to f_*(\mathscr O_X)|U$ 是双射。
\end{corollary}

\begin{proof}
事实上，由（7.8.7，(ii)）可知，该推论 (i) 中出现的整数 $m$ 必须等于 1。
\end{proof}

\begin{corollary}[7.8.9]
在（7.8.6）的假设下，存在一个开邻域 $U$（属于 $y$）、一个由唯一同构确定的
凝聚 $\mathscr O_U$-模 $\mathscr Q$，以及关于拟凝聚 $\mathscr O_U$-模
$\mathscr M$ 的函子同构
\[
  \mathrm R^1f_*(f^*(\mathscr M))
  \xrightarrow{\sim}
  \mathscr{H}\!om_{\mathscr O_U}(\mathscr Q,\mathscr M).
  \tag{7.8.9.1}\label{III.7.8.9.1.fr}
\]
\end{corollary}

\begin{proof}
事实上，该假设蕴含 $\mathscr T_0^U$ 正合，其中 $U$ 适当选取。因此只须
把（7.7.5，a)）与（7.7.5，b$'$)）的等价性应用于 $p=0$ 的情形，
并取 $\mathscr P_\bullet$ 为仅由次数 0 项 $\mathscr O_X$ 组成的复形。
\end{proof}

\begin{remarks}[7.8.10]
\textnormal{(i)} 在（7.8.6）的条件下，考虑 $f$ 的 Stein 分解（4.3.3）
\[
\xymatrix{
  X \ar[r]^{f'} & Y' \ar[r]^{g} & Y .
}
\]

\oldpage[III]{208}
其中 $Y'=\operatorname{Spec}(f_*(\mathscr O_X))$；此时有限态射 $g$
满足 $g_*(\mathscr O_{Y'})=f_*(\mathscr O_X)$ 在 $y$ 的邻域中局部自由，
并且它在 $y$ 上的纤维是一个 $\kappa(y)$-可分代数的谱（II，1.5.1）。
我们将在第 IV 章由此推出：存在开邻域 $U$，它是 $y$ 在 $Y$ 中的邻域，
使得对 $g^{-1}(U)\to U$ 这一 $g$ 的限制，\emph{每个}纤维
$g^{-1}(z)$（$z\in U$）
都是一个 $\kappa(z)$-可分代数的谱（我们将把它称为 $U$ 的一个
\emph{平展覆盖}）。于是由（7.8.7，(ii)）可知，（7.8.6）中对点 $y$
所作的假设，在 $y$ 的某个邻域的每一点也都成立。

\medskip
\textnormal{(ii)} 我们将在第 V 章看到，即使 $X$ 在 $Y$ 上是射影的
（甚至还在 $Y$ 上是“单纯的”，这一性质将在第 IV 章定义），（7.8.9）
中的 $\mathscr O_U$-模 $\mathscr Q$ 也未必局部自由；换言之，在这些条件下，
$\mathscr O_X$ 相对于 $Y$ 在点 $y$ 处未必在\emph{维数 1} 上上同调平坦。
在第 V 章中，我们将把 $\mathscr Q$ 解释为 $X$ 相对于 $Y$ 的 Picard
概形沿单位截面的 1-微分层。
\end{remarks}

\subsection{固有态射的应用：III. Euler--Poincaré 特征与 Hilbert
多项式的不变性}
\label{subsection:III.7.9.fr}
\label{III.7.9.fr}

\begin{env}[7.9.1]
\label{III.7.9.1.fr}
设 $A$ 是环，$M$ 是有限型射影 $A$-模；回忆（Bourbaki，
\emph{Alg. comm.}，第 II 章，\S~5，第 2 号），这等价于说，相伴的
$\mathscr O_X$-模 $\widetilde M$ 在
$X=\operatorname{Spec}(A)$ 上是\emph{有限型局部自由模}。对每个
$\mathfrak p\in\operatorname{Spec}(A)$，称
\emph{$M$ 在 $\mathfrak p$ 处的秩}并记作
$\operatorname{rang}_{\mathfrak p}(M)$ 的数，是自由
$A_{\mathfrak p}$-模 $M_{\mathfrak p}$ 的秩（也就是在
$\mathfrak p$ 处，局部自由 $\mathscr O_X$-模 $\widetilde M$ 的秩）。
于是
\[
  \operatorname{rang}_{\mathfrak p}M
  =\operatorname{rang}_{\mathfrak p}(M_{\mathfrak p})
  =\operatorname{rang}_{\kappa(\mathfrak p)}
  \bigl(M\otimes_A\kappa(\mathfrak p)\bigr).
  \tag{7.9.1.1}\label{III.7.9.1.1.fr}
\]
\end{env}

\begin{proposition}[7.9.2]
\label{III.7.9.2.fr}
设 $P_\bullet$ 是由有限型射影 $A$-模组成的有限复形，并对每个
$A$-模 $M$ 置
$T_\bullet(M)=H_\bullet(P_\bullet\otimes_A M)$。则对每个
$\mathfrak p\in\operatorname{Spec}(A)$，都有
\[
  \sum_i(-1)^i\operatorname{rang}_{\kappa(\mathfrak p)}
  T_i\bigl(\kappa(\mathfrak p)\bigr)
  =\sum_i(-1)^i\operatorname{rang}_{\mathfrak p}(P_i).
  \tag{7.9.2.1}\label{III.7.9.2.1.fr}
\]
\end{proposition}

\begin{proof}
事实上，按定义
$T_i(\kappa(\mathfrak p))=H_i(P_\bullet\otimes_A\kappa(\mathfrak p))$；
计及（7.9.1.1），公式（7.9.2.1）无非是有限维向量空间有限复形的
Euler--Poincaré 特征在取同调时的不变性（$0_{\mathrm I}$，11.10.2）。
\SourceCorrectionNote{法文印本及外交式来源在此误引公式（7.9.1.2）；
当前校勘法文改引实际给出该等式的（7.9.2.1），本译采用这一更正读法。}
\end{proof}

\begin{corollary}[7.9.3]
\label{III.7.9.3.fr}
函数
\[
  \mathfrak p\longmapsto
  \sum_i(-1)^i\operatorname{rang}_{\kappa(\mathfrak p)}
  T_i\bigl(\kappa(\mathfrak p)\bigr)
\]
在 $\operatorname{Spec}(A)$ 上局部常值。
\end{corollary}

\begin{theorem}[7.9.4]
\label{III.7.9.4.fr}
设 $Y$ 是局部 Noether 预概形，$f:X\to Y$ 是固有态射，
$\mathscr P_\bullet$ 是由凝聚 $\mathscr O_X$-模组成的有限复形，
且这些模在 $Y$ 上平坦。若置
\[
  \mathscr T_\bullet(\mathscr M)
  =\mathscr H^{-\bullet}
  \bigl(f,\mathscr P^\bullet\otimes_{\mathscr O_Y}\mathscr M\bigr)
  \qquad\text{（参见（7.7.1.1））}
\]
则函数

\oldpage[III]{209}
\[
  y\longmapsto
  \sum_i(-1)^i\operatorname{rang}_{\kappa(y)}
  \mathscr T_i\bigl(\kappa(y)\bigr)
  \tag{7.9.4.1}\label{III.7.9.4.1.fr}
\]
在 $Y$ 上局部常值。
\end{theorem}

\begin{proof}
可以归约到 $Y=\operatorname{Spec}(A)$ 是 Noether 环 $A$ 的仿射谱这一情形。
因为复形 $\mathscr P_\bullet$ 是有限复形，由（7.7.12，(i)）知道
\[
  \mathscr T_p(\mathscr M)
  =\mathscr H_p(\mathscr L_\bullet\otimes_{\mathscr O_Y}\mathscr M),
  \qquad \mathscr L_\bullet=\widetilde{L_\bullet},
\]
其中 $L_\bullet$ 是由有限型射影 $A$-模组成的\emph{有限}复形。
定理由（7.9.3）随即得到。
\end{proof}

\begin{env}[7.9.5]
\label{III.7.9.5.fr}
在（7.9.4）的条件下，当 $Y$ \emph{连通}时，函数（7.9.4.1）
\emph{常值}。当 $Y$ 连通且非空时，把（7.9.4.1）唯一的整数值记作
\[
  \operatorname{EP}(f,\mathscr P_\bullet)
  \quad\text{或}\quad
  \operatorname{EP}(Y,\mathscr P_\bullet),
  \quad\text{或简称}\quad
  \operatorname{EP}(\mathscr P_\bullet)
\]
（若不致混淆），并称这个整数为 $\mathscr P_\bullet$ 相对于 $f$
（或相对于 $Y$）的\emph{Euler--Poincaré 特征}。一般情形下，也把
（7.9.4.1）的右端记作
\[
  \operatorname{EP}(f,\mathscr P_\bullet;y)
  \quad\text{或}\quad
  \operatorname{EP}(Y,\mathscr P_\bullet;y)
  \quad\text{或}\quad
  \operatorname{EP}(\mathscr P_\bullet;y)
\]
\end{env}

\begin{env}[7.9.6]
\label{III.7.9.6.fr}
在（7.9.4）关于 $X$、$Y$、$f$ 的假设下，设
\[
  0\longrightarrow\mathscr P'_\bullet
  \xrightarrow{u}\mathscr P_\bullet
  \xrightarrow{v}\mathscr P''_\bullet
  \longrightarrow0
\]
是由凝聚 $\mathscr O_X$-模所组成的有限复形的正合列，这些模在
$Y$ 上平坦，
其中同态 $u$ 与 $v$ 的次数分别为\emph{偶数} $2d$、$2d'$。
因为 $\mathscr T_\bullet$ 是同调函子（7.7.1），有同调正合列
\[
  \cdots\longrightarrow
  \mathscr T_i(\mathscr P'_\bullet,\kappa(y))
  \longrightarrow
  \mathscr T_{i+2d}(\mathscr P_\bullet,\kappa(y))
  \longrightarrow
  \mathscr T_{i+2d+2d'}(\mathscr P''_\bullet,\kappa(y))
  \longrightarrow
  \mathscr T_{i-1}(\mathscr P'_\bullet,\kappa(y))
  \longrightarrow\cdots
\]
而且其中只有有限多个非零项。写出这个复形的 Euler--Poincaré
特征为零（$0_{\mathrm I}$，11.10.1），立即得到
\[
  \operatorname{EP}(\mathscr P_\bullet;y)
  =\operatorname{EP}(\mathscr P'_\bullet;y)
  +\operatorname{EP}(\mathscr P''_\bullet;y)
  \tag{7.9.6.1}\label{III.7.9.6.1.fr}
\]
对每个 $y\in Y$ 都成立。再如，若
$\mathscr P_\bullet=(\mathscr P_i)$ 且 $\mathscr P_i=0$ 对 $i<0$ 成立，
则有复形正合列
\[
\xymatrix@C=2.4em{
  \cdots \ar[r] & 0 \ar[r]\ar[d] & 0 \ar[r]\ar[d]
    & 0 \ar[r]\ar[d] & 0 \ar[r] & \cdots \\
  \cdots \ar[r] & 0 \ar[r]\ar[d] & 0 \ar[r]\ar[d]
    & \mathscr P_1 \ar[r]\ar[d] & \mathscr P_2 \ar[r] & \cdots \\
  \cdots \ar[r] & 0 \ar[r]\ar[d] & \mathscr P_0 \ar[r]\ar[d]
    & \mathscr P_1 \ar[r]\ar[d] & \mathscr P_2 \ar[r] & \cdots \\
  \cdots \ar[r] & 0 \ar[r]\ar[d] & \mathscr P_0 \ar[r]\ar[d]
    & 0 \ar[r]\ar[d] & 0 \ar[r] & \cdots \\
  \cdots \ar[r] & 0 \ar[r] & 0 \ar[r]
    & 0 \ar[r] & 0 \ar[r] & \cdots
}
\]
其中非零的竖直箭头都是恒等自同构；把（7.9.6.1）应用于这个正合列，
再对 $\mathscr P_\bullet$ 的长度作归纳，得到公式
\[
  \operatorname{EP}(\mathscr P_\bullet;y)
  =\sum_i(-1)^i\operatorname{EP}(\mathscr P_i;y).
  \tag{7.9.6.2}\label{III.7.9.6.2.fr}
\]

\oldpage[III]{210}
这里，对每个凝聚 $\mathscr O_X$-模 $\mathscr F$，若它在 $Y$ 上平坦，
以 $\operatorname{EP}(\mathscr F;y)$（或
$\operatorname{EP}(f,\mathscr F;y)$，或
$\operatorname{EP}(Y,\mathscr F;y)$）表示
$\operatorname{EP}(\mathscr Q_\bullet;y)$，其中复形
$\mathscr Q_\bullet$ 唯一的项 $\ne0$ 位于次数 $0$ 且等于 $\mathscr F$。
所以，复形的 Euler--Poincaré 特征的研究可归约为仅有一个非零项的复形。
\end{env}

\begin{proposition}[7.9.7]
\label{III.7.9.7.fr}
在（7.9.4）的假设下，设 $Y'$ 是局部 Noether 预概形，
$g:Y'\to Y$ 是态射，
\[
  X'=X\times_Y Y',
  \qquad f'=f_{(Y')}:X'\longrightarrow Y',
\]
而 $\mathscr P'_\bullet$ 是有限复形
\[
  \mathscr P_\bullet\otimes_{\mathscr O_Y}\mathscr O_{Y'}
\]
是 $\mathscr O_{X'}$-模的有限复形；$\mathscr P'_\bullet$ 由凝聚
$\mathscr O_{X'}$-模组成，且这些模在 $Y'$ 上平坦。并且对每个
$y'\in Y'$，
\[
  \operatorname{EP}(\mathscr P'_\bullet;y')
  =\operatorname{EP}(\mathscr P_\bullet;g(y')).
  \tag{7.9.7.1}\label{III.7.9.7.1.fr}
\]
\end{proposition}

\begin{proof}
$\mathscr O_{X'}$-模 $\mathscr P'_i$ 是 $\mathscr P_i$ 在投影
$X'\to X$ 下的逆像，故为凝聚模；它们在 $Y'$ 上平坦，这由
（$0_{\mathrm I}$，6.2.1）与（I，4.14.5）得到，这里问题在
$X$、$Y$、$Y'$ 上都是局部的；
最后，已知 $f'$ 固有（II，5.4.2），故（7.9.7.1）的左端有定义。
把问题如常归约到 $Y$ 与 $Y'$ 都仿射的情形，公式（7.9.7.1）便由
（6.10.4.2）、（7.7.2）及引理（7.6.7）得到。
\end{proof}

\begin{proposition}[7.9.8]
\label{III.7.9.8.fr}
假设（7.9.4）的条件成立，并且存在整数 $i_0$，使
\[
  \mathscr T_i(\kappa(y))=0
  \qquad\text{对 $i\ne i_0$ 及每个 $y\in Y$。}
\]
则
\[
  \mathscr T_{i_0}(\mathscr O_Y)
  =\mathscr H^{-i_0}(f,\mathscr P_\bullet)
\]
是局部自由 $\mathscr O_Y$-模，其在 $y\in Y$ 处的秩等于
$(-1)^{i_0}\operatorname{EP}(f,\mathscr P_\bullet;y)$。
\end{proposition}

\begin{proof}
先注意，$\mathscr T_j^{\mathscr O_y}$ 满足（7.4.4）的假设，因而
（7.4.7）可以应用；由假设与（7.5.3），
$\mathscr T_i^{\mathscr O_y}$ 对 $i\ne i_0$ 为\emph{零函子}。
于是由（7.3.3），$\mathscr T_{i_0}$ 也正合；继而由（7.8.4），
$\mathscr H^{-i_0}(f,\mathscr P_\bullet)$ 局部自由，且它在
$y\in Y$ 处的秩按定义为
\[
  \operatorname{rang}_{\kappa(y)}
  \mathscr T_{i_0}(\kappa(y))
  =(-1)^{i_0}\operatorname{EP}(f,\mathscr P_\bullet;y)
\]
因为 $\mathscr T_i(\kappa(y))=0$ 对 $i\ne i_0$ 成立。
\SourceCorrectionNote{法文印本及外交式来源在上式右端漏写因子
$(-1)^{i_0}$；命题陈述及 Euler--Poincaré 特征的定义都要求该因子，
当前校勘法文已补入，本译采用这一更正读法。}
\end{proof}

\begin{corollary}[7.9.9]
\label{III.7.9.9.fr}
设 $Y$ 是局部 Noether 预概形，$f:X\to Y$ 是固有态射，
$\mathscr F$ 是凝聚 $\mathscr O_X$-模且在 $Y$ 上平坦；假设存在整数
$i_0$，使
\[
  H^i\bigl(f^{-1}(y),
  \mathscr F\otimes_{\mathscr O_Y}\kappa(y)\bigr)=0
  \qquad\text{对每个 $i\ne i_0$ 及每个 $y\in Y$。}
\]
则 $\mathrm R^{i_0}f_*(\mathscr F)$ 是局部自由 $\mathscr O_Y$-模，
其在 $y$ 处的秩等于 $(-1)^{i_0}\operatorname{EP}(f,\mathscr F;y)$。
\end{corollary}

特别地：

\begin{corollary}[7.9.10]
\label{III.7.9.10.fr}
在（7.9.9）关于 $X$、$Y$、$\mathscr F$ 的先决条件下，假设
$\mathrm R^if_*(\mathscr F)=0$ 对所有 $i>0$ 成立。则
$f_*(\mathscr F)$ 是局部自由 $\mathscr O_Y$-模，其在 $y$ 处的秩等于
$\operatorname{EP}(f,\mathscr F;y)$。
\end{corollary}

\begin{proof}
由（7.9.9），只须证明下列引理：
\end{proof}

\begin{lemma}[7.9.10.1]
\label{III.7.9.10.1.fr}
在（7.9.10）的假设下，
\[
  H^i\bigl(f^{-1}(y),
  \mathscr F\otimes_{\mathscr O_Y}\kappa(y)\bigr)=0
\]
对所有 $i>0$ 及每个 $y\in Y$ 成立。
\end{lemma}

\begin{proof}
可以归约到 $Y=\operatorname{Spec}(A)$ 仿射的情形。沿用（7.9.4）的记号，
并令 $\mathscr P_\bullet$ 仅有次数 $0$ 的一项 $\mathscr F$，由假设
$\mathscr T_p(\mathscr O_Y)=0$ 对 $p<0$ 成立；由（7.3.7）可知
$\mathscr T_p$ 对 $p<0$ 正合，于是引理由（7.7.5，a)）与
（7.7.5，d)）的等价性得到。
\end{proof}

\oldpage[III]{211}
\begin{proposition}[7.9.11]
\label{III.7.9.11.fr}
沿用（7.9.4）的假设，设 $\mathscr L$ 是可逆 $\mathscr O_X$-模，
且相对于 $Y$ 非常丰沛；置
$\mathscr P_\bullet(n)=\mathscr P_\bullet\otimes_{\mathscr O_X}
\mathscr L^{\otimes n}$，其中 $n\in\mathbf Z$。则对每个 $y\in Y$，函数
\[
  n\longmapsto\operatorname{EP}(f,\mathscr P_\bullet(n);y)
  \tag{7.9.11.1}\label{III.7.9.11.1.fr}
\]
是系数属于 $\mathbf Q$ 的多项式，并且在 $Y$ 的同一连通分支的所有点上
都相同。
\end{proposition}

\begin{proof}
显然，$\mathscr P_\bullet(n)$ 是由 $\mathscr O_X$-模组成且在 $Y$
上平坦的复形。由（7.9.6.2），可归约到
$\mathscr P_\bullet$ 仅有次数 0 的一个非零项 $\mathscr F\ne0$；又因问题在
$Y$ 上是局部的，可假设 $Y$ 仿射且 $f$ 射影（II，5.5.3）。置
$X_y=f^{-1}(y)$，并令
$\mathscr L_y=\mathscr L\otimes_{\mathscr O_Y}\kappa(y)$，这是非常丰沛的
$\mathscr O_{X_y}$-模（II，4.4.10）。由（7.7.2），对于函子
$\mathscr T_\bullet$，即相对于复形 $\mathscr P_\bullet(n)$ 的函子，有
\[
  \mathscr T_i(\kappa(y))
  =H^{-i}(X_y,\mathscr F_y\otimes\mathscr L_y^{\otimes n})
  \qquad
  \bigl(\text{其中 }\mathscr F_y=\mathscr F\otimes_{\mathscr O_Y}\kappa(y)\bigr) ;
\]
所以 $\operatorname{EP}(f,\mathscr F(n);y)$ 正是（2.5.1）中定义的
Euler--Poincaré 特征 $\chi_{\kappa(y)}(\mathscr F_y(n))$；
（7.9.11.1）是多项式由（2.5.3）得到；此外，对每个 $n$，其值在
$Y$ 的一个连通分支上常值（7.9.4），证明完毕。
\end{proof}

\begin{env}[7.9.12]
\label{III.7.9.12.fr}
把有理系数多项式（7.9.11.1）记作
$\operatorname{PH}(f,\mathscr P_\bullet;y)$ 或
$\operatorname{PH}(\mathscr P_\bullet;y)$，并称它为
\emph{在 $y$ 处，相对于 $\mathscr P_\bullet$、$f$ 与 $\mathscr L$
的 Hilbert 多项式}（或简称\emph{在 $y$ 处 $\mathscr P_\bullet$
的 Hilbert 多项式}，或 $f$ 的 Hilbert 多项式，若不致混淆）。
当 $Y$ 连通且非空时，记号与术语中省略 $y$。所得不变量将在第 V 章
给定凝聚层的凝聚商层之“模空间”理论中起本质作用。

沿用（7.9.6）与（7.9.11）的记号，有
\[
  \operatorname{PH}(\mathscr P_\bullet;y)
  =\operatorname{PH}(\mathscr P'_\bullet;y)
  +\operatorname{PH}(\mathscr P''_\bullet;y)
  \tag{7.9.12.1}\label{III.7.9.12.1.fr}
\]
特别地，
\[
  \operatorname{PH}(\mathscr P_\bullet;y)
  =\sum_i(-1)^i\operatorname{PH}(\mathscr P_i;y) ;
  \tag{7.9.12.2}\label{III.7.9.12.2.fr}
\]
这显然由（7.9.6.1）与（7.9.6.2）得到。同样，沿用（7.9.7）的记号和
假设，有
\[
  \operatorname{PH}(\mathscr P'_\bullet;y')
  =\operatorname{PH}(\mathscr P_\bullet;g(y')).
  \tag{7.9.12.3}\label{III.7.9.12.3.fr}
\]
\end{env}

公式（7.9.12.2）把复形的 Hilbert 多项式的研究归约为单个
$\mathscr O_X$-模的 Hilbert 多项式，其中该模在 $Y$ 上平坦。后者有一个
不依赖同调考虑的
显著解释：

\begin{corollary}[7.9.13]
\label{III.7.9.13.fr}
设 $Y$ 是 Noether 预概形，$f:X\to Y$ 是固有态射，
$\mathscr L$ 是可逆 $\mathscr O_X$-模且相对于 $Y$ 非常丰沛，
$\mathscr F$ 是凝聚 $\mathscr O_X$-模且在 $Y$ 上平坦。存在整数
$n_0$，使得当 $n\geq n_0$ 时，$f_*(\mathscr F(n))$ 是局部自由
$\mathscr O_Y$-模，其在 $y\in Y$ 处的秩等于
$\operatorname{PH}(f,\mathscr F;y)(n)$。
\end{corollary}

\begin{proof}
因为态射 $f$ 射影（II，5.5.3），存在 $n_0$，使得当 $n\geq n_0$

\oldpage[III]{212}
时，有 $\mathrm R^if_*(\mathscr F(n))=0$ 对每个 $i>0$ 成立（2.2.1）；
结论因而由（7.9.10）得到。
\end{proof}

下列平坦性判据将在第 V 章凝聚层的“模空间”理论中十分重要：

\begin{proposition}[7.9.14]
\label{III.7.9.14.fr}
设 $Y$ 是 Noether 预概形，$f:X\to Y$ 是射影态射，$\mathscr L$
是可逆 $\mathscr O_X$-模且相对于 $f$ 丰沛；置
$\mathscr F(n)=\mathscr F\otimes\mathscr L^{\otimes n}$，其中任取
$\mathscr O_X$-模 $\mathscr F$ 与 $n\in\mathbf Z$。凝聚
$\mathscr O_X$-模 $\mathscr F$ 在 $Y$ 上平坦，当且仅当存在整数 $n_0$，
使得对每个 $n\geq n_0$，$f_*(\mathscr F(n))$ 都是局部自由
$\mathscr O_Y$-模。
\end{proposition}

\begin{proof}
条件的必要性与（7.9.13）中的证明相同（（2.2.1）的结论可应用于丰沛层
$\mathscr L$，因为 $f$ 射影）。为证充分性，可以归约到 $Y$ 是环
$A$ 的仿射谱这一情形；由假设与（2.2.2，(i)），诸 $A$-模
$\Gamma(X,\mathscr F(n))$ 是有限型射影模（Bourbaki，
\emph{Alg. comm.}，第 II 章，\S~5，第 2 号，定理 1）。设 $S$ 为分次环
\[
  S=\bigoplus_{n\geq0}\Gamma(X,\mathscr L^{\otimes n}) ;
\]
已知 $X$ 典范地等同于 $\operatorname{Proj}(S)$（II，4.5.2，(b) 与
5.4.4）。置
\[
  M=\bigoplus_{n\geq n_0}\Gamma(X,\mathscr F(n)) ;
\]
必要时把 $\mathscr L$ 替换为某个幂 $\mathscr L^{\otimes d}$，可以假设
$S$ 由有限多个次数 1 的元素生成（2.3.5.1）；由（II，2.7.5 与
2.7.2），于是 $\mathscr F$ 等同于 $\mathscr P\!roj_0(M)$。
对每个齐次元 $g\in S$，若其次数 $>0$，因而有
$\Gamma(X_g,\mathscr F)=M_{(g)}$；而 $M$ 是射影 $A$-模的直和，故为
平坦 $A$-模，所以 $M_g$ 也平坦（$0_{\mathrm I}$，6.3.2），从而
$M_{(g)}$ 也平坦，因为它是 $M_g$ 的直和因子（$0_{\mathrm I}$，
6.1.2）。由（I，4.14.5）可知 $\mathscr F$ 在 $Y$ 上平坦于
$X_g$ 的每一点；诸 $X_g$ 覆盖 $X$，命题得证。
\end{proof}

\bigskip
\noindent\emph{（待续。）}

\clearpage
\oldpage[III]{213}
\section*{记号索引}

\begin{flushleft}
$\mathbf{Tor}_n^A(P_\bullet,Q_\bullet)$：6.3.1。

\medskip
$\mathscr{T}\!or_n^{\mathscr O_S}(\mathscr P_\bullet,\mathscr Q_\bullet)$，
$\mathscr{T}\!or_n^S(\mathscr P_\bullet,\mathscr Q_\bullet)$：6.4.1 与 6.5.1。

\medskip
$\mathscr{T}\!or_n^S(\mathscr F,\mathscr G)$：6.5.1。

\medskip
$\mathscr{T}\!or_q^S(\mathscr P_\bullet^{(1)},\mathscr P_\bullet^{(2)},\ldots,
\mathscr P_\bullet^{(m)})$：6.5.15。

\medskip
$\mathbf H^{-n}(\mathfrak U^{(i)},\mathscr P_\bullet^{(i)})$，
$\mathbf H^{-n}(X^{(i)},\mathscr P_\bullet^{(i)})$：6.6.1。

\medskip
$\mathbf{Tor}_n^A(L_{\bullet\bullet}^{(1)},L_{\bullet\bullet}^{(2)})$，
$\mathbf{Tor}_n^S(\mathfrak U^{(1)},\mathfrak U^{(2)};
\mathscr P_\bullet^{(1)},\mathscr P_\bullet^{(2)})$：6.6.2。

\medskip
$\mathscr{T}\!or_n^S(\mathfrak U^{(1)},\mathfrak U^{(2)};
\mathscr F^{(1)},\mathscr F^{(2)})$，
$\mathscr{T}\!or_n^S(\mathfrak U^{(1)},\mathfrak U^{(2)};
\mathscr P_\bullet^{(1)},\mathscr P_\bullet^{(2)})$：6.6.2。

\medskip
$\mathscr{T}\!or_n^S(f_1,f_2;\mathscr P_\bullet^{(1)},
\mathscr P_\bullet^{(2)})$：6.6.8、6.7.1 与 6.7.3。

\medskip
$\mathscr{T}\!or_n^S((f_i)_{i\in I};(\mathscr P_\bullet^{(i)})_{i\in I})$，
$\mathscr{T}\!or_n^S(f_1,\ldots,f_m;\mathscr P_\bullet^{(1)},\ldots,
\mathscr P_\bullet^{(m)})$：6.7.11。

\medskip
$\mathbf{Ab}$，$\mathbf{Ab}_A$：7.1.1。

\medskip
$T_{(B)}$，$T^{(B)}$，$T\otimes_A B$
（$T$ 是从 $\mathbf{Ab}_A$ 到 $\mathbf{Ab}$ 的函子）：7.1.3。

\medskip
$t_M$：7.2.2。

\medskip
$T_{\mathfrak p}$：7.1.4。

\medskip
$\mathscr T_\bullet^{Y'}(\mathscr M')$，
$\mathscr T_\bullet^U(\mathscr M|U)$：7.7.2。

\medskip
$\operatorname{EP}(f,\mathscr P_\bullet;y)$，
$\operatorname{EP}(Y,\mathscr P_\bullet;y)$，
$\operatorname{EP}(\mathscr P_\bullet;y)$：7.9.5。

\medskip
$\operatorname{PH}(f,\mathscr P_\bullet;y)$，
$\operatorname{PH}(\mathscr P_\bullet;y)$：7.9.12。
\end{flushleft}

\clearpage
\oldpage[III]{214}
\section*{术语索引}

\begin{flushleft}
模复形的 Euler--Poincaré 特征：7.9.5。

\medskip
在点 $y\in Y$ 处、维数 $p$ 上上同调平坦；在 $Y$ 上、维数 $p$
上上同调平坦；在 $Y$ 上于点 $y$ 处上同调平坦；在 $Y$ 上上同调平坦
（由 $\mathscr O_X$-模组成且这些模在 $Y$ 上平坦的复形）：7.8.1。

\medskip
同伦复形：6.1.4。

\medskip
上同调维数 $\leq n$ 的赋环空间：6.5.5。

\medskip
加性协变函子的标量扩张：7.1.3。

\medskip
上同调维数 $\leq n$ 的环层：6.5.5。

\medskip
Künneth 公式：6.7.8。

\medskip
在点 $y\in Y$ 处、维数 $p$ 上同调平坦；在 $Y$ 上、维数 $p$
上同调平坦；在 $Y$ 上于点 $y$ 处同调平坦；在 $Y$ 上同调平坦
（由 $\mathscr O_X$-模组成且这些模在 $Y$ 上平坦的复形）：7.8.1。

\medskip
同伦：6.1.4。

\medskip
两个 $A$-模复形的超 Tor：6.3.1。

\medskip
两个模复形的局部超 Tor：6.4.1。

\medskip
两个模复形的整体超 Tor：6.6.2、6.6.8 与 6.7.3。

\medskip
模复形的相对 Hilbert 多项式：7.9.12。

\medskip
Künneth 谱序列：6.7.3。
\end{flushleft}

\clearpage
\oldpage[III]{215}
\section*{目录}

\begin{flushleft}
\textsc{第三章}——\textbf{凝聚层的上同调研究（续）} \dotfill 5

\medskip
\S~6. 局部与整体 Tor 函子；Künneth 公式 \dotfill 5

6.1. 引言 \dotfill 5

6.2. 预概形上模复形的超上同调 \dotfill 6

6.3. 两个模复形的超 Tor \dotfill 9

6.4. 拟凝聚模复形的局部超 Tor 函子；仿射概形的情形 \dotfill 14

6.5. 拟凝聚模复形的局部超 Tor 函子：一般情形 \dotfill 16

6.6. 拟凝聚模复形的整体超 Tor 函子与 Künneth 谱序列：
仿射基的情形 \dotfill 21

6.7. 拟凝聚模复形的整体超 Tor 函子与 Künneth 谱序列：
一般情形 \dotfill 25

6.8. 整体超 Tor 的结合性谱序列 \dotfill 32

6.9. 整体超 Tor 中的基变换谱序列 \dotfill 34

6.10. 某些上同调函子的局部结构 \dotfill 39

\medskip
\S~7. 模的协变同调函子中的基变换研究 \dotfill 43

7.1. $A$-模函子 \dotfill 43

7.2. 张量积函子的刻画 \dotfill 44

7.3. 模的同调函子的正合性判据 \dotfill 48

7.4. 函子 $\mathrm H_i(P\otimes_A M)$ 的正合性判据 \dotfill 53

7.5. Noether 局部环的情形 \dotfill 58

7.6. 正合性质的下降、半连续性定理与 Grauert 正合性判据 \dotfill 60

7.7. 固有态射的应用：I. 交换性质 \dotfill 65

7.8. 固有态射的应用：II. 上同调平坦性判据 \dotfill 72

7.9. 固有态射的应用：III. Euler--Poincaré 特征与 Hilbert
多项式的不变性 \dotfill 76

\begin{flushright}
\emph{1962 年 12 月 15 日收到。}
\end{flushright}
\end{flushleft}

% 印刷第 216 页为空白页。
\clearpage
\thispagestyle{empty}
\null
\clearpage

\oldpage[III]{217}
\section*{勘误与增补}
\begin{center}
（表 2）

\medskip
A）\emph{排印错误}
\end{center}

\noindent\textbf{$(0_{\mathrm I},\,3.2.1)$} 第 27 页第 2 行，把 $X$
替换为 $U$。

\noindent\textbf{$(0_{\mathrm I},\,3.2.6)$} 第 27 页倒数第 1 行及
第 28 页第 1 行，把 $\mathfrak S_\lambda$（共 3 处）替换为
$\mathscr H_\lambda$。

\noindent\textbf{$(0_{\mathrm I},\,3.4.5)$} 在 3.4.5) 前补一个左括号。

\noindent\textbf{$(0_{\mathrm I},\,4.2.1)$} 第 39 页第 10 行，把
$\psi_*(A)$ 替换为 $\psi_*(\mathscr A)$。

\noindent\textbf{$(0_{\mathrm I},\,5.1.3)$} 第 45 页第 16 行，把
$\mathscr F|V$ 替换为 $\mathscr F|U$。

\noindent\textbf{$(0_{\mathrm I},\,5.3.9)$} 第 48 页第 12 行，在
“使得”之前增补：在开邻域 $U$ 上方（$x$ 的邻域）。

\noindent\textbf{$(0_{\mathrm I},\,7.6.9)$} 第 73 页倒数第 3 行，把
$I_\lambda$ 替换为 $J_\lambda$。

\noindent\textbf{(I, 1.1.15)} 第 83 页第 3 行，把 $\mathfrak a$
替换为 $\mathfrak a\ne A$。

\noindent\textbf{(I, 1.4.1)} 第 90 页倒数第 2 行，把 $\widetilde N$
替换为 $N$。

\noindent\textbf{(I, 2.3.1)} 第 101 页倒数第 3 行，把 (0, 4.1.6)
替换为 (0, 4.1.7)。

\noindent\textbf{(I, 2.3.2)} 第 101 页第 11 行，把 $B$ 替换为 $B_s$，
把 $C$ 替换为 $C_t$。

\noindent\textbf{(I, 2.5.5)} 第 104 页第 19 行，把 $S$-态射替换为
$S$-预概形。

\noindent\textbf{(I, 3.3.7)} 第 109 页第 17 行，把 $Y_{(S)}$ 替换为
$Y_{(S')}$。

\noindent\textbf{(I, 3.4.5)} 第 113 页第 9 行，把 $\mathrm s$
替换为 $s$。

\noindent\textbf{(I, 3.4.8)} 第 114 页第 4 行，把 $p(x)$ 替换为
$f(x)$。

\noindent\textbf{(I, 3.5.1)} 第 115 页第 3 行，把 $1_X\times g$
替换为 $1_{X'}\times g$。

\noindent\textbf{(I, 3.7.2)} 第 119 页第 8 行，把第一个 $X$
替换为 $X'$。

\noindent\textbf{(I, 4.2.2)} 第 123 页第 8 行，在图的左侧竖直箭头中，
把 $\alpha_{\psi(y)}$ 替换为 $\varphi_{\psi(y)}$。

\noindent\textbf{(I, 4.4.3)} 第 126 页第 16 行，把
\emph{isomorphée} 替换为 \emph{isomorphes}。

\noindent\textbf{(I, 4.5.5)} 第 127 页倒数第 17 行，把 (4.2.4)
替换为 (4.2.5)。

\noindent\textbf{(I, 5.1.4)} 第 128 页倒数第 14 行，把 (2.1.7)
替换为 (2.1.8)。

\noindent\textbf{(I, 5.3.13)} 第 134 页第 20 行，把 (4.2.4)
替换为 (4.2.5)。

\noindent\textbf{(I, 6.4.2)} 第 147 页倒数第 1 行，把“覆盖”
替换为“有限覆盖”。

\noindent\textbf{(I, 9.1.13)} 第 171 页第 15 行，把
$p^{-1}(\mathscr F)$ 替换为 $p^*(\mathscr F)$。

\noindent\textbf{(I, 9.5.11)} 第 179 页第 7 行，把 $Y$ 替换为 $Y'$。

\noindent\textbf{(I, 9.6.5)} 第 180 页倒数第 17 行，把
$(0,\,4.1.4)$ 替换为 $(0,\,4.1.3)$。

\noindent\textbf{(I, 10.12.2)} 第 206 页倒数第 14 行，把
$(X_n)_{(S_n)}$ 替换为 $(X_n)_{(S_m)}$。

\clearpage
\oldpage[III]{218}
\noindent\textbf{(I)，术语索引：} 第 219 页倒数第 6 行，把
0, 4.1.4 替换为 0, 4.1.5。第 222 页倒数第 16、17、18 行，把
I, 2.1.7 替换为 I, 2.1.8；第 222 页倒数第 19 行，把 I, 2.1.8
替换为 I, 2.1.7。

\noindent\textbf{(II, 1.5.2)} 第 12 页第 17 行，在图的左侧竖直箭头旁
添上一个 $f$。

\noindent\textbf{(II, 4.2.7)} 第 75 页倒数第 16 行，把
(III, 2.1.14) 替换为 (III, 2.1.13)。

\noindent\textbf{(II, 5.2.1)} 第 97 页倒数第 16 行，把 $X-U$
替换为 $U$。第 97 页倒数第 8 行，把 $X'$ 替换为 $X_f$。

\noindent\textbf{(II, 5.3.6)} 第 100 页第 16 行，把 (4.6.18)
替换为 (4.6.17)。

\noindent\textbf{(II, 6.2.7)} 第 116 页倒数第 9 行，把第 V 章
替换为第 IV 章。

\noindent\textbf{(II, 6.6.5)} 第 132 页倒数第 15 行，把第 V 章
替换为第 IV 章。

\noindent\textbf{(II, 7.4.12)} 第 151 页第 1 行，把第 V 章
替换为第 III 章。

\noindent\textbf{(II, 8.1.4)} 第 154 页第 15 行，把 (III, 2.3.8)
替换为 (III, 2.3.7)。

\noindent\textbf{(II, 8.10.5)} 第 187 页倒数第 17 行，把
$g_{(j)x}$ 替换为 $g_{j(x)}$。

\noindent\textbf{$(0_{\mathrm{III}},\,10.2.7)$} 第 20 页第 1 行，
把 $u$-进替换为 $\mathfrak m$-进。第 20 页倒数第 13 行，把 $k$
替换为 $k_i$。

\noindent\textbf{$(0_{\mathrm{III}},\,11.1.1)$} 第 23 页第 18 行，把
$\operatorname{Im}(d_r^{p+r,q-r+1})$ 替换为
$\operatorname{Im}(d_r^{p-r,q+r-1})$。

\noindent\textbf{$(0_{\mathrm{III}},\,11.8.6)$} 第 45 页倒数第 14 行，
把箭头 $\to$ 替换为 $\rightsquigarrow$（共 4 处）。

\noindent\textbf{$(0_{\mathrm{III}},\,12.3.4)$} 第 61 页第 4 行，
把 $\rightsquigarrow$ 替换为 $\to$。

\noindent\textbf{$(0_{\mathrm{III}},\,13.2.4)$} 第 67 页倒数第 2 行，
把 $\varinjlim$ 替换为 $\varprojlim$（共 2 处）。

\noindent\textbf{(III, 1.1.7)} 第 85 页第 15 行，把 $A^r$
替换为 $\bigwedge(A^r)$（共 2 处）。

\noindent\textbf{(III, 1.2.4)} 第 87 页第 7 行，把
$\Gamma(U,\mathscr F)$ 替换为 $\Gamma(U,\mathscr O_X)$，并把
$\mathscr G$ 替换为 $\mathscr F$。第 87 页第 8 行，把
$\mathrm H^p(\mathfrak U,\mathscr G)$ 替换为
$\mathrm H^p(\mathfrak U,\mathscr F)$。

\noindent\textbf{(III, 2.5.3.1)} 第 111 页第 3 行，把 $n>0$
替换为 $n\geq0$。第 111 页第 8 行，把 $-r\leq n\leq0$
替换为 $-r\leq n<0$。

\noindent\textbf{(III, 3.1.2)} 第 116 页第 13 行，把 $\mathscr G$
替换为 $\mathscr G\in\mathbf K'$。

\noindent\textbf{(III, 4.3.6)} 第 133 页第 7 行，把
$K'=K\otimes_A\widehat A$ 替换为
$K'\supset K\otimes_A\widehat A$。第 133 页第 16 行，把 $R(X)$
替换为 $R(Y)$。

\noindent\textbf{(III, 4.4.12)} 第 138 页第 10 行，把第 V 章
替换为第 IV 章。

\noindent\textbf{(III, 4.6.6)} 第 142 页第 19 行，把“第 V 章”
替换为“在后面的一节中”。

\begin{center}
B）\emph{正文修改}\footnote{为便于检索参考文献，此后把属于第 N 章所附
勘误表的修改记作 $\mathbf{Err}_{\mathrm N}$ 后接一个编号。}
\end{center}

\noindent\textbf{$(\mathbf{Err}_{\mathrm{III}},\,1)$} 在
$(0_{\mathrm I},\,5.1.3)$ 第 45 页第 18 行，把“直和”替换为
“有限直和”。

\noindent\textbf{$(\mathbf{Err}_{\mathrm{III}},\,2)$} 在
$(0_{\mathrm I},\,5.4.3)$ 第 49 页倒数第 8 行至倒数第 4 行，把
“一般情形下……”起的文字替换为：

在一般情形下，设 $X$ 是赋环空间，并且

\oldpage[III]{219}
对每个点，$\mathscr O_x$ 都是\emph{局部}环，其中 $x\in X$。若 $\mathscr L$
是\emph{有限型} $\mathscr O_X$-模，并且存在 $\mathscr O_X$-模
$\mathscr F$，使 $\mathscr L\otimes_{\mathscr O_X}\mathscr F$
\emph{同构于} $\mathscr O_X$，则前面的论证表明，对每个 $x\in X$，
$\mathscr L_x$ 同构于 $\mathscr O_x$。由此推出 $\mathscr L$
\emph{可逆}。事实上，对每个 $x\in X$，取开邻域 $U$（属于 $x$），使
$\mathscr L|U$ 由 $n$ 个截面 $s_i\ (1\leq i\leq n)$ 在 $U$ 上
生成（5.2.1）。例如可假设 $s=s_1$ 满足 $s_x\ne0$；因为
$\mathscr L_x$ 同构于 $\mathscr O_x$，对每个 $i$，存在截面 $t_i$，
它是 $\mathscr O_X$ 在某个开邻域 $V_i\subset U$（属于 $x$）上的截面，
并使 $(t_i)_x s_x=(s_i)_x$。因而存在开邻域
$V\subset\bigcap_{i=1}^nV_i$（属于 $x$），使 $t_i s=s_i$ 在 $V$ 上成立；换言之，
$\mathscr L|V$ 由\emph{唯一}截面 $s|V$ 生成。再者，若 $z$ 是开集
$W\subset V$ 上同态
$\mathscr O_X|V\to\mathscr L|V$（5.1.1）的核的截面，其中该同态由
$s$ 定义；则 $z_y$ 零化 $\mathscr L_y$，对每个 $y\in W$ 都如此；
故由假设 $z_y=0$，从而
$z=0$；这就证明了断言。此外，考察张量积
$\mathscr L^{-1}\otimes\mathscr L\otimes\mathscr F$，立即可见
$\mathscr F$ 同构于 $\mathscr L^{-1}$。

\noindent\textbf{$(\mathbf{Err}_{\mathrm{III}},\,3)$} 在
$(0_{\mathrm I},\,7.2.4)$ 中，必须加上条件：诸幂
$\mathfrak J^n$ 是可容环 $A$ 中的\emph{闭}理想，命题才成立。
原证明不正确，修正后的陈述由 Bourbaki，\emph{Top. gén.}，第 III 章，
第 3 版，\S~3，第 5 号，命题 9 的推论 1 得到。也必须假设诸
$\mathfrak J^n$ 在 $A$ 中\emph{闭}，这适用于
$(0_{\mathrm I},\,7.2.5)$ 与 $(0_{\mathrm I},\,7.2.6)$ 的陈述。

\noindent\textbf{$(\mathbf{Err}_{\mathrm{III}},\,4)$} 在命题
(I, 2.2.5) 后增补：沿用 (I, 2.2.5) 的记号，若 $\mathfrak J$
是 $A$ 的理想，则用 $\widetilde{\mathfrak J\mathscr F}$ 表示
$\mathscr O_X$-子模 $\mathfrak J\widetilde{\mathscr F}$，后者定义于
$(0,\,4.3.5)$。

\noindent\textbf{$(\mathbf{Err}_{\mathrm{III}},\,5)$} 在
(I, 3.4.5) 中，第 112 页倒数第 1 行，把“\emph{域} $K$”替换为
“\emph{代数闭域} $K$”。第 113 页第 1 行，把“\emph{扩域}”
替换为“\emph{代数闭扩域}”。第 113 页第 1 至第 4 行，把
“$K$ 称为……”起的文字替换为：

对 $X$ 的每个取值于域 $K$ 的点，称 $K$ 为相应点的\emph{值域}；
若 $x$ 是该点的所在位置，也称该点\emph{位于} $x$。这样定义了映射
$X(K)\to X$，它把取值于 $K$ 的点送到其所在位置。

第 113 页第 7、16 行删除“几何的”；第 113 页第 10、13 行删除
“几何的”。第 115 页倒数第 10、11 行删除“\emph{几何的}”。

\noindent\textbf{$(\mathbf{Err}_{\mathrm{III}},\,6)$} 在
(I, 3.6.1) 证明末尾、第 117 页倒数第 12 行后增补：若置
$X'=X\times_Y\operatorname{Spec}(\mathscr O_y/\mathfrak a_y)$，则对
每个点 $x\in X'$（经 $p$ 等同于 $X$ 的一点），有
$\mathscr O_{X',x}=\mathscr O_{X,x}/\mathfrak a_y\mathscr O_{X,x}$。
事实上，问题在 $X$ 与 $Y$ 上是局部的，故可假设
$X=\operatorname{Spec}(B)$、$Y=\operatorname{Spec}(A)$；有
$(B/\mathfrak a_yB)_x=B_x/\mathfrak a_yB_x$，这由平坦性
$(0,\,1.3.2)$ 得到。

\noindent\textbf{$(\mathbf{Err}_{\mathrm{III}},\,7)$} 在
(I, 5.1.1) 第 127 页倒数第 3 行，把“拟凝聚 $\mathscr O_X$-模”
替换为“$\mathscr B$ 的\emph{拟凝聚理想}”。

\noindent\textbf{$(\mathbf{Err}_{\mathrm{III}},\,8)$} 在
(I, 5.1.10) 后、第 131 页第 16 行增补：

\emph{注（5.1.11）。}——稍微修改（5.1.9）的论证即可看出，结论在
\emph{不先验假设 $X$ 是预概形}的情形下仍成立，只须假设

\oldpage[III]{220}
$(X,\mathscr O_X)$ 是局部环赋环空间，$\mathscr J$ 是
$\mathscr O_X$ 的理想且 $\mathscr J^n=0$，赋环空间
$X_0=(X,\mathscr O_X/\mathscr J)$ 是仿射概形，并且理想
$\mathscr J^k/\mathscr J^{k+1}$ 是拟凝聚 $\mathscr O_{X_0}$-模。
在论证中用（1.8.1）（参见 $(\mathbf{Err}_{\mathrm{II}})$）
替代（2.2.4）即可。

\noindent\textbf{$(\mathbf{Err}_{\mathrm{III}},\,9)$} 在
(I, 5.3.1) 第 132 页第 11 行，在“或 $\Delta_X$”后插入
“或写作 $\Delta_\varphi$，如果 $\varphi:X\to S$ 是结构态射”。

\noindent\textbf{$(\mathbf{Err}_{\mathrm{III}},\,10)$} 在
(I, 5.3.9) 中，原证明不充分，因为它没有证明 $\Delta_X(X)$ 在
$X\times_SX$ 中局部闭。为得到正确证明，只须使用（4.2.4，$a$)）：
对每个 $x\in X$ 及每个仿射邻域 $U$（$x$ 在 $X$ 中的邻域），
$U\times_SU$ 是 $\Delta_X(x)$ 的仿射邻域；计及（5.3.16）
（其证明只使用定义（5.3.1.1）），问题归约为
$S=\operatorname{Spec}(B)$、$X=\operatorname{Spec}(A)$ 都是仿射概形
时证明（5.3.9）。由（5.3.1.1）显然可见，$\Delta_X$ 对应于典范同态
$A\otimes_BA\to A$，它把 $x\otimes y$ 送到 $xy$；该同态满射，
故 $\Delta_X$ 在此情形下是闭浸入（4.2.3），证明完毕。

\noindent\textbf{$(\mathbf{Err}_{\mathrm{III}},\,11)$} 在
(I, 7.1.15) 与 (I, 7.1.16) 中，第 158 页倒数第 7、15 行删除
“\emph{几何的}”。

\noindent\textbf{$(\mathbf{Err}_{\mathrm{III}},\,12)$} 把
(I, 7.4.7) 替换为：把（7.4.1）的定义（滥用语言地）扩张到下列情形：
$X$ 是\emph{既约}预概形，而且每个点都有只含\emph{有限}多个不可约分支
的开邻域。由（7.3.4）与（7.4.6）可知，对有限型拟凝聚
$\mathscr O_X$-模 $\mathscr F$，说 $\mathscr F$ 是\emph{挠层}，
等价于说 $\operatorname{Supp}(\mathscr F)$
\emph{不包含 $X$ 的任何不可约分支}。

\noindent\textbf{$(\mathbf{Err}_{\mathrm{III}},\,13)$} 在
(I, 10.11.7) 中，第 205 页第 17 至第 23 行，从“还须证明……”起的
证明末段由（10.11.6）已属多余，因为所考察的层由 $(0,\,5.3.5)$
是凝聚的。

\noindent\textbf{$(\mathbf{Err}_{\mathrm{III}},\,14)$} 在
(II, 1.7.8) 中，第 16 页倒数第 10 行后增补：当
$\mathscr E=\mathscr O_S^n$ 时，也用 $\mathbf V_S^n$ 代替
$\mathbf V(\mathscr E)$；若再有 $S=\operatorname{Spec}(A)$ 仿射，
则用 $\mathbf V_A^n$ 代替 $\mathbf V_S^n$；有
$\mathbf V_A^n=\operatorname{Spec}(A[T_1,\ldots,T_n])$，其中
$T_i$ 是不定元。

\noindent\textbf{$(\mathbf{Err}_{\mathrm{III}},\,15)$} 在
(II, 1.7.10) 中，第 17 页第 10 行删除“\emph{几何的}”；
第 17 页第 12 行及第 16--17 行删除“几何的”。

\noindent\textbf{$(\mathbf{Err}_{\mathrm{III}},\,16)$} 在
(II, 1.7.12) 中增补：同样置
\[
  \mathbf V(\mathscr O_S^n)=\mathbf V_S^n=S[T_1,\ldots,T_n].
\]

\noindent\textbf{$(\mathbf{Err}_{\mathrm{III}},\,17)$} 在
(II, 4.2.6) 中，第 75 页第 7 行删除“\emph{几何的}”；
第 75 页第 8 行删除“\emph{几何的}”。

\noindent\textbf{$(\mathbf{Err}_{\mathrm{III}},\,18)$} 在
(II, 4.6.13) 中，第 92 页第 7 行的论证需要说得更精确，因为沿用
（4.4.10）的记号，这里必须证明浸入 $\Gamma_f$ \emph{拟紧}，才能应用
(i bis)。由（4.6.4），为证明 $\mathscr L$ 相对于 $f$ 丰沛，可以归约到
$Y$ 仿射的情形。另一方面，注意在 (v) 的每组假设中，$f$ 都拟紧
（I，6.6.4）。此外，若 $g$ 分离，则

\oldpage[III]{221}
$\Gamma_f$ 是闭浸入（I，5.4.3），因而拟紧（I，6.6.4）。反之，若
$X$ 局部 Noether，则因 $Y$ 仿射而 $f$ 拟紧，$X$ 的底空间拟紧，
从而为 Noether 空间；可再次应用（I，6.6.4）证明 $\Gamma_f$ 拟紧。

\noindent\textbf{$(\mathbf{Err}_{\mathrm{III}},\,19)$}
(II, 5.2.2) 的陈述不正确：条件 $b)$、$c)$、$c')$ 并非关于
$Y$ 的局部条件，因为我们不知道，对开集 $U$（属于 $Y$），拟凝聚
$(\mathscr O_X|f^{-1}(U))$-模是否是在 $f^{-1}(U)$ 上的限制，而该限制
来自某个拟凝聚 $\mathscr O_X$-模。因此，在第 98 页第 16 行，把
“\emph{设 $f:X\to Y$ 是分离拟紧态射}”替换为：
“\emph{设 $X$、$Y$ 是两个预概形，并且 $X$ 是概形或者 $X$ 的
底空间局部 Noether；设 $f:X\to Y$ 是拟紧态射。}”在证明中应注意：
这些假设蕴含 $f$ \emph{分离}，如果假设 $X$ 是概形；这是由
（I，5.5.5 与 5.5.8）得到的；
故两种情形都可以应用（I，9.2.2）。

\noindent\textbf{$(\mathbf{Err}_{\mathrm{III}},\,20)$} 在
(II, 6.2.3) 中，第 115 页倒数第 18 行后增补：

\emph{称态射 $f:X\to Y$ 在点 $x\in X$ 处拟有限，如果存在仿射开邻域
$V$（属于 $y=f(x)$）及仿射开邻域 $U$（属于 $x$），使
$f(U)\subset V$，而态射 $U\to V$（即 $f$ 的限制）是拟有限态射。
称态射 $f:X\to Y$ 局部拟有限，如果它在 $X$ 的每一点处拟有限。}

\noindent\textbf{$(\mathbf{Err}_{\mathrm{III}},\,21)$} 在
(II, 6.4.3) 中，原证明不正确，因为有限型 $A$-子模（位于 $K$ 中）的元素
未必整于 $A$。把第 121 页最后 9 行替换为：可以假设 $E$ 无挠，
因而忠实；故 $E$ 也是忠实 $A[u]$-模（把 $A$ 嵌入 $E$ 的自同态环）。
因为 $E$ 是有限型 $A$-模，所以 $u$ \emph{整于} $A$
（Bourbaki，\emph{Alg. comm.}，第 V 章，\S~1，第 1 号，引理 1）。
立即推出 $u\otimes1$ 的特征值（在 $K$ 的一个代数闭包中）都是整于
$A$ 的元素，因而诸 $\sigma_i(u)$ 也是如此。

\noindent\textbf{$(\mathbf{Err}_{\mathrm{III}},\,22)$} 在
$(0_{\mathrm{III}},\,9.1.1)$ 中，第 12 页倒数第 18 行删除“开”。

\noindent\textbf{$(\mathbf{Err}_{\mathrm{III}},\,23)$} 在
$(0_{\mathrm{III}},\,11.7.3)$ 中，第 43 页第 10 行，在
$\mathsf C''$ 后增补：使得对每个射影对象 $P$（属于 $\mathsf C'$）
（相应地，对每个射影对象 $P'$（属于 $\mathsf C''$）），函子
$A'\rightsquigarrow T(P,A')$（相应地，
$A\rightsquigarrow T(A,P')$）在 $\mathsf C'$（相应地，
$\mathsf C$）中正合。

\noindent\textbf{$(\mathbf{Err}_{\mathrm{III}},\,24)$}\par
\noindent
$(0_{\mathrm{III}},\,13.7.7)$ 的陈述不准确，须作如下修改：
第 78 页第 6 行，删除“以及
$(\mathrm R^{n+1}T(A_k))_{k\in\mathbf Z}$”；第 78 页第 7 行，把
“$\mathrm R'^nT(A)$ 是有限型 $S$-模”替换为
“$\mathrm R'^nT(A)$ 与 $\mathrm R'^{n+1}T(A)$ 是有限型
$S$-模”；第 78 页第 12--13 行，删除“以及 $p+q=n+1$”。
在证明中，第 78 页第 23 行删除“以及对 $n+1$”。

\noindent\textbf{$(\mathbf{Err}_{\mathrm{III}},\,25)$} 在
(III, 1.4.15) 中，第 92 页倒数第 6 行，把“\emph{有限型}”
替换为“\emph{拟紧}”。

\noindent\textbf{$(\mathbf{Err}_{\mathrm{III}},\,26)$} 在
(III, 2.2.4) 中，第 102 页第 1 行，把“\emph{$\mathscr F$ 与
$\mathscr H$ 的支集在 $Y$ 上固有}”替换为“\emph{
$\operatorname{Im}(\mathscr F\to\mathscr G)$ 的支集在 $Y$ 上固有}”。
在证明中，把第 10 至第 16 行从“由此……”起的文字替换为：

\oldpage[III]{222}
置
$\mathscr G_1=\operatorname{Im}(\mathscr F\to\mathscr G)
=\operatorname{Ker}(\mathscr G\to\mathscr H)$，这是凝聚层。因为有正合列
$0\to\mathscr G_1(n)\to\mathscr G(n)\to\mathscr H(n)$，且函子 $f_*$
左正合，所以序列
$0\to f_*(\mathscr G_1(n))\to f_*(\mathscr G(n))\to f_*(\mathscr H(n))$
对每个 $n$ 都正合；因此只须证明，当 $n$ 足够大时，
$f_*(\mathscr F(n))\to f_*(\mathscr G_1(n))\to0$ 正合。换言之，
可以归约到 $\mathscr H=0$ 且 $\mathscr G$ 的支集在 $Y$ 上固有的
情形；故该支集在 $Z$ 中\emph{闭}（II，5.4.10）。
于是 $\mathscr G'=i_*(\mathscr G)$ 是凝聚 $\mathscr O_Z$-模，并满足
$\mathscr G'|X=\mathscr G$。由（I，9.4.3），存在凝聚
$\mathscr O_Z$-模 $\mathscr F'$，满足 $\mathscr F'|X=\mathscr F$。
此外，因为 $X$ 在 $Z$ 中开且
$\operatorname{Supp}(\mathscr G)\subset X$ 在 $Z$ 中闭，显然可以如下
定义层的满同态 $u':\mathscr F'\to\mathscr G'$，使 $u'|X$ 是给定的
满同态 $u:\mathscr F\to\mathscr G$：取 $u'|U=0$，只要 $U$ 是
$Z$ 中不与 $\operatorname{Supp}(\mathscr G)$ 相交的开集；并取
$u'|U=u|U$，只要 $U\subset X$ 是开集。由此，如
（2.2.2）证明开头那样，可以归约到 $Y$ 仿射的情形；于是须证明当
$n$ 足够大时，同态
$\Gamma(u):\Gamma(X,\mathscr F(n))\to\Gamma(X,\mathscr G(n))$
满射。现有交换图
\[
\xymatrix@C=35mm@R=12mm{
\Gamma(Z,\mathscr F'(n)) \ar[r]^{\Gamma(u')} \ar[d]_v &
\Gamma(Z,\mathscr G'(n)) \ar[d]^w \\
\Gamma(X,\mathscr F(n)) \ar[r]_{\Gamma(u)} & \Gamma(X,\mathscr G(n))
}
\]
其中 $v$、$w$ 是限制同态。$\mathscr G'$ 的定义还表明 $w$
\emph{双射}；另一方面，由（2.2.3），$\Gamma(u')$ 在 $n$ 足够大时
\emph{满射}，所以 $\Gamma(u)$ 也满射。

\noindent\textbf{$(\mathbf{Err}_{\mathrm{III}},\,27)$} 在
(III, 2.2.5) 中，第 102 页倒数第 7 行后增补：

(iii) 在（2.2.4）关于 $X$、$Y$、$f$、$\mathscr L$ 的假设下，
显然，若 $\mathscr H$ 是凝聚 $\mathscr O_X$-模且其支集在 $Y$ 上
\emph{固有}，则类似于（2.2.4）的论证表明，存在整数 $N$，使得对
$n\geq N$，有 $\mathrm R^if_*(\mathscr H(n))=0$ 对每个 $i>0$
成立。由上同调正合列可知，若
$u:\mathscr F\to\mathscr G$ 是凝聚 $\mathscr O_X$-模的同态，而且
$\operatorname{Ker}(u)$ 与 $\operatorname{Coker}(u)$ 的支集都在
$Y$ 上\emph{固有}，则存在 $N$，使得对 $n\geq N$，相应同态
$\mathrm R^if_*(\mathscr F(n))\to\mathrm R^if_*(\mathscr G(n))$
对每个 $i>0$ 都\emph{双射}。

\noindent\textbf{$(\mathbf{Err}_{\mathrm{III}},\,28)$} 在
(III, 4.4.9) 中，第 137 页倒数第 17 行，把“\emph{设 $Y$ 是局部
Noether 整预概形}”替换为“\emph{设 $X$、$Y$ 是两个局部 Noether
整预概形}”；第 137 页倒数第 16 行，把“\emph{有限型}”替换为
“\emph{局部有限型}”。证明本质上不变，因为 $f^{-1}(y)$ 在
$k(y)$ 上局部有限型，故仍是离散空间。

\noindent\textbf{$(\mathbf{Err}_{\mathrm{III}},\,29)$} 在
(I, 9.3.4) 中，第 173 页倒数第 19 行，把“\emph{Noether}”
替换为“\emph{拟紧}”；第 173 页倒数第 18、19 行，把
“\emph{凝聚}”替换为“\emph{有限型拟凝聚}”。在证明中归约到

\oldpage[III]{223}
$X=\operatorname{Spec}(A)$、$\mathscr F=\widetilde M$ 的情形，其中
$M$ 是有限型 $A$-模；以及 $\mathscr J=\widetilde{\mathfrak J}$
的情形，其中 $\mathfrak J$ 是 $A$ 的有限型理想；其余论证不变。

\noindent\textbf{$(\mathbf{Err}_{\mathrm{III}},\,30)$} 在
(I, 9.3.5) 中，把第 173 页倒数第 1 至第 7 行替换为：

\emph{命题（9.3.5）。——设 $X$ 是预概形，$\mathscr F$ 是有限型
拟凝聚 $\mathscr O_X$-模。则存在闭子预概形 $Y$（属于 $X$），其底空间等于
$\operatorname{Supp}(\mathscr F)$；并存在有限型拟凝聚
$\mathscr O_Y$-模 $\mathscr G$，使得若 $j:Y\to X$ 是典范注入，
则 $\mathscr F$ 同构于 $j_*(\mathscr G)$。}

只须证明理想 $\mathscr J$（$\mathscr O_X$ 的理想，即 $\mathscr F$ 的
\emph{零化理想}）是\emph{拟凝聚的}。于是取 $Y$ 为 $X$ 的闭子预概形，
它由 $\mathscr J$ 定义（I，4.1.2）；因为
$\mathscr J\mathscr F=0$，$\mathscr F$ 是
$(\mathscr O_X/\mathscr J)$-模，并可取
$\mathscr G=j^*(\mathscr F)$。为说明 $\mathscr J$ 拟凝聚，可
（因问题是局部的）归约到 $X=\operatorname{Spec}(A)$、
$\mathscr F=\widetilde M$ 的情形，其中 $M$ 是 $A$-模，由有限多个元素
$x_i\ (1\leq i\leq r)$ 生成；此时 $\mathscr J$ 是诸
$x_i$ 的零化理想之交。而 $x_i$ 的零化理想是同态
$\mathscr O_X\to\mathscr F$ 的核，该层同态对应于同态
$s\mapsto sx_i$（从 $A$ 到 $M$）；故它确是拟凝聚理想（I，4.1.1），而这类理想的任意
有限交仍是拟凝聚理想（I，1.3.10）。
\end{document}
